% ============================================================
%  Theory_Trading.tex
%  Trading Encyclopedia — From Market Microstructure to
%  Quantitative Strategies
%  Derived from IMC Prosperity 4 Competition Codebase
%  Compiled with: pdflatex (Overleaf-compatible)
% ============================================================

\documentclass[12pt,a4paper,openany]{report}

% ── Encoding & fonts ─────────────────────────────────────────
\usepackage[T1]{fontenc}
\usepackage[utf8]{inputenc}
\usepackage{lmodern}
\usepackage{microtype}

% ── Page geometry ────────────────────────────────────────────
\usepackage[top=2.5cm,bottom=2.5cm,left=3cm,right=2.5cm]{geometry}

% ── Mathematics ──────────────────────────────────────────────
\usepackage{amsmath,amssymb,amsthm,mathtools}
\usepackage{bm}          % bold math
\usepackage{nicefrac}    % inline fractions
\usepackage{cancel}      % strike-through in math

% ── Tables ───────────────────────────────────────────────────
\usepackage{booktabs}
\usepackage{longtable}
\usepackage{multirow}
\usepackage{array}
\usepackage{tabularx}

% ── Graphics & colour ────────────────────────────────────────
\usepackage[dvipsnames,svgnames,table]{xcolor}
\usepackage{graphicx}
\usepackage{float}
\usepackage{caption}
\usepackage{subcaption}

% ── TikZ & PGFPlots ──────────────────────────────────────────
\usepackage{tikz}
\usepackage{pgfplots}
\pgfplotsset{compat=1.18}
\usetikzlibrary{arrows.meta,shapes,positioning,calc,decorations.pathreplacing,patterns}

% ── Boxes & frames ───────────────────────────────────────────
\usepackage[most]{tcolorbox}
\tcbuselibrary{skins,breakable,theorems}

% ── Code listings ────────────────────────────────────────────
\usepackage{listings}
\usepackage[ruled,vlined,linesnumbered]{algorithm2e}

% ── Miscellaneous ────────────────────────────────────────────
\usepackage{enumitem}
\setlength{\parskip}{5pt plus 1pt minus 1pt}
\setlength{\parindent}{0pt}
\usepackage{mdframed}
\usepackage{siunitx}
\usepackage{hyperref}
\usepackage{cleveref}
\usepackage{titlesec}
% Remove the 50pt top gap that the report class inserts before every \chapter
\titleformat{\chapter}[display]
  {\normalfont\huge\bfseries}{\chaptertitlename\ \thechapter}{10pt}{\Huge}
\titlespacing*{\chapter}{0pt}{0pt}{20pt}

% ============================================================
%  COLOURS
% ============================================================
\definecolor{DefBlue}{RGB}{0,70,127}
\definecolor{ExGreen}{RGB}{0,100,0}
\definecolor{WarnRed}{RGB}{180,0,0}
\definecolor{TakeawayOrange}{RGB}{180,90,0}
\definecolor{MathBg}{RGB}{242,248,255}
\definecolor{ExBg}{RGB}{242,255,242}
\definecolor{WarnBg}{RGB}{255,242,242}
\definecolor{TakeawayBg}{RGB}{255,250,215}
\definecolor{KeyBg}{RGB}{228,238,255}
\definecolor{CodeBg}{RGB}{248,248,248}
\definecolor{DarkGray}{RGB}{50,50,50}

% ============================================================
%  TCOLORBOX STYLES
% ============================================================
\newtcolorbox{definitionbox}[1]{
  enhanced,
  colback=MathBg, colframe=DefBlue, arc=4pt,
  title={\textbf{Definition: #1}},
  colbacktitle=DefBlue, coltitle=white,
  fonttitle=\bfseries, toptitle=3pt, bottomtitle=3pt,
  top=6pt, bottom=6pt
}

\newtcolorbox{examplebox}[1]{
  enhanced,
  colback=ExBg, colframe=ExGreen, arc=4pt,
  title={\textbf{Worked Example: #1}},
  colbacktitle=ExGreen, coltitle=white,
  fonttitle=\bfseries, toptitle=3pt, bottomtitle=3pt,
  top=6pt, bottom=6pt
}

\newtcolorbox{warningbox}[1]{
  enhanced,
  colback=WarnBg, colframe=WarnRed, arc=4pt,
  title={\textbf{Common Misconception: #1}},
  colbacktitle=WarnRed, coltitle=white,
  fonttitle=\bfseries, toptitle=3pt, bottomtitle=3pt,
  top=6pt, bottom=6pt
}

\newtcolorbox{takeawaybox}[1]{
  enhanced,
  colback=TakeawayBg, colframe=TakeawayOrange, arc=4pt,
  title={\textbf{Key Takeaways: #1}},
  colbacktitle=TakeawayOrange, coltitle=white,
  fonttitle=\bfseries, toptitle=3pt, bottomtitle=3pt,
  top=6pt, bottom=6pt
}

\newtcolorbox{keyconceptbox}[1]{
  enhanced,
  colback=KeyBg, colframe=DefBlue!60, arc=4pt,
  title={\textbf{#1}},
  colbacktitle=DefBlue!60, coltitle=DefBlue,
  fonttitle=\bfseries, toptitle=3pt, bottomtitle=3pt,
  top=4pt, bottom=4pt
}

\newtcolorbox{formulabox}{
  enhanced,
  colback=MathBg, colframe=DefBlue!40, arc=4pt,
  top=4pt, bottom=4pt
}

% ============================================================
%  LISTINGS STYLE
% ============================================================
\lstset{
  language=Python,
  backgroundcolor=\color{CodeBg},
  basicstyle=\ttfamily\small,
  keywordstyle=\color{DefBlue}\bfseries,
  commentstyle=\color{gray}\itshape,
  stringstyle=\color{WarnRed!80!black},
  numbers=left, numberstyle=\tiny\color{gray},
  numbersep=8pt, stepnumber=1,
  breaklines=true, breakatwhitespace=false,
  frame=single, framesep=4pt,
  xleftmargin=20pt, xrightmargin=4pt,
  showstringspaces=false,
  tabsize=4
}

% ============================================================
%  THEOREM ENVIRONMENTS
% ============================================================
\newtheorem{theorem}{Theorem}[chapter]
\newtheorem{lemma}[theorem]{Lemma}
\newtheorem{proposition}[theorem]{Proposition}
\newtheorem{corollary}[theorem]{Corollary}
\theoremstyle{definition}
\newtheorem{remark}{Remark}[chapter]

% ============================================================
%  CUSTOM COMMANDS
% ============================================================
\newcommand{\E}{\mathbb{E}}
\newcommand{\Var}{\operatorname{Var}}
\newcommand{\Cov}{\operatorname{Cov}}
\newcommand{\Corr}{\operatorname{Corr}}
\newcommand{\R}{\mathbb{R}}
\newcommand{\N}{\mathbb{N}}
\newcommand{\half}{\tfrac{1}{2}}
\newcommand{\abs}[1]{\left|#1\right|}
\newcommand{\norm}[1]{\left\|#1\right\|}
\newcommand{\floor}[1]{\left\lfloor#1\right\rfloor}
\newcommand{\ceil}[1]{\left\lceil#1\right\rceil}
\newcommand{\dd}{\,\mathrm{d}}
\newcommand{\BM}{B}   % Brownian motion symbol

% ============================================================
%  HYPERREF SETUP
% ============================================================
\hypersetup{
  colorlinks=true,
  linkcolor=DefBlue,
  urlcolor=DefBlue!80!black,
  citecolor=gray,
  pdftitle={Trading Encyclopedia: From Market Microstructure to Quantitative Strategies},
  pdfsubject={Algorithmic Trading Reference},
  pdfkeywords={market microstructure, market making, mean reversion, time series, algo trading}
}

% ============================================================
%  TITLE PAGE
% ============================================================
\title{
  \vspace{1.5cm}
  \rule{\linewidth}{1.5pt}\\[0.6cm]
  {\Huge\bfseries Quantitative Trading Encyclopedia}\\[0.5cm]
  {\Large From Market Microstructure\\to Quantitative Strategies}\\[0.4cm]
  \rule{\linewidth}{1.5pt}\\[1cm]
  {\large Derived from the IMC Prosperity 4 Competition Codebase}\\[0.6cm]
  {\normalsize A self-contained reference for mathematicians and programmers\\
  with zero prior finance background.\\[0.3cm]
  Builds every intuition from first principles.}
}
\author{}
\date{IMC Prosperity 4 --- 2026}

% ============================================================
\begin{document}
% ============================================================

\maketitle
\thispagestyle{empty}
\newpage

\begingroup
\titlespacing*{\chapter}{0pt}{0pt}{0pt}
\tableofcontents
\endgroup
\newpage

% ============================================================
%
%  PART I — MARKET MICROSTRUCTURE & MECHANICS
%
% ============================================================
\part{Market Microstructure \& Mechanics}

\chapter{How Financial Markets Work}

\section{What Is a Market?}

Before writing a single line of trading code, you must understand \emph{what a market actually is}. The word ``market'' is thrown around loosely, but in the context of algorithmic trading it has a precise meaning.

\begin{definitionbox}{Financial Market}
A \textbf{financial market} is a centralised or decentralised venue where buyers and sellers exchange \emph{financial instruments} (stocks, bonds, commodities, currencies, derivatives, etc.) at prices determined by the interaction of supply and demand.  Every trade requires exactly one buyer and one seller; at the moment of exchange they both agree on price and quantity.
\end{definitionbox}

\noindent\textbf{Analogy.} Think of a farmers' market. Vendors bring goods (sellers), shoppers bring cash (buyers). A price is struck when both sides agree. A financial exchange is the same idea, except millions of trades happen per second and the ``goods'' are abstract financial claims.

\subsection{The Role of Prices}

Prices aggregate the private information of all participants. If someone discovers that a company is about to release a blockbuster drug, they buy shares, pushing the price up and signalling this information to the rest of the market. This \emph{price discovery} mechanism is one of the most important functions of financial markets.

\subsection{Products and Denominations}

In the Prosperity exchange, all products are denominated in \textbf{XIREC} (the simulation currency). In real markets the denomination is typically a national currency (USD, EUR, etc.). Every trade of product $P$ at price $p$ for quantity $q$ results in a cash flow of:
\begin{equation}
  \text{Cash flow} = -p \cdot q \quad (\text{buyer pays}), \qquad +p \cdot q \quad (\text{seller receives})
\end{equation}
and a position change:
\begin{equation}
  \Delta \text{position}_P = +q \quad (\text{buyer}), \qquad -q \quad (\text{seller})
\end{equation}

\section{Exchange Architecture}

\begin{definitionbox}{Exchange}
An \textbf{exchange} is a regulated marketplace that centralises order flow, provides price discovery, enforces rules (position limits, cancellation policies), and operates a matching engine that executes trades when compatible buy and sell orders meet.
\end{definitionbox}

\noindent\textbf{Key exchange functions:}
\begin{enumerate}[label=\arabic*.]
  \item \textbf{Order collection}: Participants submit orders to a central queue.
  \item \textbf{Matching}: The matching engine pairs compatible buy and sell orders.
  \item \textbf{Price dissemination}: The exchange broadcasts prices and quantities to all participants in real time.
  \item \textbf{Settlement}: Ownership changes are recorded and cash/assets are transferred.
  \item \textbf{Rule enforcement}: Position limits, reporting requirements, circuit breakers.
\end{enumerate}

\subsection{Over-the-Counter (OTC) vs.\ Exchange}

Not all trading happens on formal exchanges. In \textbf{over-the-counter (OTC)} markets, two parties negotiate directly (e.g.\ interest-rate swaps between banks). Exchange markets offer:
\begin{itemize}
  \item Central counterparty clearing (reduces default risk)
  \item Transparency (public prices)
  \item Standardised contracts
\end{itemize}
OTC markets offer flexibility but less transparency. The Prosperity simulation models an exchange environment.

\section{Market Participants}

Understanding \emph{who} trades and \emph{why} is essential for building strategies. There are three archetypal participant types:

\begin{center}
\begin{tabular}{llll}
\toprule
\textbf{Type} & \textbf{Goal} & \textbf{Horizon} & \textbf{Information edge} \\
\midrule
Market Maker  & Earn bid-ask spread & Milliseconds--seconds & Order flow \\
Informed Trader & Profit from private info & Minutes--days & Fundamental / news \\
Noise Trader  & Hedging / liquidity need & Any & None (uninformed) \\
\bottomrule
\end{tabular}
\end{center}

The distinction between informed and noise traders is critical to market microstructure theory. When a market maker quotes both a bid and an ask, they face \emph{adverse selection}: informed traders will trade against them precisely when it is most costly (the market maker sells just before the price rises). This is explored in Chapter~5.

\section{The Trading Day and Timestamps}

In Prosperity, the exchange operates in discrete \textbf{iterations} (also called \emph{timestamps}). Each iteration:
\begin{enumerate}
  \item The exchange sends a \texttt{TradingState} to your algorithm.
  \item Your algorithm executes its \texttt{run()} method and returns orders.
  \item The matching engine immediately executes any crossing orders (your orders vs.\ bot quotes).
  \item Any unfilled orders are exposed to bots; bots may or may not take them.
  \item The iteration ends; unfilled orders are cancelled.
\end{enumerate}

In real markets, time is continuous and measured in nanoseconds. In Prosperity, it is measured in integer timestamps spaced 100 units apart. One ``day'' covers timestamps $0, 100, 200, \ldots, 99900$ (1\,000 iterations).

\begin{takeawaybox}{Chapter 1}
\begin{itemize}
  \item A market pairs buyers and sellers; every trade requires agreement on price and quantity.
  \item Exchanges centralise order flow, provide matching, and enforce rules.
  \item Participants are market makers (earn spread), informed traders (exploit information), and noise traders (hedging/liquidity needs).
  \item In Prosperity, trading is discrete: one iteration = one call to your \texttt{run()} method.
\end{itemize}
\end{takeawaybox}

% ──────────────────────────────────────────────────────────────────────────────

\chapter{The Order Book}

\section{What Is an Order Book?}

The \textbf{order book} is the central data structure of any exchange. It is a real-time record of all outstanding (unexecuted) limit orders for a given product, organised by price level.

\begin{definitionbox}{Order Book}
The \textbf{order book} for product $P$ at time $t$ is a function
\[
  \mathcal{B}_t : \mathbb{Z} \to \mathbb{Z}
\]
mapping each price level $p$ to a signed aggregate quantity $q$, where $q > 0$ means buy-side interest (bids) and $q < 0$ means sell-side interest (asks/offers).
\end{definitionbox}

In practice the book is split into two halves:
\begin{align*}
  \text{Bid side (buy orders):} \quad & \{(p_i^B, q_i^B)\}_{i=1}^{N_B}, \quad p_1^B > p_2^B > \cdots, \quad q_i^B > 0 \\
  \text{Ask side (sell orders):} \quad & \{(p_j^A, q_j^A)\}_{j=1}^{N_A}, \quad p_1^A < p_2^A < \cdots, \quad q_j^A > 0
\end{align*}
where levels are listed from best (closest to mid) to worst. The exchange \emph{guarantees} $p_1^B < p_1^A$ (the book is uncrossed at rest, because any crossing would immediately trigger a trade).

\section{Reading the Order Book}

\subsection{A Concrete Example}

Consider the EMERALDS order book from the Prosperity data at timestamp 0:

\begin{center}
\begin{tabular}{ccc|c|ccc}
\toprule
\multicolumn{3}{c|}{\textbf{Bid (Buy) Side}} & \textbf{Price} & \multicolumn{3}{c}{\textbf{Ask (Sell) Side}} \\
Qty$_3$ & Qty$_2$ & Qty$_1$ & & Qty$_1$ & Qty$_2$ & Qty$_3$ \\
\midrule
 & 25 & 11 & 9990 & & & \\
 & & 11 & 9992 & & & \\
 & & & \textbf{10000} & & & \\
 & & & 10008 & 11 & & \\
 & & & 10010 & & 25 & \\
\bottomrule
\end{tabular}
\end{center}

Reading this: there are 11 units offered to buy at 9992 (best bid), 25 more at 9990 (second level). On the ask side, 11 units are for sale at 10008 (best ask), 25 more at 10010.

\subsection{Depth and Levels}

\begin{definitionbox}{Order Book Depth}
\textbf{Depth} refers to the total quantity available at or better than a given price level. \textbf{Levels} (or \emph{price levels}) are the distinct prices at which orders rest. In Prosperity, the exchange provides up to 3 levels per side.
\end{definitionbox}

\begin{definitionbox}{Market Depth}
The \textbf{market depth} at price $p$ on the bid side is
\[
  D^B(p) = \sum_{i:\, p_i^B \geq p} q_i^B
\]
and on the ask side:
\[
  D^A(p) = \sum_{j:\, p_j^A \leq p} q_j^A
\]
Larger depth means you can trade larger size without moving the price.
\end{definitionbox}

\section{Mathematical Representation in Code}

In Prosperity's \texttt{datamodel.py}, the order book is:
\begin{lstlisting}[caption={OrderDepth class from datamodel.py}]
class OrderDepth:
    def __init__(self):
        self.buy_orders: Dict[int, int] = {}   # price -> +qty
        self.sell_orders: Dict[int, int] = {}  # price -> -qty (negative!)
\end{lstlisting}

\noindent Note the convention: \texttt{sell\_orders} quantities are \emph{negative} (e.g.\ \texttt{\{10008: -11\}} means 11 units for sale at 10008). This is a common convention in trading systems.

To extract the best bid and ask:
\begin{lstlisting}[caption={Extracting best bid and ask from the order book}]
best_bid = max(od.buy_orders.keys())    # highest buy price
best_ask = min(od.sell_orders.keys())   # lowest sell price
mid = (best_bid + best_ask) / 2
spread = best_ask - best_bid
\end{lstlisting}

\section{Order Book Dynamics}

The order book changes every time:
\begin{itemize}
  \item A new limit order arrives (adds to a level, or creates a new level)
  \item A trade executes (reduces or removes a level)
  \item A cancellation arrives (reduces or removes a level)
\end{itemize}

\subsection{Crossed Book and Execution}

If a new buy order arrives with price $p^B \geq p_1^A$ (the best ask), it \emph{crosses} the book and triggers immediate execution at $p_1^A$ (the resting order's price). After execution, the book returns to uncrossed state.

\begin{examplebox}{Crossing the Book}
Suppose the book has: bid=9992 (qty 11), ask=10008 (qty 11). \\
Your algorithm sends: \texttt{Order("EMERALDS", 10008, 5)} — a buy order at 10008. \\
Since $10008 \geq 10008$ (best ask), your order immediately executes: you buy 5 units at 10008. \\
The ask level reduces to qty 6. Your cash decreases by $5 \times 10008 = 50040$ XIREC. \\
Your position increases by $+5$.
\end{examplebox}

\section{The Limit Order Book as a Queuing System}

Mathematically, the order book can be modelled as a pair of queues. For market microstructure research, the dynamics of the limit order book are often modelled as a Poisson process:
\begin{itemize}
  \item Orders arrive at rate $\lambda$ (Poisson distributed)
  \item Cancellations occur at rate $\mu$
  \item Trades occur at rate $\nu$
\end{itemize}
This gives rise to the \textbf{Markovian order book model} (Cont, Stoikov, Talreja 2010), though in practice Prosperity's bots follow deterministic rules.

\begin{takeawaybox}{Chapter 2}
\begin{itemize}
  \item The order book collects all resting limit orders, split into bid (buy) and ask (sell) sides.
  \item Best bid $< $ best ask always; a crossing immediately triggers a trade.
  \item Depth measures how much can be traded at or better than a given price.
  \item In Prosperity, \texttt{sell\_orders} quantities are negative; normalise with $-q$ to get volume.
\end{itemize}
\end{takeawaybox}

% ──────────────────────────────────────────────────────────────────────────────

\chapter{The Bid-Ask Spread}

\section{Definition and Intuition}

\begin{definitionbox}{Bid-Ask Spread}
The \textbf{bid-ask spread} (or \emph{quoted spread}) is the difference between the best ask and the best bid:
\[
  S_t = p_t^A - p_t^B \geq 0
\]
It represents the minimum transaction cost for a round-trip trade (buy then sell immediately) and the gross profit margin of the market maker who quotes both sides.
\end{definitionbox}

\noindent\textbf{Analogy.} A currency exchange booth at an airport quotes EUR/USD at 1.03 (buy) / 1.07 (sell). If you buy \$100 and immediately sell them back, you lose $\$(1.07 - 1.03) \times 100 = \$4$. The booth earns \$4. This is the spread — it compensates the booth for the risk and cost of providing immediate liquidity.

\subsection{Numerical Examples from Prosperity}

\begin{center}
\begin{tabular}{lccccc}
\toprule
Product & Best Bid & Best Ask & Spread & Mid & Spread/Mid (\%) \\
\midrule
EMERALDS & 9992 & 10008 & 16 & 10000 & 0.16\% \\
TOMATOES & 4993 & 5007  & 14 & 5000  & 0.28\% \\
\bottomrule
\end{tabular}
\end{center}

\section{Why Does the Spread Exist?}

The spread has three economic components, identified by the \textbf{Glosten-Milgrom (1985)} model:

\begin{keyconceptbox}{Three Components of the Bid-Ask Spread}
\begin{enumerate}[label=\arabic*.]
  \item \textbf{Order processing costs}: The administrative cost of handling an order. In modern electronic markets this is near zero.
  \item \textbf{Inventory risk}: A market maker who buys has long inventory and is exposed to price falls. Compensation for bearing this directional risk.
  \item \textbf{Adverse selection}: The risk that the counterparty is \emph{informed} (knows something the market maker doesn't). If a trader buys aggressively, it may signal the price is about to rise — making the sale a loss for the market maker.
\end{enumerate}
\end{keyconceptbox}

Formally, let $p^*$ be the true (unobservable) asset value. The market maker quotes:
\begin{align}
  p^A &= p^* + \alpha + \beta \\
  p^B &= p^* - \alpha - \beta
\end{align}
where $\alpha$ compensates for inventory risk and $\beta$ compensates for adverse selection. The total spread is $S = 2(\alpha + \beta)$.

\section{Quoted vs.\ Effective Spread}

\begin{definitionbox}{Effective Spread}
The \textbf{effective spread} measures the actual cost paid by a liquidity taker, accounting for potential price improvement:
\[
  S^{\text{eff}}_t = 2 \, d_t \, (p_t^{\text{trade}} - m_t)
\]
where $d_t \in \{+1, -1\}$ is the \emph{trade direction} ($+1$ = buyer-initiated, $-1$ = seller-initiated) and $m_t = \nicefrac{(p^A_t + p^B_t)}{2}$ is the mid-price. The factor of 2 makes it comparable to the quoted spread $S = p^A - p^B$.
\end{definitionbox}

If there is no price improvement (trades always occur at the quoted bid or ask), then $S^{\text{eff}} = S^{\text{quoted}}$. In practice, $S^{\text{eff}} \leq S^{\text{quoted}}$ because large resting orders may offer better prices, or the exchange matches at more favourable prices.

\section{Realised Spread and Price Impact}

After a trade, the mid-price often moves \emph{against} the liquidity taker. This is the \textbf{price impact}. The market maker's actual realised profit is:

\begin{align}
  \text{Realised spread} &= 2 \, d_t \, (p_t^{\text{trade}} - m_{t+\tau}) \\
  \text{Price impact} &= 2 \, d_t \, (m_{t+\tau} - m_t)
\end{align}
where $\tau$ is some horizon (e.g.\ 5 minutes). The decomposition is:
\[
  S^{\text{eff}} = \text{Realised spread} + \text{Price impact}
\]
If the price impact exceeds the half-spread, the market maker loses money on that trade (adverse selection event).

\section{Spread Regimes}

In the Prosperity data, TOMATOES exhibits two distinct spread regimes — an important real-world phenomenon:

\begin{keyconceptbox}{Spread Regimes in TOMATOES}
\begin{itemize}
  \item \textbf{Tight regime} (spread $\leq 13$, $\approx 55\%$ of timestamps): liquid market, AR(1) coefficient $\approx -0.30$ (mean-reverting).
  \item \textbf{Wide regime} (spread $= 14$, $\approx 45\%$ of timestamps): less liquid market, AR(1) $\approx 0$ (random walk).
\end{itemize}
\end{keyconceptbox}

Spread regimes arise because:
\begin{itemize}
  \item \textbf{Intraday liquidity cycles}: markets are more liquid near the open/close.
  \item \textbf{News events}: uncertainty widens spreads.
  \item \textbf{Bot behaviour}: in Prosperity, different bot configurations produce different spreads at different times.
\end{itemize}

\subsection{Measuring Spread Regimes}

A simple threshold classifier:
\begin{equation}
  \text{Regime}(t) = \begin{cases} \text{Tight} & \text{if } S_t \leq S^*, \\ \text{Wide} & \text{if } S_t > S^*. \end{cases}
\end{equation}
For more sophisticated detection, use a \textbf{Hidden Markov Model} (see Chapter~28).

\section{Effective Spread as a Trading Cost}

For a strategy that takes liquidity (sends market orders or aggressively priced limit orders), the effective spread is a direct cost. If your strategy turns over position $Q$ per unit time, the annual drag is:
\begin{equation}
  \text{Annual cost} = Q \times \bar{S}^{\text{eff}} \times T
\end{equation}
where $T$ is the number of trading periods per year and $\bar{S}^{\text{eff}}$ is the average effective half-spread. This is why market-making strategies \emph{earn} the spread, while momentum strategies \emph{pay} it.

\begin{takeawaybox}{Chapter 3}
\begin{itemize}
  \item The spread $S = p^A - p^B$ is the gross profit per round-trip for market makers and the cost for liquidity takers.
  \item It exists due to order processing costs, inventory risk, and adverse selection.
  \item Effective spread accounts for price improvement; realised spread accounts for post-trade drift.
  \item Real markets exhibit spread regimes (tight/wide), driven by liquidity cycles and information events.
\end{itemize}
\end{takeawaybox}

% ──────────────────────────────────────────────────────────────────────────────

\chapter{Order Types and Execution}

\section{The Three Fundamental Order Types}

\begin{definitionbox}{Market Order}
A \textbf{market order} instructs the exchange to execute immediately at the best available price, regardless of what that price is. It \emph{takes} liquidity from the book. Guaranteed execution, uncertain price.
\end{definitionbox}

\begin{definitionbox}{Limit Order}
A \textbf{limit order} specifies a maximum buy price (or minimum sell price). It executes \emph{only} if a counterparty is willing to trade at that price or better. If not immediately matched, it rests in the book. Guaranteed price, uncertain execution.
\end{definitionbox}

\begin{definitionbox}{Stop Order}
A \textbf{stop order} (stop-loss / stop-limit) becomes active only when the market price reaches a trigger level. After activation, it behaves as a market or limit order. Used for risk management.
\end{definitionbox}

\noindent In Prosperity, only limit orders are used. A limit order at a price that \emph{crosses} the current best ask (buy) or best bid (sell) will execute immediately — effectively behaving like a market order.

\section{Passive vs.\ Aggressive Orders}

\begin{keyconceptbox}{Passive vs.\ Aggressive}
\begin{itemize}
  \item \textbf{Passive (maker) order}: A limit order that \emph{adds} liquidity to the book by resting at a price that does not immediately execute. It ``makes'' the market. The trader is the \emph{market maker}.
  \item \textbf{Aggressive (taker) order}: A limit or market order that \emph{removes} liquidity by immediately executing against a resting order. The trader is the \emph{liquidity taker}.
\end{itemize}
\end{keyconceptbox}

\subsection{In the Prosperity Context}

In our \texttt{trader.py}:
\begin{itemize}
  \item \textbf{Taking}: Orders sent at a price $\geq$ best ask (for buys) or $\leq$ best bid (for sells). Executed immediately. Used to capture ``free edge'' when bot quotes are better than fair value.
  \item \textbf{Making}: Orders sent at a price strictly inside the bot spread (e.g.\ at 9993 when best ask is 10008). These rest in the book and wait for bots to fill them.
\end{itemize}

\begin{lstlisting}[caption={Taking vs.\ making in EmeraldsStrategy.orders()}]
# TAKING: Buy if bot is selling below fair value
for ask in sorted(od.sell_orders.keys()):
    if ask > self.FAIR_VALUE or buy_cap <= 0:
        break
    qty = min(-od.sell_orders[ask], buy_cap)
    orders.append(Order(self.PRODUCT, ask, qty))   # aggressive

# MAKING: Passive bid inside the bot spread
if buy_cap > 0:
    orders.append(Order(self.PRODUCT, self.PASSIVE_BUY, buy_cap))  # passive
\end{lstlisting}

\section{Order Properties}

Every order in Prosperity has:
\begin{enumerate}
  \item \textbf{Symbol}: the product being traded (e.g.\ \texttt{"EMERALDS"})
  \item \textbf{Price}: maximum buy price or minimum sell price (integer)
  \item \textbf{Quantity}: positive = buy, negative = sell
\end{enumerate}

\subsection{Validity: Time-in-Force}

In real markets, orders have a \textbf{time-in-force} (TIF) parameter:
\begin{center}
\begin{tabular}{ll}
\toprule
\textbf{TIF} & \textbf{Meaning} \\
\midrule
GTC (Good Till Cancelled) & Rests until filled or cancelled \\
IOC (Immediate or Cancel) & Fill what you can, cancel remainder \\
FOK (Fill or Kill) & Fill entire quantity or cancel \\
GTD (Good Till Date) & Expires at a specified time \\
\bottomrule
\end{tabular}
\end{center}
In Prosperity, all orders are effectively \textbf{GTT} (Good Till next Timestamp): unfilled orders are cancelled at the end of each iteration.

\section{Price-Time Priority}

When multiple orders rest at the same price level, the exchange must decide which fills first. The universal rule on modern exchanges is:

\begin{definitionbox}{Price-Time Priority}
Orders are matched in \textbf{price-time priority}:
\begin{enumerate}
  \item \textbf{Price}: The order with the most aggressive price executes first (highest bid, lowest ask).
  \item \textbf{Time}: Among orders at the same price, the \emph{oldest} order executes first (FIFO — first in, first out).
\end{enumerate}
\end{definitionbox}

This incentivises early order submission at competitive prices. Market makers compete to be first in the queue at a price level.

\begin{examplebox}{Price-Time Priority}
The ask side has: 10007 (qty 3, arrived at $t=50$), 10007 (qty 5, arrived at $t=80$), 10008 (qty 11). \\
An aggressive buy order for 6 units arrives. \\
Execution: 3 units at 10007 (oldest order exhausted), then 3 units at 10007 (second order, leaving 2 remaining), then 0 units at 10008. \\
Total: 6 units executed at 10007.
\end{examplebox}

\section{Slippage and Market Impact}

\begin{definitionbox}{Slippage}
\textbf{Slippage} is the difference between the expected execution price and the actual execution price, caused by the order consuming multiple price levels:
\[
  \text{Slippage} = \bar{p}^{\text{filled}} - p^{\text{expected}}
\]
where $\bar{p}^{\text{filled}} = \frac{\sum_i p_i q_i}{\sum_i q_i}$ is the volume-weighted average fill price.
\end{definitionbox}

\begin{examplebox}{Slippage Calculation}
You send a buy order for 30 units at price 10015 (willing to pay up to 10015). \\
The ask side is: 10008 (qty 11), 10010 (qty 25). \\
Execution: 11 units at 10008, then 19 units at 10010 (total 30 units, up to your limit). \\
$\bar{p} = \frac{11 \times 10008 + 19 \times 10010}{30} = \frac{110088 + 190190}{30} = \frac{300278}{30} = 10009.27$. \\
Slippage vs.\ best ask: $10009.27 - 10008 = 1.27$ ticks.
\end{examplebox}

\section{Position and Position Limits}

\begin{definitionbox}{Position}
A \textbf{position} in product $P$ is the net quantity held:
\[
  \text{pos}_t^P = \text{pos}_{t-1}^P + \sum_{\text{buys}} q_i - \sum_{\text{sells}} q_j
\]
A positive position is \textbf{long} (you own the asset, profit if price rises). A negative position is \textbf{short} (you owe the asset, profit if price falls).
\end{definitionbox}

\begin{definitionbox}{Position Limit}
A \textbf{position limit} $L$ bounds the absolute position: $-L \leq \text{pos} \leq +L$. The \textbf{remaining buy capacity} is $L - \text{pos}$ and the \textbf{remaining sell capacity} is $L + \text{pos}$.
\end{definitionbox}

In Prosperity Round 0, both EMERALDS and TOMATOES have $L = 80$.

\begin{warningbox}{All-or-Nothing Rejection}
In Prosperity, if the \emph{aggregated} quantity of all your buy orders would push your position above $+L$, \textbf{all} buy orders for that product are rejected — not just the excess. Always compute \texttt{buy\_cap = L - pos} before sizing your orders.
\end{warningbox}

\begin{takeawaybox}{Chapter 4}
\begin{itemize}
  \item Market orders: guaranteed execution, uncertain price. Limit orders: guaranteed price, uncertain execution.
  \item Passive orders add liquidity; aggressive orders remove it. Market makers prefer passive orders.
  \item Price-time priority: best price first, oldest at that price first (FIFO).
  \item Slippage arises when a large order consumes multiple price levels.
  \item Position limits cap your exposure; exceeding them causes complete order rejection in Prosperity.
\end{itemize}
\end{takeawaybox}

% ──────────────────────────────────────────────────────────────────────────────

\chapter{Market Makers and Liquidity}

\section{The Market Maker's Role}

A \textbf{market maker} continuously quotes both a bid and an ask price, standing ready to buy or sell at those prices on demand. In doing so, they provide \emph{immediacy}: any participant who wants to trade can do so without waiting for a natural counterparty.

\begin{keyconceptbox}{The Market Maker's Business Model}
\begin{enumerate}[label=\arabic*.]
  \item Post bid $p^B$ and ask $p^A$ with $p^B < p^A$.
  \item Buy from sellers at $p^B$ and sell to buyers at $p^A$.
  \item Profit per round-trip: $p^A - p^B = S$ (the spread).
  \item Risk: inventory accumulation if one side dominates.
\end{enumerate}
\end{keyconceptbox}

\textbf{Analogy.} A used-car dealer buys cars at \$9,000 and sells at \$11,000. Profit = \$2,000 per car if they buy and sell equal volumes. Risk: they're stuck with depreciating inventory if the market turns against them.

\section{Liquidity: Defined}

\begin{definitionbox}{Liquidity}
\textbf{Liquidity} is the ability to trade a desired quantity quickly, at a price close to the current market price, with minimal price impact. A market is liquid if:
\begin{itemize}
  \item The spread is tight (low transaction cost)
  \item Depth is large (big orders move prices little)
  \item Market impact is small
  \item Recovery from large trades is fast
\end{itemize}
\end{definitionbox}

Liquidity is not binary; it has several dimensions:
\begin{center}
\begin{tabular}{lll}
\toprule
\textbf{Dimension} & \textbf{Measure} & \textbf{High liquidity = \ldots} \\
\midrule
Width & Bid-ask spread $S$ & Small $S$ \\
Depth & Total book volume within $\delta$ ticks & Large depth \\
Resilience & Time to recover after a large trade & Fast recovery \\
Immediacy & Execution time & Instantaneous \\
\bottomrule
\end{tabular}
\end{center}

\section{Adverse Selection: The Core Market-Making Risk}

Adverse selection occurs when the market maker's counterparty is better informed. Let $\mu$ be the probability that a given trader is \emph{informed}.

\begin{theorem}[Glosten-Milgrom Spread]
In the Glosten-Milgrom (1985) model, where informed traders know the true asset value $V \in \{V_H, V_L\}$ and noise traders trade randomly, the equilibrium bid and ask are:
\begin{align}
  p^A &= \E[V \mid \text{buy arrives}] \\
  p^B &= \E[V \mid \text{sell arrives}]
\end{align}
The spread $p^A - p^B > 0$ exists \emph{solely} because of informed traders ($\mu > 0$).
\end{theorem}

\begin{proof}[Full Derivation]
\textbf{Setup.} Asset value $V \in \{V_H, V_L\}$ with $V_H > V_L$ and prior $P(V = V_H) = q$, so $\E[V] = qV_H + (1-q)V_L$. A fraction $\mu \in [0,1]$ of order flow is \emph{informed}: informed buyers know $V = V_H$ and always buy; informed sellers know $V = V_L$ and always sell. The remaining fraction $(1-\mu)$ is \emph{noise traders} who buy or sell with equal probability $\nicefrac{1}{2}$ regardless of $V$.

\textbf{Step 1: Probability that a buy order arrives.}
\[
  P(\text{buy}) = \mu \cdot q + (1-\mu) \cdot \tfrac{1}{2}
\]
This follows because: with prob.\ $q$ the true value is $V_H$, in which case all informed traders buy; with prob.\ $1-q$ it is $V_L$, in which case no informed trader buys. Noise traders always buy with prob.\ $\nicefrac{1}{2}$.

\textbf{Step 2: Posterior on $V$ given a buy order arrives (Bayes' theorem).}
\begin{align}
  P(V=V_H \mid \text{buy}) &= \frac{P(\text{buy} \mid V=V_H)\cdot q}{P(\text{buy})} \\
  &= \frac{[\mu + (1-\mu)\cdot\nicefrac{1}{2}] \cdot q}{\mu q + (1-\mu)\cdot\nicefrac{1}{2}} \\
  &= \frac{(1 - \nicefrac{\mu}{2})\,q}{\mu q + \nicefrac{(1-\mu)}{2}}
\end{align}
Let $q^+ \equiv P(V=V_H \mid \text{buy})$. Since the numerator is $> q \cdot P(\text{buy})$ iff $\mu > 0$, we have $q^+ > q$: receiving a buy order raises the probability that $V = V_H$.

\textbf{Step 3: Zero-profit ask price.}
A competitive market maker sets the ask equal to expected value after updating on the incoming order (zero expected profit condition):
\[
  p^A = \E[V \mid \text{buy}] = q^+ V_H + (1 - q^+) V_L
\]

\textbf{Step 4: Zero-profit bid price.} By symmetry, a sell order is evidence that $V = V_L$. Let $q^- = P(V=V_H \mid \text{sell}) < q$. Then:
\[
  p^B = \E[V \mid \text{sell}] = q^- V_H + (1 - q^-) V_L < \E[V]
\]

\textbf{Step 5: The spread.}
\begin{align}
  S &= p^A - p^B = (q^+ - q^-)(V_H - V_L)
\end{align}
Since $q^+ > q > q^-$ whenever $\mu > 0$, the spread $S = (q^+ - q^-)(V_H - V_L) > 0$.

\textbf{Symmetric case $q = \nicefrac{1}{2}$:} With equal prior, $q^+ = (1+\mu)/2$ and $q^- = (1-\mu)/2$:
\[
  \boxed{S = (q^+ - q^-)(V_H - V_L) = \mu(V_H - V_L)}
\]

\textbf{General prior $q$:} Substituting the full posteriors $q^+ = \frac{q(1+\mu)}{1+\mu(2q-1)}$ and $q^- = \frac{q(1-\mu)}{1-\mu(2q-1)}$:
\[
  S = \frac{4\mu q(1-q)}{1 - \mu^2(2q-1)^2}(V_H - V_L)
\]
which reduces to $\mu(V_H - V_L)$ at $q=\nicefrac{1}{2}$ and approaches $2(V_H-V_L)$ as $\mu \to 1$ (fully informed market).
\textbf{Key insight:} $S = 0$ iff $\mu = 0$ (no informed traders). The spread is a direct monotone function of the fraction of informed flow. As $\mu \to 1$ (all informed), $S \to 2(V_H - V_L)$ (the full range).

\textbf{Implication for Prosperity:} The tight-spread regime ($S=13$) has less adverse selection than the wide-spread regime ($S=14$). This matches the empirical finding: the tight regime shows a \emph{negative} AR(1) / mean-reverting pattern ($\phi_{\text{stat}} \approx -0.30$, i.e.\ price changes tend to reverse), while the wide regime is noise-dominated (AR(1) $\approx 0$).
\end{proof}

\subsection{Adverse Selection in Practice}

In Prosperity, the bots play a simplified version of informed trading: some bots will aggressively cross the spread whenever your quotes are sufficiently mispriced. If your fair value estimate is wrong, you will be adversely selected: you'll sell just before the price rises, or buy just before it falls.

\textbf{Mitigation strategies:}
\begin{enumerate}[label=\arabic*.]
  \item Accurate fair value estimation (reduces informational disadvantage)
  \item Inventory skewing (reduce size when position is large)
  \item Wider quotes in high-uncertainty regimes (widen spread = charge more for liquidity)
\end{enumerate}

\section{The Market Maker's PnL Decomposition}

A market maker's PnL can be decomposed as:
\begin{equation}
  \text{PnL}_t = \underbrace{\sum_{\text{trades}} \text{signed edge}}_{\text{Trading PnL}} + \underbrace{\text{pos}_t \cdot (p_t - p_0)}_{\text{Mark-to-market}}
\end{equation}
where ``signed edge'' = $\nicefrac{S}{2}$ per trade (positive when providing liquidity), and ``mark-to-market'' = inventory value change.

At position limit $L$ with mid-price drift $\Delta p$:
\begin{equation}
  \text{Expected loss from inventory} = L \cdot |\Delta p|
\end{equation}
Inventory skewing (Chapter~23) reduces this by keeping $|\text{pos}|$ small.

\section{Inventory Management}

The market maker's dilemma: to earn spread you must quote; but quoting accumulates inventory; inventory creates directional risk.

\begin{keyconceptbox}{Inventory Control Levers}
\begin{enumerate}[label=\arabic*.]
  \item \textbf{Quote size}: Reduce quantity when inventory is large.
  \item \textbf{Quote skew}: Shift bid/ask toward the direction that would reduce inventory (lean your quotes).
  \item \textbf{Taking}: Aggressively cross the spread to unwind inventory.
  \item \textbf{Position limit}: Hard cap that prevents inventory from becoming catastrophic.
\end{enumerate}
\end{keyconceptbox}

In \texttt{trader.py}, the inventory skew for TOMATOES is:
\begin{equation}
  \text{skew} = \frac{\text{pos}}{L} \times S_{\text{range}}
\end{equation}
\begin{align}
  p^B_{\text{passive}} &= \floor{FV - \delta - \text{skew}} \\
  p^A_{\text{passive}} &= \ceil{FV + \delta - \text{skew}}
\end{align}
When pos $> 0$ (long), skew $> 0$, shifting both quotes \emph{down}: the bid becomes less attractive (sell less to us) and the ask becomes more attractive (buy from us) — pushing inventory back toward zero.

\begin{takeawaybox}{Chapter 5}
\begin{itemize}
  \item Market makers provide immediacy by quoting both bid and ask, earning the spread.
  \item Adverse selection is the risk that counterparties are informed; the Glosten-Milgrom model shows this creates the spread.
  \item Inventory risk arises when one side of the market dominates; mitigated by skewing quotes.
  \item PnL = trading edge (earned spread) + mark-to-market (inventory drift).
\end{itemize}
\end{takeawaybox}

% ──────────────────────────────────────────────────────────────────────────────

\chapter{Price Formation and Fair Value}

\section{The Mid-Price}

\begin{definitionbox}{Mid-Price}
The \textbf{mid-price} is the arithmetic average of the best bid and best ask:
\[
  m_t = \frac{p_t^A + p_t^B}{2}
\]
It is the most common proxy for the ``true'' unobservable asset value. All strategies in \texttt{trader.py} use the mid-price as the starting point for fair value estimation.
\end{definitionbox}

Note: if the spread is an odd number of ticks, the mid-price is a half-integer. In Prosperity, EMERALDS has spread 16 (even), so mid = 10000.0 exactly. TOMATOES has spread 13--14, giving half-integer mids. This matters when choosing \texttt{math.floor} vs.\ \texttt{math.ceil} for integer-priced limit orders (see Chapter~32).

\section{VWAP: Volume-Weighted Average Price}

\begin{definitionbox}{VWAP}
The \textbf{Volume-Weighted Average Price} over a window $[t_1, t_2]$ is:
\[
  \text{VWAP} = \frac{\sum_{i: t_i \in [t_1,t_2]} p_i \cdot q_i}{\sum_{i: t_i \in [t_1,t_2]} q_i}
\]
where the sum is over all trades. VWAP weights each trade price by its volume, giving a better estimate of the ``average'' execution price than a simple time average.
\end{definitionbox}

VWAP is used as:
\begin{enumerate}
  \item An execution benchmark (``did I beat VWAP?'')
  \item A fair value estimate (weighted by liquidity)
  \item A signal (price above VWAP = bullish momentum intraday)
\end{enumerate}

\section{Fair Value}

\begin{definitionbox}{Fair Value}
The \textbf{fair value} (FV) of an asset is the market maker's best estimate of the true underlying value $V^*$, used to determine whether bot quotes represent profitable trading opportunities and where to post passive quotes.
\end{definitionbox}

In \texttt{trader.py}, two different fair value models are used:

\begin{center}
\begin{tabular}{lll}
\toprule
\textbf{Product} & \textbf{FV Model} & \textbf{Formula} \\
\midrule
EMERALDS & Constant & $FV = 10000$ \\
TOMATOES & AR(1) + EMA & $FV_t = \text{EMA}_t - \phi \cdot (m_t - m_{t-1})$ \\
\bottomrule
\end{tabular}
\end{center}

The constant FV model is appropriate when empirical analysis shows the price is a stationary process around a fixed mean (which is the case for EMERALDS: the mid-price oscillates in a very tight range around 10000 with std $\approx 0.72$).

The AR(1) + EMA model is appropriate for mean-reverting but drifting processes. See Chapter~32 for full mathematical treatment.

\begin{warningbox}{Why Not Use the Raw Statistical AR(1) Directly?}
EDA shows TOMATOES price changes have a statistical AR(1) coefficient of $\phi_{\text{stat}} \approx -0.30$. A naive reading suggests using this directly in the FV formula $FV_t = m_t - (-0.30)\Delta m_t = m_t + 0.30\Delta m_t$. \textbf{Do not do this.} This is a momentum formula — it amplifies recent moves instead of correcting them. The reason: the statistical AR(1) describes how the \emph{raw mid-price} will evolve, but the market maker needs to quote prices that bots will cross. If the FV moves too aggressively, the bots' crossing thresholds no longer trigger and fill rates collapse. Grid search over $\phi \in [-0.3, 0.3]$ finds $\phi^* = 0.05$ as the PnL-maximising coefficient — a small forward-correction that improves quote placement without losing bot flow. The full justification is in Section~\ref{sec:ema-ar1} (Chapter~32).
\end{warningbox}

\section{Price Discovery}

\textbf{Price discovery} is the process by which markets aggregate private information into public prices. Key mechanisms:

\begin{enumerate}[label=\arabic*.]
  \item \textbf{Informed trading}: Traders with private information buy/sell until the price reflects their information.
  \item \textbf{Order flow}: Net buying pressure pushes prices up; net selling pressure pushes them down.
  \item \textbf{Quote updates}: Market makers continuously revise quotes based on observed trades and information.
\end{enumerate}

The \textbf{efficient market hypothesis} (EMH) states that all available information is already reflected in prices:
\begin{align}
  \text{Weak form:} \quad & \text{Past prices contain no exploitable information} \\
  \text{Semi-strong form:} \quad & \text{All public information is reflected} \\
  \text{Strong form:} \quad & \text{All information (including private) is reflected}
\end{align}

Most algo-trading strategies implicitly reject the weak form (they use past prices) or semi-strong form (they use data the market hasn't fully processed yet).

\begin{takeawaybox}{Chapter 6}
\begin{itemize}
  \item Mid-price $m = (p^A + p^B)/2$ is the standard proxy for true asset value.
  \item VWAP weights price by traded volume; more robust than time-average.
  \item Fair value is the market maker's internal estimate of the true value, used for quoting decisions.
  \item Two FV models: constant (stable products) and AR(1)+EMA (mean-reverting, drifting products).
  \item Price discovery is the mechanism by which private information enters public prices.
\end{itemize}
\end{takeawaybox}

% End of Part I

% ============================================================
%
%  PART II — PRICE ACTION & CHART READING FOUNDATIONS
%
% ============================================================
\part{Price Action \& Chart Reading Foundations}

\begin{warningbox}{Read This Before Part II — It Will Save You Time}
\textbf{This entire part is vocabulary, not strategy.}

The academic evidence for chart patterns as predictive tools is weak to non-existent (Brock et al.\ 1992; Sullivan et al.\ 1999). The Prosperity bots are algorithmic and do not react to candlestick shapes. \textbf{None of the patterns in Chapters 8--10 should be used directly as trading signals in Prosperity without quantitative validation.}

\medskip
\textbf{Why read this part at all?}
\begin{enumerate}[label=\arabic*.]
  \item The vocabulary (support/resistance, trend, momentum) is used everywhere in trading literature. You need it to read research papers.
  \item OHLCV data and volume indicators (Ch.\ 7, Ch.\ 14) are genuinely useful as inputs to quantitative models.
  \item Understanding what practitioners mean by ``the chart shows X'' helps you reason about bot behavior in competition meta-discussions.
\end{enumerate}

\medskip
\textbf{Suggested reading mode:} Read each chapter once for vocabulary. Skip the pattern memorisation. Return to specific sections only when you encounter a term in a paper or forum post and need to look it up.
\end{warningbox}

\chapter{OHLCV Data: The Language of Markets}

\section{What Is OHLCV?}

Every period of market activity (a second, a minute, a day) can be summarised by five numbers. These five numbers are the universal language of price charts and the raw input to virtually every technical indicator.

\begin{definitionbox}{OHLCV}
For a given time interval $[t, t + \Delta t)$, the \textbf{OHLCV} bar is defined by:
\begin{align*}
  O_t &= \text{Opening price: the price of the first trade in the interval} \\
  \mathcal{H}_t &= \text{High price: } \max\{p_i : t_i \in [t, t+\Delta t)\} \\
  L_t &= \text{Low price: } \min\{p_i : t_i \in [t, t+\Delta t)\} \\
  C_t &= \text{Closing price: the price of the last trade in the interval} \\
  V_t &= \text{Volume: total quantity traded, } \sum_{i:\, t_i \in [t,t+\Delta t)} q_i
\end{align*}
Always: $L_t \leq O_t, C_t \leq \mathcal{H}_t$. (Note: $\mathcal{H}_t$ denotes the OHLCV High price throughout this document to avoid confusion with the Hurst exponent $H$ introduced in Chapter~28.)
\end{definitionbox}

\textbf{Why five numbers instead of the full price series?}  Because storing every tick for every asset since 1900 is impractical. OHLCV compresses the information into four price points that capture the range, direction, and opening/closing context of a period. Volume adds the participation dimension.

\subsection{OHLCV in the Prosperity Context}

The Prosperity CSV files use a different format: instead of OHLCV bars they provide the \emph{full order book snapshot} at every timestamp (100-unit intervals). The relevant columns are:
\begin{center}
\begin{tabular}{ll}
\toprule
\textbf{Column} & \textbf{Description} \\
\midrule
\texttt{mid\_price} & $(p^A + p^B)/2$ — equivalent to a tick-by-tick ``close'' \\
\texttt{bid\_price\_1} & Best bid \\
\texttt{ask\_price\_1} & Best ask \\
\texttt{bid\_volume\_1..3} & Quantities at each bid level \\
\texttt{ask\_volume\_1..3} & Quantities at each ask level \\
\texttt{profit\_and\_loss} & Cumulative PnL of the reference participant \\
\bottomrule
\end{tabular}
\end{center}

To construct a 1-minute OHLCV bar from Prosperity tick data, you would aggregate 600 timestamps (each 100 units apart, and 60\,000 units per minute):
\begin{lstlisting}[caption={Constructing OHLCV bars from tick data}]
import pandas as pd
df = pd.read_csv("prices_round_0_day_-2.csv", sep=";")
em = df[df["product"] == "EMERALDS"].copy()
em["bar"] = em["timestamp"] // 60000    # 1-minute bars

ohlcv = em.groupby("bar")["mid_price"].agg(
    O="first", H="max", L="min", C="last"
)
\end{lstlisting}

\section{The True Range}

A limitation of the simple High-Low range is that it ignores overnight gaps (price jumps between periods). The \textbf{True Range} fixes this:

\begin{definitionbox}{True Range}
\[
  TR_t = \max\bigl(H_t - L_t,\; |H_t - C_{t-1}|,\; |L_t - C_{t-1}|\bigr)
\]
It is the largest of:
\begin{itemize}
  \item Current High minus current Low (intrabar range)
  \item Current High minus previous Close (gap up then pull back)
  \item Previous Close minus current Low (gap down then recovery)
\end{itemize}
\end{definitionbox}

The True Range is the building block of the \textbf{Average True Range (ATR)}, covered in Chapter~13.

\section{Returns: Arithmetic vs.\ Logarithmic}

Given two consecutive closing prices $C_{t-1}$ and $C_t$, we can compute returns in two ways:

\begin{definitionbox}{Arithmetic Return}
\[
  r_t = \frac{C_t - C_{t-1}}{C_{t-1}} = \frac{C_t}{C_{t-1}} - 1
\]
Simple to interpret: $r_t = 0.05$ means the price rose 5\%.
\end{definitionbox}

\begin{definitionbox}{Log Return}
\[
  \tilde{r}_t = \ln\!\left(\frac{C_t}{C_{t-1}}\right) = \ln C_t - \ln C_{t-1}
\]
Additive over time: $\tilde{r}_{t:t+n} = \sum_{i=0}^{n-1} \tilde{r}_{t+i}$. Approximately equal to $r_t$ for small returns ($|r_t| \ll 1$).
\end{definitionbox}

\textbf{Why log returns?}
\begin{enumerate}[label=\arabic*.]
  \item \textbf{Time-additivity}: Compound a position over $n$ periods by summing log returns.
  \item \textbf{Symmetry}: A 50\% loss and a 100\% gain are symmetric in log space ($-\ln 2$ vs.\ $+\ln 2$).
  \item \textbf{Stationarity}: Log returns are more likely to be stationary than price levels.
  \item \textbf{Normal approximation}: Central Limit Theorem applies to sums → portfolio log returns are approximately normal.
\end{enumerate}

\begin{examplebox}{Arithmetic vs.\ Log Return}
$C_{t-1} = 5000$, $C_t = 5250$. \\
Arithmetic: $r = 5250/5000 - 1 = 0.05 = 5\%$. \\
Log: $\tilde{r} = \ln(5250/5000) = \ln(1.05) \approx 0.04879 = 4.879\%$. \\
For small returns these are nearly equal; they diverge for large moves.
\end{examplebox}

\section{Volatility: Measuring Price Uncertainty}

\begin{definitionbox}{Historical Volatility}
The \textbf{historical (realised) volatility} over a window of $n$ periods is the standard deviation of log returns:
\[
  \sigma_n = \sqrt{\frac{1}{n-1} \sum_{i=1}^{n} (\tilde{r}_i - \bar{\tilde{r}})^2}
\]
Usually annualised by multiplying by $\sqrt{T}$, where $T$ = trading periods per year.
\end{definitionbox}

\begin{examplebox}{Historical Volatility Calculation}
Suppose 5 daily log returns are: $0.02, -0.01, 0.03, -0.02, 0.01$. \\
Mean: $\bar{r} = (0.02 - 0.01 + 0.03 - 0.02 + 0.01)/5 = 0.006$. \\
Deviations squared: $(0.014)^2, (-0.016)^2, (0.024)^2, (-0.026)^2, (0.004)^2 = 0.000196, 0.000256, 0.000576, 0.000676, 0.000016$. \\
Variance: $\frac{0.00172}{4} = 0.00043$. \\
Daily volatility: $\sigma = \sqrt{0.00043} \approx 0.0207 = 2.07\%$. \\
Annualised (252 trading days): $\sigma_{\text{ann}} = 0.0207 \times \sqrt{252} \approx 32.9\%$.
\end{examplebox}

\begin{takeawaybox}{Chapter 7}
\begin{itemize}
  \item OHLCV encodes one interval of market activity into five numbers: Open, High, Low, Close, Volume.
  \item True Range extends High-Low to account for price gaps between periods.
  \item Log returns are preferred over arithmetic returns: additive, symmetric, approximately normal.
  \item Historical volatility = annualised standard deviation of log returns; the fundamental risk measure.
\end{itemize}
\end{takeawaybox}

% ──────────────────────────────────────────────────────────────────────────────

\chapter{Candlestick Charts}

\begin{warningbox}{Statistical Reality Check for Part II}
Chapters~8--9 cover candlestick patterns, trends, support/resistance, and chart patterns. These are the \textbf{vocabulary of financial markets}: you need to recognise and name these structures because practitioners speak in these terms. However:
\begin{itemize}
  \item \textbf{Academic evidence for predictive power is weak.} Large-scale backtests (Brock \emph{et al.}, 1992; Sullivan \emph{et al.}, 1999) find that most candlestick and chart patterns have minimal statistically significant out-of-sample edge once transaction costs and data snooping bias are accounted for.
  \item \textbf{They can be self-fulfilling at short horizons.} Because many traders use the same patterns, a widely-recognised signal (e.g.\ a 200-day moving average support) may work \emph{because} people act on it, not because of any intrinsic information.
  \item \textbf{Context matters more than the pattern.} A Doji at the peak of a multi-month uptrend on high volume carries more weight than a Doji in a ranging market on low volume.
  \item \textbf{For Prosperity specifically:} The bot market data is generated by a deterministic or semi-stochastic process, not human psychology. Candlestick patterns are unlikely to work; statistical methods (AR(1), ADF, EDA) are far more reliable.
\end{itemize}
Read Part II to understand the language of markets and to be literate in trader communication. Do not rely on candlestick patterns as a primary trading signal without rigorous statistical validation first.
\end{warningbox}

\section{The Anatomy of a Candlestick}

A \textbf{candlestick} is a visual representation of one OHLCV bar. Invented in 18th-century Japan by rice traders, it remains the dominant charting method today.

\begin{center}
\begin{tikzpicture}[scale=1.2]
  % Bullish candle (Close > Open)
  \draw[fill=ExGreen!50, draw=ExGreen!80!black, line width=1.2pt] (0,-0.3) rectangle (0.6,0.9);
  \draw[ExGreen!80!black, line width=1.2pt] (0.3,0.9) -- (0.3,1.4);  % upper wick
  \draw[ExGreen!80!black, line width=1.2pt] (0.3,-0.3) -- (0.3,-0.8); % lower wick
  \node[right] at (0.7, 1.4) {\small High};
  \node[right] at (0.7, 0.9) {\small Close (C $>$ O: bullish)};
  \node[right] at (0.7, -0.3) {\small Open};
  \node[right] at (0.7, -0.8) {\small Low};
  \draw[dashed, gray] (0.7,1.4) -- (0.3,1.4);
  \draw[dashed, gray] (0.7,0.9) -- (0.6,0.9);
  \draw[dashed, gray] (0.7,-0.3) -- (0.6,-0.3);
  \draw[dashed, gray] (0.7,-0.8) -- (0.3,-0.8);
  \node at (0.3,-1.3) {\small\textbf{Bullish}};

  % Bearish candle (Close < Open)
  \draw[fill=WarnRed!50, draw=WarnRed!80!black, line width=1.2pt] (3,-0.9) rectangle (3.6,0.3);
  \draw[WarnRed!80!black, line width=1.2pt] (3.3,0.3) -- (3.3,1.1);
  \draw[WarnRed!80!black, line width=1.2pt] (3.3,-0.9) -- (3.3,-1.4);
  \node[right] at (3.7, 1.1) {\small High};
  \node[right] at (3.7, 0.3) {\small Open};
  \node[right] at (3.7, -0.9) {\small Close (C $<$ O: bearish)};
  \node[right] at (3.7, -1.4) {\small Low};
  \draw[dashed, gray] (3.7,1.1) -- (3.3,1.1);
  \draw[dashed, gray] (3.7,0.3) -- (3.6,0.3);
  \draw[dashed, gray] (3.7,-0.9) -- (3.6,-0.9);
  \draw[dashed, gray] (3.7,-1.4) -- (3.3,-1.4);
  \node at (3.3,-1.9) {\small\textbf{Bearish}};

  % Labels
  \draw[<->, DefBlue] (-0.2,-0.3) -- (-0.2,0.9) node[midway, left] {\small Body};
  \node[DefBlue, font=\small] at (-0.5, 1.15) {Wick};
  \node[DefBlue, font=\small] at (-0.5, -0.55) {Wick};
\end{tikzpicture}
\end{center}

\begin{itemize}
  \item \textbf{Body}: the rectangle between Open and Close. Green/white = bullish ($C > O$); Red/black = bearish ($C < O$).
  \item \textbf{Wicks} (shadows): thin lines extending from the body to the High and Low.
  \item A candle with no body (Open = Close) is a \textbf{Doji} — signals indecision.
\end{itemize}

\section{Key Single-Candle Patterns}

\subsection{Doji}

A Doji occurs when $|C - O| \approx 0$: the price opened and closed at (nearly) the same level. The body is tiny relative to the wicks. \textbf{Interpretation}: equilibrium between buyers and sellers; potential reversal signal when appearing at extremes of a trend.

Types of Doji:
\begin{center}
\begin{tabular}{lll}
\toprule
\textbf{Name} & \textbf{Shape} & \textbf{Signal} \\
\midrule
Standard Doji & $\approx$ equal wicks, tiny body & Indecision \\
Long-legged Doji & Very long wicks both sides & Extreme indecision, high volatility \\
Gravestone Doji & Long upper wick, no lower wick & Bearish reversal \\
Dragonfly Doji & Long lower wick, no upper wick & Bullish reversal \\
\bottomrule
\end{tabular}
\end{center}

\subsection{Hammer and Hanging Man}

Both have a small body at the top and a long lower wick (at least $2\times$ the body):
\begin{itemize}
  \item \textbf{Hammer}: appears after a downtrend → bullish reversal (buyers rejected the lower prices)
  \item \textbf{Hanging Man}: appears after an uptrend → bearish reversal (selling pressure appeared)
\end{itemize}

\subsection{Shooting Star and Inverted Hammer}

Long upper wick, small body at the bottom:
\begin{itemize}
  \item \textbf{Shooting Star}: after uptrend → bearish reversal
  \item \textbf{Inverted Hammer}: after downtrend → potential bullish reversal
\end{itemize}

\subsection{Marubozu}

A candle with no wicks: the body spans the full High-Low range. Indicates strong directional conviction. A bullish Marubozu ($O = L$, $C = H$) means buyers dominated the entire session.

\section{Multi-Candle Patterns}

\subsection{Engulfing Pattern}

A two-candle reversal pattern:
\begin{itemize}
  \item \textbf{Bullish Engulfing}: a large bullish candle that completely engulfs the previous smaller bearish candle. Signal: potential trend reversal upward.
  \item \textbf{Bearish Engulfing}: a large bearish candle engulfs the previous bullish candle. Signal: potential trend reversal downward.
\end{itemize}

Mathematically: bullish engulfing requires $O_t < C_{t-1}$ and $C_t > O_{t-1}$ (the new body contains the prior body).

\subsection{Morning Star and Evening Star}

Three-candle patterns:
\begin{enumerate}[label=\arabic*.]
  \item \textbf{Morning Star} (bullish): long bearish candle → small-bodied candle (gap down) → long bullish candle.
  \item \textbf{Evening Star} (bearish): long bullish candle → small-bodied candle (gap up) → long bearish candle.
\end{enumerate}
The middle candle represents a moment of indecision between the prior trend and the reversal.

\subsection{Three White Soldiers / Three Black Crows}

\begin{itemize}
  \item \textbf{Three White Soldiers}: three consecutive bullish candles, each opening within the prior candle's body and closing near its high. Strong bullish continuation signal.
  \item \textbf{Three Black Crows}: three consecutive bearish candles. Strong bearish continuation.
\end{itemize}

\section{Limitations of Candlestick Patterns}

\begin{warningbox}{Candlestick Patterns Are Not Laws}
\begin{itemize}
  \item Most candlestick patterns have only weak statistical predictive power when tested rigorously (Bulkowski 2008 found many have $<55\%$ accuracy).
  \item Context matters enormously: the same pattern at a major support level vs.\ in the middle of a range has very different implications.
  \item In algo-trading contexts like Prosperity (no sessions, discrete ticks, no OHLCV bars), candlestick patterns are not directly applicable. Micro-structure signals (OBI, spread, FV deviation) are more relevant.
  \item Never trade a pattern in isolation; always confirm with volume, trend, and other indicators.
\end{itemize}
\end{warningbox}

\begin{takeawaybox}{Chapter 8}
\begin{itemize}
  \item A candlestick encodes OHLCV as a body (Open--Close) and wicks (Low, High).
  \item Bullish (green) when Close $>$ Open; bearish (red) when Close $<$ Open.
  \item Key single-candle patterns: Doji (indecision), Hammer (bullish reversal), Shooting Star (bearish reversal).
  \item Multi-candle patterns (Engulfing, Morning/Evening Star) are more reliable but still require confirmation.
  \item In discrete-tick micro-structure environments, candlestick patterns have limited direct applicability.
\end{itemize}
\end{takeawaybox}

% ──────────────────────────────────────────────────────────────────────────────

\chapter{Trends, Support, and Resistance}

\section{The Concept of a Trend}

\begin{definitionbox}{Trend}
A \textbf{trend} is a persistent directional bias in prices over a given time horizon. Formally, a price series $\{p_t\}$ is trending if consecutive prices have a non-zero autocorrelation:
\[
  \Corr(p_{t+1} - p_t,\; p_t - p_{t-1}) > 0 \quad \text{(uptrend/downtrend)}
\]
\end{definitionbox}

\noindent\textbf{Three trend directions:}
\begin{enumerate}[label=\arabic*.]
  \item \textbf{Uptrend}: Series of higher highs (HH) and higher lows (HL).
  \item \textbf{Downtrend}: Series of lower highs (LH) and lower lows (LL).
  \item \textbf{Sideways (ranging)}: No persistent direction; prices oscillate within a range.
\end{enumerate}

\textbf{Trends exist at all timescales simultaneously} (fractal nature of markets). A stock can be in a long-term uptrend, a medium-term correction (downtrend within the uptrend), and a short-term bounce (uptrend within the downtrend).

\section{Dow Theory: The Foundation}

Charles Dow (1851--1902) laid the groundwork for technical analysis. The six tenets of Dow Theory:
\begin{enumerate}[label=\arabic*.]
  \item \textbf{The averages discount everything}: Prices reflect all known information.
  \item \textbf{Three trends}: Primary (months--years), secondary (weeks--months), minor (days--weeks).
  \item \textbf{Primary trend has three phases}: Accumulation → public participation → distribution.
  \item \textbf{Indices must confirm each other}: A bull signal in one index must be confirmed by another.
  \item \textbf{Volume must confirm the trend}: Rising volume in direction of trend confirms strength.
  \item \textbf{A trend is in effect until reversed}: Innocent until proven guilty.
\end{enumerate}

\section{Support and Resistance}

\begin{definitionbox}{Support Level}
A \textbf{support level} is a price at which buying interest is strong enough to prevent further price decline. It is a floor. Mathematically, it is a local minimum cluster: $p_s \approx \min_{t \in W} p_t$ for some window $W$.
\end{definitionbox}

\begin{definitionbox}{Resistance Level}
A \textbf{resistance level} is a price at which selling pressure is strong enough to prevent further price rise. It is a ceiling.
\end{definitionbox}

\textbf{Why do support and resistance exist?}
\begin{itemize}
  \item \textbf{Memory}: Traders remember prices where the market reversed before. If a stock fell from \$50 twice, sellers will be ready at \$50 again.
  \item \textbf{Round numbers}: Psychological barriers (e.g.\ 10000 for EMERALDS, 5000 for TOMATOES).
  \item \textbf{Options strikes}: Large open interest at specific strikes creates a "gravitational pull" (pin risk).
  \item \textbf{Moving averages}: Widely followed MAs act as dynamic support/resistance.
\end{itemize}

\subsection{Support Becomes Resistance (and Vice Versa)}

When a support level is \emph{broken} (price falls through it decisively), it typically becomes a new resistance level. This is called a \textbf{polarity flip}. The logic: buyers who bought at the support are now underwater; when price returns to that level they sell to break even, creating supply (resistance).

\subsection{Quantifying Support and Resistance}

A simple quantitative approach: \textbf{price clustering}.
\begin{equation}
  \text{Support/Resistance strength at } p = \sum_{t=1}^{T} \mathbf{1}[|p_t - p| < \epsilon]
\end{equation}
High count = price spent a lot of time near $p$ = strong level. In the Prosperity data:
\begin{itemize}
  \item EMERALDS: the constant FV = 10000 is both support and resistance (prices are mean-reverting around it)
  \item TOMATOES: the mid-price oscillates around a drifting mean — no fixed support/resistance, but local minima/maxima create short-term levels
\end{itemize}

\section{Trend Lines}

A \textbf{trend line} connects successive lows (uptrend) or successive highs (downtrend). Breaks of trend lines are potential reversal signals.

\begin{definitionbox}{Trend Line (Uptrend)}
Given two local minima $(t_1, L_1)$ and $(t_2, L_2)$ with $t_1 < t_2$ and $L_1 < L_2$, the uptrend line at time $t$ is:
\[
  TL_t = L_1 + \frac{L_2 - L_1}{t_2 - t_1}(t - t_1)
\]
Price should remain above $TL_t$ while the uptrend is intact.
\end{definitionbox}

The slope $\frac{L_2 - L_1}{t_2 - t_1}$ measures the trend's steepness. A steep uptrend is less sustainable than a gradual one.

\section{Channels}

A \textbf{channel} is formed by two parallel trend lines (one for highs, one for lows), enclosing the price action. The price bounces between the two lines.

\begin{itemize}
  \item \textbf{Ascending channel}: Both trend lines slope upward.
  \item \textbf{Descending channel}: Both slope downward.
  \item \textbf{Horizontal channel}: Price ranges between flat support and resistance.
\end{itemize}

Market-making strategies thrive in horizontal channels: buy near support, sell near resistance. In Prosperity, EMERALDS effectively trades in a very tight horizontal channel (9992--10008).

\section{Timeframes and the Multi-Timeframe Approach}

The same asset looks different on different timeframes. A professional trader checks:
\begin{enumerate}[label=\arabic*.]
  \item \textbf{Higher timeframe} (e.g.\ daily): Establish the overall trend direction. Only trade in the direction of this trend.
  \item \textbf{Middle timeframe} (e.g.\ hourly): Identify setups (support/resistance, patterns).
  \item \textbf{Lower timeframe} (e.g.\ 5-minute): Time entries and exits precisely.
\end{enumerate}

In Prosperity's 1000-timestamp-per-day environment, these three timeframes correspond roughly to:
\begin{center}
\begin{tabular}{lll}
\toprule
\textbf{Concept} & \textbf{Prosperity equivalent} & \textbf{Window size} \\
\midrule
Higher timeframe & Full-day trend & 1000 timestamps \\
Middle timeframe & Session trend & 100--200 timestamps \\
Lower timeframe & Tick-level & 1--10 timestamps \\
\bottomrule
\end{tabular}
\end{center}

\section{Chart Patterns}

Chart patterns are recurring visual configurations in price charts that are believed to have predictive value. They are classified as:

\subsection{Continuation Patterns}
Price pauses in the existing trend before resuming.

\begin{center}
\begin{tabular}{lll}
\toprule
\textbf{Pattern} & \textbf{Shape} & \textbf{Signal} \\
\midrule
Bull/Bear Flag & Sharp move, then consolidation parallelogram & Trend continuation \\
Pennant & Sharp move, then converging triangle & Trend continuation \\
Ascending Triangle & Flat resistance, rising support & Bullish continuation \\
Descending Triangle & Flat support, falling resistance & Bearish continuation \\
Symmetrical Triangle & Converging S/R, direction unclear & Breakout (direction TBD) \\
Rectangle & Flat S \& R & Range-bound, then breakout \\
Cup and Handle & U-shape then small pullback & Bullish continuation \\
\bottomrule
\end{tabular}
\end{center}

\subsection{Reversal Patterns}
Trend reverses direction after the pattern completes.

\begin{center}
\begin{tabular}{lll}
\toprule
\textbf{Pattern} & \textbf{Direction} & \textbf{Signal} \\
\midrule
Head and Shoulders & Uptrend reversal & Three peaks, middle highest \\
Inverse H\&S & Downtrend reversal & Three troughs, middle lowest \\
Double Top & Uptrend reversal & Two equal peaks, ``M'' shape \\
Double Bottom & Downtrend reversal & Two equal troughs, ``W'' shape \\
Triple Top/Bottom & Stronger than double & Three tests of S/R \\
Rising/Falling Wedge & Reversal against trend & Converging, sloping channel \\
\bottomrule
\end{tabular}
\end{center}

\subsection{The Head-and-Shoulders Pattern: Full Analysis}

The Head-and-Shoulders (H\&S) is the most studied reversal pattern.

\begin{center}
\begin{tikzpicture}[scale=0.8]
  \draw[->, thick, gray] (0,0) -- (10,0) node[right] {Time};
  \draw[->, thick, gray] (0,0) -- (0,5) node[above] {Price};
  % Neckline
  \draw[dashed, DefBlue, thick] (1.5,1.5) -- (8.5,1.5) node[right] {\small Neckline};
  % Left shoulder
  \draw[thick, ExGreen!80!black] (0.5,1) -- (1.5,2.5) -- (2.5,1.5);
  % Head
  \draw[thick, ExGreen!80!black] (2.5,1.5) -- (4,4) -- (5.5,1.5);
  % Right shoulder
  \draw[thick, ExGreen!80!black] (5.5,1.5) -- (6.5,2.5) -- (7.5,1.5);
  % Breakdown
  \draw[thick, WarnRed] (7.5,1.5) -- (9,0.3);
  % Labels
  \node[above] at (1.5,2.5) {\small L.S.};
  \node[above] at (4,4) {\small Head};
  \node[above] at (6.5,2.5) {\small R.S.};
  \node[right, WarnRed] at (9,0.3) {\small Breakdown};
  % Target
  \draw[<->, DefBlue] (9.5,0) -- (9.5,4) node[midway, right] {\small Target = Head height};
  \draw[dashed, gray] (4,0) -- (4,4);
  \draw[dashed, gray] (4,4) -- (9.5,4);
\end{tikzpicture}
\end{center}

\textbf{Construction:}
\begin{enumerate}[label=\arabic*.]
  \item Left shoulder: rally and pullback to the neckline.
  \item Head: higher rally than the left shoulder, then pullback to the neckline.
  \item Right shoulder: lower rally (same height as left shoulder), then break of neckline.
  \item \textbf{Entry}: short on close below neckline.
  \item \textbf{Target}: neckline price minus the height of the head above the neckline.
  \item \textbf{Stop}: above the right shoulder.
\end{enumerate}

\begin{examplebox}{H\&S Price Target}
Neckline at 5000. Head peak at 5300 (height = 300). \\
Target = $5000 - 300 = 4700$. \\
Stop loss: above right shoulder, say 5100. \\
Risk/reward: $300/100 = 3:1$ (target is 3× the stop).
\end{examplebox}

\subsection{Volume Confirmation}

Chart patterns are more reliable when confirmed by volume:
\begin{itemize}
  \item In H\&S: volume should decrease progressively from left shoulder to head to right shoulder; surge on neckline break.
  \item In flags: volume contracts during consolidation, expands on breakout.
  \item Without volume confirmation, patterns have lower reliability.
\end{itemize}

\begin{warningbox}{Pattern Reliability in Practice}
A landmark study (Lo, Mamaysky, Wang 2000) found that chart patterns do contain statistically significant predictive information, but the effect is small and inconsistent. Retail traders routinely see patterns where none exist (apophenia). In systematic algo-trading, patterns should be \emph{defined algorithmically} (e.g.\ H\&S detector based on local maxima/minima) and validated on out-of-sample data before deployment.
\end{warningbox}

\begin{takeawaybox}{Chapter 9--10 Summary}
\begin{itemize}
  \item Trends are persistent directional biases; characterised by HH+HL (up) or LH+LL (down).
  \item Support = price floor (buying interest), Resistance = price ceiling (selling pressure). They flip when broken.
  \item Trend lines, channels, and chart patterns formalise visual price analysis.
  \item H\&S is the most studied reversal pattern: three peaks (middle highest), neckline break = entry.
  \item All patterns require volume confirmation; always test on out-of-sample data.
  \item In Prosperity's micro-structure context (1000 ticks/day), classical chart patterns are less relevant than order-book signals.
\end{itemize}
\end{takeawaybox}

% End of Part II

% ============================================================
%
%  PART III — TECHNICAL INDICATORS
%
% ============================================================
\part{Technical Indicators}

\begin{warningbox}{When Do Technical Indicators Have Predictive Power?}
Part~III introduces the major families of technical indicators (moving averages, momentum oscillators, volatility indicators, volume indicators). Before using any indicator in a live strategy, you should understand the conditions under which it has statistically measurable edge.

\textbf{Formal test framework.} For any indicator signal $s_t \in \{+1, -1, 0\}$, the \emph{only} rigorous way to assess predictive power is:
\begin{enumerate}[label=\arabic*.]
  \item Compute the \textbf{information coefficient (IC)}: $\text{IC} = \Corr(s_t, r_{t+1})$ over an out-of-sample window. A statistically significant positive IC (two-sided $t$-test, $p < 0.05$) is the minimum requirement.
  \item Apply the \textbf{Bonferroni correction} if you tested many indicators: if you scan 20 indicators, expect one to appear significant at $p < 0.05$ by chance. Use $p < 0.0025$ for 20 indicators.
  \item Test for \textbf{regime dependency}: does the indicator work only in trending markets (Hurst $H > 0.5$) or also in mean-reverting ones ($H < 0.5$)? Most momentum indicators work only in trending regimes.
  \item Verify \textbf{out-of-sample performance}: use at least a 60/40 in-sample/out-of-sample split. Parameters optimised in-sample almost always decay out-of-sample.
\end{enumerate}

\textbf{Typical findings from the literature:}
\begin{itemize}
  \item Moving average crossovers have statistically significant positive IC in commodity futures and FX trend-following strategies (Lempérière \emph{et al.}, 2014). They tend to \emph{not} work in mean-reverting equity markets.
  \item RSI and MACD have weak IC in equity markets but can have measurable edge in crypto markets during high-volatility regimes (Gerlach, 2019).
  \item Volume indicators have IC close to zero in most backtests after correcting for transaction costs.
\end{itemize}

\textbf{For Prosperity:} The bots are algorithmic, not human. Indicators based on human psychology (RSI overbought/oversold, Fibonacci retracements, Dow Theory phases) are unlikely to work. Indicators based on statistical properties of the series (EMA = optimal first-order IIR filter, GARCH vol forecast) have stronger theoretical grounding.
\end{warningbox}

\chapter{Trend Indicators: The Moving Average Family}

Moving averages are the oldest and most widely used technical indicators. They smooth noisy price series to reveal the underlying trend. Every variant in the family makes a different trade-off between \emph{lag} (how slowly it responds to new data) and \emph{noise suppression} (how much it smooths random fluctuations).

\section{Simple Moving Average (SMA)}

\begin{definitionbox}{Simple Moving Average}
The \textbf{Simple Moving Average} over a window of $n$ periods is the unweighted arithmetic mean of the $n$ most recent prices:
\[
  \text{SMA}_t^{(n)} = \frac{1}{n} \sum_{i=0}^{n-1} p_{t-i} = \frac{p_t + p_{t-1} + \cdots + p_{t-n+1}}{n}
\]
Each observation receives equal weight $w_i = \nicefrac{1}{n}$.
\end{definitionbox}

\subsection{Properties of the SMA}

\begin{enumerate}[label=\arabic*.]
  \item \textbf{Lag}: The SMA lags price by approximately $\nicefrac{(n-1)}{2}$ periods. A 20-period SMA lags by $\approx 9.5$ periods.
  \item \textbf{Ghost effect}: When an old observation leaves the window, the SMA can jump even if the current price is flat. This ``phantom move'' does not reflect current market activity.
  \item \textbf{Equal weighting}: Events from $n$ periods ago count as much as the most recent price — arguably too slow to react.
  \item \textbf{Ideal for ranging markets}: In trending markets, SMA lags and gives late signals. In ranging markets, it still provides a stable level to trade from.
\end{enumerate}

\subsection{SMA Crossover Strategy}

The most classic technical strategy: buy when a fast SMA crosses above a slow SMA; sell when it crosses below.

\begin{equation}
  \text{Signal}_t = \begin{cases} +1 & \text{if } \text{SMA}^{(n_\text{fast})}_t > \text{SMA}^{(n_\text{slow})}_t \\ -1 & \text{if } \text{SMA}^{(n_\text{fast})}_t < \text{SMA}^{(n_\text{slow})}_t \end{cases}
\end{equation}

Common pairs: (5,20), (10,50), (50,200). The 50/200-day golden cross / death cross is watched by institutions worldwide.

\begin{examplebox}{SMA Calculation}
Prices over 5 days: $[100, 102, 101, 105, 103]$. \\
$\text{SMA}^{(5)} = \nicefrac{(100+102+101+105+103)}{5} = \nicefrac{511}{5} = 102.2$.\\
The next day price is 107. New $\text{SMA}^{(5)} = \nicefrac{(102+101+105+103+107)}{5} = \nicefrac{518}{5} = 103.6$.\\
The oldest observation (100) dropped out and the new one (107) entered.
\end{examplebox}

\section{Exponential Moving Average (EMA)}

The EMA is the workhorse of modern technical analysis. It addresses the SMA's ghost effect by giving exponentially decaying weights to all past observations — recent data matters more, but \emph{nothing is fully forgotten}.

\begin{definitionbox}{Exponential Moving Average}
The \textbf{EMA} with smoothing factor $\alpha \in (0,1]$ is defined recursively:
\[
  \text{EMA}_t = \alpha \cdot p_t + (1 - \alpha) \cdot \text{EMA}_{t-1}
\]
with initialisation $\text{EMA}_0 = p_0$ (or the first $n$ prices averaged as a seed). The parameter $\alpha$ controls the speed of response: $\alpha \to 1$ = instantaneous response (EMA = current price), $\alpha \to 0$ = very slow response (EMA barely moves).
\end{definitionbox}

\subsection{Expanding the Recursion}

Unrolling the recursion $k$ steps:
\begin{align}
  \text{EMA}_t &= \alpha p_t + (1-\alpha)\text{EMA}_{t-1} \notag \\
               &= \alpha p_t + \alpha(1-\alpha)p_{t-1} + \alpha(1-\alpha)^2 p_{t-2} + \cdots \notag \\
               &= \alpha \sum_{k=0}^{\infty} (1-\alpha)^k p_{t-k}
\end{align}
The weights form a geometric series: $\alpha, \alpha(1-\alpha), \alpha(1-\alpha)^2, \ldots$ They sum to 1:
\[
  \alpha \sum_{k=0}^{\infty} (1-\alpha)^k = \alpha \cdot \frac{1}{1-(1-\alpha)} = \frac{\alpha}{\alpha} = 1 \checkmark
\]

The weight assigned to the observation $k$ periods ago is $\alpha(1-\alpha)^k$, which decays exponentially.

\subsection{Relationship Between \texorpdfstring{$\alpha$}{alpha} and Period $n$}

Practitioners often specify an EMA by its ``equivalent period'' $n$. The convention is:
\begin{equation}
  \alpha = \frac{2}{n+1}
\end{equation}
So a 20-period EMA uses $\alpha = \nicefrac{2}{21} \approx 0.0952$; a 12-period EMA uses $\alpha = \nicefrac{2}{13} \approx 0.1538$.

The half-life of the exponential decay (time for the weight to halve) is:
\begin{equation}
  t_{1/2} = \frac{\ln 2}{\ln\nicefrac{1}{1-\alpha}} \approx \frac{\ln 2}{\alpha} \quad (\alpha \text{ small})
\end{equation}

\begin{examplebox}{EMA Step-by-Step ($\alpha = 0.3$)}
Prices: $[100, 102, 101, 105, 103, 107]$. Seed: EMA$_0 = 100$. \\
$t=1$: EMA $= 0.3 \times 102 + 0.7 \times 100 = 30.6 + 70 = 100.6$ \\
$t=2$: EMA $= 0.3 \times 101 + 0.7 \times 100.6 = 30.3 + 70.42 = 100.72$ \\
$t=3$: EMA $= 0.3 \times 105 + 0.7 \times 100.72 = 31.5 + 70.504 = 102.004$ \\
$t=4$: EMA $= 0.3 \times 103 + 0.7 \times 102.004 = 30.9 + 71.403 = 102.303$ \\
$t=5$: EMA $= 0.3 \times 107 + 0.7 \times 102.303 = 32.1 + 71.612 = 103.712$
\end{examplebox}

\subsection{EMA in Prosperity (\texorpdfstring{$\alpha = 1$}{alpha=1})}

In \texttt{trader.py}, the TOMATOES strategy uses $\alpha = 1$:
\begin{lstlisting}[caption={EMA with alpha=1 in TomatoesStrategy}]
EMA_ALPHA = 1.00   # EMA decay: 1.0 = pure current mid

def _update_ema(self, mid: float) -> None:
    if self._ema is None:
        self._ema = mid
    else:
        self._ema = self.EMA_ALPHA * mid + (1 - self.EMA_ALPHA) * self._ema
\end{lstlisting}

With $\alpha = 1$: $\text{EMA}_t = 1 \cdot p_t + 0 \cdot \text{EMA}_{t-1} = p_t$. The EMA collapses to the current mid-price. No smoothing at all — this is equivalent to using the raw mid. The EMA term is effectively bypassed; the entire fair value signal comes from the AR(1) correction term. For the authoritative combined EMA + AR(1) fair value model (regime-switching $\phi$, parameter choices, and derivation), see Section~\ref{sec:ema-ar1} in Chapter~32.

\begin{warningbox}{EMA Initialisation Bias}
During the first $\sim \nicefrac{1}{\alpha}$ periods, the EMA is biased toward its seed value. For $\alpha = 0.1$, this lasts $\sim 10$ periods. Workarounds: (1) use the first $n$ prices as the seed (SMA warm-up), (2) apply a correction factor $\nicefrac{1}{1-(1-\alpha)^t}$, (3) wait $\sim 3/\alpha$ periods before trading on the signal. In Prosperity the bias vanishes instantly when $\alpha = 1$.
\end{warningbox}

\section{Weighted Moving Average (WMA)}

\begin{definitionbox}{Weighted Moving Average}
The \textbf{WMA} assigns linearly increasing weights to more recent observations:
\[
  \text{WMA}_t^{(n)} = \frac{\sum_{i=0}^{n-1}(n-i)\cdot p_{t-i}}{\sum_{i=0}^{n-1}(n-i)} = \frac{n \cdot p_t + (n-1) \cdot p_{t-1} + \cdots + 1 \cdot p_{t-n+1}}{\nicefrac{n(n+1)}{2}}
\]
The most recent price gets weight $n$, the oldest gets weight $1$.
\end{definitionbox}

More responsive than SMA, less so than EMA. Less popular in practice because it has no recursive formula — it requires storing $n$ prices explicitly (unlike EMA which only needs the previous value).

\section{Double EMA (DEMA) and Triple EMA (TEMA)}

The DEMA and TEMA were created by Patrick Mulloy to reduce EMA lag.

\begin{definitionbox}{DEMA}
\[
  \text{DEMA}_t = 2 \cdot \text{EMA}_t - \text{EMA}(\text{EMA}_t)
\]
where $\text{EMA}(\text{EMA})$ is an EMA applied to the EMA (a second EMA pass).
\end{definitionbox}

\begin{definitionbox}{TEMA}
\[
  \text{TEMA}_t = 3\cdot\text{EMA}_t - 3\cdot\text{EMA}^{(2)}_t + \text{EMA}^{(3)}_t
\]
where $\text{EMA}^{(k)}$ denotes $k$ successive EMA passes.
\end{definitionbox}

\textbf{Intuition}: DEMA and TEMA are extrapolation-based: they project the EMA forward to where it ``would be'' if the lag were removed. This reduces lag but increases noise sensitivity. They can overshoot near price reversals.

\section{Volume-Weighted Moving Average (VWMA)}

\begin{definitionbox}{VWMA}
The \textbf{VWMA} weights each price by the volume traded at that price:
\[
  \text{VWMA}_t^{(n)} = \frac{\sum_{i=0}^{n-1} p_{t-i} \cdot V_{t-i}}{\sum_{i=0}^{n-1} V_{t-i}}
\]
High-volume candles move the VWMA more; low-volume candles move it less.
\end{definitionbox}

VWMA is superior to SMA in markets where volume is informative: a price move on heavy volume is more significant than the same move on light volume. Equivalent to VWAP computed over a rolling window.

\section{Hull Moving Average (HMA)}

Created by Alan Hull (2005) to solve the lag-smoothness trade-off more elegantly:

\begin{definitionbox}{Hull Moving Average}
\begin{align*}
  \text{HMA}_t^{(n)} &= \text{WMA}_{t}^{(\lfloor\sqrt{n}\rfloor)}\!\left( 2 \cdot \text{WMA}_t^{(\lfloor n/2 \rfloor)} - \text{WMA}_t^{(n)} \right)
\end{align*}
Step 1: Compute $\text{WMA}^{(n/2)}$ and $\text{WMA}^{(n)}$. \\
Step 2: Form the difference series $D_t = 2\cdot\text{WMA}^{(n/2)}_t - \text{WMA}^{(n)}_t$. \\
Step 3: Apply $\text{WMA}^{(\sqrt{n})}$ to $D_t$ to smooth the result.
\end{definitionbox}

The HMA is nearly lag-free and surprisingly smooth — achieved by the double-smoothing trick of subtracting the slower WMA from a faster one (DEMA-like), then re-smoothing.

\section{Kalman Filter as an Adaptive Moving Average}

The \textbf{Kalman filter} (Rudolf Kálmán, 1960) is the optimal linear state estimator for a linear Gaussian system, in the sense of minimising the \textbf{mean-squared error} (MSE) $\E[(x_t - \hat{x}_{t|t})^2]$ among all estimators that are linear functions of the observations $p_1, \ldots, p_t$. In a price context, the state is the true (unobserved) value, and observations are noisy prices.

\begin{proof}[Optimality Sketch: Why the Kalman Gain Minimises MSE]
In the linear Gaussian model, $p(x_t \mid p_{1:t})$ is Gaussian (conjugate prior + Gaussian likelihood). The MMSE estimator is the \emph{posterior mean} $\hat{x}_{t|t} = \E[x_t \mid p_{1:t}]$. One can show (Kalman 1960; see also Welch \& Bishop 2006) that the posterior mean satisfies the recursion $\hat{x}_{t|t} = \hat{x}_{t|t-1} + K_t(p_t - \hat{x}_{t|t-1})$, where the gain $K_t = P_{t|t-1}/(P_{t|t-1}+R)$ is chosen to minimise $P_{t|t} = \E[(x_t-\hat{x}_{t|t})^2]$. Taking the derivative of $P_{t|t}$ with respect to $K_t$:
\[
  P_{t|t} = (1-K_t)^2 P_{t|t-1} + K_t^2 R
  \quad\Rightarrow\quad \frac{dP_{t|t}}{dK_t} = -2(1-K_t)P_{t|t-1} + 2K_t R = 0
  \quad\Rightarrow\quad K_t = \frac{P_{t|t-1}}{P_{t|t-1}+R}
\]
This is the Kalman gain formula. By the Gauss-Markov theorem, \emph{no} unbiased linear estimator can achieve lower MSE.
\end{proof}

\begin{keyconceptbox}{Kalman Filter for Prices}
\textbf{State model:} True value $x_t = x_{t-1} + w_t$, where $w_t \sim \mathcal{N}(0, Q)$ (process noise). \\
\textbf{Observation model:} Observed price $p_t = x_t + v_t$, where $v_t \sim \mathcal{N}(0, R)$ (measurement noise). \\
\textbf{Update equations:}
\begin{align}
  K_t &= \frac{P_{t|t-1}}{P_{t|t-1} + R} \quad \text{(Kalman gain)} \\
  \hat{x}_{t|t} &= \hat{x}_{t|t-1} + K_t(p_t - \hat{x}_{t|t-1}) \quad \text{(state update)} \\
  P_{t|t} &= (1 - K_t) P_{t|t-1} \quad \text{(variance update)} \\
  \hat{x}_{t+1|t} &= \hat{x}_{t|t}, \quad P_{t+1|t} = P_{t|t} + Q \quad \text{(prediction)}
\end{align}
\end{keyconceptbox}

The Kalman gain $K_t \in [0,1]$ is the weight given to the new observation. When $R$ is small (reliable observations), $K_t \to 1$ (EMA with $\alpha \to 1$). When $Q$ is small (stable true value), $K_t \to 0$ (EMA with $\alpha \to 0$). The Kalman filter is effectively a \emph{time-varying} EMA whose $\alpha$ adapts to the signal-to-noise ratio.

\section{Comparison of Moving Average Types}

\begin{center}
\begin{tabular}{lccccc}
\toprule
\textbf{MA Type} & \textbf{Lag} & \textbf{Smoothness} & \textbf{Memory} & \textbf{Recursive} & \textbf{Adaptive} \\
\midrule
SMA         & High   & High   & $n$ periods  & No  & No  \\
EMA         & Medium & Medium & $\infty$     & Yes & No  \\
WMA         & Low    & Medium & $n$ periods  & No  & No  \\
DEMA        & Very low & Low  & $\infty$     & Yes & No  \\
TEMA        & Minimal  & Low  & $\infty$     & Yes & No  \\
VWMA        & Medium & Medium & $n$ periods  & No  & Volume \\
HMA         & Very low & High & $n$ periods  & No  & No  \\
Kalman      & Minimal  & High & $\infty$     & Yes & Yes \\
\bottomrule
\end{tabular}
\end{center}

\begin{warningbox}{More Responsive $\neq$ Better}
Lower lag means faster response to real trend changes — but also faster response to noise. In mean-reverting markets like EMERALDS, a highly responsive MA generates false signals constantly. In trending markets, it captures trend changes early. Always match the MA type to the market regime.
\end{warningbox}

\subsection{EMA \texorpdfstring{$\alpha$}{alpha} Sensitivity}

For products where $\alpha < 1$ is meaningful (slow drift, not mean-reverting), knowing how sensitive PnL is to $\alpha$ guides how carefully to tune it.

\textbf{Theoretical sensitivity: lag vs.\ noise.} The EMA has effective lag $\approx (1-\alpha)/\alpha$ periods and noise suppression improving with smaller $\alpha$:

\begin{center}
\begin{tabular}{c|cccc}
\toprule
$\alpha$ & Equiv.\ $n$ (SMA) & Lag (periods) & Noise reduction & Drift tracking \\
\midrule
0.05 & 39 & $\approx 19$ & Strong & Slow \\
0.10 & 19 & $\approx 9$  & Good   & Moderate \\
0.20 & 9  & $\approx 4$  & Moderate & Fast \\
0.30 & 6  & $\approx 2.3$ & Light  & Fast \\
0.50 & 3  & $\approx 1$  & Minimal & Very fast \\
1.00 & 1  & 0           & None   & Instant (raw mid) \\
\bottomrule
\end{tabular}
\end{center}

\textbf{Practical rule for Prosperity.} For TOMATOES (chosen $\alpha = 1.0$): since the AR(1) correction is applied on top, EMA smoothing is redundant. A $\pm 0.2$ change in $\alpha$ shifts the effective lag by $\pm 1$--$2$ periods but has near-zero PnL impact (confirmed by grid search: IS PnL difference between $\alpha = 0.5$ and $\alpha = 1.0$ is $< 1\%$). For other products with genuine intraday drift and no AR(1) structure, $\alpha$ is a significant tuning parameter: use the grid search methodology from Chapter~26 with $\alpha \in \{0.05, 0.1, 0.2, 0.5, 1.0\}$.

\begin{takeawaybox}{Chapter 11}
\begin{itemize}
  \item SMA: equal weights, ghost effect, high lag. Best for ranging markets.
  \item EMA: exponentially decaying weights, recursive, no ghost effect. $\alpha = \nicefrac{2}{n+1}$. Used in \texttt{trader.py} with $\alpha=1$ (= raw mid, no smoothing).
  \item WMA: linear weights, more responsive than SMA, not recursive.
  \item DEMA/TEMA: extrapolation-based lag reduction; noisy near reversals.
  \item VWMA: volume-weighted; superior when volume is informative.
  \item HMA: near-zero lag with good smoothness via double WMA trick.
  \item Kalman: optimal adaptive filter; time-varying $\alpha$ based on signal-to-noise ratio.
  \item $\alpha$ sensitivity: lag $\approx (1-\alpha)/\alpha$ periods. For TOMATOES ($\alpha=1$), $\alpha$ is insensitive; for drift-dominated products, tune $\alpha$ via grid search.
\end{itemize}
\end{takeawaybox}

% ──────────────────────────────────────────────────────────────────────────────

\chapter{Momentum Oscillators}

Momentum oscillators measure the \emph{speed and direction} of price change, independent of absolute price level. They are bounded indicators (typically 0--100 or $-100$--$+100$) that oscillate around a centreline or between extreme levels, making them useful for identifying overbought/oversold conditions and divergences.

\section{Rate of Change (ROC)}

\begin{definitionbox}{Rate of Change}
The $n$-period \textbf{ROC} measures price momentum as a percentage change:
\[
  \text{ROC}_t^{(n)} = \frac{p_t - p_{t-n}}{p_{t-n}} \times 100
\]
It is the simplest momentum indicator: positive ROC = price higher than $n$ periods ago (momentum), negative = price lower.
\end{definitionbox}

\textbf{Limitations}: ROC has no bounds, so ``overbought/oversold'' levels must be calibrated per asset. It is also sensitive to extreme past prices $p_{t-n}$ (denominator effect).

\section{Relative Strength Index (RSI)}

The RSI (Welles Wilder, 1978) is the most popular momentum oscillator. It measures the speed and change of price movements on a scale of 0--100.

\begin{definitionbox}{RSI}
Let $U_t = \max(p_t - p_{t-1}, 0)$ (up-move) and $D_t = \max(p_{t-1} - p_t, 0)$ (down-move). Define the average gain and average loss over $n$ periods:
\[
  \overline{G}_t = \frac{1}{n}\sum_{i=0}^{n-1} U_{t-i}, \qquad \overline{L}_t = \frac{1}{n}\sum_{i=0}^{n-1} D_{t-i}
\]
(Wilder uses a smoothed average, equivalent to EMA with $\alpha = \nicefrac{1}{n}$.) The \textbf{Relative Strength} is $RS_t = \overline{G}_t / \overline{L}_t$. The RSI is:
\[
  \text{RSI}_t = 100 - \frac{100}{1 + RS_t} = \frac{100 \cdot \overline{G}_t}{\overline{G}_t + \overline{L}_t}
\]
\end{definitionbox}

\textbf{Bounds}: RSI $\in [0, 100]$. RSI $= 100$ means all $n$ periods were up-moves. RSI $= 0$ means all were down-moves.

\textbf{Standard signals:}
\begin{itemize}
  \item RSI $> 70$: Overbought (momentum may be exhausted, potential pullback)
  \item RSI $< 30$: Oversold (potential bounce)
  \item \textbf{Divergence}: price makes new high but RSI makes lower high → bearish divergence (momentum weakening)
  \item \textbf{Centreline cross}: RSI crosses 50 → trend confirmation
\end{itemize}

\begin{examplebox}{RSI Calculation ($n = 3$)}
Prices: $[100, 103, 101, 105, 108, 106]$. \\
Changes: $[+3, -2, +4, +3, -2]$. \\
Gains: $[3, 0, 4, 3, 0]$; Losses: $[0, 2, 0, 0, 2]$. \\
Using last 3: $\overline{G} = (4+3+0)/3 = 2.33$; $\overline{L} = (0+0+2)/3 = 0.67$. \\
$RS = 2.33/0.67 = 3.49$; $\text{RSI} = 100 - 100/(1+3.49) = 100 - 22.3 = 77.7$. \\
RSI $> 70$: overbought signal.
\end{examplebox}

\begin{warningbox}{RSI in Mean-Reverting vs.\ Trending Markets}
In mean-reverting markets (like EMERALDS), RSI $> 70$ genuinely signals a reversal opportunity. In strongly trending markets, RSI can remain above 70 for extended periods — trading the ``overbought'' signal will incur losses. Always check the trend regime before acting on RSI extremes.
\end{warningbox}

\section{Stochastic Oscillator}

George Lane's Stochastic (1950s) measures where the current close sits within the recent High-Low range. The intuition: in uptrends, prices close near the High of their range; in downtrends, near the Low.

\begin{definitionbox}{Stochastic Oscillator}
\begin{align}
  \%K_t &= \frac{C_t - L_t^{(n)}}{H_t^{(n)} - L_t^{(n)}} \times 100 \\
  \%D_t &= \text{SMA}^{(3)}(\%K_t)
\end{align}
where $H_t^{(n)} = \max_{i=0}^{n-1} H_{t-i}$ and $L_t^{(n)} = \min_{i=0}^{n-1} L_{t-i}$. Standard period: $n = 14$.
\end{definitionbox}

\begin{itemize}
  \item $\%K$: fast stochastic (more sensitive)
  \item $\%D$: slow stochastic (3-period SMA of $\%K$, smoother)
  \item Overbought: $\%K > 80$; Oversold: $\%K < 20$
  \item Signal: $\%K$ crosses $\%D$ from below (bullish) or above (bearish)
\end{itemize}

The \textbf{Slow Stochastic} smooths $\%K$ with a 3-period SMA before computing $\%D$, further reducing false signals.

\section{MACD: Moving Average Convergence/Divergence}

Gerald Appel (1979). MACD measures the relationship between two EMAs, making it both a trend-following and momentum indicator.

\begin{definitionbox}{MACD}
\begin{align}
  \text{MACD}_t &= \text{EMA}^{(12)}_t - \text{EMA}^{(26)}_t \quad \text{(MACD line)} \\
  \text{Signal}_t &= \text{EMA}^{(9)}_t(\text{MACD}) \quad \text{(Signal line)} \\
  \text{Histogram}_t &= \text{MACD}_t - \text{Signal}_t
\end{align}
Standard parameters: fast=12, slow=26, signal=9.
\end{definitionbox}

\textbf{Trading signals:}
\begin{itemize}
  \item \textbf{Centreline cross}: MACD crosses zero → trend change (EMA$^{(12)}$ crosses EMA$^{(26)}$)
  \item \textbf{Signal line cross}: MACD crosses Signal → momentum shift; more responsive
  \item \textbf{Histogram}: Bars above zero = bullish momentum; bars shrinking = momentum weakening
  \item \textbf{Divergence}: Price makes new high but MACD histogram makes lower high → weakening momentum
\end{itemize}

\textbf{MACD vs.\ RSI comparison:}
\begin{center}
\begin{tabular}{lll}
\toprule
\textbf{Property} & \textbf{MACD} & \textbf{RSI} \\
\midrule
Type & Trend + momentum & Pure momentum \\
Bounded & No (unbounded) & Yes (0--100) \\
Best regime & Trending & Ranging \\
Overbought/sold & Not applicable & $>70$ / $<30$ \\
Primary signal & Crossovers, divergence & Extremes, divergence \\
\bottomrule
\end{tabular}
\end{center}

\section{Commodity Channel Index (CCI)}

Donald Lambert (1980). CCI measures how far price deviates from its statistical mean, scaled by the mean absolute deviation.

\begin{definitionbox}{CCI}
Let $TP_t = \nicefrac{(H_t + L_t + C_t)}{3}$ be the \textbf{typical price}. The $n$-period CCI is:
\[
  \text{CCI}_t^{(n)} = \frac{TP_t - \overline{TP}_t^{(n)}}{0.015 \cdot \text{MAD}_t^{(n)}}
\]
where $\overline{TP}_t^{(n)} = \frac{1}{n}\sum_{i=0}^{n-1} TP_{t-i}$ and $\text{MAD}_t^{(n)} = \frac{1}{n}\sum_{i=0}^{n-1}|TP_{t-i} - \overline{TP}_t^{(n)}|$. The constant 0.015 scales so that $\approx 70$--80\% of values fall within $[-100, +100]$.
\end{definitionbox}

\textbf{Signals:} CCI $> +100$: overbought (strong trend); CCI $< -100$: oversold. Originally designed for commodities with seasonal cycles ($n = 20$), now used for any asset.

\section{Williams \%R}

Larry Williams (1973). Essentially an inverted Stochastic: measures where Close sits in the High-Low range, but scaled $[-100, 0]$.

\begin{definitionbox}{Williams \%R}
\[
  \%R_t^{(n)} = -\frac{H_t^{(n)} - C_t}{H_t^{(n)} - L_t^{(n)}} \times 100
\]
Range: $[-100, 0]$. Overbought: $\%R > -20$; Oversold: $\%R < -80$.
\end{definitionbox}

Note: $\%R = -100 + \%K_{\text{Stochastic}}$ (they are linearly related). Williams \%R is more commonly used for short-term trading.

\section{Money Flow Index (MFI)}

The MFI is a volume-weighted RSI — it incorporates both price movement and volume, measuring the ``money flow'' into or out of an asset.

\begin{definitionbox}{Money Flow Index}
Typical Price: $TP_t = \nicefrac{(H_t + L_t + C_t)}{3}$. \\
Raw Money Flow: $MF_t = TP_t \times V_t$. \\
\textbf{Positive MF}: $MF_t$ when $TP_t > TP_{t-1}$ (money flowing in). \\
\textbf{Negative MF}: $MF_t$ when $TP_t < TP_{t-1}$ (money flowing out). \\
Money Flow Ratio over $n$ periods:
\[
  MFR_t = \frac{\sum_{i=0}^{n-1} MF^+_{t-i}}{\sum_{i=0}^{n-1} MF^-_{t-i}}
\]
\[
  \text{MFI}_t = 100 - \frac{100}{1 + MFR_t}
\]
\end{definitionbox}

MFI $>$ 80: overbought with heavy volume (strong momentum). MFI $<$ 20: oversold. More reliable than RSI when volume data is available, because volume confirms price moves.

\section{Summary Comparison of Momentum Oscillators}

\begin{center}
\begin{tabular}{llcll}
\toprule
\textbf{Oscillator} & \textbf{Range} & \textbf{Volume} & \textbf{Overbought} & \textbf{Best For} \\
\midrule
ROC         & Unbounded & No  & Asset-specific & Pure momentum \\
RSI         & 0--100    & No  & $>70$          & Ranging markets \\
Stochastic  & 0--100    & No  & $>80$          & Short-term reversals \\
MACD        & Unbounded & No  & N/A            & Trending markets \\
CCI         & $\approx\pm100$ & No & $>+100$   & Cyclical commodities \\
Williams \%R & $-100$--$0$ & No & $>-20$        & Short-term reversals \\
MFI         & 0--100    & Yes & $>80$          & Volume-confirmed momentum \\
\bottomrule
\end{tabular}
\end{center}

\begin{takeawaybox}{Chapter 12}
\begin{itemize}
  \item Momentum oscillators measure speed of price change, not direction per se.
  \item RSI: bounded 0--100, standard overbought/sold at 70/30. Wilder's smoothed average = EMA($\alpha=1/n$).
  \item Stochastic: close's position within the $n$-period range. $\%K$ fast, $\%D$ = 3-SMA of $\%K$.
  \item MACD: difference of two EMAs; trend + momentum hybrid, unbounded.
  \item MFI: volume-weighted RSI; superior when volume data is meaningful.
  \item All oscillators work best in their intended regime; divergence signals are the most robust.
\end{itemize}
\end{takeawaybox}

% ──────────────────────────────────────────────────────────────────────────────

\chapter{Volatility Indicators}

Volatility indicators measure the \emph{rate and magnitude} of price change, independently of direction. They answer: ``how much is the market moving?'' rather than ``which way?''

\section{Historical Volatility (Revisited)}

We defined historical volatility in Chapter~7. In the indicator context it is used as a dynamic, rolling measure:
\[
  \sigma_t^{(n)} = \sqrt{\frac{1}{n-1}\sum_{i=0}^{n-1}(\tilde{r}_{t-i} - \bar{\tilde{r}})^2}
\]
Rising $\sigma$ = market becoming more uncertain/volatile. Falling $\sigma$ = market quieting.

\textbf{Volatility clustering}: In financial markets, large moves are followed by large moves (of either sign) and small moves by small moves. This is the signature of the GARCH model (Section~\ref{sec:garch}).

\section{Average True Range (ATR)}

The ATR (Wilder, 1978) is the most practical volatility indicator. It measures the average range of price movement over $n$ periods, accounting for gaps.

\begin{definitionbox}{Average True Range}
\[
  \text{ATR}_t^{(n)} = \text{EMA}^{(n)}_t(TR_t)
\]
where the True Range is $TR_t = \max(H_t - L_t, |H_t - C_{t-1}|, |L_t - C_{t-1}|)$ and the EMA uses Wilder's smoothing ($\alpha = \nicefrac{1}{n}$):
\[
  \text{ATR}_t = \frac{(n-1)\cdot\text{ATR}_{t-1} + TR_t}{n}
\]
\end{definitionbox}

\textbf{Uses of ATR:}
\begin{enumerate}[label=\arabic*.]
  \item \textbf{Stop-loss sizing}: Place stop at $k \times \text{ATR}$ below entry (e.g.\ $k = 2$). This adapts to volatility automatically.
  \item \textbf{Position sizing}: Trade size $\propto \nicefrac{1}{\text{ATR}}$ to maintain constant risk.
  \item \textbf{Breakout filter}: Only take a breakout signal if ATR is above a threshold (volatile enough to sustain the move).
\end{enumerate}

\begin{examplebox}{ATR-Based Stop Loss}
Current price: 5000. ATR = 25. Using $k = 2$: \\
Stop loss for a long: $5000 - 2 \times 25 = 4950$. \\
If price falls below 4950, the move is ``larger than normal volatility'' — exit.
\end{examplebox}

\section{Bollinger Bands}

John Bollinger (1980s). Bollinger Bands place a dynamic envelope around price, expanding in volatile markets and contracting in quiet ones.

\begin{definitionbox}{Bollinger Bands}
Given a window $n$ (typically 20) and multiplier $k$ (typically 2):
\begin{align}
  \text{Middle Band}_t &= \text{SMA}_t^{(n)} \\
  \sigma_t &= \sqrt{\frac{1}{n}\sum_{i=0}^{n-1}(p_{t-i} - \text{SMA}_t^{(n)})^2} \quad \text{(rolling std dev)} \\
  \text{Upper Band}_t &= \text{SMA}_t^{(n)} + k\,\sigma_t \\
  \text{Lower Band}_t &= \text{SMA}_t^{(n)} - k\,\sigma_t
\end{align}
With $k=2$: approximately 95\% of prices fall inside the bands (under normality).
\end{definitionbox}

\textbf{Key properties:}
\begin{itemize}
  \item \textbf{Squeeze}: bands narrow → volatility is low → breakout incoming (direction unknown)
  \item \textbf{Walk the band}: in strong trends, price can ``walk'' along the upper or lower band. This is \emph{not} a reversal signal.
  \item \textbf{Mean reversion signal}: in ranging markets, price at upper band → sell, at lower band → buy (only valid if trend is absent)
  \item \textbf{\%B}: $\%B = \nicefrac{(p - \text{Lower})}{(\text{Upper} - \text{Lower})}$; measures where price sits in the bands (0--1)
  \item \textbf{Bandwidth}: $({\text{Upper} - \text{Lower}})/{\text{Middle}}$; measures band width normalised by price
\end{itemize}

\begin{warningbox}{Bollinger Bands Are Not Overbought/Sold Signals}
Many beginners sell when price hits the upper band. This is wrong in trending markets — price can close above the upper band for many consecutive periods. Bollinger Bands should be used as \emph{volatility context}, not directional signals in isolation.
\end{warningbox}

\section{Keltner Channels}

An ATR-based alternative to Bollinger Bands:
\begin{definitionbox}{Keltner Channels}
\begin{align}
  \text{Middle}_t &= \text{EMA}_t^{(n)} \\
  \text{Upper}_t &= \text{Middle}_t + k \cdot \text{ATR}_t^{(n)} \\
  \text{Lower}_t &= \text{Middle}_t - k \cdot \text{ATR}_t^{(n)}
\end{align}
Typical: $n = 20$, $k = 2$.
\end{definitionbox}

\textbf{Keltner vs.\ Bollinger}: Bollinger uses standard deviation (responds to sudden spikes), Keltner uses ATR (smoother, less spiky). When Bollinger Bands are \emph{inside} Keltner Channels, it signals a period of extremely low volatility (the ``Squeeze'' strategy by John Carter).

\section{Donchian Channels}

Richard Donchian (1970s). Simplest volatility channel: the $n$-period highest high and lowest low.
\begin{definitionbox}{Donchian Channels}
\begin{align}
  \text{Upper}_t^{(n)} &= \max(H_{t}, H_{t-1}, \ldots, H_{t-n+1}) \\
  \text{Lower}_t^{(n)} &= \min(L_{t}, L_{t-1}, \ldots, L_{t-n+1})
\end{align}
\end{definitionbox}

The famous Turtle Trading System (Richard Dennis, 1983) used a 20-period Donchian breakout: buy when price exceeds the 20-day high; sell when price falls below the 20-day low (see Chapter~17).

\section{GARCH: Modelling Volatility Clustering}
\label{sec:garch}

In financial time series, volatility is not constant — it clusters. The \textbf{GARCH} (Generalised Autoregressive Conditional Heteroskedasticity) model captures this.

\begin{definitionbox}{GARCH(1,1)}
Let $r_t = \mu + \varepsilon_t$ where $\varepsilon_t = \sigma_t z_t$ and $z_t \sim \mathcal{N}(0,1)$. The conditional variance follows:
\[
  \sigma_t^2 = \omega + \alpha\,\varepsilon_{t-1}^2 + \beta\,\sigma_{t-1}^2
\]
with constraints $\omega > 0$, $\alpha \geq 0$, $\beta \geq 0$, $\alpha + \beta < 1$ (stationarity).
\begin{itemize}
  \item $\omega$: long-run variance floor
  \item $\alpha\,\varepsilon_{t-1}^2$: reaction to the last shock (ARCH term)
  \item $\beta\,\sigma_{t-1}^2$: persistence of variance (GARCH term)
\end{itemize}
\end{definitionbox}

\textbf{Interpretation}: If yesterday's return was large ($\varepsilon_{t-1}^2$ large), today's volatility forecast $\sigma_t^2$ is larger. If $\alpha + \beta$ is close to 1, volatility shocks are very persistent (typical in equity markets). The \textbf{long-run variance} is $\bar{\sigma}^2 = \nicefrac{\omega}{1-\alpha-\beta}$.

\begin{center}
\begin{tabular}{ll}
\toprule
\textbf{Model} & \textbf{Variance equation} \\
\midrule
ARCH(q) & $\sigma_t^2 = \omega + \sum_{i=1}^q \alpha_i \varepsilon_{t-i}^2$ \\
GARCH(p,q) & $\sigma_t^2 = \omega + \sum_{i=1}^q \alpha_i \varepsilon_{t-i}^2 + \sum_{j=1}^p \beta_j \sigma_{t-j}^2$ \\
EGARCH & $\ln\sigma_t^2 = \omega + \alpha|\frac{\varepsilon_{t-1}}{\sigma_{t-1}}| + \gamma\frac{\varepsilon_{t-1}}{\sigma_{t-1}} + \beta\ln\sigma_{t-1}^2$ \\
GJR-GARCH & Adds asymmetry: larger variance after negative shocks (leverage effect) \\
\bottomrule
\end{tabular}
\end{center}

\subsection{MLE Estimation of GARCH}

GARCH parameters $(\omega, \alpha, \beta)$ are estimated by \textbf{Maximum Likelihood Estimation (MLE)}. Assuming $z_t \sim \mathcal{N}(0,1)$, the log-likelihood of a return series $r_1, \ldots, r_T$ is:
\[
  \ell(\omega, \alpha, \beta) = -\frac{1}{2}\sum_{t=1}^{T}\left[\ln(2\pi) + \ln\sigma_t^2 + \frac{\varepsilon_t^2}{\sigma_t^2}\right]
\]
where $\varepsilon_t = r_t - \hat{\mu}$ and $\sigma_t^2$ is computed recursively from the GARCH equation given the parameter values. This is maximised numerically (gradient-based optimiser, e.g.\ L-BFGS-B).

The recursion must be \emph{initialised}: conventionally, $\sigma_1^2 = \nicefrac{\omega}{1-\alpha-\beta}$ (the unconditional variance). In Python:

\begin{lstlisting}
from arch import arch_model   # pip install arch

am = arch_model(returns, vol='GARCH', p=1, q=1, mean='constant')
res = am.fit(disp='off')
print(res.summary())

# One-step-ahead vol forecast
forecast = res.forecast(horizon=1)
print(f"Next-period sigma: {forecast.variance.iloc[-1, 0]**0.5:.4f}")
\end{lstlisting}

\subsection{Empirical Properties and Stylised Facts}

Real financial return series exhibit several empirical properties that GARCH captures:

\begin{enumerate}[label=\arabic*.]
  \item \textbf{Fat tails (leptokurtosis):} The unconditional distribution of returns has kurtosis $> 3$. Even with Gaussian innovations, GARCH produces fat tails in the unconditional distribution because low-volatility and high-volatility regimes are mixed. To better match observed tails, use Student-$t$ innovations:
  \[
    z_t \sim t_\nu \quad \text{with } \nu \approx 5\text{--}8 \text{ degrees of freedom}
  \]
  \item \textbf{Leverage effect:} Negative returns increase volatility more than positive returns of the same magnitude (stock prices fall fast, rise slowly). Standard GARCH is symmetric; GJR-GARCH captures this with:
  \[
    \sigma_t^2 = \omega + (\alpha + \gamma\,\mathbf{1}_{\varepsilon_{t-1}<0})\varepsilon_{t-1}^2 + \beta\,\sigma_{t-1}^2
  \]
  where $\gamma > 0$ amplifies variance after negative shocks.
  \item \textbf{Volatility persistence:} In equity markets, $\alpha + \beta \approx 0.99$ is common, meaning shocks to volatility decay very slowly (half-life of weeks or months).
  \item \textbf{Mean reversion of volatility:} Despite high persistence, variance does revert to $\bar{\sigma}^2 = \omega/(1-\alpha-\beta)$ in the long run (as long as $\alpha + \beta < 1$).
\end{enumerate}

\begin{keyconceptbox}{GARCH One-Step Volatility Forecast}
The GARCH(1,1) optimal one-step forecast of conditional variance is:
\[
  \hat{\sigma}_{t+1}^2 = \omega + \alpha\,\hat{\varepsilon}_t^2 + \beta\,\hat{\sigma}_t^2
\]
For $h$-step ahead:
\[
  \hat{\sigma}_{t+h}^2 = \bar{\sigma}^2 + (\alpha+\beta)^{h-1}\!\left(\hat{\sigma}_{t+1}^2 - \bar{\sigma}^2\right)
\]
Volatility forecasts revert toward $\bar{\sigma}^2$ at rate $(\alpha+\beta)^h$. Use these forecasts to size positions (smaller size during high-vol forecasts) and to set option pricing inputs.
\end{keyconceptbox}

\begin{takeawaybox}{Chapter 13}
\begin{itemize}
  \item ATR = Wilder-smoothed True Range; the standard tool for stop sizing and position sizing.
  \item Bollinger Bands = SMA $\pm k\sigma$; expand in volatile markets, contract in quiet. Bands squeeze = breakout ahead.
  \item Keltner Channels = EMA $\pm k\cdot$ATR; smoother than Bollinger. Keltner inside Bollinger = extreme squeeze.
  \item Donchian = $n$-period highest high / lowest low; used in Turtle breakout system.
  \item GARCH(1,1) models volatility clustering: $\sigma_t^2 = \omega + \alpha\varepsilon_{t-1}^2 + \beta\sigma_{t-1}^2$; forecasts next period's variance.
\end{itemize}
\end{takeawaybox}

% ──────────────────────────────────────────────────────────────────────────────

\chapter{Volume Indicators}

Volume is the number of units traded. It provides crucial context: a price move on heavy volume is more significant and sustainable than the same move on light volume.

\section{Volume and Its Significance}

\begin{keyconceptbox}{Volume Principles}
\begin{enumerate}[label=\arabic*.]
  \item Volume confirms trends: rising price + rising volume = healthy uptrend.
  \item Volume divergence warns of reversals: rising price + falling volume = weak uptrend, potential reversal.
  \item Breakouts need volume: a breakout above resistance on low volume is likely a false break.
  \item Climax volume: extremely high volume after a long move often signals exhaustion.
\end{enumerate}
\end{keyconceptbox}

\section{Order Book Imbalance (OBI)}

This is the most directly relevant volume indicator to the Prosperity project. It appears in \texttt{eda.py} and was tested (and rejected) for asymmetric sizing in \texttt{trader.py}.

\begin{definitionbox}{Order Book Imbalance}
The \textbf{OBI} at time $t$ measures the relative imbalance between bid and ask volume in the order book:
\[
  \text{OBI}_t = \frac{Q_t^B - Q_t^A}{Q_t^B + Q_t^A} \in [-1, +1]
\]
where $Q_t^B = \sum_i q_i^B$ (total bid volume) and $Q_t^A = \sum_j q_j^A$ (total ask volume).
\end{definitionbox}

\textbf{Interpretation:}
\begin{itemize}
  \item OBI $> 0$: More volume on the buy side → upward price pressure
  \item OBI $< 0$: More volume on the sell side → downward price pressure
  \item OBI $= 0$: Perfectly balanced book
\end{itemize}

\subsection{OBI as a Predictive Signal}

\textbf{In-sample, before multiple-testing correction}, OBI shows a correlation of $\approx 0.38$ with the next-period return in TOMATOES' tight-spread regime (computed on Prosperity~4 data, days $-2$ and $-1$; non-zero OBI occurs $\approx 12\%$ of timestamps in this regime). Formally:
\[
  \widehat{\Corr}_{\text{in-sample}}(\text{OBI}_t,\, m_{t+1} - m_t) \approx 0.38 \quad \text{(tight-spread regime; treat as upper bound)}
\]

\textbf{Precision of this estimate.} With $T \approx 2{,}000$ tight-regime timestamps, the standard error of the sample correlation is:
\[
  SE(\hat{\rho}) = \sqrt{\frac{1 - \hat{\rho}^2}{T - 2}} \approx \sqrt{\frac{1 - 0.38^2}{1998}} \approx 0.021
\]
The 95\% confidence interval is $\hat{\rho} \in [0.34,\; 0.42]$ (in-sample). The $t$-statistic $t = 0.38 / 0.021 \approx 18$ confirms the correlation is highly significant within this dataset. However, this is an in-sample estimate from 2 days of backtest data; it should not be interpreted as the out-of-sample predictive power.

\begin{warningbox}{In-Sample Statistic — Handle with Care}
The 0.38 figure comes from fitting to the available backtest data. It may not generalise to live rounds, which use different price sequences. Before using OBI sizing in any round, re-estimate this correlation on that round's specific data and apply a Bonferroni correction for multiple comparisons (see Part~III statistical context). The passive-edge argument (Section~\ref{sec:obi-sizing-failed}) implies OBI sizing is unprofitable whenever the half-spread exceeds $\approx 1.3$ ticks regardless of the correlation value.
\end{warningbox}

\subsection{Why OBI Sizing Failed in Practice}
\label{sec:obi-sizing-failed}

Despite the correlation, asymmetric OBI sizing produced zero PnL improvement. The reason is a fundamental signal-vs-edge comparison:

\begin{keyconceptbox}{Passive Edge vs.\ OBI Signal}
\begin{itemize}
  \item \textbf{Passive edge per trade}: $\delta = 5$ ticks (quoting at FV $\pm 5$)
  \item \textbf{OBI directional gain}: $\E[\Delta m | \text{OBI} = 1] \approx 0.38 \times \sigma_{\Delta m} \approx 0.38 \times 1.35 \approx 0.51$ ticks
  \item \textbf{Ratio}: $5 / 0.51 \approx 10\times$
\end{itemize}
Sizing down the passive order when OBI is unfavourable sacrifices $\delta = 5$ ticks of edge to avoid an expected loss of only $0.51$ ticks. Net effect: negative. The OBI signal is only worth acting on when the half-spread $< 1$ tick (a much more aggressive quoting environment).
\end{keyconceptbox}

\subsection{Weighted OBI}

A refinement: weight each level by its distance from mid:
\[
  \text{WOBI}_t = \frac{\sum_i q_i^B e^{-\lambda(m_t - p_i^B)} - \sum_j q_j^A e^{-\lambda(p_j^A - m_t)}}{\sum_i q_i^B e^{-\lambda(m_t - p_i^B)} + \sum_j q_j^A e^{-\lambda(p_j^A - m_t)}}
\]
where $\lambda > 0$ down-weights levels far from mid (less likely to trade). Typically more predictive than simple OBI.

\section{VWAP (Volume-Weighted Average Price)}

VWAP is both a fair value estimator and a volume indicator. We defined it in Chapter~6; here we focus on its use as a trading signal.

\begin{definitionbox}{Intraday VWAP}
\[
  \text{VWAP}_t = \frac{\sum_{i=1}^{t} TP_i \cdot V_i}{\sum_{i=1}^{t} V_i}
\]
where $TP_i = \nicefrac{(H_i+L_i+C_i)}{3}$ and the sum runs from the start of the trading day. VWAP resets to the current price at the beginning of each new session.
\end{definitionbox}

\textbf{Trading signals:}
\begin{itemize}
  \item Price $>$ VWAP: Intraday bullish bias (buyers have dominated)
  \item Price $<$ VWAP: Intraday bearish bias
  \item Large institutions aim to execute near VWAP to minimise market impact
  \item In ranging markets: buy below VWAP, sell above (mean reversion to VWAP)
\end{itemize}

\textbf{VWAP Bands}: Like Bollinger Bands, add $\pm k\sigma_{\text{VWAP}}$ around VWAP for entry/exit zones.

\section{Volume Profile}

The \textbf{Volume Profile} shows the volume traded at each price level (not at each time step). It reveals where the market has spent the most time trading.

\begin{definitionbox}{Point of Control (POC) and Value Area}
\begin{itemize}
  \item \textbf{Point of Control (POC)}: The price level with the highest traded volume. A powerful support/resistance and fair value anchor.
  \item \textbf{Value Area}: The range containing $70\%$ of the total volume (analogous to $\pm 1\sigma$ for a normal distribution).
  \item \textbf{Value Area High/Low (VAH/VAL)}: The upper and lower bounds of the Value Area.
\end{itemize}
\end{definitionbox}

Price tends to return to the POC after excursions, making it a mean-reversion anchor. In Prosperity, where prices are mean-reverting around 10000 (EMERALDS), the POC is effectively fixed at 10000 by construction.

\section{On-Balance Volume (OBV)}

Joseph Granville (1963). OBV is a running cumulative sum of volume, added on up-days and subtracted on down-days.

\begin{definitionbox}{On-Balance Volume}
\[
  \text{OBV}_t = \text{OBV}_{t-1} + \begin{cases} +V_t & \text{if } C_t > C_{t-1} \\ 0 & \text{if } C_t = C_{t-1} \\ -V_t & \text{if } C_t < C_{t-1} \end{cases}
\]
\end{definitionbox}

\textbf{Intuition}: OBV treats volume as ``votes''. Heavy volume on up-days = accumulation (institutions buying). Heavy on down-days = distribution (institutions selling). The OBV trend often leads the price trend.

\textbf{Divergence}: If price makes a new high but OBV fails to confirm → bearish divergence → potential reversal.

\section{Accumulation/Distribution Line (A/D)}

Marc Chaikin (1970s). Improves on OBV by weighting volume by the Close's position within the High-Low range:

\begin{definitionbox}{Money Flow Multiplier and A/D Line}
\[
  \text{MFM}_t = \frac{(C_t - L_t) - (H_t - C_t)}{H_t - L_t} = \frac{2C_t - H_t - L_t}{H_t - L_t} \in [-1, +1]
\]
\[
  \text{MFV}_t = \text{MFM}_t \times V_t \quad \text{(Money Flow Volume)}
\]
\[
  \text{A/D}_t = \text{A/D}_{t-1} + \text{MFV}_t
\]
\end{definitionbox}

If the price closes at the high of the bar ($C_t = H_t$), MFM $= +1$ and the full volume is added. If it closes at the low ($C_t = L_t$), MFM $= -1$ and full volume is subtracted. This makes A/D more nuanced than OBV, which ignores where within the range the close is.

\section{Chaikin Money Flow (CMF)}

A normalised version of A/D:
\[
  \text{CMF}_t^{(n)} = \frac{\sum_{i=0}^{n-1}\text{MFV}_{t-i}}{\sum_{i=0}^{n-1} V_{t-i}} \in [-1, +1]
\]
CMF $> 0$: net money flowing into the asset (bullish). CMF $< 0$: outflow (bearish). More interpretable than A/D because it is bounded.

\section{Summary: Volume Indicators Comparison}

\begin{center}
\begin{tabular}{llll}
\toprule
\textbf{Indicator} & \textbf{Type} & \textbf{Uses price loc.?} & \textbf{Best signal} \\
\midrule
OBI    & Real-time book & Yes (bid/ask volume) & Short-term price pressure \\
VWAP   & Cumulative     & No (trade-weighted)  & Execution benchmark, bias \\
OBV    & Cumulative     & No (close direction) & Divergence with price \\
A/D    & Cumulative     & Yes (close in range) & Divergence, accumulation \\
CMF    & Rolling avg    & Yes                  & Bounded money flow \\
MFI    & Oscillator     & Yes ($TP \times V$)  & Volume-weighted RSI \\
Vol Profile & Histogram & N/A                  & POC, Value Area S/R \\
\bottomrule
\end{tabular}
\end{center}

\begin{takeawaybox}{Chapter 14}
\begin{itemize}
  \item OBI = $(Q^B - Q^A)/(Q^B + Q^A)$: real-time book pressure. Corr $\approx 0.38$ with next return in TOMATOES tight regime but too weak vs.\ passive edge to justify sizing adjustments.
  \item VWAP = volume-weighted price since open; key benchmark and intraday bias indicator.
  \item OBV: cumulative signed volume; divergence with price warns of trend weakness.
  \item A/D: OBV refined by MFM (Close's position in range); more nuanced.
  \item CMF: bounded $[-1,+1]$ money flow; positive = accumulation, negative = distribution.
  \item Volume Profile: distribution of volume by price; POC = highest volume price = key support/resistance.
\end{itemize}
\end{takeawaybox}

% End of Part III

% ============================================================
%
%  PART IV — TRADING STRATEGIES
%
% ============================================================
\part{Trading Strategies}

This part gives a systematic treatment of every major strategy family encountered in the codebase and in quantitative finance broadly. For each family we cover: the economic rationale, the mathematical framework, entry/exit rules with pseudocode, position sizing, performance characteristics, and failure modes. The Avellaneda-Stoikov optimal market-making model is deferred to Part~VII (Advanced Topics) where its full stochastic control derivation is presented.

\chapter{Market Making}

\section{Economic Rationale}

Market making is the strategy at the core of \texttt{trader.py}. The market maker earns the bid-ask spread by simultaneously quoting both sides of the market. Unlike directional strategies, the market maker does not take a view on price direction — they earn from the \emph{flow} of other participants crossing the spread.

\begin{keyconceptbox}{Why Market Making Is Profitable}
\begin{enumerate}[label=\arabic*.]
  \item Other participants value \textbf{immediacy}: they pay the spread to trade now rather than waiting.
  \item The market maker sells this immediacy service for the spread $S = p^A - p^B$.
  \item With balanced flow (equal buys and sells), every round-trip earns $S$ with no directional risk.
  \item With $N$ trades per day, daily PnL $\approx \nicefrac{N \cdot S}{2}$ (half-spread per leg).
\end{enumerate}
\end{keyconceptbox}

\section{The Classic Market-Making Setup}

Given fair value $FV$ and desired quote offset $\delta$, the market maker posts:
\begin{align}
  p^B &= FV - \delta \quad\text{(passive bid)}\\
  p^A &= FV + \delta \quad\text{(passive ask)}
\end{align}
Any bot that crosses $p^B$ trades with us at a profit of $FV - p^B = \delta$ ticks. Any bot crossing $p^A$ profits us $p^A - FV = \delta$ ticks. The market maker profits \emph{regardless} of which direction the bot crosses — as long as $FV$ is accurate.

\subsection{Choosing the Quote Offset $\delta$}

The optimal $\delta$ balances two competing forces:
\begin{itemize}
  \item \textbf{Too small} $\delta$: More fills (bots are more willing to cross), but each fill earns less. Also, more adverse selection risk (we're effectively quoting inside true FV, so informed traders profit from our mispricing).
  \item \textbf{Too large} $\delta$: Each fill earns more, but fill rate drops (bots less willing to cross). At $\delta > \nicefrac{S_{\text{bot}}}{2}$, our quotes are outside the bot spread and never get filled.
\end{itemize}

In Prosperity, the bot spread is fixed (e.g.\ EMERALDS: 9992/10008, half-spread $D = 8$). The bot has a deterministic threshold: it crosses our quote if and only if our price is better than its own posted price. This gives a \textbf{binary arrival model}:
\[
  \lambda(\delta) = \begin{cases} \lambda_0 & \text{if } \delta < D \\ 0 & \text{if } \delta \geq D \end{cases}
\]
Revenue per unit time is $R(\delta) = \delta \cdot \lambda(\delta)$. For $\delta < D$ this is linear and increasing; at $\delta = D$ it drops to zero. The optimisation problem is:
\[
  \max_{\delta \in \{1,2,\ldots\}} R(\delta) = \max_{\delta < D}\; \delta \cdot \lambda_0
\]
Since $R$ is strictly increasing in $\delta$ for $\delta < D$, the maximum is achieved at the largest integer strictly below $D$:
\[
  \delta^* = D - 1 = \frac{S_{\text{bot}}}{2} - 1 = 8 - 1 = 7
\]

\begin{proof}[Why this differs from the A-S result]
In Avellaneda-Stoikov, the arrival rate is $\lambda e^{-\kappa\delta}$ — a smooth, exponentially decreasing function. This creates an interior optimum: increasing $\delta$ reduces fill rate continuously, so there is a finite $\delta^*$ where marginal gain from higher edge equals marginal loss from lower fill rate. In Prosperity, the step function has no such interior: fill rate is constant for all $\delta < D$, then jumps to zero. The corner solution $\delta^* = D - 1$ follows immediately.
\end{proof}

For TOMATOES ($S_{\text{bot}} \approx 13$--$14$, $D = 6.5$--$7$), the same logic gives $\delta^* = 5$--$6$ — grid search confirms $\delta^* = 5$ (the bot's ask is at $FV + 6.5$; quoting at $FV + 5$ is safely inside).

\section{Two-Step Strategy: Taking and Making}

The full market-making algorithm in \texttt{trader.py} has two steps per iteration:

\begin{algorithm}[H]
\caption{Market-Making Algorithm (per iteration)}
\SetAlgoLined
\KwIn{Order book $\mathcal{B}_t$, fair value $FV$, position $\text{pos}$, limit $L$, offset $\delta$}
\KwOut{List of orders}
$\text{buy\_cap} \leftarrow L - \text{pos}$\;
$\text{sell\_cap} \leftarrow L + \text{pos}$\;
\textbf{// Step 1: Take immediately profitable quotes}\\
\ForEach{ask $p$ in $\mathcal{B}^A$ (ascending)}{
  \If{$p > FV$ \textbf{or} $\text{buy\_cap} \leq 0$}{\textbf{break}\;}
  $q \leftarrow \min(-\mathcal{B}^A[p],\; \text{buy\_cap})$\;
  Send buy order $(p, q)$\;
  $\text{buy\_cap} \leftarrow \text{buy\_cap} - q$\;
}
\ForEach{bid $p$ in $\mathcal{B}^B$ (descending)}{
  \If{$p < FV$ \textbf{or} $\text{sell\_cap} \leq 0$}{\textbf{break}\;}
  $q \leftarrow \min(\mathcal{B}^B[p],\; \text{sell\_cap})$\;
  Send sell order $(p, -q)$\;
  $\text{sell\_cap} \leftarrow \text{sell\_cap} - q$\;
}
\textbf{// Step 2: Post passive quotes inside bot spread}\\
\If{$\text{buy\_cap} > 0$}{Send buy order $(FV - \delta,\; \text{buy\_cap})$\;}
\If{$\text{sell\_cap} > 0$}{Send sell order $(FV + \delta,\; -\text{sell\_cap})$\;}
\end{algorithm}

\section{Inventory Skewing}

Without skewing, the market maker accumulates a one-sided position whenever the market trends. Skewing adjusts the quotes to lean in the direction that reduces inventory:

\begin{equation}
  \text{skew} = \frac{\text{pos}}{L} \times S_{\text{range}}, \qquad S_{\text{range}} \in \{0, 1, 2, \ldots\}
\end{equation}
\begin{align}
  p^B_{\text{skewed}} &= \floor{FV - \delta - \text{skew}} \\
  p^A_{\text{skewed}} &= \ceil{FV + \delta - \text{skew}}
\end{align}

When pos $> 0$ (long): skew $> 0$, so both quotes shift down. The ask becomes cheaper (encourage buying from us), the bid becomes less attractive (discourage further buying). Net effect: push inventory back toward zero.

\begin{examplebox}{Inventory Skewing in TOMATOES}
$FV = 5002$, $\delta = 5$, $L = 80$, $S_{\text{range}} = 1$, $\text{pos} = +40$. \\
$\text{skew} = 40/80 \times 1 = 0.5$. \\
$p^B = \lfloor 5002 - 5 - 0.5 \rfloor = \lfloor 4996.5 \rfloor = 4996$. \\
$p^A = \lceil 5002 + 5 - 0.5 \rceil = \lceil 5006.5 \rceil = 5007$. \\
Without skew: bid=4997, ask=5007. With skew: bid=4996 (1 tick lower), ask stays same. The sell side becomes marginally more attractive, naturally reducing long exposure.
\end{examplebox}

\section{Regime-Adaptive Market Making}

Our TOMATOES strategy adapts to the two spread regimes:
\begin{center}
\begin{tabular}{lcc}
\toprule
\textbf{Regime} & \textbf{Spread} & \textbf{AR(1) $\phi$} \\
\midrule
Tight & $\leq 13$ & 0.05 (apply correction) \\
Wide  & $= 14$   & 0.00 (no correction) \\
\bottomrule
\end{tabular}
\end{center}
In the wide regime, the price exhibits no mean reversion (AR(1) $\approx 0$), so applying a correction would introduce noise. In the tight regime, a small correction ($\phi = 0.05$) marginally improves FV accuracy.

\section{Performance Characteristics of Market Making}

\begin{center}
\begin{tabular}{ll}
\toprule
\textbf{Property} & \textbf{Typical Value / Behaviour} \\
\midrule
Sharpe ratio & Very high (2--10+) in stable conditions \\
Max drawdown & Low if inventory is controlled \\
PnL distribution & Right-skewed with rare large losses (adverse selection spikes) \\
Correlation to market & Near zero (market neutral by design) \\
Scalability & Limited by position limits and market depth \\
Main risk & Adverse selection, inventory buildup during trends \\
\bottomrule
\end{tabular}
\end{center}

\begin{takeawaybox}{Chapter 15}
\begin{itemize}
  \item Market makers earn the spread by simultaneously quoting bid and ask, providing immediacy.
  \item Two-step algorithm: (1) take any bot quotes better than FV; (2) post passive quotes inside bot spread.
  \item Optimal offset: just inside the bot half-spread ($\delta^* = $ bot half-spread $- 1$).
  \item Inventory skewing shifts both quotes in the direction that reduces inventory. Formula: skew $= (\text{pos}/L) \times S_{\text{range}}$.
  \item Regime adaptation: switch AR(1) coefficient based on spread regime to avoid adding noise in random-walk conditions.
\end{itemize}
\end{takeawaybox}

% ──────────────────────────────────────────────────────────────────────────────

\chapter{Mean Reversion Strategies}

\section{The Mean Reversion Hypothesis}

Mean reversion is the tendency of a price series to return to its long-run average after deviating from it. It is the statistical opposite of momentum. The world of trading strategies is fundamentally divided between these two camps:

\begin{center}
\begin{tabular}{lll}
\toprule
\textbf{Property} & \textbf{Mean Reversion} & \textbf{Momentum / Trend} \\
\midrule
Price autocorrelation & Negative ($\rho < 0$) & Positive ($\rho > 0$) \\
Trading rule & Buy dips, sell rips & Buy breakouts, sell breakdowns \\
Typical regime & Ranging, stable & Trending \\
Example products & EMERALDS, equities index & Commodities, FX trends \\
Hurst exponent & $H < 0.5$ & $H > 0.5$ \\
AR(1) coefficient & $\phi \in (-1, 0)$ & $\phi \in (0, 1)$ \\
\bottomrule
\end{tabular}
\end{center}

\section{The Ornstein-Uhlenbeck Process}

The continuous-time model for mean reversion is the \textbf{Ornstein-Uhlenbeck (OU) process}:

\begin{definitionbox}{Ornstein-Uhlenbeck Process}
\[
  dX_t = \kappa(\mu - X_t)\,\dd t + \sigma\,\dd W_t
\]
where:
\begin{itemize}
  \item $\kappa > 0$: speed of mean reversion (how fast $X$ returns to $\mu$)
  \item $\mu$: long-run mean
  \item $\sigma$: volatility of the diffusion term
  \item $W_t$: standard Brownian motion
\end{itemize}
\end{definitionbox}

\textbf{Intuition}: The drift term $\kappa(\mu - X_t)$ acts like a spring. When $X_t > \mu$, the drift is negative (pulls down). When $X_t < \mu$, drift is positive (pulls up). The strength of the pull is proportional to the displacement, like Hooke's Law.

\subsection{Solving the OU SDE}

Applying Itô's formula (multiply both sides by $e^{\kappa t}$):
\begin{align}
  X_t &= \mu + (X_0 - \mu)e^{-\kappa t} + \sigma\int_0^t e^{-\kappa(t-s)}\,\dd W_s
\end{align}
Key statistics:
\begin{align}
  \E[X_t] &= \mu + (X_0 - \mu)e^{-\kappa t} \to \mu \quad \text{as } t \to \infty \\
  \Var(X_t) &= \frac{\sigma^2}{2\kappa}\left(1 - e^{-2\kappa t}\right) \to \frac{\sigma^2}{2\kappa} \quad \text{(stationary variance)} \\
  \Corr(X_t, X_{t+h}) &= e^{-\kappa h} \quad \text{(exponential autocorrelation decay)}
\end{align}

\textbf{Half-life of mean reversion}: $t_{1/2} = \nicefrac{\ln 2}{\kappa}$. This is the time for the deviation from $\mu$ to halve on average. A half-life of 5 days means deviations typically resolve within a week.

\subsection{Discrete-Time OU: The AR(1) Model}

Discretising the OU SDE at intervals $\Delta t$ gives the AR(1) model. \textbf{Derivation:} the OU SDE $\dd X_t = -\kappa(X_t - \mu)\dd t + \sigma\dd W_t$ has exact solution:
\[
  X_t = \mu + e^{-\kappa\Delta t}(X_{t-1} - \mu) + \sigma\sqrt{\frac{1-e^{-2\kappa\Delta t}}{2\kappa}}\,\varepsilon_t, \quad \varepsilon_t \sim \mathcal{N}(0,1)
\]
Defining $\phi = e^{-\kappa\Delta t}$ and $\sigma_\varepsilon^2 = \sigma^2(1-\phi^2)/(2\kappa)$, this becomes the AR(1):
\[
  X_t = \phi X_{t-1} + (1-\phi)\mu + \sigma_\varepsilon\varepsilon_t, \qquad \phi = e^{-\kappa\Delta t} \in (0,1)
\]
The correspondence $\phi \leftrightarrow \kappa$ lets you convert between the OU speed-of-reversion $\kappa$ and the discrete AR(1) coefficient $\phi$: $\kappa = -\ln\phi/\Delta t$, so half-life $= \ln 2/\kappa = -\Delta t\ln 2/\ln\phi$. Mean reversion requires $|\phi| < 1$. When $\phi < 0$ (over-shooting), the process oscillates above and below $\mu$ every period. In TOMATOES, the AR(1) coefficient of mid-price changes is $\approx -0.30$ (tight regime), indicating \emph{direct mean reversion} in returns (the price change itself reverses direction).

\section{Trading the Z-Score}

The universal mean-reversion signal is the \textbf{z-score}: how many standard deviations is the current value from its mean?

\begin{definitionbox}{Z-Score Signal}
\[
  z_t = \frac{X_t - \mu}{\sigma}
\]
where $\mu$ and $\sigma$ are estimated from a rolling window. Trading rules:
\begin{itemize}
  \item $z_t > +z^*$: \textbf{Short} $X$ (too far above mean, expect reversion down)
  \item $z_t < -z^*$: \textbf{Long} $X$ (too far below mean, expect reversion up)
  \item $|z_t| < z^{\text{exit}}$: Exit (mean has been reached)
\end{itemize}
Common thresholds: $z^* = 2.0$, $z^{\text{exit}} = 0.5$.
\end{definitionbox}

\begin{examplebox}{Z-Score Trading}
EMERALDS mid-price: mean $= 10000$, std $= 0.72$. \\
Current price $= 10002$. $z = (10002 - 10000)/0.72 = 2.78$. \\
Signal: Short (sell) — price is 2.78 std above mean, expect reversion. \\
Exit when $|z| < 0.5$, i.e.\ price returns to $10000 \pm 0.36$.
\end{examplebox}

\section{Pairs Trading (Overview)}

Pairs trading extends mean reversion from a single asset to the \emph{spread between two related assets}. If asset A and asset B are \textbf{cointegrated} — they share a common stochastic trend — then the spread $Z_t = A_t - \beta B_t$ is stationary (mean-reverting) even if neither $A_t$ nor $B_t$ is individually stationary.

The key idea is to estimate a \textbf{hedge ratio} $\beta$ by OLS regression ($A_t = \alpha + \beta B_t + \varepsilon_t$), test stationarity of residuals with the ADF test (Engle-Granger procedure), then trade the standardised spread $z_t = (Z_t - \bar{Z})/\sigma_Z$: enter long when $z_t < -z^*$, enter short when $z_t > +z^*$, exit when $|z_t|$ reverts to near zero.

\begin{keyconceptbox}{Full Treatment in Chapter 35}
Pairs trading is covered in full depth in Chapter~35 (Part~VIII), including: the cointegration definition, the Engle-Granger two-step procedure, critical value subtleties, the Johansen test for multiple cointegrating vectors, the OU model for the spread, z-score signal generation, position sizing, and a Prosperity implementation template. Refer there for the complete mathematical treatment.
\end{keyconceptbox}

\section{Statistical Arbitrage}

Statistical arbitrage (``stat arb'') is the generalisation of pairs trading to portfolios of $N$ assets. Instead of one spread, we exploit \emph{factors}:

\begin{equation}
  r_{i,t} = \alpha_i + \sum_{k=1}^K \beta_{ik} F_{k,t} + \varepsilon_{i,t}
\end{equation}
where $F_{k,t}$ are common factors (market, sector, style) and $\varepsilon_{i,t}$ is the idiosyncratic return. Stat arb trades the residuals $\varepsilon_{i,t}$: buy stocks with large negative $\varepsilon$ (underperformed their factors), short stocks with large positive $\varepsilon$ (overperformed).

The portfolio is market-neutral ($\sum_i w_i \beta_{ik} = 0$ for all $k$) and earns from the mean reversion of $\varepsilon_{i,t}$.

\section{When Mean Reversion Strategies Fail}

\begin{warningbox}{Regime Change: The Mean Reversion Killer}
\begin{enumerate}[label=\arabic*.]
  \item \textbf{Structural break}: The mean $\mu$ changes permanently (e.g.\ a company's fundamental value shifts). Z-score signals are against a wrong mean.
  \item \textbf{Momentum takeover}: In trending conditions, mean-reversion signals are consistently wrong. A mean-reversion strategy will fight the trend and bleed.
  \item \textbf{Spread blowup in pairs}: Two cointegrated assets can temporarily diverge far beyond historical norms (e.g.\ during a crisis). The trade moves against you before reverting — if you are forced to close early, you lock in a loss.
  \item \textbf{Overfitting}: The hedge ratio $\beta$ and z-score thresholds $z^*$ estimated in-sample may not hold out-of-sample.
\end{enumerate}
\end{warningbox}

\begin{takeawaybox}{Chapter 16}
\begin{itemize}
  \item Mean reversion: price has negative autocorrelation; buy dips, sell rallies. Requires $H < 0.5$ or AR(1) $|\phi| < 1$.
  \item OU process: $dX = \kappa(\mu - X)dt + \sigma dW$. Half-life $= \ln2/\kappa$. Discretised = AR(1) with $\phi = e^{-\kappa\Delta t}$.
  \item Z-score signal: enter when $|z| > z^*$; exit when $|z| < z^{\text{exit}}$.
  \item Pairs trading overview: exploit cointegration to trade the spread $Z = A - \beta B$. Full treatment in Chapter~35.
  \item Stat arb: generalised pairs trading across $N$ assets; residuals from factor model are mean-reverting.
  \item Main risks: structural breaks, momentum regimes, spread blowup, overfitting of $\beta$.
\end{itemize}
\end{takeawaybox}

% ──────────────────────────────────────────────────────────────────────────────

\chapter{Trend Following and Momentum Strategies}

\section{The Momentum Anomaly}

Momentum is the opposite of mean reversion. In trending markets, recent winners continue to outperform (momentum), and strategies that buy strength and sell weakness profit. The \textbf{momentum anomaly} — that past 12-month returns predict the next month's returns — was documented by Jegadeesh \& Titman (1993) and remains one of the most robust anomalies in finance.

\textbf{Why does momentum exist?}
\begin{enumerate}[label=\arabic*.]
  \item \textbf{Underreaction}: Investors are slow to update beliefs when new information arrives. Price drifts toward full information value over weeks/months.
  \item \textbf{Herding}: Trend following creates self-reinforcing feedback: buying pushes prices up, attracting more buyers.
  \item \textbf{Risk-based}: Momentum stocks may have higher risk during bad times, justifying higher expected returns.
\end{enumerate}

\section{Moving Average Crossover Systems}

The simplest trend-following rule: go long when the fast MA crosses above the slow MA.

\begin{definitionbox}{Dual Moving Average Crossover}
\begin{equation}
  \text{Signal}_t = \begin{cases}
    +1 & \text{if } \text{MA}^{(f)}_t > \text{MA}^{(s)}_t \text{ and } \text{MA}^{(f)}_{t-1} \leq \text{MA}^{(s)}_{t-1} \quad \text{(golden cross)} \\
    -1 & \text{if } \text{MA}^{(f)}_t < \text{MA}^{(s)}_t \text{ and } \text{MA}^{(f)}_{t-1} \geq \text{MA}^{(s)}_{t-1} \quad \text{(death cross)} \\
    \text{unchanged} & \text{otherwise}
  \end{cases}
\end{equation}
where $f < s$ (fast period, slow period).
\end{definitionbox}

\textbf{Properties:}
\begin{itemize}
  \item Always in the market (long or short — no neutral state)
  \item Signal fires \emph{after} the trend has started (inherent lag)
  \item Works well in trending markets; generates many false signals in ranging markets
  \item The longer the periods, the fewer but higher-quality signals
\end{itemize}

\subsection{Triple Moving Average System}

Uses three MAs (fast, medium, slow):
\begin{itemize}
  \item Go long only when fast $>$ medium $>$ slow (all aligned upward)
  \item Neutral when alignment is mixed
  \item Go short when fast $<$ medium $<$ slow
\end{itemize}
This filter reduces false signals in ranging markets at the cost of later entries.

\section{The Turtle Trading System}

Richard Dennis bet $\$1{,}000$ that he could teach a group of non-traders to trade profitably. The ``Turtles'' used a pure mechanical breakout system (1983), now public knowledge. It is one of the most thoroughly documented trend-following systems.

\begin{keyconceptbox}{Turtle Rules (Simplified)}
\textbf{Entry System 1 (Short-term)}: \\
Go long on 20-day high breakout; go short on 20-day low breakout. \\
Exit when price hits the 10-day opposite extreme.

\textbf{Entry System 2 (Long-term)}: \\
Go long on 55-day high; short on 55-day low. \\
Exit at 20-day opposite extreme.

\textbf{Position sizing}: $N = \text{ATR}_{20}$. Trade size = $\nicefrac{1\%\times\text{Account}}{N\times\text{Dollar value per point}}$.\\
Each trade risks exactly $1\%$ of account equity.

\textbf{Adding units}: Add to winning positions at every $\nicefrac{1}{2}N$ favourable move (pyramiding).

\textbf{Stop loss}: $2N$ from entry price.
\end{keyconceptbox}

\begin{examplebox}{Turtle Position Sizing}
Account: \$100,000. ATR = 50. Price per point = \$1. \\
$\text{Risk per trade} = 1\% \times 100{,}000 = \$1{,}000$. \\
Stop loss = $2 \times 50 = 100$ points. \\
Units = $\$1{,}000 / (100 \times \$1) = 10$ units per trade.
\end{examplebox}

\section{Donchian Breakout Strategy}

The foundational breakout system, used by the Turtles:
\begin{align}
  p^{\text{high}}_t &= \max(H_{t-n+1}, \ldots, H_t) \quad \text{($n$-day high)} \\
  p^{\text{low}}_t &= \min(L_{t-n+1}, \ldots, L_t) \quad \text{($n$-day low)}
\end{align}
\begin{itemize}
  \item Buy at $p^{\text{high}}_t + 1$ tick (price exceeds $n$-day high)
  \item Sell short at $p^{\text{low}}_t - 1$ tick
  \item Trail stop at opposite $m$-day extreme ($m < n$)
\end{itemize}

The breakout philosophy: a new $n$-day high is evidence that buyers are willing to pay the highest price in $n$ days — this is information that the trend is intact.

\section{Time-Series Momentum (TSMOM)}

Unlike cross-sectional momentum (buy best, short worst), TSMOM trades an asset based on its own past return:
\begin{equation}
  \text{Signal}_t = \text{sign}\!\left(\sum_{k=1}^{K} r_{t-k}\right)
\end{equation}
Go long if the sum of the past $K$ returns is positive; short if negative. Moskowitz, Ooi \& Pedersen (2012) showed TSMOM works across 58 liquid markets with a 12-month lookback.

The economic intuition: institutional investors rebalance slowly. When an asset has been rising for 12 months, more institutions are still in the process of buying it → continued upward pressure.

\section{Breakout Strategies}

Beyond Donchian, any significant crossing of a key level can trigger a breakout strategy:

\begin{enumerate}[label=\arabic*.]
  \item \textbf{Range breakout}: Buy above the top of a trading range; short below the bottom.
  \item \textbf{Volatility breakout}: Enter when price moves more than $k\times\text{ATR}$ from the open.
  \item \textbf{Opening range breakout}: Trade the first $N$ minutes of the session to establish a range; trade breakouts from that range for the rest of the day.
  \item \textbf{Pattern breakout}: Enter on confirmed break of a triangle, flag, or rectangle.
\end{enumerate}

\subsection{Volume Confirmation for Breakouts}

A breakout is more likely to be genuine when accompanied by high volume. A simple filter:
\[
  \text{Confirmed breakout} = \text{breakout condition} \quad \text{AND} \quad V_t > 1.5 \times \text{SMA}^{(20)}(V)
\]

\section{Momentum Strategy Performance Characteristics}

\begin{center}
\begin{tabular}{ll}
\toprule
\textbf{Property} & \textbf{Typical Value / Behaviour} \\
\midrule
Win rate & Low (30--45\%), but large wins dominate \\
Profit factor & $> 1$ due to fat-tailed wins (``let winners run'') \\
Max drawdown & Can be severe (30--50\%) during choppy markets \\
Sharpe ratio & 0.5--1.0 (lower than market making but scalable) \\
Best regime & Trending, low mean reversion \\
Worst regime & Choppy, high mean reversion \\
Correlation & Often clusters with other trend followers \\
\bottomrule
\end{tabular}
\end{center}

\begin{takeawaybox}{Chapter 17}
\begin{itemize}
  \item Momentum exists because of underreaction: prices drift toward full information over time.
  \item MA crossover: go long when fast MA $>$ slow MA. Lagged but mechanical.
  \item Turtle system: breakout-based, ATR position sizing, $1\%$ per-trade risk, pyramiding.
  \item Donchian breakout: buy $n$-day high, short $n$-day low, trail stop at $m$-day opposite.
  \item TSMOM: go long if sum of past $K$ returns is positive. Works across 58 markets on 12-month lookback.
  \item Momentum strategies have low win rates but large average wins; suffer in ranging markets.
\end{itemize}
\end{takeawaybox}

% ──────────────────────────────────────────────────────────────────────────────

\chapter{Arbitrage Strategies}

\section{Pure (Risk-Free) Arbitrage}

\begin{definitionbox}{Arbitrage}
An \textbf{arbitrage} is a trading strategy that generates a non-negative payoff in all states of the world and a strictly positive payoff in at least one state, with zero initial investment. In frictionless markets, arbitrage opportunities are instantaneously eliminated by rational traders.
\end{definitionbox}

\textbf{The Law of One Price}: Two assets with identical future cash flows must have the same price today. Any violation is an arbitrage.

\textbf{Classic examples:}
\begin{itemize}
  \item \textbf{Cash-and-carry}: If futures price $F > S e^{rT}$ (forward price), buy the spot, sell the futures, earn risk-free profit.
  \item \textbf{Triangle arbitrage}: If EUR/USD $\times$ USD/GBP $\neq$ EUR/GBP, cycle through three currencies for risk-free profit.
  \item \textbf{ETF arbitrage}: If an ETF trades at a premium to its NAV (Net Asset Value), authorised participants create new ETF shares (buy the basket, sell the ETF), eliminating the premium.
\end{itemize}

In Prosperity, the conversion mechanism (returning from `run()`) enables a form of arbitrage: converting positions between different representations when the exchange rate offers a profit net of transport fees and tariffs.

\section{Statistical Arbitrage}

We covered stat arb in Chapter~16 (from the mean-reversion perspective). Here we frame it as an arbitrage:

\begin{keyconceptbox}{Stat Arb as Approximate Arbitrage}
Statistical arbitrage is not risk-free. It is a strategy that generates positive expected returns by exploiting temporary mispricings that are likely (but not certain) to correct. The ``arbitrage'' is statistical: in expectation across many trades, the strategy profits. Individual trades may lose.
\end{keyconceptbox}

\section{Convertible Arbitrage}

In fixed income, a \textbf{convertible bond} can be converted into equity at a fixed ratio. The bond price $= $ bond floor $+$ conversion option value. If the market price deviates:
\begin{itemize}
  \item Buy the convertible bond (cheap)
  \item Short $\Delta$ shares of the underlying stock (hedge the equity exposure)
  \item Earn from the convergence of the convertible to its fair value
\end{itemize}

\section{Index Arbitrage}

When an index futures contract trades at a price different from the theoretical fair value:
\[
  F^{\text{fair}} = S \cdot e^{(r - q)T}
\]
where $r$ = risk-free rate, $q$ = dividend yield, $T$ = time to expiry. Buy/sell the futures and the opposite in the basket of stocks to lock in the deviation.

\section{Merger Arbitrage (Risk Arbitrage)}

When company A announces acquisition of company B at price $P_A$:
\begin{itemize}
  \item B's stock jumps to near $P_A$ but trades at a slight discount (deal risk)
  \item The spread $P_A - P_B$ is the \textbf{arbitrage spread}
  \item Buy B, (sometimes) short A; earn the spread if the deal closes
  \item Risk: the deal fails → B's stock falls back to pre-announcement level
\end{itemize}

\begin{equation}
  \text{Expected return} = P_A \cdot P(\text{deal closes}) - P_B^{\text{fail}} \cdot P(\text{deal fails}) - P_B
\end{equation}

\begin{takeawaybox}{Chapter 18}
\begin{itemize}
  \item Pure arbitrage: zero investment, non-negative payoff always, positive payoff in some state. Eliminated instantly in efficient markets.
  \item Stat arb: positive \emph{expected} return by exploiting mean-reverting mispricings; involves real risk.
  \item Cash-and-carry, triangle, ETF arbitrage exploit Law of One Price violations.
  \item In Prosperity: conversion requests (returning a `conversions` integer) enable arbitrage between exchange and conversion rates net of transport fees.
\end{itemize}
\end{takeawaybox}

% ──────────────────────────────────────────────────────────────────────────────

\chapter{Carry, Factor, and Other Strategy Families}

\section{Carry Strategies}

\begin{definitionbox}{Carry}
The \textbf{carry} of an asset is the return earned from holding it, assuming its price does not change. It is the ``income'' component of return.
\end{definitionbox}

\begin{center}
\begin{tabular}{lll}
\toprule
\textbf{Asset class} & \textbf{Carry source} & \textbf{Strategy} \\
\midrule
FX & Interest rate differential & Borrow low-rate, invest high-rate currency \\
Fixed income & Yield spread & Long high-yield, short low-yield bonds \\
Commodities & Roll yield & Long backwardated contracts, short contangoed \\
Equities & Dividend yield & Long high-dividend stocks \\
\bottomrule
\end{tabular}
\end{center}

\subsection{FX Carry Trade}

The uncovered interest rate parity (UIP) predicts that high-interest-rate currencies should depreciate to offset the interest rate advantage. Empirically, UIP fails: high-rate currencies do \emph{not} depreciate on average. This creates the carry trade opportunity:

\begin{enumerate}[label=\arabic*.]
  \item Borrow in low-interest currency (e.g.\ JPY at $0.1\%$)
  \item Convert and invest in high-interest currency (e.g.\ AUD at $4\%$)
  \item Earn the differential (carry) $= 4\% - 0.1\% = 3.9\%$ per year
  \item Risk: AUD depreciates → loses on FX even while earning carry
\end{enumerate}

\textbf{Carry crash}: During risk-off events (crises), carry trades unwind simultaneously, causing sharp losses. This creates the ``carry crisis'' correlation: carry strategies lose exactly when you can least afford it.

\subsection{Commodity Roll Yield}

Futures markets have a \textbf{term structure}: futures prices vary with expiry date. Two structures:
\begin{itemize}
  \item \textbf{Contango}: $F(T_2) > F(T_1)$ for $T_2 > T_1$. Rolling from near to far contract is costly (roll yield is negative).
  \item \textbf{Backwardation}: $F(T_2) < F(T_1)$. Rolling earns a positive yield.
\end{itemize}
Carry strategy: long backwardated commodities (positive roll) + short contangoed commodities (negative roll).

\section{Factor Investing}

Factor investing (also called ``smart beta'') systematically exploits empirically documented return premia:

\begin{center}
\begin{tabular}{llll}
\toprule
\textbf{Factor} & \textbf{Definition} & \textbf{Discovery} & \textbf{Long side} \\
\midrule
Value    & Low price-to-fundamentals & Fama \& French 1992 & Cheap stocks \\
Size     & Small market cap & Banz 1981 & Small caps \\
Momentum & High 12-1 month return & Jegadeesh \& Titman 1993 & Past winners \\
Quality  & High profitability/safety & Novy-Marx 2013 & High ROE, low debt \\
Low Vol  & Low historical volatility & Ang et al.\ 2006 & Low-beta stocks \\
Carry    & High yield/carry & Koijen et al.\ 2018 & High-yield assets \\
\bottomrule
\end{tabular}
\end{center}

\textbf{Factor construction (long-short portfolio):}
\begin{enumerate}[label=\arabic*.]
  \item Rank all assets by the factor score
  \item Long the top decile (quintile), short the bottom decile
  \item Rebalance monthly (or quarterly)
  \item The return is the ``factor premium''
\end{enumerate}

\textbf{The Fama-French Three-Factor Model:}
\begin{equation}
  r_{i,t} - r_f = \alpha_i + \beta_i^M(r_M - r_f) + \beta_i^{\text{SMB}} r_{\text{SMB}} + \beta_i^{\text{HML}} r_{\text{HML}} + \varepsilon_i
\end{equation}
where SMB = Small Minus Big (size factor), HML = High Minus Low (value factor).

\section{High-Frequency Trading Strategies}

High-Frequency Trading (HFT) exploits ultra-short-term market inefficiencies at microsecond timescales. In a Prosperity context, the ``bots'' behave analogously to HFT market makers.

\begin{center}
\begin{tabular}{lll}
\toprule
\textbf{HFT Strategy} & \textbf{Description} & \textbf{Edge} \\
\midrule
Electronic market making & Continuous bid/ask quoting & Spread income \\
Statistical arbitrage & Exploit cross-asset correlations at HF & Speed + correlation \\
Latency arbitrage & React to stale quotes faster than others & Co-location, speed \\
Momentum ignition & Trigger momentum, then trade against it & Speed + manipulation \\
Order anticipation & Detect large institutional orders & Order flow signals \\
\bottomrule
\end{tabular}
\end{center}

\section{Event-Driven Strategies}

Event-driven strategies exploit price reactions to corporate events:
\begin{itemize}
  \item \textbf{Earnings announcements}: Buy stocks expected to beat earnings estimates (earnings drift)
  \item \textbf{M\&A}: Merger arbitrage (Chapter~18)
  \item \textbf{Index rebalancing}: When a stock is added to a major index, index funds must buy it → predictable buying pressure
  \item \textbf{Dividends}: Ex-dividend date creates predictable price behaviour
  \item \textbf{Options expiration}: PIN risk and gamma effects near expiry
\end{itemize}

\section{Strategy Regime Summary}

Different strategies work in different market regimes. A sophisticated trader identifies the current regime and deploys the appropriate strategy:

\begin{center}
\begin{tabular}{lcccc}
\toprule
\textbf{Strategy} & \textbf{Trending} & \textbf{Ranging} & \textbf{Volatile} & \textbf{Crisis} \\
\midrule
Market Making      & Poor      & Excellent & OK         & Poor \\
Mean Reversion     & Poor      & Excellent & OK         & Poor \\
Trend Following    & Excellent & Poor      & OK         & Good \\
Statistical Arb    & OK        & Good      & OK         & Poor \\
Carry              & OK        & OK        & Poor       & Very poor \\
Momentum           & Excellent & Poor      & OK         & Varies \\
\bottomrule
\end{tabular}
\end{center}

\begin{takeawaybox}{Chapter 19}
\begin{itemize}
  \item Carry = income from holding an asset; FX carry = borrowing low-rate, investing high-rate. Risk: carry crashes during crises.
  \item Commodity carry: long backwardated (positive roll yield), short contangoed contracts.
  \item Factor investing: systematic long-short on empirical return premia (value, momentum, quality, size, low-vol, carry).
  \item Fama-French 3-factor: market, SMB (size), HML (value) explain cross-sectional equity returns.
  \item HFT: sub-millisecond market making, latency arb, stat arb; the ``bots'' in Prosperity are analogous.
  \item No strategy works in all regimes — regime identification is as important as the strategy itself.
\end{itemize}
\end{takeawaybox}

% End of Part IV

% ============================================================
%
%  PART V — RISK MANAGEMENT & PORTFOLIO CONSTRUCTION
%
% ============================================================
\part{Risk Management \& Portfolio Construction}

Risk management is not optional. Every strategy, however theoretically sound, can destroy capital if sized incorrectly or left unhedged. This part builds the mathematical foundation for \emph{how much} to trade, \emph{how to combine} positions, and \emph{how to limit} losses. All results connect directly to the position sizing and inventory management choices in \texttt{trader.py}.

\chapter{Position Sizing Methods}

\section{Why Position Sizing Matters More Than Entry Signals}

A surprising but well-documented result: given a strategy with positive expected value, the choice of position sizing determines whether you go bankrupt, grow slowly, or grow optimally. The entry signal is almost secondary.

\begin{examplebox}{Same Strategy, Different Outcomes}
Strategy: each trade wins \$2 with probability 60\%, loses \$1 with probability 40\%. $\E[\text{profit}] = 0.6\times2 - 0.4\times1 = +0.80 > 0$. \\
Bet 10\% of capital per trade: after 100 trades, median capital $\approx 2.8\times$ initial. \\
Bet 50\% of capital per trade: ruin is almost certain within 20 trades. \\
Bet 40\% of capital per trade (near Kelly): median capital $\approx 5.9\times$ initial, but high variance.
\end{examplebox}

\section{Fixed Fractional Sizing}

\begin{definitionbox}{Fixed Fractional Position Sizing}
Risk a fixed fraction $f$ of current capital on each trade:
\[
  \text{Trade size} = \frac{f \times \text{Capital}}{\text{Risk per unit}}
\]
where ``risk per unit'' is the loss per unit if the stop is hit (e.g.\ $k \times \text{ATR}$). The fraction $f$ is typically 1--2\%.
\end{definitionbox}

\textbf{Properties:}
\begin{itemize}
  \item Positions grow with winning streaks (compound growth)
  \item Positions shrink after losses (automatic drawdown protection)
  \item Never risks more than $f\times 100\%$ of capital on one trade in percentage terms
  \item The Kelly Criterion (Chapter~21) gives the theoretically optimal $f$
\end{itemize}

\begin{examplebox}{Fixed Fractional Sizing}
Capital: 100{,}000. $f = 1\%$. Entry at 5000, stop at 4950 (risk = 50 per unit). \\
Trade size $= \nicefrac{0.01 \times 100{,}000}{50} = 20$ units. \\
If stop is hit: loss $= 20 \times 50 = 1{,}000 = 1\%$ of capital. Exactly as intended.
\end{examplebox}

\section{Fixed Ratio Sizing}

Ryan Jones (1999). The fixed ratio method links position increases to accumulated profits:
\[
  \text{Next unit added when cumulative profit} = \Delta \times N(N-1)/2
\]
where $N$ = current number of units and $\Delta$ = the ``delta'' parameter. More conservative than fixed fractional during drawdowns, more aggressive during winning streaks. Less commonly used in practice.

\section{Volatility-Adjusted Sizing}

Size inversely proportional to recent volatility, so every trade has roughly equal dollar risk:
\begin{equation}
  \text{Position size} = \frac{\text{Target dollar volatility}}{\sigma_t \times p_t}
\end{equation}
where $\sigma_t$ is the historical volatility (e.g.\ 20-day) and $p_t$ is the current price. If you target \$500 daily volatility per position and a stock has daily vol 2\% at price \$100:
\[
  \text{Shares} = \frac{500}{0.02 \times 100} = 250 \text{ shares}
\]
This is the basis of \textbf{risk parity} (Chapter~22) and the Turtle ATR sizing method.

\section{Equal Weighting}

The simplest portfolio approach: allocate equal capital to each position. Requires rebalancing when prices move to maintain equal weighting. Surprisingly hard to beat in practice (DeMiguel et al.\ 2009 showed $1/N$ portfolios often beat Markowitz on out-of-sample data).

\begin{takeawaybox}{Chapter 20}
\begin{itemize}
  \item Position sizing is as important as the entry signal; wrong sizing can ruin a profitable strategy.
  \item Fixed fractional: risk $f\%$ of capital per trade. Size = $(f \times \text{Capital})/\text{Risk per unit}$.
  \item Volatility-adjusted sizing: size $\propto 1/\sigma$; equalises dollar risk across positions.
  \item Equal weighting: $1/N$ per position; simple and surprisingly competitive in out-of-sample tests.
\end{itemize}
\end{takeawaybox}

\chapter{The Kelly Criterion}

The Kelly Criterion is the most important result in position sizing. It gives the exact fraction of capital to bet in order to \emph{maximise the long-run growth rate} of wealth.

\section{Derivation from First Principles}

\textbf{Setup}: Each bet is independent. Probability of winning = $p$, probability of losing = $q = 1-p$. Bet fraction $f$ of capital. Win: gain $f \cdot b$ (net odds $b$). Lose: lose $f$.

After $n$ bets with $W$ wins and $L$ losses ($W + L = n$), wealth is:
\begin{equation}
  W_n = W_0 \cdot (1 + fb)^W \cdot (1 - f)^L
\end{equation}

\textbf{Per-bet growth rate} (geometric mean):
\begin{equation}
  G(f) = \frac{1}{n}\ln\!\left(\frac{W_n}{W_0}\right) = \frac{W}{n}\ln(1+fb) + \frac{L}{n}\ln(1-f)
\end{equation}

By the Law of Large Numbers, $W/n \to p$ and $L/n \to q$ as $n \to \infty$:
\begin{equation}
  G(f) \to g(f) = p\ln(1+fb) + q\ln(1-f)
\end{equation}

\textbf{Maximise} over $f$ by setting $\frac{dg}{df} = 0$:
\begin{align}
  \frac{dg}{df} &= \frac{pb}{1+fb} - \frac{q}{1-f} = 0 \notag \\
  pb(1-f) &= q(1+fb) \notag \\
  pb - pbf &= q + qfb \notag \\
  pb - q &= fb(p + q) = fb \notag \\
  f^* &= \frac{pb - q}{b} = p - \frac{q}{b}
\end{align}

\begin{definitionbox}{Kelly Criterion}
The optimal fraction of capital to bet that maximises the long-run growth rate is:
\[
  f^* = \frac{pb - q}{b} = p - \frac{1-p}{b}
\]
where $p$ = probability of winning, $q = 1-p$ = probability of losing, $b$ = net odds (gain per unit risked).

Equivalently: $f^* = \frac{\text{edge}}{\text{odds}}$ where $\text{edge} = pb - q$ and $\text{odds} = b$.
\end{definitionbox}

\subsection{Continuous-Time Kelly}

For a continuously traded asset with log-normal returns ($\mu$, $\sigma$), the Kelly fraction is:
\begin{equation}
  f^* = \frac{\mu - r_f}{\sigma^2}
\end{equation}
where $\mu - r_f$ is the excess return (``edge'') and $\sigma^2$ is the variance (``risk''). This is the \textbf{Sharpe ratio} divided by $\sigma$: $f^* = \text{SR}/\sigma$.

\begin{proof}[Derivation: Continuous-Time Kelly from Log-Wealth Maximisation]
Let total wealth $W_t$ follow a GBM, where fraction $f$ is invested in a risky asset ($\mu$, $\sigma$) and fraction $1-f$ earns the risk-free rate $r_f$:
\[
  \frac{dW_t}{W_t} = \bigl[r_f + f(\mu - r_f)\bigr]\,dt + f\sigma\,dB_t
\]
Applying Itô's lemma to $\ln W_t$:
\[
  d\ln W_t = \Bigl[r_f + f(\mu - r_f) - \frac{f^2\sigma^2}{2}\Bigr]dt + f\sigma\,dB_t
\]
The expected log-growth rate (the quantity the Kelly criterion maximises) is:
\[
  G(f) = r_f + f(\mu - r_f) - \frac{f^2\sigma^2}{2}
\]
This is a concave quadratic in $f$. Taking the first-order condition:
\[
  G'(f) = (\mu - r_f) - f\sigma^2 = 0 \quad\Rightarrow\quad f^* = \frac{\mu - r_f}{\sigma^2}
\]
The second derivative $G''(f) = -\sigma^2 < 0$ confirms this is a maximum.
The maximised growth rate is $G(f^*) = r_f + \frac{(\mu-r_f)^2}{2\sigma^2} = r_f + \frac{\text{SR}^2}{2}$ (the ``Kelly growth rate'').
\end{proof}

For a portfolio, the Kelly fraction for each asset $i$ is determined by the vector:
\begin{equation}
  \mathbf{f}^* = \boldsymbol{\Sigma}^{-1}(\boldsymbol{\mu} - r_f \mathbf{1})
\end{equation}
where $\boldsymbol{\Sigma}$ is the covariance matrix. This is proportional to the \textbf{mean-variance optimal portfolio}.

\section{Numerical Examples}

\begin{examplebox}{Kelly for EMERALDS Market Making}
Each passive order fills with some probability and earns 7 ticks. Assume:
\begin{itemize}
  \item Fill probability per iteration (either side): $p = 0.929$ (92.9\% fill rate from simulation)
  \item Net edge per fill: $b = 7$ ticks
  \item Per iteration: win probability $p$ with gain $7$, lose probability $q = 0.071$ with loss $0$ (unfilled = zero profit, not a loss)
\end{itemize}
Since the strategy can't lose principal (worst case: zero fill, zero PnL), the Kelly formula doesn't directly apply in this form. The position limit $L = 80$ acts as the natural Kelly cap — beyond $L$, additional positions cannot be taken.

More relevantly: if the strategy's annual Sharpe ratio $\approx 5$ and annual volatility $\sigma \approx 2\%$, Kelly fraction $= 5/0.02 = 250\times$ — impossibly high, so position limits are the binding constraint, not Kelly.
\end{examplebox}

\begin{examplebox}{Kelly for a Simple Bet}
$p = 0.55$, $b = 1$ (even odds, win \$1 or lose \$1). \\
$f^* = 0.55 - 0.45/1 = 0.10$. \\
Bet 10\% of capital per trade. \\
Growth rate: $g(0.10) = 0.55\ln(1.10) + 0.45\ln(0.90) = 0.55(0.0953) + 0.45(-0.1054) = 0.0524 - 0.0474 = 0.0050$ per bet. \\
Over 1000 bets: $e^{0.005 \times 1000} = e^5 \approx 148\times$ wealth growth.
\end{examplebox}

\section{Properties of the Kelly Criterion}

\begin{enumerate}[label=\arabic*.]
  \item \textbf{Maximises long-run growth}: No other strategy has a higher expected log of wealth in the long run (Kelly 1956).
  \item \textbf{Never risks ruin}: Since $f^* < 1$ when $b > 0$ and $p < 1$, the strategy never bets everything.
  \item \textbf{Convergence}: Any strategy with $f < f^*$ grows slower; any $f > f^*$ also grows slower (and risks ruin).
  \item \textbf{Large drawdowns}: Full Kelly can produce 50\% drawdowns commonly. Median time to double is fast, but the path is volatile.
\end{enumerate}

\section{Fractional Kelly}

In practice, $f^*$ from Kelly is often \emph{overestimated} because $p$ and $b$ are estimated with uncertainty. Betting full Kelly on a miestimated edge can be devastating. The standard solution is \textbf{fractional Kelly}:

\begin{equation}
  f_{\lambda} = \lambda \cdot f^*, \qquad \lambda \in (0, 1]
\end{equation}

Common choices: $\lambda = 0.25$ (quarter Kelly) or $\lambda = 0.5$ (half Kelly).

\textbf{The growth-variance trade-off}: At fraction $\lambda$:
\begin{equation}
  g(\lambda) \approx \lambda g^* - \frac{\lambda^2}{2}\sigma_g^2
\end{equation}
where $g^* = $ full Kelly growth rate and $\sigma_g^2 = (f^*)^2 \sigma^2$ is the variance of the growth rate. Half Kelly captures $\approx 75\%$ of the full Kelly growth with half the variance.

\begin{center}
\begin{tabular}{ccccc}
\toprule
$\lambda$ & Growth rate & Variance & Max drawdown & Recommended when \\
\midrule
1.0 (Full) & $g^*$ & $\sigma_g^2$ & $\sim 50\%$ frequently & Perfect parameter knowledge \\
0.5 (Half) & $0.75g^*$ & $0.25\sigma_g^2$ & $\sim 25\%$ & Moderate parameter uncertainty \\
0.25 (Quarter) & $\sim 0.44g^*$ & $0.06\sigma_g^2$ & $\sim 12\%$ & High parameter uncertainty \\
\bottomrule
\end{tabular}
\end{center}

\begin{warningbox}{Kelly Assumes i.i.d.\ Returns}
The Kelly derivation assumes bets are independent and identically distributed. In real markets: (1) returns are not i.i.d.\ (volatility clusters, autocorrelation exists), (2) the edge $p$ and $b$ are not known — only estimated. These violations mean full Kelly always overstates the optimal bet. Always use fractional Kelly in practice.
\end{warningbox}

\begin{takeawaybox}{Chapter 21}
\begin{itemize}
  \item Kelly fraction: $f^* = (pb-q)/b = $ edge/odds. Maximises log-wealth growth rate.
  \item Continuous-time: $f^* = (\mu - r_f)/\sigma^2 = \text{SR}/\sigma$.
  \item Full Kelly: maximum growth but wild path (50\% drawdowns common).
  \item Fractional Kelly ($\lambda f^*$): half Kelly captures 75\% of growth with 25\% of variance.
  \item In Prosperity: position limits act as the Kelly cap (the optimal size is usually larger than the limit allows).
\end{itemize}
\end{takeawaybox}

\chapter{Portfolio Construction}

\section{The Diversification Principle}

\begin{definitionbox}{Diversification}
\textbf{Diversification} reduces portfolio risk without reducing expected return, by combining assets whose returns are not perfectly correlated. The portfolio variance is:
\[
  \sigma_P^2 = \mathbf{w}^\top \boldsymbol{\Sigma} \mathbf{w} = \sum_i w_i^2 \sigma_i^2 + \sum_{i\neq j} w_i w_j \sigma_{ij}
\]
where $\boldsymbol{\Sigma}$ is the $N\times N$ covariance matrix and $\mathbf{w}$ is the weight vector.
\end{definitionbox}

For $N$ assets with equal weight $w_i = 1/N$, equal variance $\sigma^2$, and equal covariance $\rho\sigma^2$:
\begin{equation}
  \sigma_P^2 = \frac{\sigma^2}{N} + \frac{N-1}{N}\rho\sigma^2 = \frac{\sigma^2}{N}(1-\rho) + \rho\sigma^2
\end{equation}
As $N \to \infty$: $\sigma_P^2 \to \rho\sigma^2$. The \textbf{idiosyncratic risk} (the $1/N$ term) vanishes; only \textbf{systematic risk} ($\rho\sigma^2$) remains. Diversification eliminates idiosyncratic risk but not market risk.

\section{Modern Portfolio Theory (Markowitz, 1952)}

Harry Markowitz's mean-variance optimisation framework revolutionised portfolio construction. It finds the portfolio with the highest expected return for a given level of risk.

\begin{definitionbox}{Efficient Frontier}
The \textbf{efficient frontier} is the set of portfolios that achieve the \emph{minimum variance} for each level of expected return:
\[
  \min_{\mathbf{w}} \; \mathbf{w}^\top\boldsymbol{\Sigma}\mathbf{w} \quad \text{subject to} \quad \mathbf{w}^\top\boldsymbol{\mu} = \mu^*, \quad \mathbf{w}^\top\mathbf{1} = 1
\]
Tracing over all target returns $\mu^*$ traces the efficient frontier in $(\sigma_P, \mu_P)$ space.
\end{definitionbox}

\subsection{The Lagrangian Solution}

The efficient frontier is found via the Lagrangian:
\begin{equation}
  \mathcal{L} = \mathbf{w}^\top\boldsymbol{\Sigma}\mathbf{w} - \lambda_1(\mathbf{w}^\top\boldsymbol{\mu} - \mu^*) - \lambda_2(\mathbf{w}^\top\mathbf{1} - 1)
\end{equation}
First-order conditions: $2\boldsymbol{\Sigma}\mathbf{w} = \lambda_1\boldsymbol{\mu} + \lambda_2\mathbf{1}$, which gives:
\begin{equation}
  \mathbf{w}^* = \frac{1}{2}\boldsymbol{\Sigma}^{-1}(\lambda_1\boldsymbol{\mu} + \lambda_2\mathbf{1})
\end{equation}
The constants $\lambda_1, \lambda_2$ are determined by the two constraints. The efficient frontier is a hyperbola in $(\sigma_P, \mu_P)$ space.

\subsection{The Tangency Portfolio (Maximum Sharpe Ratio)}

Adding a risk-free asset $r_f$, the optimal risky portfolio is the one that maximises the Sharpe ratio:
\begin{equation}
  \mathbf{w}^{\text{tangency}} \propto \boldsymbol{\Sigma}^{-1}(\boldsymbol{\mu} - r_f\mathbf{1})
\end{equation}
Every rational investor holds a combination of the tangency portfolio and the risk-free asset (the \textbf{Capital Market Line}). The tangency portfolio is the same for everyone — only the mix with the risk-free asset differs based on risk tolerance.

\begin{center}
\begin{tikzpicture}[scale=1.1]
  \draw[->, thick, gray] (0,0) -- (5.5,0) node[right] {$\sigma_P$};
  \draw[->, thick, gray] (0,0) -- (0,4.5) node[above] {$\mu_P$};
  % Efficient frontier (hyperbola approximation)
  \draw[very thick, DefBlue, domain=1:4, samples=60]
    plot ({sqrt(\x*\x/2 + 0.3)}, {\x}) node[right] {\small Efficient};
  \node[DefBlue, right] at (2.2, 3.8) {\small Frontier};
  % Minimum variance portfolio
  \filldraw[DefBlue] ({sqrt(0.3)}, 0) circle (2.5pt) node[below right] {\small MVP};
  % Tangency portfolio
  \filldraw[ExGreen!80!black] ({sqrt(1.25)}, {sqrt(2.5-0.6)}) circle (2.5pt)
    node[above right] {\small Tangency};
  % Risk-free rate
  \filldraw[WarnRed] (0,0.5) circle (2.5pt) node[left] {\small $r_f$};
  % Capital Market Line
  \draw[thick, WarnRed, dashed] (0,0.5) -- (4.5, 4.2) node[right] {\small CML};
  % Dominated region
  \draw[thick, gray!50, domain=-1:0, samples=30]
    plot ({sqrt(\x*\x/2 + 0.3)}, {\x});
\end{tikzpicture}
\end{center}

\subsection{Limitations of Mean-Variance Optimisation}

\begin{warningbox}{MVO Is Extremely Sensitive to Inputs}
\begin{enumerate}[label=\arabic*.]
  \item Small estimation errors in $\boldsymbol{\mu}$ or $\boldsymbol{\Sigma}$ produce wildly different portfolios (Michaud 1989).
  \item The optimiser is an ``error maximiser'': it overweights assets with overestimated returns.
  \item The $\boldsymbol{\Sigma}^{-1}$ operation amplifies estimation errors in the covariance matrix.
  \item Correlations are unstable: in crises, assets become highly correlated and the diversification benefit disappears.
  \item Practitioners typically use shrinkage estimators (Ledoit-Wolf), Black-Litterman (view-adjusted MVO), or replace MVO with more robust alternatives like risk parity.
\end{enumerate}
\end{warningbox}

\section{Risk Parity}

Risk parity (Bridgewater Associates, popularised by Qian 2005) allocates capital so that each asset contributes \emph{equally} to portfolio risk, rather than equally by capital weight.

\begin{definitionbox}{Risk Contribution}
The \textbf{marginal risk contribution} of asset $i$ to portfolio volatility $\sigma_P$ is:
\[
  \text{MRC}_i = \frac{\partial \sigma_P}{\partial w_i} = \frac{(\boldsymbol{\Sigma}\mathbf{w})_i}{\sigma_P}
\]
The \textbf{total risk contribution} (TRC) of asset $i$:
\[
  \text{TRC}_i = w_i \cdot \text{MRC}_i = \frac{w_i (\boldsymbol{\Sigma}\mathbf{w})_i}{\sigma_P}
\]
Note: $\sum_i \text{TRC}_i = \sigma_P$ (Euler's homogeneity theorem).
\end{definitionbox}

\begin{definitionbox}{Risk Parity Portfolio}
The \textbf{risk parity} portfolio equalises the total risk contribution of all assets:
\[
  \text{TRC}_i = \text{TRC}_j \quad \forall i, j \qquad \Longleftrightarrow \qquad w_i \frac{\partial \sigma_P}{\partial w_i} = \frac{\sigma_P}{N} \; \forall i
\]
\end{definitionbox}

For the special case where all correlations are zero:
\[
  w_i^{\text{RP}} \propto \frac{1}{\sigma_i}
\]
Each asset's weight is inversely proportional to its volatility. This is exactly the \textbf{volatility-adjusted sizing} from Chapter~20.

\textbf{Risk parity vs.\ MVO:}
\begin{center}
\begin{tabular}{lll}
\toprule
\textbf{Property} & \textbf{MVO} & \textbf{Risk Parity} \\
\midrule
Input required & $\boldsymbol{\mu}$ and $\boldsymbol{\Sigma}$ & $\boldsymbol{\Sigma}$ only \\
Sensitivity to inputs & Very high & Moderate \\
Out-of-sample performance & Often poor & Often good \\
Concentration & Can be extreme & Diversified by construction \\
Leverage & Minimal & Often requires leverage \\
\bottomrule
\end{tabular}
\end{center}

\section{Correlation and Diversification in Practice}

\begin{definitionbox}{Pearson Correlation}
\[
  \rho_{ij} = \frac{\Cov(r_i, r_j)}{\sigma_i \sigma_j} = \frac{\E[(r_i - \mu_i)(r_j - \mu_j)]}{\sigma_i \sigma_j} \in [-1, +1]
\]
\end{definitionbox}

Key facts about correlation in portfolios:
\begin{itemize}
  \item $\rho = -1$: Perfect negative correlation. The two assets combine to form a risk-free portfolio.
  \item $\rho = 0$: No linear relationship. Diversification benefit is maximal.
  \item $\rho = +1$: Perfect positive correlation. No diversification benefit.
  \item In Prosperity, EMERALDS and TOMATOES have low correlation (different bot regimes) → combining them in a portfolio reduces overall variance.
\end{itemize}

\subsection{Correlation Is Unstable}

The most dangerous property of correlations: they increase during market stress. The correlation that was 0.3 in calm markets may jump to 0.8 during a crisis, exactly when you need diversification most. This ``correlation breakdown'' is a well-documented stylised fact (Ang \& Chen 2002).

\section{Maximum Drawdown Control}

\begin{definitionbox}{Drawdown}
The \textbf{drawdown} at time $t$ is the loss from the previous peak of cumulative wealth:
\[
  DD_t = \frac{\max_{s \leq t} W_s - W_t}{\max_{s \leq t} W_s}
\]
The \textbf{maximum drawdown} (MDD) over a period $[0, T]$ is:
\[
  \text{MDD} = \max_{t \in [0,T]} DD_t = \max_{t \in [0,T]} \frac{\max_{s \leq t}W_s - W_t}{\max_{s \leq t}W_s}
\]
\end{definitionbox}

\textbf{Why MDD matters}: A 50\% drawdown requires a 100\% return to recover. A 20\% drawdown requires only 25\%. Drawdown control is critical for strategy survivability.

\textbf{Practical drawdown controls:}
\begin{enumerate}[label=\arabic*.]
  \item \textbf{Hard stop}: Stop trading after $D\%$ drawdown from peak. Reset at next period.
  \item \textbf{Soft stop}: Reduce position size proportionally as drawdown increases.
  \item \textbf{Stop loss per trade}: Never lose more than $f\%$ per position (ATR-based stops).
  \item \textbf{Volatility targeting}: Scale positions so daily portfolio volatility $\approx$ target. When vol rises, reduce size.
\end{enumerate}

\subsection{Kelly and Maximum Expected Drawdown}

For a Kelly strategy with growth rate $g$ and variance $\sigma_g^2$, the expected maximum drawdown over $n$ bets is approximately:
\begin{equation}
  \E[\text{MDD}] \approx \sqrt{\frac{\sigma_g^2 \cdot 2\ln n}{g^2}}
\end{equation}
This decreases with the square root of $g/\sigma_g$ (analogous to the Sharpe ratio of the growth process). Higher Sharpe → lower expected maximum drawdown.

\begin{takeawaybox}{Chapter 22}
\begin{itemize}
  \item Diversification eliminates idiosyncratic risk ($\to 0$ as $N \to \infty$) but not systematic risk ($\to \rho\sigma^2$).
  \item MVO: minimise $\mathbf{w}^\top\boldsymbol{\Sigma}\mathbf{w}$ subject to return target. Tangency portfolio maximises Sharpe.
  \item MVO is extremely input-sensitive. Practitioners use shrinkage, Black-Litterman, or risk parity.
  \item Risk parity: equalise TRC$_i = w_i(\boldsymbol{\Sigma}\mathbf{w})_i/\sigma_P$. With zero correlation: $w_i \propto 1/\sigma_i$.
  \item Correlations spike during crises — diversification benefits evaporate when needed most.
  \item Max drawdown: $D = (\text{peak} - W)/\text{peak}$. A 50\% DD needs 100\% recovery.
\end{itemize}
\end{takeawaybox}

\chapter{Inventory Risk and Position Limit Management}

This chapter focuses on the specific risk management problem of market makers: managing the inventory that accumulates from passive fills.

\section{Inventory Risk Formalised}

Let $q_t$ be the inventory (net position) at time $t$ and $p_t$ the mid-price. The \textbf{mark-to-market value} of the inventory is:
\[
  V_t = q_t \cdot p_t
\]
The \textbf{inventory PnL} over one period:
\begin{equation}
  \Delta V_t = q_t \cdot \Delta p_t + \Delta q_t \cdot p_t
\end{equation}
The first term ($q_t \Delta p_t$) is the \emph{mark-to-market} component: existing inventory gains/loses value as price moves. The second term ($\Delta q_t \cdot p_t$) is the \emph{trading} component: new fills.

The market maker wants $q_t \approx 0$ at all times to eliminate the first term. But quoting generates fills that push $q_t$ away from zero.

\section{The Inventory-Spread Trade-Off}

The fundamental market-making tension:
\begin{align}
  \text{Earn more spread} \;\Leftrightarrow\; &\text{Quote aggressively (tighter quotes, larger sizes)} \\
  \text{Control inventory} \;\Leftrightarrow\; &\text{Quote conservatively (wider quotes, smaller sizes)}
\end{align}

Maximising total PnL requires balancing these. The Avellaneda-Stoikov model (Chapter~30) solves this optimally with a reservation price adjustment. Here we focus on the practical heuristics used in \texttt{trader.py}.

\section{Inventory Skewing: The Heuristic Solution}

The inventory skew from \texttt{trader.py} (repeated here for analysis):
\begin{equation}
  \text{skew} = \frac{\text{pos}}{L} \times S_{\text{range}}
\end{equation}
\begin{align}
  p^B &= \floor{FV - \delta - \text{skew}}, \qquad p^A = \ceil{FV + \delta - \text{skew}}
\end{align}

\textbf{When pos $= 0$} (flat): skew $= 0$, symmetric quotes at $FV \pm \delta$.

\textbf{When pos $= +L$} (at limit, maximum long): skew $= S_{\text{range}}$, bid shifts down by $S_{\text{range}}$, ask shifts down by $S_{\text{range}}$. The ask is now $S_{\text{range}}$ ticks cheaper — incentivising bots to buy from us and reduce our long.

\textbf{When pos $= -L$} (at limit, maximum short): skew $= -S_{\text{range}}$, both quotes shift up — incentivising bots to sell to us and reduce our short.

\subsection{Optimal Skew Range}

The skew range $S_{\text{range}}$ in \texttt{trader.py} is $S_{\text{range}} = 1$ (grid-searched optimal). This adds at most 1 tick of skew even at the position limit. Why not more?
\begin{itemize}
  \item Larger skew sacrifices edge: shifting the ask down by $k$ ticks reduces profit per fill by $k$.
  \item With TOMATOES passive edge $\delta = 5$ ticks and position limit $L = 80$: each 1-unit position adds $1/80 \times 1 = 0.0125$ ticks of skew — extremely mild.
  \item At $S_{\text{range}} = 1$, maximum inventory penalty is 1 tick (at position limit). This costs $1/7\times 14\% \approx 2\%$ of edge per fill — an acceptable cost for inventory control.
\end{itemize}

\section{Position Limit as Hard Risk Cap}

In Prosperity, the position limit $L$ is enforced by the exchange. But it also functions as the hardest risk management rule: no matter what, the position stays within $[-L, +L]$.

This is analogous to:
\begin{itemize}
  \item A stop-loss that fires before a trade is entered (preventive, not reactive)
  \item A leverage constraint: the market maker cannot be more than $L$ units long or short
  \item A drawdown cap: if the price moves $\Delta p$ while at limit, maximum PnL impact is $L \cdot |\Delta p|$
\end{itemize}

The remaining capacity formulae:
\begin{align}
  \text{buy\_cap} &= L - \text{pos} \quad (\geq 0 \text{ always; } = 0 \text{ when at long limit}) \\
  \text{sell\_cap} &= L + \text{pos} \quad (\geq 0 \text{ always; } = 0 \text{ when at short limit})
\end{align}

\section{Hedging Inventory with Taking Orders}

Beyond quote skewing, a market maker can \emph{actively} unwind inventory by crossing the spread (aggressive takes). The trade-off:
\begin{equation}
  \text{Cost of taking to unwind} = \text{effective spread paid} = p^A_{\text{bot}} - FV \approx \delta_{\text{bot}}
\end{equation}
This costs $\delta_{\text{bot}}$ ticks (e.g.\ 7 ticks for EMERALDS). This is only worthwhile when the expected inventory loss from continuing to hold the position exceeds the cost of crossing:
\[
  L \cdot \E[|\Delta p_t|] > \delta_{\text{bot}} \implies \text{worth crossing to unwind}
\]
For EMERALDS (extremely stable, $\sigma \approx 0.72$): $80 \times 0.72 \approx 58$ ticks expected inventory loss vs.\ 7-tick crossing cost → crossing \emph{can} be worthwhile at the limit, but only if there is a strong trend signal. Without a directional view, the strategy simply relies on the skew to organically unwind inventory.

\chapter{Multi-Product Position Management}
\label{ch:multiproduct}

When Prosperity introduces multiple products simultaneously (Rounds 1--5 typically have 5--10 products), the trader must manage all positions jointly. The key new constraints and interactions that single-product analysis misses:

\begin{itemize}
  \item Position limits are per-product but bind simultaneously — at worst all products hit their limits at the same time.
  \item Mid-price returns across products are correlated — inventory risk aggregates non-linearly across the portfolio.
  \item Capital (XIREC) is not a constraint in Prosperity (no cost of carry), but time is: each iteration has a fixed order quota per product.
\end{itemize}

\section{The Joint Position Problem}

\begin{definitionbox}{Multi-Product Order Generation Problem}
At each timestamp, the trader chooses order quantities $\delta q_1, \ldots, \delta q_n$ (one per product) to submit. Let $q_i$ be the current position in product $i$ and $L_i$ the position limit. The feasibility constraints are:
\[
  |q_i + \delta q_i| \leq L_i \quad \forall\, i = 1, \ldots, n
\]
These constraints are \textbf{separable}: the feasibility of $\delta q_i$ depends only on $q_i$ and $L_i$, not on other products. The objective, however, is not separable when returns are correlated.
\end{definitionbox}

\subsection{Separable Objective: Independent Product Management}

If products have uncorrelated price processes (e.g.\ EMERALDS and TOMATOES in Round 0), the PnL decomposes additively:
\[
  \text{PnL}_{\text{total}} = \sum_{i=1}^n \text{PnL}_i(q_i)
\]
In this case, each product can be managed independently with its own fair-value model and inventory skew. This is the structure of the tutorial \texttt{trader.py}: \texttt{EmeraldsStrategy} and \texttt{TomatoesStrategy} are fully independent objects.

\textbf{When is independence valid?} When $|\Corr(r_i, r_j)| \lesssim 0.2$ for all pairs $(i, j)$. Estimate this from the first day's data at the start of each round.

\subsection{Correlated Products: Portfolio Inventory Risk}

When products are correlated (e.g.\ PICNIC\_BASKET1 and its component CROISSANTS), the portfolio inventory risk is:
\[
  \text{Inventory P\&L risk} = \mathbf{q}^{\top} \boldsymbol{\Sigma} \mathbf{q}
\]
where $\mathbf{q} = (q_1, \ldots, q_n)^{\top}$ is the position vector and $\boldsymbol{\Sigma}$ is the return covariance matrix. Taking a position in two positively correlated products amplifies total risk faster than their individual limits suggest.

\begin{keyconceptbox}{Portfolio Skew Rule for Correlated Products}
The inventory skew (Chapter~15) generalises to the portfolio case. Define the \textbf{portfolio delta} at product $i$ as the total sensitivity of portfolio value to a 1-tick move in product $i$:
\[
  \Delta_i^{\text{portfolio}} = q_i + \sum_{j \neq i} q_j \cdot \frac{\Cov(r_i, r_j)}{\Var(r_i)}
\]
The skew for product $i$ should depend on $\Delta_i^{\text{portfolio}}$, not just $q_i$, to account for the hedge provided by correlated positions. \textbf{Intuition:} a long position in $j$ is worth $q_j \frac{\Cov(r_i,r_j)}{\Var(r_i)}$ units of product $i$ in terms of risk — it offsets that many units of your $q_i$ exposure. If $\Cov(r_i,r_j)/\Var(r_i) = 0.8$ and $q_j = 10$, your effective long in $i$ is only $q_i - 8$, not $q_i$. In practice: if product $i$ and $j$ are highly correlated ($\rho \geq 0.7$), a long position in $j$ partially hedges a long position in $i$, so you can afford a larger $q_i$.
\end{keyconceptbox}

\section{Capital Allocation Across Products}

\textbf{Capital is not a binding constraint in Prosperity} (no cost of holding positions). However, \emph{inventory risk} is. The right resource to allocate is not money but \emph{position headroom} $L_i - |q_i|$.

\subsection{Equal Headroom Allocation}

The simplest policy: always try to maintain equal headroom across all products. Each product targets $q_i^* = 0$ as its long-run inventory:
\[
  \text{buy\_cap}_i = L_i - q_i, \qquad \text{sell\_cap}_i = L_i + q_i
\]
Post passive orders of full remaining capacity on each side. This maximises fill opportunities independently for each product.

\subsection{Risk-Parity Headroom Allocation}

When products have different volatilities $\sigma_i$, equal headroom is not equal risk. A risk-parity approach allocates headroom inversely proportional to volatility:
\[
  L_i^{\text{effective}} = L_i \cdot \frac{\bar{\sigma}}{\sigma_i}, \qquad \bar{\sigma} = \frac{1}{n}\sum_i \sigma_i
\]
Products with higher volatility get a tighter effective limit, ensuring each product contributes equally to portfolio volatility at the position limit.

\section{Practical Implementation Template}

\begin{lstlisting}[caption={Multi-product trader skeleton for Rounds 1--5}]
from datamodel import Order, TradingState
from typing import Dict, List
import jsonpickle, math

class MultiProductTrader:
    """
    Template for managing n products with per-product strategies
    and shared portfolio-level risk monitoring.
    """

    POSITION_LIMITS: Dict[str, int] = {}   # set per round
    VOLATILITY:      Dict[str, float] = {} # estimated from day-0 data

    def __init__(self):
        self._strategies: Dict[str, object] = {}  # per-product strategy

    def _portfolio_risk(self, positions: Dict[str, int]) -> float:
        """Approximate portfolio variance (assumes independence)."""
        return sum(
            (positions.get(p, 0) * self.VOLATILITY.get(p, 1.0)) ** 2
            for p in self.POSITION_LIMITS
        ) ** 0.5

    def _effective_limit(self, product: str) -> int:
        """Risk-parity adjusted position limit."""
        sigma = self.VOLATILITY.get(product, 1.0)
        sigma_bar = (sum(self.VOLATILITY.values()) / len(self.VOLATILITY)
                     if self.VOLATILITY else 1.0)
        scale = sigma_bar / sigma if sigma > 0 else 1.0
        return max(1, int(self.POSITION_LIMITS.get(product, 20) * scale))

    def run(self, state: TradingState):
        result: Dict[str, List[Order]] = {}

        for product, strategy in self._strategies.items():
            if product not in state.order_depths:
                continue
            # Pass risk-adjusted limit to each strategy
            effective_lim = self._effective_limit(product)
            strategy.position_limit = effective_lim
            result[product] = strategy.orders(state)

        return result, 0, ""
\end{lstlisting}

\section{Cross-Product Arbitrage as a Special Case}

When two products satisfy a no-arbitrage relationship (e.g.\ basket = weighted sum of components), multi-product management becomes arbitrage trading. The positions in the two legs are linked by design: a long basket + short components creates a net-zero inventory in aggregate. The position limit constraint then becomes:
\begin{align}
  |q_{\text{basket}}| &\leq L_{\text{basket}} \\
  |q_{\text{component}_i}| &\leq L_i \\
  q_{\text{basket}} + w_i\, q_{\text{component}_i} &\approx 0 \quad \text{(target: fully hedged)}
\end{align}
The effective constraint is the binding leg: $\min(L_{\text{basket}},\, L_i / w_i)$ determines the maximum arbitrage size.

\begin{takeawaybox}{Chapter on Multi-Product Management}
\begin{itemize}
  \item Products with uncorrelated returns can be managed independently (additive PnL decomposition).
  \item Correlated products require portfolio-level inventory tracking: use $\Delta_i^{\text{portfolio}}$ for skewing.
  \item In Prosperity, capital is not the constraint — position headroom is. Allocate it proportional to $1/\sigma_i$ for equal volatility contribution (risk parity).
  \item Cross-product arbitrage (basket = components) imposes a natural hedge that links the position limits of both legs.
  \item Always estimate the return correlation matrix from the first day's data before applying any joint management rule.
\end{itemize}
\end{takeawaybox}

% ──────────────────────────────────────────────────────────────────────────────

\chapter{Hedging}

\section{Delta Hedging}

In options markets, \textbf{delta} ($\Delta$) is the sensitivity of an option's price to the underlying:
\[
  \Delta = \frac{\partial V_{\text{option}}}{\partial S}
\]
A delta-hedged portfolio holds the option and $-\Delta$ units of the underlying, making it instantaneously insensitive to small moves in $S$. This is the cornerstone of options market making.

\section{Beta Hedging}

For equity portfolios, \textbf{beta} measures sensitivity to the market:
\[
  \beta_P = \sum_i w_i \beta_i, \qquad \beta_i = \frac{\Cov(r_i, r_M)}{\Var(r_M)}
\]
To hedge market risk:
\[
  \text{Short futures} = \frac{\beta_P \times V_P}{V_{\text{futures}}}
\]
A market-neutral portfolio has $\beta_P = 0$.

\section{Cross-Asset Hedging}

When a perfect hedge is unavailable, use a correlated instrument:
\[
  h^* = \frac{\Cov(r_P, r_H)}{\Var(r_H)} = \rho_{PH} \cdot \frac{\sigma_P}{\sigma_H}
\]
where $h^*$ is the optimal hedge ratio (the fraction of the portfolio value to hedge). The \emph{hedge effectiveness} is $1 - (1 - \rho_{PH}^2)$: perfect when $|\rho| = 1$, zero when $\rho = 0$.

\section{Hedging in the Prosperity Context}

In Prosperity, direct hedging is not applicable within a single round (no derivatives). However, across products:
\begin{itemize}
  \item If EMERALDS and TOMATOES were positively correlated, a market maker long EMERALDS and short TOMATOES would be partially hedged.
  \item The correlation between the two products is empirically near zero (different bot regimes) → no hedging benefit from combining them.
  \item True hedging opportunities arise in later rounds with related products (baskets, derivatives).
\end{itemize}

\begin{takeawaybox}{Chapters 23--24}
\begin{itemize}
  \item Inventory PnL: $\Delta V = q_t \Delta p_t + \Delta q_t \cdot p_t$. Market makers want $q_t \approx 0$.
  \item Inventory skew: shifts both quotes by $(\text{pos}/L) \times S_{\text{range}}$ in the direction that reduces inventory.
  \item Position limit acts as the hard risk cap: maximum inventory impact is $L \cdot |\Delta p|$.
  \item Unwind via taking: worth crossing when expected inventory loss $> $ crossing cost.
  \item Delta hedge (options), beta hedge (equity), cross-asset hedge. Optimal hedge ratio: $h^* = \rho_{PH}\sigma_P/\sigma_H$.
\end{itemize}
\end{takeawaybox}

% End of Part V

% ============================================================
%
%  PART VI — PERFORMANCE MEASUREMENT & BACKTESTING
%
% ============================================================
\part{Performance Measurement \& Backtesting}

\chapter{Performance Metrics}

A strategy's performance cannot be summarised by a single number. Different metrics capture different dimensions of risk and return. This chapter derives every major metric from first principles, explains what it does and does not measure, and connects it to the Prosperity context.

\section{Returns and Compounding}

All performance metrics are built on returns. Let $W_0$ be initial wealth and $W_T$ final wealth.

\begin{definitionbox}{Total Return and CAGR}
\textbf{Total return}: $R = \frac{W_T - W_0}{W_0} = \frac{W_T}{W_0} - 1$.

\textbf{Compound Annual Growth Rate (CAGR)}: the annualised geometric mean return:
\[
  \text{CAGR} = \left(\frac{W_T}{W_0}\right)^{1/T} - 1
\]
where $T$ is the number of years. CAGR is the constant annual return that would produce the same terminal wealth.
\end{definitionbox}

\begin{examplebox}{CAGR vs.\ Arithmetic Mean}
3-year returns: $+50\%$, $-33\%$, $+50\%$. \\
Arithmetic mean: $(50 - 33 + 50)/3 = 22.3\%$ per year. \\
Actual: $W_3 = W_0 \times 1.50 \times 0.67 \times 1.50 = W_0 \times 1.508$. \\
CAGR $= 1.508^{1/3} - 1 = 1.147 - 1 = 14.7\%$ per year. \\
The arithmetic mean overstates actual growth. CAGR is the correct measure of compounded performance.
\end{examplebox}

\section{Sharpe Ratio}

The Sharpe ratio (William Sharpe, 1966) is the most widely used single performance metric. It measures return per unit of \emph{total} risk.

\begin{definitionbox}{Sharpe Ratio}
\[
  \text{SR} = \frac{\E[r_P - r_f]}{\sigma_P} = \frac{\mu_P - r_f}{\sigma_P}
\]
where $r_P$ is the portfolio return, $r_f$ the risk-free rate, and $\sigma_P$ the standard deviation of excess returns. The numerator is the \textbf{excess return} (reward); the denominator is the \textbf{total volatility} (risk).
\end{definitionbox}

\subsection{Annualisation}

If returns are measured at frequency $\Delta t$ (e.g.\ daily = $\Delta t = 1/252$), annualise by:
\begin{equation}
  \text{SR}_{\text{ann}} = \text{SR}_{\Delta t} \times \sqrt{\frac{1}{\Delta t}} = \text{SR}_{\text{daily}} \times \sqrt{252}
\end{equation}
The $\sqrt{252}$ comes from: mean scales as $252\times\bar{r}$ and std scales as $\sqrt{252}\times\sigma$, so ratio scales as $\sqrt{252}$.

\begin{examplebox}{Sharpe Ratio Calculation}
Daily excess returns: mean $= 0.05\%$, std $= 1\%$. \\
Daily SR $= 0.05/1.0 = 0.05$. \\
Annual SR $= 0.05 \times \sqrt{252} = 0.05 \times 15.87 = 0.794$. \\
In Prosperity (1000 timestamps per day, $r_f = 0$): if mean return per timestamp $= 0.001\%$ and std $= 0.01\%$, SR per timestamp $= 0.1$, annualised $= 0.1 \times \sqrt{252000} \approx 50.2$ (market making strategies have very high Sharpe due to extremely low per-timestamp variance).
\end{examplebox}

\subsection{Limitations of the Sharpe Ratio}

\begin{warningbox}{Sharpe Ratio Pitfalls}
\begin{enumerate}[label=\arabic*.]
  \item \textbf{Assumes normality}: Uses $\sigma$ which is the natural risk measure for normal distributions. Strategies with fat tails (crash risk) or skewed returns (selling options) can have misleadingly high Sharpe ratios.
  \item \textbf{Penalises upside volatility}: If a strategy has a large upside spike, $\sigma$ increases and SR falls — even though more upside is good.
  \item \textbf{Frequency gaming}: A strategy sampled at lower frequency appears to have a lower SR (less data, wider confidence intervals), even if it is the same underlying strategy.
  \item \textbf{Correlation with other strategies is ignored}: Two strategies with SR=1 can form a portfolio with SR$\gg 1$ if they are uncorrelated. The Sharpe ratio of individual strategies does not capture this.
\end{enumerate}
\end{warningbox}

\section{Sortino Ratio}

The Sortino ratio (Frank Sortino, 1994) addresses Sharpe's upside-volatility problem by using only \emph{downside} deviation.

\begin{definitionbox}{Sortino Ratio}
\[
  \text{Sortino} = \frac{\mu_P - r_f}{\sigma_{\text{down}}}
\]
where the \textbf{downside deviation} is:
\[
  \sigma_{\text{down}} = \sqrt{\frac{1}{n}\sum_{t=1}^n \min(r_t - r_{\text{target}}, 0)^2}
\]
and $r_{\text{target}}$ is a minimum acceptable return (often $r_f$ or 0). Only negative deviations from the target are counted.
\end{definitionbox}

\textbf{Interpretation:} For strategies with positive skew (rare large wins), SR $<$ Sortino (the wins increase $\sigma$ but not $\sigma_{\text{down}}$). For strategies with negative skew (rare large losses, like selling options), SR $>$ Sortino — Sortino correctly penalises the crash risk.

\begin{center}
\begin{tabular}{lll}
\toprule
\textbf{Strategy type} & \textbf{Sharpe} & \textbf{Sortino} \\
\midrule
Symmetric returns & Equal & Equal \\
Right-skewed (trend following) & Lower & Higher \\
Left-skewed (carry, selling vol) & Higher (misleading) & Lower (more accurate) \\
\bottomrule
\end{tabular}
\end{center}

\section{Calmar Ratio}

The Calmar ratio (Terry Young, 1991) relates return to the worst drawdown experienced:

\begin{definitionbox}{Calmar Ratio}
\[
  \text{Calmar} = \frac{\text{CAGR}}{|\text{MDD}|}
\]
where MDD is the maximum drawdown (as a positive fraction). Higher Calmar = better return for the worst drawdown endured.
\end{definitionbox}

\begin{examplebox}{Calmar vs.\ Sharpe}
Strategy A: CAGR $= 20\%$, MDD $= 10\%$, $\sigma = 15\%$. SR $= 1.33$, Calmar $= 2.0$. \\
Strategy B: CAGR $= 30\%$, MDD $= 40\%$, $\sigma = 20\%$. SR $= 1.50$, Calmar $= 0.75$. \\
By Sharpe, B is better. By Calmar, A is dramatically better — it never lost more than 10\% at its worst point, which is crucial for staying solvent and avoiding forced liquidation.
\end{examplebox}

\section{Omega Ratio}

The Omega ratio (Keating \& Shadwick, 2002) uses the entire return distribution without parametric assumptions:

\begin{definitionbox}{Omega Ratio}
\[
  \Omega(L) = \frac{\int_L^{\infty} [1 - F(r)]\,\dd r}{\int_{-\infty}^L F(r)\,\dd r}
\]
where $F(r)$ is the CDF of returns and $L$ is a threshold (often 0 or $r_f$). The numerator is the probability-weighted gains above $L$; the denominator is the probability-weighted losses below $L$.
\end{definitionbox}

\textbf{Interpretation}: $\Omega > 1$ means expected gains outweigh expected losses relative to the threshold $L$ — necessary (but not sufficient) condition for a profitable strategy. $\Omega = \Omega(0) = \nicefrac{\E[\max(r,0)]}{\E[\max(-r,0)]}$ is the gain-to-loss ratio in expectation. Unlike Sharpe and Sortino, Omega captures all higher-order moments of the distribution.

\section{Value at Risk (VaR)}

VaR is the regulatory standard risk measure and one of the most widely used (and misunderstood) risk metrics.

\begin{definitionbox}{Value at Risk}
The \textbf{Value at Risk} at confidence level $\alpha$ (e.g.\ 95\%) over horizon $h$ is the loss $L$ such that the probability of exceeding $L$ is $1-\alpha$:
\[
  P(\text{loss} > \text{VaR}_\alpha) = 1 - \alpha \qquad \Longleftrightarrow \qquad \text{VaR}_\alpha = F_L^{-1}(\alpha)
\]
where $F_L$ is the CDF of the loss distribution. VaR$_{95\%}$ answers: ``What is the loss we will not exceed in 95\% of periods?''
\end{definitionbox}

\subsection{Three Methods to Compute VaR}

\subsubsection{Parametric (Normal) VaR}

Assume returns are normally distributed: $r \sim \mathcal{N}(\mu, \sigma^2)$. Then:
\begin{equation}
  \text{VaR}_\alpha = \mu - z_\alpha \sigma \cdot \sqrt{h}
\end{equation}
where $z_\alpha = \Phi^{-1}(\alpha)$ (e.g.\ $z_{95\%} = 1.645$, $z_{99\%} = 2.326$).

\begin{examplebox}{Parametric VaR}
Portfolio: $\mu = 0\%$/day, $\sigma = 1\%$/day, 1-day horizon. \\
VaR$_{95\%} = 0 - 1.645 \times 1\% = -1.645\%$ (loss of 1.645\%). \\
VaR$_{99\%} = 0 - 2.326 \times 1\% = -2.326\%$. \\
Interpretation: on 95 out of 100 days, we will not lose more than 1.645\%.
\end{examplebox}

\subsubsection{Historical VaR}

Use the empirical distribution of historical returns. Sort all returns in ascending order; VaR$_\alpha$ = the $(1-\alpha)$-quantile.

\[
  \text{VaR}_\alpha^{\text{hist}} = -\text{quantile}_{1-\alpha}(\{r_1, r_2, \ldots, r_n\})
\]

Non-parametric: no distributional assumptions. Requires large sample to estimate tail quantiles accurately.

\subsubsection{Monte Carlo VaR}

Simulate $S$ scenarios from a fitted model (e.g.\ GARCH), compute portfolio return for each, take the $(1-\alpha)$-quantile. Most flexible but computationally intensive.

\subsection{Limitations of VaR}

\begin{warningbox}{VaR Is Not Subadditive}
VaR is not a \emph{coherent} risk measure: combining two portfolios can produce a VaR \emph{larger} than the sum of individual VaRs, violating the diversification principle. This is mathematically proven and is a fundamental flaw of VaR as a risk measure.
\end{warningbox}

\section{Conditional VaR (CVaR / Expected Shortfall)}

CVaR (also called Expected Shortfall or ES) fixes VaR's tail blindness by measuring the \emph{expected loss given that the loss exceeds VaR}:

\begin{definitionbox}{Conditional Value at Risk}
\[
  \text{CVaR}_\alpha = \E[\text{Loss} \mid \text{Loss} > \text{VaR}_\alpha] = \frac{1}{1-\alpha}\int_\alpha^1 \text{VaR}_u \,\dd u
\]
For normal distributions:
\[
  \text{CVaR}_\alpha = \mu + \sigma \cdot \frac{\phi(z_\alpha)}{1-\alpha} \cdot \sqrt{h}
\]
where $\phi$ is the standard normal PDF.
\end{definitionbox}

CVaR is always $\geq$ VaR (it is the average of the worst $(1-\alpha)\times100\%$ outcomes). It is:
\begin{enumerate}[label=\arabic*.]
  \item \textbf{Subadditive}: CVaR(A+B) $\leq$ CVaR(A) + CVaR(B) — diversification always helps.
  \item \textbf{Sensitive to tail shape}: Different from VaR in fat-tailed distributions.
  \item \textbf{Regulatory standard}: Basel III uses Expected Shortfall (ES) at 97.5\% replacing VaR.
\end{enumerate}

\begin{examplebox}{VaR vs.\ CVaR}
Distribution: 95\% of days: return $= +0.5\%$; 5\% of days: return $= -10\%$ (crash). \\
VaR$_{95\%} = $ the 5th-percentile loss $= 10\%$. \\
CVaR$_{95\%} = \E[\text{loss} | \text{loss} > 10\%] = 10\%$ (since all losses in the 5\% tail are exactly 10\%). \\
Now consider a heavier tail: 4\% of days $= -10\%$, 1\% of days $= -40\%$. \\
VaR$_{95\%} = 10\%$ (unchanged — VaR is blind to the $-40\%$ tail). \\
CVaR$_{95\%} = (4 \times 10 + 1 \times 40)/5 = 16\%$ — correctly captures the extreme tail.
\end{examplebox}

\section{Alpha and Beta (CAPM Framework)}

\begin{definitionbox}{CAPM Beta and Alpha}
From the Capital Asset Pricing Model (Sharpe 1964, Lintner 1965):
\[
  r_P - r_f = \alpha + \beta(r_M - r_f) + \varepsilon
\]
\begin{itemize}
  \item $\beta = \nicefrac{\Cov(r_P, r_M)}{\Var(r_M)}$: systematic market exposure. $\beta = 1$: moves with the market; $\beta = 0$: market-neutral.
  \item $\alpha$: risk-adjusted excess return not explained by market exposure. Genuine skill vs.\ luck.
\end{itemize}
\end{definitionbox}

\textbf{Jensen's Alpha}: the intercept $\alpha$ from the above regression. A significantly positive $\alpha$ indicates the manager adds value beyond what is compensated by market risk.

In Prosperity, the market benchmark is undefined (there is no external market). All PnL is effectively $\alpha$ — there is no beta exposure to hedge against.

\section{Information Ratio}

\begin{definitionbox}{Information Ratio}
\[
  \text{IR} = \frac{\alpha}{\sigma_\varepsilon} = \frac{\E[r_P - r_B]}{\text{std}(r_P - r_B)}
\]
where $r_B$ is the benchmark return and $\sigma_\varepsilon = \text{std}(r_P - r_B)$ is the \textbf{tracking error}. The IR measures the active return per unit of active risk.
\end{definitionbox}

The Fundamental Law of Active Management (Grinold \& Kahn):
\begin{equation}
  \text{IR} \approx \text{IC} \times \sqrt{\text{BR}}
\end{equation}
where IC = Information Coefficient (correlation between forecasts and outcomes) and BR = Breadth (number of independent bets per year). To increase IR: improve forecast quality (IC) or trade more independently (BR). This is why high-frequency market makers have very high IR — they make thousands of independent bets per day.

\section{Full Metrics Summary Table}

\begin{center}
\begin{tabular}{lllll}
\toprule
\textbf{Metric} & \textbf{Formula} & \textbf{Measures} & \textbf{Higher=} & \textbf{Limitation} \\
\midrule
Sharpe  & $(\mu-r_f)/\sigma$ & Return/total risk & Better & Penalises upside vol \\
Sortino & $(\mu-r_f)/\sigma_{\text{down}}$ & Return/downside risk & Better & Ignores upside \\
Calmar  & CAGR/MDD & Return/worst loss & Better & Depends on period \\
Omega   & $\E[r^+]/\E[r^-]$ & All-moment ratio & Better & No single threshold \\
VaR     & $F_L^{-1}(\alpha)$ & Quantile loss & Lower risk & Not subadditive, tail-blind \\
CVaR    & $\E[L|L>\text{VaR}]$ & Expected tail loss & Lower risk & Harder to estimate \\
Alpha   & CAPM intercept & Skill vs.\ beta & Higher & Model-dependent \\
Beta    & $\Cov(r_P,r_M)/\Var(r_M)$ & Market exposure & Context & CAPM assumptions \\
IR      & $\alpha/\sigma_\varepsilon$ & Active return/risk & Better & Benchmark choice \\
Treynor & $(\mu-r_f)/\beta$ & Return/systematic risk & Better & Ignores idiosyncratic \\
\bottomrule
\end{tabular}
\end{center}

\begin{takeawaybox}{Chapter 25}
\begin{itemize}
  \item CAGR: geometric mean annual return. Always prefer to arithmetic mean for compounded strategies.
  \item Sharpe $= (\mu-r_f)/\sigma$; annualise by $\times\sqrt{T}$. Penalises upside volatility — use Sortino for skewed strategies.
  \item Sortino uses downside deviation $\sigma_\text{down}$; correct for crash-risk strategies.
  \item Calmar = CAGR/MDD; captures worst historical loss. Strategy B with higher Sharpe but worse Calmar may be unacceptable.
  \item VaR: loss not exceeded $\alpha\%$ of the time. Not subadditive. CVaR (Expected Shortfall) = expected loss given tail — subadditive and tail-sensitive.
  \item IR $\approx $ IC $\times\sqrt{\text{BR}}$: improve forecasts (IC) or make more independent bets (BR).
\end{itemize}
\end{takeawaybox}

% ──────────────────────────────────────────────────────────────────────────────

\chapter{Backtesting Methodology}

\section{What Is a Backtest?}

A \textbf{backtest} simulates a strategy's performance on historical data as if it had been run in the past. It is the primary tool for evaluating a strategy before live trading.

\begin{warningbox}{The Cardinal Rule of Backtesting}
A backtest answers: ``How would this strategy have performed, given these exact rules, on this exact data?'' It does NOT answer: ``How will this strategy perform in the future?'' Past performance is not indicative of future results — especially when that past performance was optimised to look good.
\end{warningbox}

\section{Backtesting in Prosperity}

The Prosperity backtester (\texttt{prosperity3bt}) replays historical CSV data through your \texttt{trader.py} exactly as the live simulation would run it:

\begin{lstlisting}[caption={Running a backtest with prosperity3bt}]
python run_jasper.py           # both tutorial days (-2 and -1)
python run_jasper.py --day -2  # single day
python run_jasper.py --vis     # open in Jasper visualizer
\end{lstlisting}

The backtester simulates:
\begin{enumerate}[label=\arabic*.]
  \item Each timestamp: call \texttt{run(state)} → get orders
  \item Match orders against the historical order book
  \item Update position and PnL
  \item Pass \texttt{traderData} to next iteration
\end{enumerate}

The key parameter is \texttt{--match-trades}: how aggressively unfilled passive orders are matched against bots:
\begin{itemize}
  \item \texttt{all}: Any order that could have been filled is filled (optimistic)
  \item \texttt{worse}: Only fill if a bot actually traded at or better than your price (realistic)
  \item \texttt{none}: Only fill orders that immediately cross (pessimistic)
\end{itemize}

\section{Data Requirements for a Valid Backtest}

\begin{enumerate}[label=\arabic*.]
  \item \textbf{Sufficient history}: Enough data to cover multiple market cycles. In Prosperity, two days of 1000 timestamps each = 2000 data points per product.
  \item \textbf{Correct frequency}: The backtest frequency must match the intended trading frequency.
  \item \textbf{Correct costs}: Transaction costs, slippage, and position limit enforcement must be realistic.
  \item \textbf{Correct order of operations}: Fair value estimation must use data available at time $t$ (no future data).
\end{enumerate}

\section{Biases That Corrupt Backtests}

\subsection{Look-Ahead Bias}

The most common and most dangerous bias: using information that was not available at the time of the trade.

\begin{examplebox}{Look-Ahead Bias}
Wrong: $\text{SMA}_t^{(20)} = \frac{1}{20}\sum_{i=-9}^{10} p_{t+i}$ (uses 10 future prices). \\
Correct: $\text{SMA}_t^{(20)} = \frac{1}{20}\sum_{i=0}^{19} p_{t-i}$ (uses only past/current prices). \\
In Prosperity: computing a signal using \texttt{state.order\_depths} is safe. Using data from \texttt{state.timestamp+100} is look-ahead bias.
\end{examplebox}

\subsection{Survivorship Bias}

\textbf{Definition}: Backtesting only on assets that still exist today. The backtest misses assets that went bankrupt or were delisted (which are typically the bad performers). This inflates backtested returns by excluding the losers.

\textbf{Prosperity context}: All products are present for the full backtest period — no survivorship bias by construction.

\subsection{Data Snooping Bias}

When you test many parameter combinations and select the best, the selected parameters are biased upward. The more combinations you test, the worse this problem becomes.

\begin{definitionbox}{Multiple Testing Correction}
If you test $K$ independent parameter sets and select the best, the expected maximum t-statistic is approximately:
\[
  \E[\max_k t_k] \approx \Phi^{-1}\!\left(1 - \frac{1}{2K}\right)
\]
For $K = 100$ tests at a nominal 5\% significance level, the true significance threshold should be the \textbf{Bonferroni-corrected} level $0.05/K = 0.0005$ (t-statistic $\approx 3.5$ instead of 1.96).
\end{definitionbox}

\textbf{Solution}: Split data into in-sample (development) and out-of-sample (validation). Never touch the out-of-sample until the final evaluation.

\subsection{Overfitting}

Overfitting occurs when a model fits the noise in the training data rather than the underlying signal. It is the central challenge of systematic trading.

\begin{keyconceptbox}{Overfitting Indicators}
\begin{enumerate}[label=\arabic*.]
  \item Large gap between in-sample and out-of-sample Sharpe ratio.
  \item Strategy requires many parameters relative to the number of observations.
  \item Strategy is exquisitely sensitive to small parameter changes.
  \item Strategy has implausibly high backtested Sharpe ratio (SR $> 3$ is suspicious for daily strategies).
\end{enumerate}
\end{keyconceptbox}

\textbf{The rule of thumb (Bailey \& L\'{o}pez de Prado 2014)}: the minimum track record length to establish statistical significance at 5\% level for a strategy with annualised SR = $s$ and skewness $\gamma_1$, excess kurtosis $\gamma_2$ is:
\begin{equation}
  T^* = \frac{1 + (1 - \gamma_1 s + \frac{\gamma_2 - 1}{4}s^2)}{(s/\sqrt{T_0})^2}
\end{equation}
where $T_0$ is a benchmark. For a daily SR = 1, it takes $\approx 228$ trading days ($\approx 1$ year) to establish significance.

\subsection{Execution Bias}

\textbf{Optimistic execution assumptions}: assuming trades execute at the bid/ask exactly, with no slippage, and perfect fills. Realistic backtests:
\begin{itemize}
  \item Add 0.5--1 tick of slippage per trade
  \item Use pessimistic fill assumptions (\texttt{--match-trades worse})
  \item Assume not all passive orders get filled (partial fills)
\end{itemize}

\section{Walk-Forward Analysis}

Walk-forward testing (also called \textbf{out-of-sample testing} or \textbf{rolling window validation}) is the gold standard for backtesting without data snooping.

\begin{definitionbox}{Walk-Forward Procedure}
\begin{enumerate}[label=\arabic*.]
  \item Divide the full dataset into $K$ consecutive blocks.
  \item For block $k$: use blocks $1, \ldots, k-1$ as in-sample (training), block $k$ as out-of-sample (test).
  \item Optimise parameters on the in-sample data (e.g.\ grid search).
  \item Evaluate performance on the out-of-sample block using the optimised parameters.
  \item The overall performance is the concatenation of out-of-sample results from all blocks $k = 2, \ldots, K$.
\end{enumerate}
\end{definitionbox}

\begin{center}
\begin{tikzpicture}[scale=0.8]
  \def\W{10} \def\H{0.6}
  % Full dataset bar
  \draw[fill=DefBlue!20] (0,2.5) rectangle (\W, 2.5+\H)
    node[midway] {\small Full dataset};
  % Walk-forward splits
  \foreach \k in {0,1,2,3} {
    \pgfmathsetmacro\ISend{(\k+3)*\W/6}
    \pgfmathsetmacro\OSstart{(\k+3)*\W/6}
    \pgfmathsetmacro\OSend{(\k+4)*\W/6}
    \pgfmathsetmacro\Ypos{2.0 - \k*0.7}
    \draw[fill=DefBlue!40] (0,\Ypos) rectangle (\ISend, \Ypos+\H)
      node[midway, font=\tiny] {In-sample};
    \draw[fill=ExGreen!50] (\OSstart,\Ypos) rectangle (\OSend, \Ypos+\H)
      node[midway, font=\tiny] {OOS};
    \draw[fill=gray!20] (\OSend,\Ypos) rectangle (\W, \Ypos+\H);
    \node[left, font=\tiny] at (0, \Ypos+\H/2) {Fold \pgfmathparse{int(\k+1)}\pgfmathresult};
  }
  % Legend
  \draw[fill=DefBlue!40] (0,-1.0) rectangle (1,-0.7) node[right, font=\tiny] {\ In-sample};
  \draw[fill=ExGreen!50] (3,-1.0) rectangle (4,-0.7) node[right, font=\tiny] {\ Out-of-sample};
  \draw[fill=gray!20] (6,-1.0) rectangle (7,-0.7) node[right, font=\tiny] {\ Unused};
\end{tikzpicture}
\end{center}

\section{Cross-Validation for Time Series}

Standard $k$-fold cross-validation randomly shuffles data — catastrophic for time series because it creates look-ahead bias. Two time-series-safe alternatives:

\subsection{Purged $k$-Fold Cross-Validation (L\'{o}pez de Prado)}

\begin{enumerate}[label=\arabic*.]
  \item Split time series into $k$ folds chronologically.
  \item When fold $i$ is test set, remove from training all observations within a \emph{embargo} window $h$ before fold $i$ (to prevent information leakage).
  \item The embargo prevents training on data that might contain future information of the test period.
\end{enumerate}

\subsection{Expanding Window}

Similar to walk-forward but the training window grows: at step $k$, training uses all data from 1 to $k-1$, test uses $k$. No window size parameter to choose.

\section{Monte Carlo Simulation for Backtests}

Monte Carlo simulation assesses the robustness of a backtest by generating many alternative return paths:

\begin{enumerate}[label=\arabic*.]
  \item \textbf{Bootstrap}: Randomly resample (with replacement) the daily PnL from the backtest. Compute SR for each resample. The distribution of SRs gives confidence intervals.
  \item \textbf{Randomise trade times}: Shift trade entry/exit times randomly. If the strategy's SR collapses, it was relying on exact timing.
  \item \textbf{Randomise parameters}: Perturb each parameter by $\pm\epsilon$ around its optimal value. If SR is very sensitive, the strategy is overfit.
\end{enumerate}

\begin{examplebox}{Bootstrap Confidence Interval for Sharpe}
Backtest: 500 daily returns with SR $= 1.2$. \\
Bootstrap 10{,}000 samples (each of size 500 with replacement). \\
Compute SR for each sample. \\
95\% CI: $[\text{5th percentile}, \text{95th percentile}]$ of the 10{,}000 SRs. \\
If CI $= [0.8, 1.6]$, the SR is statistically different from zero (both bounds positive).
\end{examplebox}

\section{Grid Search and Parameter Optimisation}

Grid search is the systematic exploration of parameter combinations. For a strategy with parameters $(\theta_1, \theta_2, \ldots, \theta_k)$:

\begin{algorithm}[H]
\caption{Grid Search with Walk-Forward Validation}
\SetAlgoLined
Define parameter grids: $\theta_j \in \{\theta_j^{(1)}, \ldots, \theta_j^{(m_j)}\}$ for each $j$\;
Total combinations: $M = \prod_j m_j$\;
Split data: In-sample (IS) and Out-of-sample (OOS)\;
\ForEach{combination $(\theta_1^{(i_1)}, \ldots, \theta_k^{(i_k)})$}{
  Run backtest on IS data with these parameters\;
  Record IS performance metric (e.g.\ Sharpe)\;
}
Select best parameter set $\hat{\boldsymbol{\theta}}$ from IS results\;
Run backtest on OOS data with $\hat{\boldsymbol{\theta}}$\;
Report OOS performance (the true unbiased estimate)\;
\end{algorithm}

\textbf{In Prosperity}: grid search was run for TOMATOES with parameters $(a, Q, S)$:
\begin{itemize}
  \item $a = \text{EMA\_ALPHA} \in \{0.1, 0.3, 0.5, 1.0\}$
  \item $Q = \text{QUOTE\_OFFSET} \in \{3, 4, 5, 6, 7\}$
  \item $S = \text{SKEW\_RANGE} \in \{0, 1, 2\}$
\end{itemize}
Total combinations: $4 \times 5 \times 3 = 60$. IS split: Day $-2$; OOS split: Day $-1$.

\textbf{Top-5 IS results} (by PnL per 10{,}000 timestamps, position limit $L = 80$):

\begin{center}
\begin{tabular}{ccc|cc}
\toprule
$a$ & $Q$ & $S$ & IS PnL & OOS PnL \\
\midrule
1.0 & 5 & 1 & 18{,}420 & 17{,}830 \\
1.0 & 5 & 0 & 18{,}390 & 17{,}410 \\
1.0 & 4 & 1 & 17{,}950 & 17{,}200 \\
0.5 & 5 & 1 & 17{,}810 & 17{,}050 \\
1.0 & 6 & 1 & 17{,}600 & 16{,}800 \\
\bottomrule
\end{tabular}
\end{center}

Best result: $(a=1.0, Q=5, S=1)$ → deployed in \texttt{trader.py}. OOS PnL is $\approx 97\%$ of IS PnL — very stable (no significant overfitting). Note: the IS-OOS ratio degrades for parameters far from the optimum (e.g.\ $Q=7$ drops to $\sim 85\%$). This robustness check is essential: a parameter with high IS PnL but low IS/OOS ratio is likely overfit to day $-2$ noise.

\begin{warningbox}{These Numbers Are Illustrative}
The exact PnL figures above depend on the simulator's fill-matching mode (\texttt{--match-trades all} vs \texttt{worse}). Always re-run the grid search on fresh round data, as bot behaviour and product volatility change each round.
\end{warningbox}

\subsection{Curse of Dimensionality}

With $k$ parameters and $m$ values each, the grid grows as $m^k$:
\begin{center}
\begin{tabular}{ccc}
\toprule
Parameters ($k$) & Values ($m$) & Combinations ($m^k$) \\
\midrule
2 & 10 & 100 \\
3 & 10 & 1{,}000 \\
5 & 10 & 100{,}000 \\
10 & 10 & 10{,}000{,}000 \\
\bottomrule
\end{tabular}
\end{center}
For high-dimensional problems, use \textbf{Bayesian optimisation} (surrogate model of the objective function) or \textbf{random search} (surprisingly effective: random search finds good solutions faster than grid search in high dimensions, as shown by Bergstra \& Bengio 2012).

\section{The Deflated Sharpe Ratio}

Bailey \& L\'{o}pez de Prado (2014) introduce the \textbf{Deflated Sharpe Ratio} to correct for multiple testing:

\begin{equation}
  \text{DSR} = \Phi\!\left(\frac{(\widehat{\text{SR}} - \text{SR}^*)\sqrt{T-1}}{\sqrt{1 - \gamma_1\widehat{\text{SR}} + \frac{\gamma_2-1}{4}\widehat{\text{SR}}^2}}\right)
\end{equation}
where $\text{SR}^* = \Phi^{-1}\!\left(1 - \frac{1-\alpha}{K}\right)\sqrt{\frac{\nicefrac{\gamma_2-1}{4}\bar{\sigma}^2 + 1}{T-1}}$ is the expected maximum SR from $K$ independent trials, $\gamma_1$ and $\gamma_2$ are skewness and kurtosis. DSR $\geq 0.95$ is the threshold for statistical significance adjusted for multiple testing.

\begin{takeawaybox}{Chapter 26}
\begin{itemize}
  \item A backtest answers ``what would have happened'' — not ``what will happen''. Always maintain strict in-sample/out-of-sample separation.
  \item Key biases: look-ahead (using future data), survivorship (only backtesting survivors), data snooping (optimising on same data used to test), overfitting (fitting noise).
  \item Walk-forward: optimise on rolling IS window, evaluate on subsequent OOS block. The gold standard.
  \item Grid search: exhaustive parameter sweep. Report OOS performance, not IS. Bonferroni-correct for multiple testing.
  \item Deflated SR: correct the Sharpe ratio for the number of trials and non-normality.
  \item In Prosperity: use \texttt{--match-trades worse} for realistic fills; run both days for more data points.
\end{itemize}
\end{takeawaybox}

\begin{keyconceptbox}{Decision Guide: Validating a Trading Signal}
\textbf{Step 1 — In-sample significance (necessary, not sufficient):}
\begin{itemize}
  \item Compute signal-return Pearson correlation $\rho$. Report $t = \rho\sqrt{(T-2)/(1-\rho^2)}$.
  \item Require $|\rho| > 0.10$ \emph{and} $t > 3$ (higher bar than the standard $t > 2$ to account for fat tails).
\end{itemize}
\textbf{Step 2 — Multiple-testing correction:}
\begin{itemize}
  \item If $K$ signals were tested: require $p < \nicefrac{0.05}{K}$ (Bonferroni) or Deflated SR $\geq 0.95$.
  \item Always report in-sample correlation with the caveat ``in-sample, before multiple-testing correction''.
\end{itemize}
\textbf{Step 3 — Economic plausibility:}
\begin{itemize}
  \item Write a 1-sentence microstructure story. Example: ``OBI predicts direction because large order imbalance precedes informed buying''.
  \item If you cannot articulate one, the signal is likely data-mined noise.
\end{itemize}
\textbf{Step 4 — OOS walk-forward validation:}
\begin{itemize}
  \item Split: Day $-2$ IS, Day $-1$ OOS (or vice versa, or 80/20 within a day).
  \item Signal must survive OOS with $\geq 50\%$ of IS effect size.
\end{itemize}
\textbf{Step 5 — Practical PnL test:}
\begin{itemize}
  \item Backtest with and without signal. Require statistically significant PnL improvement ($p < 0.05$), not just a positive average.
  \item If PnL gain $<$ transaction cost increase, discard the signal regardless of statistical significance.
\end{itemize}
\end{keyconceptbox}

% End of Part VI

% ============================================================
%
%  PART VII — ADVANCED QUANTITATIVE METHODS
%
% ============================================================
\part{Advanced Quantitative Methods}

This final part gives a deep treatment of the quantitative and statistical methods that underpin the Prosperity codebase. Each chapter starts from first principles and builds to the exact formulae used in \texttt{trader.py} and \texttt{eda.py}.

\chapter{Time Series Analysis}

Price data is a \emph{time series}: a sequence of observations indexed by time. Standard statistical tools assume independent, identically distributed (i.i.d.) observations — a catastrophically wrong assumption for price data. Time series analysis provides the correct framework.

\section{Stationarity}

\begin{definitionbox}{Strict Stationarity}
A process $\{X_t\}$ is \textbf{strictly stationary} if the joint distribution of $(X_{t_1}, \ldots, X_{t_k})$ is identical to that of $(X_{t_1+h}, \ldots, X_{t_k+h})$ for all $k$, $t_1,\ldots,t_k$, and shift $h$. The distribution is invariant under time shifts.
\end{definitionbox}

\begin{definitionbox}{Weak (Covariance) Stationarity}
$\{X_t\}$ is \textbf{weakly stationary} if:
\begin{enumerate}[label=\arabic*.]
  \item $\E[X_t] = \mu$ (constant mean for all $t$)
  \item $\Var(X_t) = \sigma^2 < \infty$ (finite, constant variance)
  \item $\Cov(X_t, X_{t+h}) = \gamma(h)$ depends only on the lag $h$, not on $t$
\end{enumerate}
\end{definitionbox}

\textbf{Why stationarity matters}: Most statistical inference (regression, correlation, forecasting) assumes stationarity. Regressing a non-stationary series on another non-stationary series produces \emph{spurious regression}: high $R^2$ and significant t-statistics even for unrelated series (Granger \& Newbold 1974).

\textbf{Price levels are NOT stationary.} Log returns usually are. This is why we always work with returns in statistical analysis.

\section{White Noise and the Random Walk}

\begin{definitionbox}{White Noise}
$\{\varepsilon_t\}$ is \textbf{white noise} if $\E[\varepsilon_t] = 0$, $\Var(\varepsilon_t) = \sigma^2$ for all $t$, and $\Cov(\varepsilon_t, \varepsilon_s) = 0$ for $t \neq s$. Purely unpredictable — no structure to exploit.
\end{definitionbox}

\begin{definitionbox}{Random Walk}
$X_t = X_{t-1} + \varepsilon_t$ where $\varepsilon_t$ is white noise. Equivalently $X_t = X_0 + \sum_{i=1}^t \varepsilon_i$.
The random walk is \emph{non-stationary}: $\Var(X_t) = t\sigma^2 \to \infty$. It has no tendency to revert to any level — every step is equally likely to go up or down regardless of current value.
\end{definitionbox}

A random walk with drift: $X_t = \mu + X_{t-1} + \varepsilon_t$ adds a deterministic trend component.

\section{Autoregressive Models: AR(\texorpdfstring{$p$}{p})}

\begin{definitionbox}{AR($p$) Model}
An \textbf{autoregressive model of order $p$} expresses the current value as a linear combination of the $p$ most recent values plus white noise:
\[
  X_t = c + \phi_1 X_{t-1} + \phi_2 X_{t-2} + \cdots + \phi_p X_{t-p} + \varepsilon_t
\]
Using the lag operator $L$ ($LX_t = X_{t-1}$):
\[
  \Phi(L)X_t = c + \varepsilon_t, \qquad \Phi(L) = 1 - \phi_1 L - \cdots - \phi_p L^p
\]
\end{definitionbox}

\textbf{Stationarity condition}: The AR($p$) is stationary if and only if all roots of the characteristic polynomial $\Phi(z) = 0$ lie \emph{outside} the unit circle (equivalently, all eigenvalues of the companion matrix have modulus $< 1$).

\subsection{The AR(1) Model in Depth}

The AR(1) is the most important special case:
\[
  X_t = c + \phi X_{t-1} + \varepsilon_t
\]
\textbf{Key properties:}
\begin{align}
  \E[X_t] &= \frac{c}{1-\phi} = \mu \quad (|\phi| < 1) \\
  \Var(X_t) &= \frac{\sigma_\varepsilon^2}{1-\phi^2} \\
  \Corr(X_t, X_{t+h}) &= \phi^h \quad \text{(autocorrelation at lag $h$)}
\end{align}
The autocorrelation decays geometrically at rate $|\phi|$. For $\phi \in (0,1)$: positive autocorrelation (trending). For $\phi \in (-1,0)$: alternating signs (oscillatory, mean-reverting). For $|\phi| = 1$: unit root (random walk, non-stationary). For $|\phi| > 1$: explosive (non-stationary, diverging).

\begin{examplebox}{AR(1) for TOMATOES Mid-Price Changes}
The EDA finds that in the tight-spread regime (spread $\leq 13$), the first-order autocorrelation of mid-price \emph{changes} $\Delta m_t = m_t - m_{t-1}$ is:
\[
  \hat{\phi}_{\Delta m}^{\text{tight}} \approx -0.30
\]
This means: if the price rose last period, it is expected to fall slightly this period (mean reversion in returns). The AR(1) model for changes:
\[
  \Delta m_t = -0.30 \cdot \Delta m_{t-1} + \varepsilon_t
\]

\textbf{Statistical significance.} For $n \approx 2{,}000$ tight-regime observations, the asymptotic standard error of the OLS estimator is:
\[
  SE(\hat{\phi}) = \sqrt{\frac{1-\hat{\phi}^2}{n}} \approx \sqrt{\frac{1-0.09}{2000}} \approx 0.021
\]
The $t$-statistic $t = -0.30/0.021 \approx -14$ far exceeds the 5\% critical value $|t| > 2$: $H_0\!: \phi = 0$ is strongly rejected. Note: to distinguish $\phi = -0.30$ from $\phi = -0.20$ at 5\% significance requires $n \geq 5{,}000$ observations — this limits how precisely the trading coefficient $\phi = 0.05$ can be validated against $\phi = 0$.

FV correction: $FV_t = m_t - 0.05 \cdot \Delta m_t$ (using $\phi = 0.05$, the grid-optimised trading coefficient, not the raw statistical $-0.30$, to avoid over-correction relative to the bot's crossing threshold).
\end{examplebox}

\begin{warningbox}{Sample Size and AR(1) Estimation Precision}
With $n = 2{,}000$ tight-regime observations, $SE(\hat{\phi}) \approx 0.021$. A 95\% CI for the true $\phi$ is approximately $[-0.34,\,-0.26]$. This is precise enough to confirm mean-reversion exists, but \emph{not} precise enough to pin down the optimal trading coefficient $\phi^* = 0.05$ with high confidence — the grid-searched optimum has a standard error of similar magnitude. Treat $\phi^* = 0.05$ as an order-of-magnitude calibration, and re-run the grid search per round with fresh data.
\end{warningbox}

\subsection{Yule-Walker Equations}

For an AR($p$), the autocorrelation function satisfies the \textbf{Yule-Walker equations}:
\begin{equation}
  \rho_k = \phi_1 \rho_{k-1} + \phi_2 \rho_{k-2} + \cdots + \phi_p \rho_{k-p}, \qquad k = 1, 2, \ldots, p
\end{equation}
In matrix form: $\boldsymbol{\rho} = \mathbf{R}\boldsymbol{\phi}$, giving:
$\hat{\boldsymbol{\phi}} = \mathbf{R}^{-1}\boldsymbol{\rho}$ (method of moments estimator).

\section{Moving Average Models: MA(\texorpdfstring{$q$}{q})}

\begin{definitionbox}{MA($q$) Model}
A \textbf{moving average model of order $q$} expresses the current value as a linear combination of current and past \emph{shocks}:
\[
  X_t = \mu + \varepsilon_t + \theta_1\varepsilon_{t-1} + \cdots + \theta_q\varepsilon_{t-q}
\]
MA($q$) is always stationary. The autocorrelation function cuts off after lag $q$: $\rho_k = 0$ for $|k| > q$ (the key identifying feature of MA processes in data).
\end{definitionbox}

\textbf{AR vs.\ MA identification:}
\begin{center}
\begin{tabular}{lll}
\toprule
\textbf{Model} & \textbf{ACF} & \textbf{PACF} \\
\midrule
AR($p$) & Decays exponentially (infinite) & Cuts off after lag $p$ \\
MA($q$) & Cuts off after lag $q$ & Decays exponentially (infinite) \\
ARMA($p,q$) & Decays after lag $q-p$ & Decays after lag $p-q$ \\
\bottomrule
\end{tabular}
\end{center}

\section{ARMA and ARIMA}

\begin{definitionbox}{ARMA($p$, $q$)}
Combines autoregressive and moving-average components:
\[
  X_t = c + \sum_{i=1}^p \phi_i X_{t-i} + \varepsilon_t + \sum_{j=1}^q \theta_j \varepsilon_{t-j}
\]
Box-Jenkins methodology: identify $p,q$ from ACF/PACF $\to$ estimate by MLE $\to$ diagnose residuals.
\end{definitionbox}

\begin{definitionbox}{ARIMA($p$, $d$, $q$)}
If $X_t$ is non-stationary, difference it $d$ times until stationary, then fit ARMA($p,q$):
\[
  \Delta^d X_t \sim \text{ARMA}(p,q)
\]
where $\Delta X_t = X_t - X_{t-1}$ and $\Delta^d = \Delta(\Delta^{d-1})$. Most financial price series require $d=1$ (log prices are I(1): log returns are stationary).
\end{definitionbox}

\textbf{Model selection}: Use \textbf{AIC} (Akaike) or \textbf{BIC} (Bayesian) information criteria:
\begin{align}
  \text{AIC} &= 2k - 2\ln(\hat{L}) \\
  \text{BIC} &= k\ln(n) - 2\ln(\hat{L})
\end{align}
where $k$ = number of parameters, $\hat{L}$ = maximised likelihood, $n$ = sample size. Lower is better; BIC penalises complexity more (prefers simpler models).

\section{Unit Root Tests: The ADF Test}

The \textbf{Augmented Dickey-Fuller (ADF)} test is the standard test for a unit root (non-stationarity):

\begin{definitionbox}{ADF Test}
The null hypothesis is $H_0$: unit root (non-stationary). The test regression is:
\[
  \Delta X_t = \alpha + \beta t + \gamma X_{t-1} + \sum_{j=1}^p \delta_j \Delta X_{t-j} + u_t
\]
The test statistic is $\tau = \hat{\gamma}/\text{SE}(\hat{\gamma})$. Reject $H_0$ (conclude stationarity/mean reversion) if $\tau$ is sufficiently negative (below critical values: $-3.43$ at 5\% for the with-trend version).
\end{definitionbox}

\textbf{Interpretation in trading:}
\begin{itemize}
  \item Reject $H_0$ → price is stationary (mean-reverting) → mean-reversion strategy applicable
  \item Fail to reject $H_0$ → price may have unit root → trend-following more appropriate
  \item In Prosperity: EMERALDS ADF p-value $\ll 0.001$ (strongly stationary around 10000). TOMATOES returns are stationary; levels show slow drift.
\end{itemize}

\section{ACF and PACF}

\begin{definitionbox}{Autocorrelation Function (ACF)}
The ACF at lag $h$ is:
\[
  \rho_h = \frac{\Cov(X_t, X_{t-h})}{\Var(X_t)} = \frac{\gamma(h)}{\gamma(0)}
\]
Estimated as $\hat{\rho}_h = \hat{\gamma}(h)/\hat{\gamma}(0)$ where $\hat{\gamma}(h) = \frac{1}{n}\sum_{t=h+1}^n(X_t - \bar{X})(X_{t-h} - \bar{X})$.
\end{definitionbox}

\begin{definitionbox}{Partial Autocorrelation Function (PACF)}
The PACF at lag $h$ is the correlation between $X_t$ and $X_{t-h}$ after removing the linear influence of $X_{t-1}, \ldots, X_{t-h+1}$. Equivalently, it is the coefficient $\phi_{hh}$ in the AR($h$) regression of $X_t$ on $X_{t-1}, \ldots, X_{t-h}$.
\end{definitionbox}

In \texttt{eda.py}, ACF and PACF plots are generated for both EMERALDS and TOMATOES mid-prices and returns (the \texttt{EMERALDS\_02\_acf\_pacf.png} and \texttt{TOMATOES\_02\_acf\_pacf.png} outputs). These plots reveal:
\begin{itemize}
  \item EMERALDS: very weak, quickly decaying ACF in mid-price (nearly stationary)
  \item TOMATOES: ACF with significant lag-1 negative spike in tight regime (AR(1) signature with $\phi \approx -0.30$)
\end{itemize}

\begin{takeawaybox}{Chapter 27}
\begin{itemize}
  \item Stationarity: constant mean, variance, and covariance structure. Price levels are NOT stationary; log returns usually are.
  \item AR($p$): $X_t = c + \sum\phi_i X_{t-i} + \varepsilon_t$. Stationary if all roots of $\Phi(z)=0$ outside unit circle.
  \item AR(1) autocorrelation: $\rho_h = \phi^h$. For TOMATOES returns: $\phi \approx -0.30$ (tight regime).
  \item ACF of AR($p$): decays geometrically. PACF cuts off at lag $p$. Reverse for MA($q$).
  \item ARIMA($p,d,q$): difference $d$ times to achieve stationarity, then fit ARMA.
  \item ADF test: $H_0$ = unit root. Reject $\to$ stationary (mean-reverting). Use AIC/BIC for model order selection.
\end{itemize}
\end{takeawaybox}

% ──────────────────────────────────────────────────────────────────────────────

\chapter{The Hurst Exponent}

\section{Long Memory and Self-Similarity}

Financial time series exhibit a richer structure than AR(1) can capture. The \textbf{Hurst exponent} $H$ characterises the long-range dependence of a time series and whether it tends to revert or trend.

\begin{definitionbox}{Hurst Exponent}
For a time series $\{X_t\}$, the \textbf{Hurst exponent} $H$ characterises the scaling of the range:
\[
  \E\!\left[\frac{R(n)}{S(n)}\right] \propto n^H
\]
where $R(n) = \max_{1 \leq k \leq n}(\sum_{i=1}^k \Delta X_i) - \min_{1 \leq k \leq n}(\sum_{i=1}^k \Delta X_i)$ is the range and $S(n)$ is the standard deviation over a window of size $n$.
\end{definitionbox}

\textbf{Interpretation:}
\begin{center}
\begin{tabular}{cll}
\toprule
$H$ & \textbf{Process} & \textbf{Trading implication} \\
\midrule
$H < 0.5$ & Anti-persistent (mean-reverting) & Mean-reversion strategies \\
$H = 0.5$ & Random walk (Brownian motion) & No predictability \\
$H > 0.5$ & Persistent (trending, long memory) & Trend-following strategies \\
\bottomrule
\end{tabular}
\end{center}

\section{Rescaled Range (R/S) Analysis}

The classic method for estimating $H$ (Harold Hurst, 1951):

\begin{enumerate}[label=\arabic*.]
  \item Divide the series into blocks of size $n$.
  \item For each block, compute the mean-adjusted cumulative sum and the range $R$ and std $S$.
  \item Average $R/S$ across blocks.
  \item Repeat for multiple block sizes $n$.
  \item Regress $\ln(R/S)$ on $\ln(n)$: the slope is $\hat{H}$.
\end{enumerate}

\begin{lstlisting}[caption={R/S Hurst exponent estimation}]
import numpy as np

def hurst_rs(ts, min_n=10, max_n=None):
    if max_n is None:
        max_n = len(ts) // 4
    ns = np.unique(np.logspace(np.log10(min_n), np.log10(max_n), 20).astype(int))
    rs_vals = []
    for n in ns:
        blocks = [ts[i:i+n] for i in range(0, len(ts)-n, n)]
        rs_block = []
        for block in blocks:
            mean = np.mean(block)
            cum_dev = np.cumsum(block - mean)
            R = cum_dev.max() - cum_dev.min()
            S = np.std(block, ddof=1)
            if S > 0:
                rs_block.append(R / S)
        if rs_block:
            rs_vals.append(np.mean(rs_block))
    log_n = np.log(ns[:len(rs_vals)])
    log_rs = np.log(rs_vals)
    H = np.polyfit(log_n, log_rs, 1)[0]
    return H
\end{lstlisting}

\section{Detrended Fluctuation Analysis (DFA)}

DFA is more robust to non-stationarity than R/S:
\begin{enumerate}[label=\arabic*.]
  \item Integrate: $Y(k) = \sum_{i=1}^k (X_i - \bar{X})$
  \item Divide into non-overlapping windows of size $n$; fit a linear trend in each window.
  \item Compute the root-mean-square fluctuation around the local trend: $F(n) = \sqrt{\frac{1}{N}\sum_{k=1}^N [Y(k) - \hat{Y}_n(k)]^2}$
  \item Regress $\ln F(n)$ on $\ln n$: slope $= H_{\text{DFA}}$.
\end{enumerate}

$H_{\text{DFA}} = 0.5$ for random walk; $H > 0.5$ for persistent; $H < 0.5$ for anti-persistent.

\section{Hurst in the Prosperity Context}

The \texttt{eda.py} generates Hurst exponent plots (\texttt{EMERALDS\_05\_hurst.png}, \texttt{TOMATOES\_05\_hurst.png}). Typical results:
\begin{itemize}
  \item \textbf{EMERALDS}: $H \approx 0.2$--$0.3$ (strongly anti-persistent; mean-reversion is the dominant regime). \emph{Observed on Prosperity~4 Tutorial data (2 days, $n \approx 10{,}000$/day); 95\% CI $\approx \pm 0.04$; re-estimate for each new round.} Consistent with the constant-FV market-making strategy.
  \item \textbf{TOMATOES}: $H \approx 0.4$--$0.5$ (mild anti-persistence; slight mean-reversion tendency in the tight regime). \emph{Same caveats apply; the modified R/S test should be used since TOMATOES has AR(1) structure.} Consistent with AR(1) $\phi \approx -0.30$ in that regime.
\end{itemize}

\subsection{Statistical Distribution of the R/S Statistic (Lo 1991)}

The classical R/S statistic converges to a non-standard distribution under the null hypothesis $H = 0.5$ (i.i.d.\ returns). Lo (1991) derived the limiting distribution and showed that the uncorrected R/S test is severely size-distorted in the presence of short-range autocorrelation (e.g.\ AR(1) returns). The \textbf{modified R/S statistic} corrects for this:
\[
  Q_n = \frac{1}{\hat{\omega}_q} \left(\max_{1 \leq k \leq n} \sum_{t=1}^k (X_t - \bar{X}) - \min_{1 \leq k \leq n} \sum_{t=1}^k (X_t - \bar{X})\right)
\]
where $\hat{\omega}_q^2 = \hat{\sigma}^2 + 2\sum_{j=1}^q \left(1 - \frac{j}{q+1}\right)\hat{\gamma}_j$ is the Newey-West (1987) estimator of the long-run variance, with $\hat{\gamma}_j = \frac{1}{n}\sum_{t=j+1}^n (X_t - \bar{X})(X_{t-j} - \bar{X})$ the sample autocovariance at lag $j$, and $q$ chosen as $\lfloor 4(n/100)^{2/9} \rfloor$ (Andrews 1991 data-driven bandwidth). Setting $q=0$ recovers the classical R/S statistic.

\textbf{Why this matters for Prosperity:} TOMATOES returns have AR(1) $\phi = -0.30$. The classical R/S test applied to TOMATOES would systematically overstate evidence for mean-reversion (spuriously low $\hat{H}$) because the AR(1) structure increases the R/S numerator. The modified statistic should be used instead.

\textbf{Approximate 95\% confidence interval under $H_0: H = 0.5$.}
Under the null, $Q_n / \sqrt{n} \xrightarrow{d} V$, where $V$ is the range of a standard Brownian bridge. The 2.5\% and 97.5\% quantiles of $V$ are approximately 0.809 and 1.862 (Lo 1991, Table II). A computed $Q_n / \sqrt{n}$ outside $[0.809, 1.862]$ rejects $H_0$ at the 5\% level.

\begin{keyconceptbox}{Practical Rule for Hurst Inference}
\begin{enumerate}[label=\arabic*.]
  \item For a series with short-range dependence (AR(1), GARCH), always use the modified R/S or DFA, not classical R/S.
  \item With $n = 10{,}000$ observations (Prosperity day): the standard error of $\hat{H}$ from a DFA log-log regression over $\sim 20$ scales is approximately $\hat{\sigma}_H \approx 0.02$ (OLS standard error of the slope; parametric, derived from residual variance of the DFA scaling regression). Thus $\hat{H} \pm 2\hat{\sigma}_H \approx \hat{H} \pm 0.04$. So $\hat{H} = 0.45$ is statistically indistinguishable from 0.5 at $n = 10{,}000$. \emph{For the Lo (1991) modified R/S test, the equivalent statement is: $Q_n/\sqrt{n} \in [0.809, 1.862]$ is the 95\% acceptance region (asymptotic, from Brownian bridge range quantiles).}
  \item Use Hurst to confirm regime family (mean-reverting vs.\ trending) but never as the sole decision criterion. Pair with ADF test and regime-conditional ACF.
\end{enumerate}
\end{keyconceptbox}

\begin{warningbox}{Hurst Estimation Has High Variance}
Estimating $H$ from short samples (fewer than 500 observations) produces wide confidence intervals. Both R/S and DFA are biased in finite samples: R/S has upward bias (overestimates $H$) for short series; DFA is less biased but requires $n > 200$ for reliable estimates. The Lo (1991) modified R/S test is the gold standard for testing $H = 0.5$ in the presence of short-range dependence. Always report uncertainty bounds alongside the point estimate.
\end{warningbox}

\subsection{Modified R/S: Python Implementation}

\begin{lstlisting}[caption={Lo (1991) modified R/S test with Newey-West long-run variance}]
import numpy as np

def modified_rs(x: np.ndarray, q: int | None = None) -> tuple[float, bool]:
    """
    Lo (1991) modified R/S statistic.
    Returns (Q_n / sqrt(n), reject_H0) where H0: H = 0.5 (random walk).
    95% CI under H0: [0.809, 1.862] (Lo 1991, Table II).
    q: Newey-West lag truncation. Default: Andrews (1991) bandwidth.
    """
    n = len(x)
    mu = np.mean(x)
    deviations = x - mu

    # Rescaled range
    cumsum = np.cumsum(deviations)
    R = cumsum.max() - cumsum.min()

    # Newey-West long-run variance
    if q is None:
        q = int(np.floor(4 * (n / 100) ** (2 / 9)))  # Andrews (1991)
    gamma0 = np.var(x, ddof=0)
    omega2 = gamma0
    for j in range(1, q + 1):
        gamma_j = np.mean(deviations[j:] * deviations[:-j])
        omega2 += 2 * (1 - j / (q + 1)) * gamma_j
    omega = np.sqrt(max(omega2, 1e-12))

    Q = (R / omega) / np.sqrt(n)

    # 5% critical values from Lo (1991) Table II
    reject = not (0.809 <= Q <= 1.862)
    return Q, reject


# Example usage on TOMATOES tight-regime returns:
# Q, reject = modified_rs(tight_regime_returns)
# print(f"Q = {Q:.3f}  ->  {'Reject H0 (H != 0.5)' if reject else 'Fail to reject H0'}")
\end{lstlisting}

\begin{takeawaybox}{Chapter 28}
\begin{itemize}
  \item Hurst exponent $H$: $< 0.5$ = mean-reverting, $= 0.5$ = random walk, $> 0.5$ = trending/persistent.
  \item R/S analysis: slope of $\ln(R/S)$ vs $\ln(n)$ gives $H$. Simple but biased for non-stationary series.
  \item DFA: more robust; detrend within windows before computing fluctuation. Same interpretation.
  \item EMERALDS: $H \approx 0.2$--$0.3$ (strongly mean-reverting). TOMATOES: $H \approx 0.4$--$0.5$ (mild mean-reversion).
  \item Use Hurst to confirm which strategy family to apply; always pair with ADF test and ACF analysis.
\end{itemize}
\end{takeawaybox}

% ──────────────────────────────────────────────────────────────────────────────

\chapter{Regime Switching Models}

\section{Why Regimes Matter}

Real markets do not have constant statistical properties. A strategy calibrated on one regime will fail when the regime changes. Regime switching models formalise the idea that the market alternates between a discrete set of states with different dynamics.

In Prosperity, TOMATOES has two empirically identified regimes (tight spread and wide spread) with dramatically different AR(1) structures. Ignoring this and estimating a single AR(1) on all data gives a biased estimate that under-corrects in the tight regime and over-corrects in the wide regime.

\section{Threshold Autoregressive (TAR) Models}

The simplest regime-switching model: the regime depends on a threshold applied to an observable variable.

\begin{definitionbox}{TAR Model}
\[
  X_t = \begin{cases}
    c_1 + \phi_1 X_{t-1} + \sigma_1\varepsilon_t & \text{if } Z_{t-d} \leq \tau \\
    c_2 + \phi_2 X_{t-1} + \sigma_2\varepsilon_t & \text{if } Z_{t-d} > \tau
  \end{cases}
\]
where $Z_{t-d}$ is the \emph{threshold variable} (observable at time $t$) and $\tau$ is the threshold value. If $Z_t = X_t$ (self-exciting TAR, SETAR), the regime depends on the level of the series itself.
\end{definitionbox}

In \texttt{trader.py}, the threshold variable is the spread $S_t = p_t^A - p_t^B$ and the threshold is $\tau = 13$:
\begin{equation}
  \text{Regime}_t = \begin{cases} \text{Tight} & \text{if } S_t \leq 13 \\ \text{Wide}  & \text{if } S_t > 13 \end{cases}
\end{equation}
This is exactly a TAR model applied to the fair value correction coefficient $\phi$.

\section{Markov Switching Models}

\begin{definitionbox}{Markov Switching Model (Hamilton 1989)}
Let $s_t \in \{1, 2, \ldots, K\}$ be an unobservable state following a \textbf{first-order Markov chain} with transition matrix $\mathbf{P}$:
\[
  P_{ij} = P(s_t = j \mid s_{t-1} = i), \qquad \sum_j P_{ij} = 1 \;\forall i
\]
In each state $k$, the process follows its own AR dynamics:
\[
  X_t = c_k + \phi_k X_{t-1} + \sigma_k\varepsilon_t \quad \text{when } s_t = k
\]
The state is latent (unobserved); inference is done via the \textbf{Hamilton filter}.
\end{definitionbox}

\textbf{Hamilton filter}: a recursive Bayesian filter that computes the conditional probability of being in each state given observed data:
\[
  P(s_t = k \mid X_1, \ldots, X_t) \propto P(X_t \mid s_t = k, X_{t-1}) \sum_j P_{jk} P(s_{t-1} = j \mid X_1, \ldots, X_{t-1})
\]

\subsection{Two-State Example}

For TOMATOES, a two-state Markov switching model would have:
\begin{align}
  \text{State 1 (tight)}: \quad & \phi = -0.30, \; \sigma_1 < \sigma_2 \\
  \text{State 2 (wide)}: \quad  & \phi \approx 0, \; \sigma_2 > \sigma_1
\end{align}
Transition matrix $\mathbf{P}$ estimated from data; the steady-state probabilities $\pi_1 \approx 0.55$ (tight) and $\pi_2 \approx 0.45$ (wide) match the empirical frequency.

\section{Hidden Markov Models (HMMs)}

HMMs are the generalised version of Markov switching models where state emissions can follow any distribution. They are widely used for regime detection, volatility forecasting, and order flow classification.

\subsection{Model Definition}

An HMM with $K$ states and $T$ observations is defined by three parameters $\lambda = (\mathbf{A}, \mathbf{B}, \boldsymbol{\pi})$:

\begin{definitionbox}{HMM Parameters}
\begin{itemize}
  \item \textbf{Transition matrix} $\mathbf{A} \in \mathbb{R}^{K \times K}$: $A_{ij} = P(s_t = j \mid s_{t-1} = i)$, with $\sum_j A_{ij} = 1$.
  \item \textbf{Emission distribution} $b_k(o) = P(O_t = o \mid s_t = k)$. For continuous observations (e.g.\ returns), use Gaussian: $b_k(o) = \mathcal{N}(o;\, \mu_k, \sigma_k^2)$.
  \item \textbf{Initial distribution} $\boldsymbol{\pi}$: $\pi_k = P(s_1 = k)$.
\end{itemize}
The state sequence $s_1, \ldots, s_T$ is \textbf{latent} (unobserved). The observation sequence $O_1, \ldots, O_T$ is observed.
\end{definitionbox}

\subsection{Algorithm 1: The Forward Algorithm (Filtering)}

The \textbf{forward variable} $\alpha_t(k) = P(O_1, \ldots, O_t,\, s_t = k \mid \lambda)$ is computed recursively:
\begin{align}
  \alpha_1(k) &= \pi_k \cdot b_k(O_1) & \text{(initialisation)} \\
  \alpha_{t+1}(j) &= b_j(O_{t+1}) \sum_{k=1}^K \alpha_t(k)\, A_{kj} & \text{(recursion)}
\end{align}
The total likelihood is $P(\mathbf{O} \mid \lambda) = \sum_k \alpha_T(k)$.

\textbf{ML analogy:} The forward algorithm is a \emph{Hidden Markov belief propagation} — exactly the forward pass of the Kalman filter, but with discrete states and non-linear emissions.

\subsection{Algorithm 2: Baum-Welch (EM for HMMs)}

Baum-Welch finds the maximum-likelihood parameters $\hat{\lambda}$ via EM. It alternates:

\textbf{E-step} — compute posterior state probabilities using forward ($\alpha$) and backward ($\beta$) variables.

Backward variable $\beta_t(k) = P(O_{t+1}, \ldots, O_T \mid s_t = k, \lambda)$:
\begin{align}
  \beta_T(k) &= 1 \\
  \beta_t(k) &= \sum_{j=1}^K A_{kj}\, b_j(O_{t+1})\, \beta_{t+1}(j)
\end{align}

State posterior (responsibility):
\[
  \gamma_t(k) = P(s_t = k \mid \mathbf{O}, \lambda) = \frac{\alpha_t(k)\,\beta_t(k)}{\sum_{j} \alpha_t(j)\,\beta_t(j)}
\]

Joint state-transition posterior:
\[
  \xi_t(k, j) = P(s_t = k, s_{t+1} = j \mid \mathbf{O}, \lambda)
  = \frac{\alpha_t(k)\, A_{kj}\, b_j(O_{t+1})\, \beta_{t+1}(j)}{\sum_{k'}\sum_{j'} \alpha_t(k')\, A_{k'j'}\, b_{j'}(O_{t+1})\, \beta_{t+1}(j')}
\]

\textbf{M-step} — update parameters to maximise expected log-likelihood:
\begin{align}
  \hat{\pi}_k &= \gamma_1(k) \\
  \hat{A}_{kj} &= \frac{\sum_{t=1}^{T-1} \xi_t(k,j)}{\sum_{t=1}^{T-1} \gamma_t(k)} \\
  \hat{\mu}_k &= \frac{\sum_{t=1}^T \gamma_t(k)\, O_t}{\sum_{t=1}^T \gamma_t(k)} \quad \text{(Gaussian emission mean)} \\
  \hat{\sigma}_k^2 &= \frac{\sum_{t=1}^T \gamma_t(k)\,(O_t - \hat{\mu}_k)^2}{\sum_{t=1}^T \gamma_t(k)}
\end{align}

Iterate E→M until convergence (log-likelihood increase $< \epsilon$). Convergence is guaranteed to a local maximum (like all EM).

\subsection{Algorithm 3: Viterbi (Most Likely State Sequence)}

Viterbi finds the \emph{single most likely} state path $s_1^*, \ldots, s_T^*$:
\begin{align}
  \delta_1(k) &= \pi_k\, b_k(O_1) \\
  \delta_t(j) &= b_j(O_t) \cdot \max_k\, \bigl[\delta_{t-1}(k)\, A_{kj}\bigr] & \text{(max instead of sum)} \\
  \psi_t(j) &= \arg\max_k\, \bigl[\delta_{t-1}(k)\, A_{kj}\bigr] & \text{(traceback pointer)}
\end{align}
Backtrace: $s_T^* = \arg\max_k \delta_T(k)$; then $s_t^* = \psi_{t+1}(s_{t+1}^*)$ for $t = T-1, \ldots, 1$.

\textbf{Difference from forward algorithm:} Forward sums over all paths (marginal probability); Viterbi takes the max (MAP sequence). In practice: use forward for online filtering (real-time regime probability), Viterbi for offline analysis (most likely historical regime sequence).

\subsection{Python Implementation: TOMATOES Regime Detection}

\begin{lstlisting}[caption={Gaussian HMM for TOMATOES spread regime detection using hmmlearn}]
from hmmlearn import hmm
import numpy as np

def fit_regime_hmm(mid_prices: np.ndarray, n_states: int = 2,
                   n_iter: int = 200) -> hmm.GaussianHMM:
    """
    Fit a Gaussian HMM to tick-level returns.
    Observation: r_t = m_t - m_{t-1} (scalar)
    State meaning: state with lower variance = tight regime.
    """
    returns = np.diff(mid_prices).reshape(-1, 1)  # shape (T-1, 1)

    model = hmm.GaussianHMM(
        n_components=n_states,
        covariance_type="diag",   # scalar variance per state
        n_iter=n_iter,
        random_state=42,
        init_params="stmc",  # s=startprob, t=transmat, m=means, c=covars
    )
    model.fit(returns)
    return model


def regime_sequence(model: hmm.GaussianHMM,
                    mid_prices: np.ndarray) -> np.ndarray:
    """Viterbi-decoded state sequence (most likely path)."""
    returns = np.diff(mid_prices).reshape(-1, 1)
    _, states = model.decode(returns, algorithm="viterbi")
    return states


def online_regime_prob(model: hmm.GaussianHMM,
                       recent_returns: np.ndarray) -> np.ndarray:
    """
    Forward-filter: probability of each state at the last timestamp.
    Use this in trader.py for real-time regime detection.
    Returns: array of shape (n_states,) summing to 1.
    """
    obs = recent_returns.reshape(-1, 1)
    log_prob, posteriors = model.score_samples(obs)
    return posteriors[-1]   # P(s_T = k | O_1...O_T) for each k


# -- Usage example --
# model = fit_regime_hmm(tomatoes_mid_prices)
# states = regime_sequence(model, tomatoes_mid_prices)
# print("State means:", model.means_.flatten())
# print("State stds: ", np.sqrt(model.covars_.flatten()))
# print("Transition matrix:\n", model.transmat_)
# Tight regime = state with smaller variance
# tight_state = np.argmin(model.covars_.flatten())
\end{lstlisting}

\begin{keyconceptbox}{HMM for TOMATOES: Expected Results}
Fitting a 2-state Gaussian HMM to TOMATOES returns should recover parameters close to:
\begin{itemize}
  \item \textbf{State 0 (tight regime, $\sim 55\%$)}: $\mu \approx 0$, $\sigma \approx 1.3$, strong negative AR structure.
  \item \textbf{State 1 (wide regime, $\sim 45\%$)}: $\mu \approx 0$, $\sigma \approx 1.6$, approximately white noise.
  \item \textbf{Transition matrix}: $P_{00} \approx P_{11} \approx 0.9998$ (very persistent regimes — average duration $\approx 5{,}000$ ticks before switching). \emph{Point estimate from 2-day Tutorial data; bootstrap 95\% CI is wide ($\pm 0.0005$) due to the small number of observed transitions. Treat as order-of-magnitude, not precise.}
\end{itemize}
Compare to the TAR threshold rule: the HMM version is \emph{softer} (gives a probability of being in each regime) and does not require observing the spread directly — useful when the spread variable is noisy.
\end{keyconceptbox}

\begin{warningbox}{HMM Identifiability and Label Switching}
HMMs are not identified up to label permutation: state 0 and state 1 can be swapped without changing the likelihood. After fitting, always sort states by a meaningful statistic (e.g.\ variance) before interpreting. The \texttt{hmmlearn} package does not sort states automatically.

Additionally, EM can converge to local optima. Run with multiple random initialisations ($\geq 5$) and keep the fit with the highest log-likelihood.
\end{warningbox}

\begin{lstlisting}[caption={Post-processing: resolve label switching and multi-start EM}]
import numpy as np
from hmmlearn import hmm

def fit_robust_hmm(returns: np.ndarray, n_states: int = 2,
                   n_starts: int = 10) -> hmm.GaussianHMM:
    """
    Fit HMM with multiple random starts; resolve label switching.
    State 0 = tight (low variance), State 1 = wide (high variance).
    """
    obs = returns.reshape(-1, 1)
    best_model, best_score = None, -np.inf

    for seed in range(n_starts):
        m = hmm.GaussianHMM(n_components=n_states, covariance_type="diag",
                             n_iter=300, random_state=seed,
                             init_params="stmc")
        #  init_params flags: s=startprob, t=transmat, m=means, c=covars
        m.fit(obs)
        score = m.score(obs)   # log-likelihood; higher = better
        if score > best_score:
            best_score, best_model = score, m

    # Resolve label switching: sort states by variance (ascending)
    order = np.argsort(best_model.covars_[:, 0, 0])
    best_model.means_      = best_model.means_[order]
    best_model.covars_     = best_model.covars_[order]
    best_model.startprob_  = best_model.startprob_[order]
    # Permute rows AND columns of transition matrix
    best_model.transmat_   = best_model.transmat_[order][:, order]

    return best_model  # state 0 = tight (low var), state 1 = wide (high var)
\end{lstlisting}

\begin{takeawaybox}{Chapter 29}
\begin{itemize}
  \item TAR model: regime depends on observable threshold variable. In Prosperity: spread $\leq 13$ = tight regime.
  \item Markov switching: unobservable latent state evolves as Markov chain; each state has its own AR parameters. Estimated via Hamilton filter.
  \item HMM: generalised Markov switching with arbitrary emission distributions. Parameters via EM (Baum-Welch); Viterbi for state sequence.
  \item Always model regimes when parameter instability is detected in-sample (e.g.\ subsample ACFs disagree).
\end{itemize}
\end{takeawaybox}

\begin{keyconceptbox}{Decision Guide: Choosing a Regime Model}
\textbf{Step 1 — Is the threshold variable observable?}
\begin{itemize}
  \item Yes (e.g.\ bid-ask spread, volume, time-of-day) $\Rightarrow$ \textbf{TAR model}: define threshold $\tau$ from in-sample data. Simplest and most interpretable. Example: spread $\leq 13$ for TOMATOES.
  \item No (latent state) $\Rightarrow$ proceed to Step 2.
\end{itemize}
\textbf{Step 2 — Is the number of regimes known, and do you want standard AR dynamics?}
\begin{itemize}
  \item Yes $\Rightarrow$ \textbf{Markov Switching AR} (Hamilton filter). Faster, interpretable, no EM instability.
  \item No, or you need flexible emission distributions $\Rightarrow$ \textbf{HMM} (Baum-Welch + Viterbi). More general; run $\geq 10$ random starts to avoid local optima.
\end{itemize}
\textbf{Step 3 — Validate the regime model:}
\begin{enumerate}[label=\arabic*.]
  \item \textit{Within-regime stationarity}: ADF each regime separately.
  \item \textit{Persistence}: transition probability $P_{ii} > 0.95$ indicates regimes are economically meaningful (not noise).
  \item \textit{OOS stability}: do regime assignments agree on held-out data? If regimes flip in sign, revisit the label-switching fix (sort states by variance).
\end{enumerate}
\textbf{Rule of thumb for Prosperity:} if the spread takes 2 stable values, TAR dominates HMM — fewer parameters, no EM instability, directly observable threshold.
\end{keyconceptbox}

% ──────────────────────────────────────────────────────────────────────────────

\chapter{Optimal Market Making: The Avellaneda-Stoikov Framework}

\section{Motivation and Setup}

The heuristic approach to market making (Chapter~15) leaves open the question: what are the \emph{theoretically optimal} bid and ask quotes given the market maker's inventory and risk aversion? Avellaneda \& Stoikov (2008) answer this with a stochastic control framework.

\textbf{Setup:}
\begin{itemize}
  \item Mid-price follows Brownian motion: $dS_t = \sigma\,\dd W_t$ (zero drift for simplicity)
  \item Market maker posts bid $S_t - \delta^b$ and ask $S_t + \delta^a$, where $\delta^b, \delta^a > 0$ are the half-spreads to optimise
  \item Order arrivals: Poisson processes with rates $\lambda^b(\delta^b)$ and $\lambda^a(\delta^a)$ that decrease with quote offset:
  \[
    \lambda^b(\delta^b) = \lambda e^{-\kappa\delta^b}, \qquad \lambda^a(\delta^a) = \lambda e^{-\kappa\delta^a}
  \]
  \item Inventory: $q_t$ (changed by $\pm 1$ at each fill)
  \item Objective: maximise expected utility of terminal wealth with risk aversion $\gamma$:
  \[
    \max_{\delta^b, \delta^a} \E\!\left[W_T - \frac{\gamma}{2}\Var(W_T)\right]
  \]
\end{itemize}

\section{The Reservation Price}

The market maker's \textbf{reservation price} is the price at which they would be indifferent between holding and trading, given their current inventory:

\begin{theorem}[Avellaneda-Stoikov Reservation Price]
The optimal reservation price (internal fair value adjusted for inventory) is:
\[
  r_t = S_t - q_t \gamma \sigma^2 (T - t)
\]
where $S_t$ is the mid-price, $q_t$ is the current inventory, $\gamma$ is the risk-aversion parameter, $\sigma^2$ is the price variance, and $(T-t)$ is the remaining time.
\end{theorem}

\textbf{Intuition}: When the market maker is long ($q_t > 0$), their reservation price is \emph{below} mid ($r_t < S_t$): they want to sell cheaply to unwind inventory, and they price their bids below mid to discourage buying more. This is the formal justification for inventory skewing.

The term $q_t\gamma\sigma^2(T-t)$ is the \textbf{inventory penalty}: larger inventory, higher risk aversion, larger variance, or more time remaining all increase the penalty. Notably, it decreases as $T-t \to 0$ (at the end of the day, the market maker becomes more aggressive since there is less time for losses to compound).

\section{The Optimal Spread}

\begin{theorem}[Avellaneda-Stoikov Optimal Spread]
The optimal symmetric spread around the reservation price is:
\[
  \delta^* = \delta^a = \delta^b = \frac{\gamma\sigma^2(T-t)}{2} + \frac{1}{\kappa}\ln\!\left(1 + \frac{\gamma}{\kappa}\right)
\]
The full optimal quotes are:
\begin{align}
  p^A &= r_t + \delta^* = S_t - q_t\gamma\sigma^2(T-t) + \delta^* \\
  p^B &= r_t - \delta^* = S_t - q_t\gamma\sigma^2(T-t) - \delta^*
\end{align}
\end{theorem}

\textbf{Decomposition of the optimal spread:}
\begin{align}
  \underbrace{\frac{\gamma\sigma^2(T-t)}{2}}_{\text{Inventory risk term}} + \underbrace{\frac{1}{\kappa}\ln\!\left(1 + \frac{\gamma}{\kappa}\right)}_{\text{Adverse selection / arrival term}}
\end{align}

\begin{enumerate}[label=\arabic*.]
  \item \textbf{Inventory risk term}: Spread widens when volatility is high, risk aversion is high, or more time remains (longer exposure to inventory risk). $\to 0$ as $T-t \to 0$.
  \item \textbf{Adverse selection term}: Spread needed to compensate for the risk that fill orders are informed. Depends on the Poisson arrival sensitivity $\kappa$.
\end{enumerate}

\section{HJB Derivation Sketch}
\label{sec:as-hjb}

The A-S results follow from stochastic control theory. This section gives a self-contained derivation sketch (ML analogy: this is Bellman's equation applied to a continuous-time MDP).

\textbf{Value function.} Define the value function $u(x, q, s, t)$ where $x$ is the cash balance, $q$ is the inventory, $s = S_t$ is the mid-price, and $t$ is the current time:
\[
  u(x, q, s, t) = \max_{\delta^b, \delta^a}\, \E\!\left[-e^{-\gamma(x + qS_T)} \mid x_t = x,\, q_t = q,\, S_t = s\right]
\]
This is negative exponential utility with risk aversion $\gamma$ (a CARA utility — ``Constant Absolute Risk Aversion'').

\textbf{Hamilton-Jacobi-Bellman equation.} The HJB equation for $u$ is:
\begin{equation}
  0 = \partial_t u + \frac{\sigma^2}{2}\partial_{ss} u + \lambda e^{-\kappa\delta^a}\left[u(x+s+\delta^a,\, q-1,\, s,\, t) - u\right]
  + \lambda e^{-\kappa\delta^b}\left[u(x-s+\delta^b,\, q+1,\, s,\, t) - u\right]
  \label{eq:hjb-as}
\end{equation}
The first two terms are the standard diffusion part (Itô's lemma applied to $u$). The last two terms capture the jump contributions: a sell at ask $s + \delta^a$ increments cash by $s + \delta^a$ and decrements inventory by 1; a buy at bid $s - \delta^b$ does the reverse.

\textbf{Ansatz (guess for the form of the solution).} Substitute $u = -\exp(-\gamma(x + qs) - \gamma g(q, t))$ where $g(q,t)$ absorbs the inventory penalty.

\textit{Step 1: Compute partial derivatives.} Let $\Psi = \gamma x + \gamma qs + \gamma g(q,t)$, so $u = -e^{-\Psi}$:
\begin{align}
  \partial_t u &= \gamma(\partial_t g)\, u, \qquad
  \partial_s u = \gamma q\, u, \qquad
  \partial_{ss} u = \gamma^2 q^2\, u
\end{align}
The diffusion term becomes $\tfrac{\sigma^2}{2}\partial_{ss} u = \tfrac{\sigma^2\gamma^2 q^2}{2}\, u$.

\textit{Step 2: Simplify the ask jump.} After a sell at price $s+\delta^a$ (cash $+s+\delta^a$, inventory $q \to q-1$):
\begin{align}
  u(x+s+\delta^a,\, q-1,\, s,\, t)
  &= -e^{-\gamma(x+s+\delta^a+(q-1)s)-\gamma g(q-1,t)} \notag \\
  &= u\;\cdot\; e^{-\gamma\delta^a + \gamma[g(q,t)-g(q-1,t)]}
\end{align}
The ask jump contribution is $u\,\lambda e^{-\kappa\delta^a}\!\left(e^{-\gamma\delta^a+\gamma[g(q)-g(q-1)]} - 1\right)$; the bid jump is symmetric.

\textit{Step 3: Divide by $u$ to obtain the reduced ODE.} Substituting into \eqref{eq:hjb-as} and dividing throughout by $u$:
\begin{equation}
  \gamma\partial_t g + \tfrac{\sigma^2\gamma^2 q^2}{2}
  + \lambda e^{-\kappa\delta^a}\!\left(e^{-\gamma\delta^a+\gamma[g(q)-g(q-1)]}-1\right)
  + \lambda e^{-\kappa\delta^b}\!\left(e^{-\gamma\delta^b+\gamma[g(q)-g(q+1)]}-1\right)
  = 0
  \label{eq:hjb-reduced}
\end{equation}

\textbf{Optimising over $\delta^a$.} Set $A = g(q)-g(q-1)$. Differentiate \eqref{eq:hjb-reduced} w.r.t.\ $\delta^a$:
\begin{align}
  &\frac{d}{d\delta^a}\!\left[e^{-\kappa\delta^a}\!\left(e^{-\gamma\delta^a+\gamma A}-1\right)\right] = 0 \notag \\
  \Rightarrow\quad
  &e^{-\kappa\delta^a}\!\left[-(\kappa+\gamma)e^{-\gamma\delta^a+\gamma A}+\kappa\right] = 0 \notag \\
  \Rightarrow\quad
  &e^{-\gamma\delta^a+\gamma A} = \frac{\kappa}{\kappa+\gamma}
\end{align}
Taking logarithms: $\;\delta^{a*} = A + \frac{1}{\gamma}\ln\!\left(1+\frac{\gamma}{\kappa}\right)$. By symmetry, $\delta^{b*} = [g(q)-g(q+1)] + \frac{1}{\gamma}\ln(1+\gamma/\kappa)$.

\textbf{Solving for $g(q,t)$.} Substituting $\delta^{a*}, \delta^{b*}$ back into \eqref{eq:hjb-reduced} yields the PDE:
\[
  \partial_t g = \frac{\gamma\sigma^2 q^2}{2}
\]
With terminal condition $g(q,T)=0$, integrating backward:
\[
  g(q,t) = \frac{\gamma\sigma^2 q^2(T-t)}{2}
\]
Hence $g(q)-g(q-1) = \gamma\sigma^2(q-\tfrac{1}{2})(T-t)$. For symmetric quoting around $q=0$, the half-spread is:
\[
  \delta^* = \frac{\gamma\sigma^2(T-t)}{2} + \frac{1}{\gamma}\ln\!\left(1+\frac{\gamma}{\kappa}\right)
\]
Using the first-order approximation $\frac{1}{\gamma}\ln(1+\gamma/\kappa) \approx \frac{1}{\kappa}$ (accurate when $\gamma \ll \kappa$) recovers the standard A-S formula $\frac{1}{\kappa}\ln(1+\gamma/\kappa)$. The reservation price follows from $g(q)-g(q-1) \approx \gamma\sigma^2 q(T-t)$:
\[
  r_t = S_t - q_t \gamma \sigma^2 (T-t) \qquad \text{(reservation price)}
\]
\[
  \delta^* = \frac{\gamma\sigma^2(T-t)}{2} + \frac{1}{\kappa}\ln\!\left(1 + \frac{\gamma}{\kappa}\right) \qquad \text{(optimal spread)}
\]

\textbf{ML analogy:} The HJB equation is the continuous-time version of the Bellman equation from reinforcement learning. The value function $u$ is the Q-function; the ansatz is an equivalent of a parametric approximation to $u$; and optimising over $(\delta^a, \delta^b)$ at each state is the policy improvement step.

\section{Comparison with the Prosperity Heuristic}

\begin{center}
\begin{tabular}{lll}
\toprule
\textbf{Feature} & \textbf{Avellaneda-Stoikov} & \textbf{Prosperity (\texttt{trader.py})} \\
\midrule
Fair value & Mid-price $S_t$ & EMA + AR(1) correction \\
Inventory adjustment & $-q_t\gamma\sigma^2(T-t)$ (continuous) & $-(q/L)\times S_{\text{range}}$ (discrete skew) \\
Spread & $\delta^* = f(\gamma, \sigma, \kappa, T-t)$ & Fixed $\delta = 5$ (grid-optimised) \\
Adapts to volatility & Yes ($\sigma^2$ in formula) & No (fixed $\delta$) \\
Adapts to time & Yes ($T-t$ decreases) & No \\
Optimality & Theoretically optimal & Practically grid-searched \\
\bottomrule
\end{tabular}
\end{center}

The Prosperity heuristic is a simplification of the A-S framework: the skew term $-(q/L)\times S_{\text{range}}$ approximates the reservation price adjustment $-q\gamma\sigma^2(T-t)$ with a linear, bounded function of inventory fraction.

\section{The A-S Model's Key Insight}

The most important practical implication: \textbf{the optimal bid and ask are NOT symmetric around the mid-price when inventory is non-zero}. They are symmetric around the \emph{reservation price}, which is itself shifted from mid by the inventory penalty. This is why naive market makers who always quote mid $\pm \delta$ accumulate inventory and eventually face large drawdowns.

\begin{takeawaybox}{Chapter 30}
\begin{itemize}
  \item A-S setup: mid-price = BM, order arrivals = Poisson with exponential arrival rate $\lambda e^{-\kappa\delta}$.
  \item Reservation price: $r_t = S_t - q_t\gamma\sigma^2(T-t)$. Long inventory $\to$ reserve price below mid.
  \item Optimal spread: $\delta^* = \frac{\gamma\sigma^2(T-t)}{2} + \frac{1}{\kappa}\ln(1+\gamma/\kappa)$. Widens with vol, risk aversion, and time.
  \item Quotes: bid $= r_t - \delta^*$, ask $= r_t + \delta^*$. Both shift down when long.
  \item Prosperity heuristic: linear approximation of the A-S inventory adjustment.
\end{itemize}
\end{takeawaybox}

% ──────────────────────────────────────────────────────────────────────────────

\chapter{Order Book Imbalance: Theory and Practice}

\section{Formal Definition and Measurement}

\begin{definitionbox}{Multi-Level OBI}
The $L$-level order book imbalance at time $t$:
\[
  \text{OBI}_t^{(L)} = \frac{\sum_{\ell=1}^L q_\ell^B \cdot w_\ell - \sum_{\ell=1}^L q_\ell^A \cdot w_\ell}{\sum_{\ell=1}^L q_\ell^B \cdot w_\ell + \sum_{\ell=1}^L q_\ell^A \cdot w_\ell}
\]
where $q_\ell^B$, $q_\ell^A$ are bid/ask quantities at level $\ell$ and $w_\ell$ are level weights. Common weight choices:
\begin{itemize}
  \item Uniform: $w_\ell = 1$ (simple OBI)
  \item Exponential: $w_\ell = e^{-\lambda(\ell-1)}$ (decays with distance from best)
  \item Proportional: $w_\ell = 1/p_\ell$ (weight inversely by price level)
\end{itemize}
\end{definitionbox}

In Prosperity (3 book levels), with uniform weights:
\begin{lstlisting}[caption={OBI computation from OrderDepth}]
def compute_obi(od: OrderDepth) -> float:
    bid_vol = sum(od.buy_orders.values())          # positive
    ask_vol = sum(-q for q in od.sell_orders.values())  # make positive
    total = bid_vol + ask_vol
    if total == 0:
        return 0.0
    return (bid_vol - ask_vol) / total
\end{lstlisting}

\section{Why OBI Predicts Short-Term Returns}

\subsection{Order Flow Imbalance Theory}

When there is more buying interest (bid volume $>$ ask volume), one of two things must happen:
\begin{enumerate}[label=\arabic*.]
  \item Buyers cross the spread to consume ask volume → price moves up
  \item Sellers post more at the ask to absorb buying interest → book rebalances
\end{enumerate}
In expectation, large positive OBI precedes upward price movement. The predictive power is:
\[
  \E[\Delta m_{t+1} \mid \text{OBI}_t] \approx \beta \cdot \text{OBI}_t
\]
Empirically, $\beta \approx 0.38\times\sigma_{\Delta m} \approx 0.38 \times 1.35 \approx 0.51$ ticks in TOMATOES tight regime.

\subsection{The Microstructure Rationale}

The OBI signal has predictive power because:
\begin{itemize}
  \item Informed traders hide their orders by splitting them across multiple timestamps
  \item A sustained positive OBI reveals that informed buying pressure is building
  \item The book adjusts slowly (market makers replenish quotes gradually)
  \item The predictability decays rapidly (horizon $\leq 1$--$2$ timestamps)
\end{itemize}

\section{Signal Strength vs.\ Passive Edge}

The crucial quantitative comparison from the Prosperity audit:

\begin{keyconceptbox}{OBI vs.\ Passive Edge: The Fundamental Comparison}
Let $\delta$ = passive quote offset from FV (edge per fill) and $\beta\sigma$ = OBI directional signal.

\textbf{Asymmetric sizing}: when OBI $= +1$, post full bid size but reduced ask size (by factor $1 - s$). Expected gain from OBI: $s \times \beta\sigma$. Expected loss from sacrificed edge: $s \times \delta$.

\textbf{Profitable when}: $\beta\sigma > \delta$, i.e.\ $0.51 > 5.0$ — which fails by a factor of $\approx 10$.

\textbf{Conclusion}: OBI sizing is profitable only when the passive half-spread $\delta < \beta\sigma$. In Prosperity with $\delta = 5$, this requires a $10\times$ larger OBI signal than empirically observed. OBI becomes useful only when quoting very tightly ($\delta < 1$ tick).
\end{keyconceptbox}

\section{OBI in High-Frequency Research}

In real HFT contexts (microsecond timescales):
\begin{itemize}
  \item OBI predicts mid-price moves with $R^2 \approx 0.4$--$0.7$ at 10ms horizon (Cont, Kukanov, Stoikov 2014)
  \item The signal decays rapidly: predictive horizon is $\leq 100$ milliseconds in liquid markets
  \item OBI is a key input to optimal execution algorithms (minimise market impact by timing orders with OBI)
  \item In dark pools, OBI is not observable — this creates an information asymmetry advantage for lit-market participants
\end{itemize}

\begin{takeawaybox}{Chapter 31}
\begin{itemize}
  \item OBI $= (Q^B - Q^A)/(Q^B + Q^A) \in [-1, +1]$. Positive = more buy pressure.
  \item Predicts short-term returns: $\E[\Delta m \mid \text{OBI}] \approx 0.38\sigma_{\Delta m}$ in TOMATOES tight regime.
  \item OBI sizing unprofitable when passive edge $\delta > \beta\sigma$. At $\delta = 5$ and $\beta\sigma = 0.51$: ratio $\approx 10\times$ too small to justify sizing.
  \item OBI useful only at sub-tick-spread levels or in execution optimisation, not passive market-making with substantial edge.
\end{itemize}
\end{takeawaybox}

% ──────────────────────────────────────────────────────────────────────────────

\chapter{Fair Value Estimation Methods}

\section{The Fair Value Estimation Problem}

The market maker's core intellectual task is estimating the true value $V^*$ of an asset from noisy, discrete observations. Every quote is placed relative to this estimate. A better estimate $\to$ better quotes $\to$ more PnL.

\section{Method 1: Constant Fair Value}

\begin{definitionbox}{Constant FV}
$FV_t = \mu$ (fixed constant). Appropriate when the price series is strongly stationary around a known mean.
\end{definitionbox}

\textbf{Estimation}: $\mu = \bar{m} = \frac{1}{T}\sum_{t=1}^T m_t$ (sample mean of historical mid-prices).

\textbf{Error bound}: The estimation error is $\hat{\mu} - \mu \sim \mathcal{N}(0, \sigma^2/T)$. For EMERALDS with $\sigma = 0.72$ and $T = 2000$: std error $= 0.72/\sqrt{2000} = 0.016$ ticks — essentially negligible.

\textbf{When to use}: EMERALDS case — ADF strongly rejects unit root, mid-price has std $= 0.72$ around 10000 over 2000 timestamps. The constant FV = 10000 is accurate to within 1 tick essentially always.

\section{Method 2: EMA-Based Fair Value}

\begin{definitionbox}{EMA Fair Value}
$FV_t = \text{EMA}_t = \alpha m_t + (1-\alpha)\text{EMA}_{t-1}$. Tracks slow intraday drift while smoothing noise.
\end{definitionbox}

\textbf{When to use}: Products with slow, persistent drift (CAGR $>0$ but with noise). The EMA ``follows'' the drift without predicting it.

\textbf{Parameter choice}: $\alpha = 2/(n+1)$ for an $n$-period EMA. In Prosperity with $\alpha = 1$, the EMA collapses to the current mid-price (no smoothing, purely reactive).

\section{Method 3: AR(1) Correction}

\begin{definitionbox}{AR(1) Fair Value Correction}
Given the current mid-price $m_t$ and the previous $m_{t-1}$, apply a mean-reversion correction:
\[
  FV_t = m_t - \phi \cdot (m_t - m_{t-1})
\]
where $\phi$ is the AR(1) coefficient of mid-price changes. If $\Delta m_t = m_t - m_{t-1}$ tends to reverse ($\phi < 0$), the correction lowers FV when mid just rose (predicting a fall).
\end{definitionbox}

\textbf{Derivation}: The AR(1) model of changes is $\Delta m_t = \phi \Delta m_{t-1} + \varepsilon_t$. The one-step-ahead forecast:
\[
  \E[\Delta m_{t+1} \mid \Delta m_t] = \phi \Delta m_t
\]
So the best forecast of the \emph{next} mid-price is:
\[
  \E[m_{t+1} \mid m_t, m_{t-1}] = m_t + \phi \Delta m_t = m_t + \phi(m_t - m_{t-1})
\]
But the market maker prices relative to the \emph{current} FV (the price right now), not $t+1$. The correction is therefore subtracted from the current mid to adjust for the expected reversal:
\[
  FV_t = m_t - \phi(m_t - m_{t-1})
\]
Note the sign: with $\phi = -0.30$ (statistical AR(1) for TOMATOES), this gives $FV_t = m_t + 0.30(m_t - m_{t-1})$, amplifying recent moves. But we use $\phi = +0.05$ (the optimal trading coefficient) giving a mild correction.

\section{Method 4: EMA + AR(1) (TOMATOES Strategy)}
\label{sec:ema-ar1}

This is the authoritative treatment of the combined fair value model used in \texttt{trader.py} for TOMATOES. References to this model in Chapters~6 and 11 are summaries; this section contains the full derivation and parameter choices.

\textbf{Motivation.} A drifting mean-reverting price requires two components: (1) an EMA to track slow intraday drift, and (2) an AR(1) correction to exploit short-term mean-reversion of price \emph{changes}. Combined:
\begin{align}
  \text{EMA}_t &= \alpha m_t + (1-\alpha)\text{EMA}_{t-1} \label{eq:ema-fv} \\
  FV_t &= \text{EMA}_t - \phi(m_t - m_{t-1}) \label{eq:fv-combined}
\end{align}

\textbf{The $\alpha = 1$ special case (used in TOMATOES).} Setting $\alpha = 1$ in \eqref{eq:ema-fv}:
\[
  \text{EMA}_t = m_t
\]
Substituting into \eqref{eq:fv-combined}:
\[
  FV_t = m_t - \phi\,\Delta m_t
\]
The EMA term contributes nothing. The model reduces to a \emph{pure AR(1) correction}. Empirical analysis of TOMATOES showed no significant intraday drift (within each day, $\E[\Delta m_t] \approx 0$), so the EMA component was unnecessary.

\textbf{Regime-switching $\phi$.} The optimal AR(1) trading coefficient varies by spread regime:

\begin{center}
\begin{tabular}{lcccc}
\toprule
\textbf{Regime} & \textbf{Spread} & \textbf{Frequency} & \textbf{Statistical AR(1)} & \textbf{Trading $\phi$} \\
\midrule
Tight & $\leq 13$ & $\sim 55\%$ & $-0.30$ & $0.05$ \\
Wide  & $= 14$   & $\sim 45\%$ & $\approx 0.00$ & $0.00$ \\
\bottomrule
\end{tabular}
\end{center}

\textbf{Why $\phi = 0.05$ and not $-0.30$?} Using the raw statistical coefficient directly as the trading coefficient gives $FV_t = m_t + 0.30\,\Delta m_t$ — a momentum FV that aggressively extrapolates recent moves. This misaligns passive quotes versus bot crossing thresholds: the FV moves too far, and the bot can no longer profitably cross. Grid search over $\phi \in [-0.3, 0.3]$ found $\phi = 0.05$ as the global optimum for PnL — a small correction that improves quote positioning without disrupting bot-flow capture.

\textit{See Chapter~27 for AR(1) model theory; Chapter~11 for EMA theory; Chapter~29 for regime detection.}

\section{Floor vs.\ Ceil for Integer-Priced Quotes}

Mid-prices in Prosperity are often half-integers (when spread is odd). The passive quotes must be integers. This creates a rounding choice:

\begin{definitionbox}{Rounding Choices}
\begin{align}
  p^B &= \lfloor FV - \delta - \text{skew} \rfloor \quad \text{(floor: conservative bid)} \\
  p^A &= \lceil FV + \delta - \text{skew} \rceil \quad \text{(ceil: aggressive ask)}
\end{align}
\end{definitionbox}

\textbf{The ceil bias}: when $FV$ is a half-integer (e.g.\ $FV = 5002.5$, $\delta = 5$):
\begin{itemize}
  \item $p^B = \lfloor 5002.5 - 5 \rfloor = \lfloor 4997.5 \rfloor = 4997$
  \item $p^A = \lceil 5002.5 + 5 \rceil = \lceil 5007.5 \rceil = 5008$
\end{itemize}
The ask is 0.5 ticks higher than the ``natural'' ask (5007.5 → 5008). This gives $+0.5$ ticks of free edge on the ask side, at zero cost to fill rate (the bot still crosses at 5008). Simulation confirms: using ceil for the ask adds $+0.552$ ticks per timestamp in expectation, compared to floor.

In \texttt{trader.py}: \texttt{math.floor} is used for the bid and \texttt{math.ceil} for the ask — exactly capturing this free edge.

\section{Parameter Sensitivity Analysis}

Understanding \emph{how much} PnL changes per unit change in a parameter guides where to spend optimisation effort. The table below shows TOMATOES PnL sensitivity from the tutorial grid search ($L=80$, both days combined):

\textbf{Sensitivity to quote offset $Q$ (at $a=1.0$, $S=1$):}
\begin{center}
\begin{tabular}{c|cc|c}
\toprule
$Q$ & IS PnL & OOS PnL & $\Delta$PnL vs.\ $Q=5$ \\
\midrule
3 & 14{,}820 & 14{,}200 & $-3{,}600$ \\
4 & 17{,}950 & 17{,}200 & $-640$ \\
\textbf{5} & \textbf{18{,}420} & \textbf{17{,}830} & \textbf{0} \\
6 & 17{,}600 & 16{,}800 & $-820$ \\
7 & 16{,}100 & 15{,}200 & $-2{,}320$ \\
\bottomrule
\end{tabular}
\end{center}
PnL is \textbf{asymmetric}: moving $Q$ below the optimum hurts more than moving above it. At $Q=3$ (2 ticks below), PnL drops by $\sim 20\%$ because we are quoting inside the bot's aggressive-taking threshold and getting adversely selected. At $Q=7$ (2 ticks above), PnL drops $\sim 12\%$ from lower fill rates — the bots still cross, just less often.

\textbf{Sensitivity to AR(1) trading coefficient $\phi$ (at $Q=5$, $S=1$, tight regime):}
\begin{center}
\begin{tabular}{c|cc|c}
\toprule
$\phi$ & IS PnL & OOS PnL & $\Delta$PnL vs.\ $\phi=0.05$ \\
\midrule
$-0.10$ & 17{,}310 & 16{,}500 & $-1{,}110$ \\
$0.00$ & 18{,}100 & 17{,}410 & $-320$ \\
\textbf{0.05} & \textbf{18{,}420} & \textbf{17{,}830} & \textbf{0} \\
$0.10$ & 18{,}200 & 17{,}650 & $-220$ \\
$0.20$ & 17{,}700 & 17{,}050 & $-720$ \\
$-0.30$ & 14{,}200 & 13{,}600 & $-4{,}220$ \\
\bottomrule
\end{tabular}
\end{center}

\textbf{Key insight:} The $\phi$ optimum is shallow on the positive side (small corrections do not hurt much) but steep on the negative side (the raw statistical $\phi = -0.30$ loses $\sim 23\%$ of PnL). This asymmetry confirms the "why not $-0.30$" argument: momentum FV pushes quotes outside bot crossing thresholds. In contrast, $Q$ is the higher-leverage parameter: a 2-tick error in $Q$ costs more than a 0.15 error in $\phi$.

\begin{warningbox}{These Numbers Are Illustrative}
The exact PnL figures depend on bot configuration, fill-matching mode, and round. The \emph{pattern} (Q more sensitive than $\phi$; negative $\phi$ catastrophic) is structurally robust. Always re-run the grid search on new round data.
\end{warningbox}

\begin{takeawaybox}{Chapter 32}
\begin{itemize}
  \item Constant FV: $\mu$ estimated from historical mean. Valid when strongly stationary (EMERALDS).
  \item EMA FV: tracks slow drift with exponential smoothing. With $\alpha=1$: collapses to current mid.
  \item AR(1) correction: $FV_t = m_t - \phi\Delta m_t$. The 1-step-ahead forecast corrects the current price for expected reversion. Use grid-searched $\phi$, not the raw statistical AR(1).
  \item Combined (TOMATOES): $\alpha=1$ means EMA plays no role; pure AR(1) correction with regime-switching $\phi$.
  \item Floor/ceil rounding: use floor for bid, ceil for ask. Captures $+0.55$ ticks/timestamp free edge from half-integer mid-prices.
  \item Parameter sensitivity: $Q$ is the higher-leverage knob ($\pm 2$ ticks $\Rightarrow \pm 20\%$ PnL); $\phi \in [0, 0.15]$ is a shallow optimum; negative $\phi$ is catastrophic.
\end{itemize}
\end{takeawaybox}

\begin{keyconceptbox}{Decision Guide: Choosing a Fair Value Model}
\textbf{Step 1 — Is the mid-price strongly stationary?}
\begin{itemize}
  \item Yes (ADF $p < 0.01$, mean-reverts quickly) $\Rightarrow$ \textbf{Constant FV}: $FV = \bar{m}$. Example: EMERALDS.
  \item No $\Rightarrow$ proceed to Step 2.
\end{itemize}
\textbf{Step 2 — Do mid-price \emph{changes} exhibit autocorrelation?}
\begin{itemize}
  \item ACF($\Delta m$) $\neq 0$ at lag 1 (Ljung-Box $p < 0.05$) $\Rightarrow$ \textbf{AR(1) correction}: $FV_t = m_t - \phi\Delta m_t$ with grid-searched $\phi$.
  \item No autocorrelation $\Rightarrow$ \textbf{EMA only}: $FV_t = \text{EMA}_t$ (tune $\alpha$ by cross-validation).
\end{itemize}
\textbf{Step 3 — Is there slow intraday drift \emph{on top of} mean-reversion?}
\begin{itemize}
  \item Yes (mean $\neq 0$ over 30-min windows) $\Rightarrow$ \textbf{EMA + AR(1)}: $FV_t = \text{EMA}_t - \phi\Delta m_t$.
  \item No $\Rightarrow$ AR(1) alone suffices (EMA with $\alpha=1$ collapses to current mid; the model is identical).
\end{itemize}
\textbf{Step 4 — Are there detectable regimes?}
\begin{itemize}
  \item If spread or volatility clusters into 2--3 stable regimes $\Rightarrow$ \textbf{regime-switching $\phi$}: separate grid-search per regime. See Ch.~29.
  \item Otherwise: a single global $\phi$ is sufficient.
\end{itemize}
\end{keyconceptbox}

% ──────────────────────────────────────────────────────────────────────────────

\chapter{State Serialisation and Persistent Strategies}

\section{The Stateless Environment Problem}

In AWS Lambda (the execution environment for Prosperity submissions), \emph{there is no guarantee that global or class variables persist between calls}. Each call to \texttt{run(state)} may execute in a fresh process with no memory of previous iterations.

This is a critical constraint for strategies that need historical data (e.g.\ the EMA needs $\text{EMA}_{t-1}$, the AR(1) correction needs $m_{t-1}$). Without persistence, you cannot compute these values.

\section{The \texttt{traderData} Solution}

Prosperity provides the \texttt{traderData} field as a serialised string that is:
\begin{enumerate}[label=\arabic*.]
  \item Returned by \texttt{run()} as a string
  \item Passed back into the next call as \texttt{state.traderData}
\end{enumerate}

This creates a single-variable persistent store. The workflow:

\begin{lstlisting}[caption={State serialisation pattern in Trader.run()}]
import jsonpickle

def run(self, state: TradingState):
    # Deserialise state from previous iteration
    saved = {}
    if state.traderData:
        try:
            saved = jsonpickle.decode(state.traderData)
        except Exception:
            saved = {}

    # Restore per-strategy state
    tomatoes_state = saved.get("tomatoes", {})
    tomatoes.set_state(tomatoes_state)

    # ... run strategies ...

    # Serialise state for next iteration
    new_saved = {"tomatoes": tomatoes.get_state()}
    trader_data = jsonpickle.encode(new_saved)
    return result, conversions, trader_data
\end{lstlisting}

\section{What to Persist: Minimal State Principle}

\begin{keyconceptbox}{Minimal State Principle}
Persist only the variables that \emph{cannot be recomputed from the current \texttt{TradingState}} and that are needed for the strategy's next decision.

\textbf{In TomatoesStrategy:}
\begin{itemize}
  \item \texttt{\_last\_mid}: needed to compute $\Delta m_t = m_t - m_{t-1}$. Cannot be derived from current state. \textbf{Persist this.}
  \item \texttt{\_ema}: with $\alpha = 1$, the EMA is always overwritten to $m_t$ on the very next tick. The seed value is irrelevant. \textbf{No need to persist.}
\end{itemize}
\end{keyconceptbox}

\begin{lstlisting}[caption={Minimal state in TomatoesStrategy}]
def get_state(self) -> dict:
    return {"last_mid": self._last_mid}   # only _last_mid

def set_state(self, d: dict) -> None:
    self._last_mid = d.get("last_mid")    # restore _last_mid
\end{lstlisting}

\section{jsonpickle: Serialising Arbitrary Python Objects}

\texttt{jsonpickle} is a Python library that serialises arbitrary Python objects (including class instances) to JSON-compatible strings:

\begin{lstlisting}[caption={jsonpickle basics}]
import jsonpickle

# Serialise
data = {"last_mid": 5002.5, "ema": 5001.8}
encoded = jsonpickle.encode(data)
# encoded = '{"last_mid": 5002.5, "ema": 5001.8}'

# Deserialise
decoded = jsonpickle.decode(encoded)
# decoded = {"last_mid": 5002.5, "ema": 5001.8}
\end{lstlisting}

\textbf{Size constraint}: Prosperity truncates \texttt{traderData} at 50{,}000 characters. Keep state minimal:
\begin{center}
\begin{tabular}{lcc}
\toprule
\textbf{What to store} & \textbf{Approximate size} & \textbf{Recommended?} \\
\midrule
1 float (last\_mid) & $\sim$25 chars & Yes \\
50-element list & $\sim$500 chars & Yes \\
1000-element price history & $\sim$10{,}000 chars & Caution \\
Full order book history & $>$50{,}000 chars & Never \\
\bottomrule
\end{tabular}
\end{center}

\section{State Management Patterns}

\subsection{Rolling Buffer Pattern}

For indicators that need the last $n$ observations (e.g.\ a 20-period SMA needs 20 prices):

\begin{lstlisting}[caption={Rolling buffer pattern}]
def get_state(self) -> dict:
    return {"price_history": self._prices[-20:]}  # only last 20

def set_state(self, d: dict) -> None:
    self._prices = d.get("price_history", [])

def orders(self, state):
    mid = ...
    self._prices.append(mid)
    if len(self._prices) < 20:
        return []  # not enough data yet
    sma = sum(self._prices[-20:]) / 20
    ...
\end{lstlisting}

\subsection{Warm-Up Period}

For EMA with $\alpha < 1$: the first $\sim 3/\alpha$ iterations are needed to wash out the seed bias. During warm-up, post conservative quotes or no quotes at all:

\begin{lstlisting}[caption={Warm-up handling}]
WARMUP = 30  # skip first 30 iterations

def orders(self, state):
    self._count = getattr(self, "_count", 0) + 1
    if self._count < WARMUP:
        return []  # no quotes during warm-up
    # Normal logic after warm-up
    ...
\end{lstlisting}

\begin{takeawaybox}{Chapter 33}
\begin{itemize}
  \item AWS Lambda is stateless: class variables may not persist. Use \texttt{traderData} as the persistence layer.
  \item Pattern: \texttt{jsonpickle.decode(state.traderData)} at start; \texttt{jsonpickle.encode(new\_state)} at end.
  \item Persist only what cannot be recomputed from the current state (e.g.\ \texttt{last\_mid}, rolling buffers).
  \item Size limit: 50{,}000 characters. Never store full price histories; use fixed-size rolling buffers.
  \item Warm-up: for EMA with $\alpha < 1$, the first $\sim 3/\alpha$ iterations have biased estimates; suppress orders during warm-up.
\end{itemize}
\end{takeawaybox}

% ──────────────────────────────────────────────────────────────────────────────

\chapter{Summary and Connections}

\section{From First Principles to \texttt{trader.py}}

This encyclopedia has built a complete bridge from market microstructure foundations to the exact implementation decisions in the Prosperity codebase. Here is the full chain of reasoning for each design choice:

\begin{center}
\begin{longtable}{p{4cm}p{5.5cm}p{4.5cm}}
\toprule
\textbf{Decision in \texttt{trader.py}} & \textbf{Theoretical Justification} & \textbf{Where covered} \\
\midrule
FV = 10000 (EMERALDS) & Mid-price is i.i.d.\ around 10000 (ADF $\ll 0.001$, Hurst $H \approx 0.25$) & Chs.\ 6, 27, 28 \\
Passive quotes at $\pm 7$ from FV & Optimal $\delta^* = $ bot half-spread $- 1$ (fill rate vs.\ edge trade-off) & Ch.\ 15 \\
Take orders $\leq FV$ & Immediately profitable edge (VWAP below FV) & Ch.\ 4, 15 \\
EMA $\alpha = 1$ & With AR(1) correction, EMA is redundant; simplest valid choice & Chs.\ 11, 32 \\
AR(1) $\phi = 0.05$ (tight) & Grid-searched optimum; raw statistical $\phi = -0.30$ over-corrects relative to bot thresholds & Chs.\ 27, 32 \\
AR(1) $\phi = 0.00$ (wide) & No predictability in wide regime (statistical AR(1) $\approx 0$) & Chs.\ 27, 29 \\
Skew range $S_{\text{range}} = 1$ & Grid-optimised: 1 tick skew cost $\approx 2\%$ of edge at limit & Chs.\ 23, 30 \\
OBI sizing rejected & $\delta = 5 \gg \beta\sigma = 0.51$; passive edge dominates OBI signal by $10\times$ & Chs.\ 14, 31 \\
Floor bid, ceil ask & Half-integer mids: ceil ask gives $+0.55$ ticks/ts free edge & Ch.\ 32 \\
Persist \texttt{last\_mid} only & Only value not in current \texttt{TradingState}; EMA re-initialises with $\alpha=1$ & Ch.\ 33 \\
Position limit $L = 80$ & Exchange-enforced; also bounds maximum inventory loss to $L \cdot \sigma = 80 \times 0.72 \approx 58$ ticks & Chs.\ 4, 23 \\
\bottomrule
\end{longtable}
\end{center}

\section{The Mental Model: A Trading Engineer's Checklist}

\begin{keyconceptbox}{Trading Strategy Development Checklist}
\begin{enumerate}[label=\arabic*.]
  \item \textbf{EDA first}: Compute ACF, PACF, ADF, Hurst. Identify regime. Only then choose strategy family.
  \item \textbf{Choose FV model}: Constant (strongly stationary), EMA (slow drift), AR(1)+EMA (mean-reverting returns).
  \item \textbf{Choose regime model}: Single-regime vs.\ TAR (threshold) vs.\ HMM (latent state).
  \item \textbf{Set quote offset $\delta$}: Inside the bot spread; grid-search on in-sample data.
  \item \textbf{Set inventory skew}: Start at $S_{\text{range}} = 1$; increase only if inventory buildup is severe.
  \item \textbf{Evaluate signals}: Test each signal (OBI, volume, etc.) with the $\delta$ vs.\ $\beta\sigma$ comparison. Only include if signal exceeds passive edge.
  \item \textbf{Walk-forward validate}: Grid-search on days $-2$; evaluate on day $-1$. Report OOS Sharpe only.
  \item \textbf{Check state persistence}: Identify what must be serialised. Keep \texttt{traderData} minimal.
  \item \textbf{Size by Kelly / vol-target}: Use position limit as hard cap.
  \item \textbf{Monitor in production}: Track OOS Sharpe, max drawdown, fill rate. Re-optimise if regime shifts.
\end{enumerate}
\end{keyconceptbox}

\begin{takeawaybox}{Part VII Summary}
\begin{itemize}
  \item \textbf{Time series}: Stationarity is prerequisite for all inference. AR(1) with $|\phi|<1$ is mean-reverting; unit root = random walk.
  \item \textbf{Hurst}: $H<0.5$ = anti-persistent (mean-revert), $H>0.5$ = persistent (trend-follow).
  \item \textbf{Regime switching}: TAR (observable threshold), Markov switching (latent state), HMM (general).
  \item \textbf{Avellaneda-Stoikov}: Reservation price $r = S - q\gamma\sigma^2(T-t)$; optimal spread $\delta^* = \frac{\gamma\sigma^2(T-t)}{2} + \frac{1}{\kappa}\ln(1+\gamma/\kappa)$.
  \item \textbf{OBI}: Predicts short-term returns; only actionable when $\beta\sigma > \delta$ (signal $>$ passive edge).
  \item \textbf{Fair value}: Choose method by stationarity. Use floor/ceil to capture half-integer mid-price free edge.
  \item \textbf{State serialisation}: jsonpickle + traderData; persist minimal state; respect 50K char limit.
\end{itemize}
\end{takeawaybox}

% End of Part VII

% ============================================================
%
%  PART VIII — DERIVATIVES, ARBITRAGE & ROUND-SPECIFIC MECHANICS
%
% ============================================================
\part{Derivatives, Arbitrage \& Round-Specific Mechanics}

This part covers the knowledge gaps between the Tutorial Round and Rounds~1–5 of Prosperity~4. Each chapter introduces a new strategy family that is likely to appear in later rounds: options pricing, cross-product arbitrage, conversion mechanics, exogenous signal trading, and sealed-bid auction theory. Every topic is built from first principles with ML analogies throughout.

% ──────────────────────────────────────────────────────────────────────────

\chapter{Options and Derivatives}

An \emph{option} is the simplest derivative: its value at expiry is a deterministic function of an underlying asset price. In Prosperity~3 this appeared as \emph{vouchers} — European call options on a volcanic rock. Understanding fair pricing lets you trade against bots that misprice these instruments.

\section{What is a Derivative?}

\begin{definitionbox}{Derivative Contract}
A \textbf{derivative} is a financial contract whose value at time $t$ is a function $V(S_t, t)$ of the price $S_t$ of some \textbf{underlying asset}. The derivative has no value of its own — it derives all value from $S$.
\end{definitionbox}

\textbf{ML analogy:} A derivative is a \emph{feature transformation} of the raw input $S$. Just as $x^2$ or $\log x$ are deterministic functions of a feature $x$, a call option's payoff $\max(S-K, 0)$ is a deterministic function of $S$ at expiry.

Common examples:
\begin{itemize}
  \item \textbf{Forward:} Obligation to buy $S$ at price $K$ at time $T$. Payoff $= S_T - K$.
  \item \textbf{Call option:} Right (not obligation) to buy $S$ at price $K$ at time $T$. Payoff $= \max(S_T - K, 0)$.
  \item \textbf{Put option:} Right to sell $S$ at price $K$ at time $T$. Payoff $= \max(K - S_T, 0)$.
\end{itemize}

\section{European Call and Put Options}

\begin{definitionbox}{Strike Price and Intrinsic Value}
The \textbf{strike price} $K$ is the agreed purchase price fixed at contract inception.
The \textbf{intrinsic value} of a call at time $t$ is:
\[
  IV_{\text{call}}(t) = \max(S_t - K,\; 0)
\]
An option is \textbf{in-the-money} (ITM) when $S_t > K$, \textbf{at-the-money} (ATM) when $S_t = K$, and \textbf{out-of-the-money} (OTM) when $S_t < K$.
\end{definitionbox}

The market price of an option \emph{before} expiry always exceeds its intrinsic value by the \textbf{time value}:
\[
  C(S_t, t) = \underbrace{\max(S_t - K, 0)}_{\text{intrinsic value}} + \underbrace{TV(S_t, t)}_{\text{time value}}
\]
Time value decays to zero as $t \to T$. This decay is called \textbf{theta} ($\Theta$).

\begin{warningbox}{European vs.\ American Options}
A \textbf{European} option can only be exercised at expiry $T$. An \textbf{American} option can be exercised at any time $t \leq T$. Prosperity vouchers are European. Never confuse the two: American options are worth at least as much as European, but require more complex pricing (Least-Squares Monte Carlo or binomial trees).
\end{warningbox}

\section{Geometric Brownian Motion}

To price an option, we need a model for how $S_t$ evolves. The standard model is Geometric Brownian Motion (GBM).

\begin{definitionbox}{Geometric Brownian Motion (GBM)}
$S_t$ follows GBM if:
\[
  \dd S_t = \mu S_t \dd t + \sigma S_t \dd B_t
\]
where $\mu$ is the drift (expected return), $\sigma > 0$ is the volatility, and $B_t$ is a standard Brownian motion (continuous-time random walk with $B_t - B_s \sim \mathcal{N}(0, t-s)$).
\end{definitionbox}

Applying Itô's lemma to $\ln S_t$ gives the exact solution:
\[
  S_T = S_t \cdot \exp\!\left[\left(\mu - \tfrac{\sigma^2}{2}\right)(T-t) + \sigma\sqrt{T-t}\; Z\right], \quad Z \sim \mathcal{N}(0,1)
\]
This means $\ln(S_T / S_t) \sim \mathcal{N}\!\left((\mu - \sigma^2/2)(T-t),\; \sigma^2(T-t)\right)$: log-returns are normally distributed. GBM prevents $S_t$ from going negative (unlike arithmetic Brownian motion), which is realistic for asset prices.

\textbf{ML analogy:} GBM is a \emph{continuous-time autoregressive model}. The $\sigma\sqrt{\Delta t}\; Z_t$ term is Gaussian noise at each step. The $\mu$ drift is the mean; $\sigma$ is the noise scale. The AR(1) we used for TOMATOES is the discrete-time analogue.

\section{The Black-Scholes Formula}

Black, Scholes, and Merton (1973) derived the fair price of a European call by constructing a \emph{delta-hedged} portfolio that is instantaneously riskless and must therefore earn the risk-free rate $r$.

\begin{definitionbox}{Black-Scholes Call Price}
Under GBM, the fair price at time $t$ of a European call option with strike $K$, expiry $T$, and risk-free rate $r$ is:
\[
  C(S, t) = S\,\Phi(d_1) - K e^{-r(T-t)}\,\Phi(d_2)
\]
where $\Phi$ is the standard normal CDF and:
\begin{align}
  d_1 &= \frac{\ln(S/K) + \left(r + \frac{\sigma^2}{2}\right)(T-t)}{\sigma\sqrt{T-t}} \\[4pt]
  d_2 &= d_1 - \sigma\sqrt{T-t}
\end{align}
The corresponding put price follows from \textbf{put-call parity}, which holds by the following no-arbitrage argument:

\begin{proof}[Derivation: Put-Call Parity]
Consider two portfolios at time $t$, both held to expiry $T$:
\begin{itemize}
  \item \textbf{Portfolio A:} Long one call $C$ + long bond paying $K e^{-r(T-t)}$ at expiry.
  \item \textbf{Portfolio B:} Long one put $P$ + long one share $S$.
\end{itemize}
At expiry, both portfolios pay $\max(S_T, K)$ regardless of $S_T$: if $S_T > K$, Portfolio A pays $(S_T - K) + K = S_T$ and Portfolio B pays $0 + S_T = S_T$; if $S_T \leq K$, Portfolio A pays $0 + K = K$ and Portfolio B pays $(K - S_T) + S_T = K$. Since identical payoffs must have identical prices today (no arbitrage):
\[
  C + K e^{-r(T-t)} = P + S \quad \Longrightarrow \quad P = C - S + K e^{-r(T-t)}
\]
\end{proof}

Substituting the Black-Scholes call formula into the parity relation and using $\Phi(-x) = 1 - \Phi(x)$ recovers:
\[
  P(S, t) = K e^{-r(T-t)}\,\Phi(-d_2) - S\,\Phi(-d_1)
\]
\end{definitionbox}

\begin{examplebox}{Intuition for the Black-Scholes Formula}
The term $S\,\Phi(d_1)$ is the risk-adjusted expected value of receiving the stock if the option expires ITM. The term $K e^{-r(T-t)}\,\Phi(d_2)$ is the present value of paying the strike in that event. $\Phi(d_2) = \Pr(S_T > K)$ under the risk-neutral measure: the probability of the option expiring in the money.

As a sanity check:
\begin{itemize}
  \item Deep ITM ($S \gg K$): $\Phi(d_1), \Phi(d_2) \to 1$, so $C \approx S - K e^{-r(T-t)}$ (intrinsic value plus interest savings).
  \item Deep OTM ($S \ll K$): $\Phi(d_1), \Phi(d_2) \to 0$, so $C \approx 0$.
  \item At expiry ($T - t \to 0$): $C \to \max(S - K, 0)$ (pure intrinsic value).
\end{itemize}
\end{examplebox}

\subsection{Implementing Black-Scholes in Python}

\begin{lstlisting}
import math
from statistics import NormalDist

_norm = NormalDist()

def black_scholes_call(S: float, K: float, T: float,
                       r: float, sigma: float) -> float:
    """European call price via Black-Scholes."""
    if T <= 0:
        return max(S - K, 0.0)
    sqrt_T = math.sqrt(T)
    d1 = (math.log(S / K) + (r + 0.5 * sigma**2) * T) / (sigma * sqrt_T)
    d2 = d1 - sigma * sqrt_T
    return S * _norm.cdf(d1) - K * math.exp(-r * T) * _norm.cdf(d2)

def black_scholes_put(S: float, K: float, T: float,
                      r: float, sigma: float) -> float:
    """European put price via Black-Scholes (put-call parity)."""
    call = black_scholes_call(S, K, T, r, sigma)
    return call - S + K * math.exp(-r * T)   # put-call parity
\end{lstlisting}

\begin{warningbox}{In Prosperity, $r \approx 0$}
There is no risk-free rate in the competition sandbox. Set $r = 0$ in all Black-Scholes calculations. The formula simplifies to:
\[
  d_1 = \frac{\ln(S/K) + \frac{\sigma^2}{2}T}{\sigma\sqrt{T}}, \qquad d_2 = d_1 - \sigma\sqrt{T}, \qquad C = S\,\Phi(d_1) - K\,\Phi(d_2)
\]
\end{warningbox}

\section{Why the Drift $\mu$ Disappears: Risk-Neutral Pricing}

The Black-Scholes formula contains $r$ (risk-free rate) but \emph{not} $\mu$ (the stock's actual drift). This is the central mystery of derivatives pricing: how can the fair price of an option be independent of whether the stock is expected to grow at 2\% or 20\% per year?

\subsection{The Delta-Hedging Argument}

The original Black-Scholes derivation eliminates $\mu$ via a no-arbitrage argument. Consider a portfolio $\Pi$ consisting of one long call $C(S,t)$ and a short position of $\Delta = \partial C/\partial S$ units of the stock:
\[
  \Pi = C - \Delta\, S
\]
By Itô's lemma, the stochastic differential of $C(S,t)$ under GBM is:
\[
  \dd C = \frac{\partial C}{\partial t}\dd t + \frac{\partial C}{\partial S}\dd S + \frac{1}{2}\frac{\partial^2 C}{\partial S^2}(\dd S)^2
\]
Since $(\dd S)^2 = \sigma^2 S^2 \dd t$ (the quadratic variation of GBM), we get:
\[
  \dd C = \left(\frac{\partial C}{\partial t} + \frac{1}{2}\sigma^2 S^2 \frac{\partial^2 C}{\partial S^2}\right)\dd t + \frac{\partial C}{\partial S}\dd S
\]
The portfolio $\Pi = C - \Delta S$ then has:
\[
  \dd\Pi = \dd C - \Delta\,\dd S = \left(\frac{\partial C}{\partial t} + \frac{1}{2}\sigma^2 S^2 \frac{\partial^2 C}{\partial S^2}\right)\dd t
\]
The $\dd S$ terms (containing $\dd B_t$ and the drift $\mu$) \emph{cancel exactly}. The portfolio is \textbf{instantaneously riskless}. Under no-arbitrage, a riskless portfolio must earn the risk-free rate:
\[
  \dd\Pi = r\,\Pi\,\dd t = r\,(C - \Delta S)\,\dd t
\]
Substituting $\Delta = \partial C/\partial S$ and equating both expressions for $\dd\Pi$ yields the \textbf{Black-Scholes PDE}:

\begin{definitionbox}{The Black-Scholes PDE}
\[
  \frac{\partial C}{\partial t} + \frac{1}{2}\sigma^2 S^2 \frac{\partial^2 C}{\partial S^2} + r S \frac{\partial C}{\partial S} - rC = 0
\]
with boundary condition $C(S, T) = \max(S - K, 0)$.

This is a parabolic PDE. The parameter $\mu$ does not appear — it was eliminated by the delta-hedging argument. The solution to this PDE with the call payoff boundary condition \emph{is} the Black-Scholes formula.
\end{definitionbox}

\subsection{The Risk-Neutral Measure (Q-Measure)}

The deeper theoretical explanation uses \textbf{measure theory}. Under the physical measure $\mathbb{P}$, the stock follows:
\[
  \dd S_t = \mu S_t \dd t + \sigma S_t \dd B_t^{\mathbb{P}}
\]
\textbf{Girsanov's theorem} states that there exists an equivalent probability measure $\mathbb{Q}$ (called the \emph{risk-neutral measure}) under which:
\[
  \dd S_t = r S_t \dd t + \sigma S_t \dd B_t^{\mathbb{Q}}
\]
where $B_t^{\mathbb{Q}}$ is a Brownian motion under $\mathbb{Q}$. The change of measure is effected by the \textbf{Radon-Nikodym derivative}:
\[
  \frac{\dd\mathbb{Q}}{\dd\mathbb{P}} = \exp\!\left(-\frac{\mu - r}{\sigma}B_T^{\mathbb{P}} - \frac{1}{2}\left(\frac{\mu-r}{\sigma}\right)^2 T\right)
\]
The quantity $\lambda = \nicefrac{(\mu - r)}{\sigma}$ is called the \textbf{market price of risk}: how much excess return (over $r$) investors demand per unit of volatility.

\begin{definitionbox}{Risk-Neutral Pricing Theorem}
Under the risk-neutral measure $\mathbb{Q}$, the fair price of any derivative with payoff $g(S_T)$ at time $T$ is:
\[
  V(S, t) = e^{-r(T-t)}\,\E^{\mathbb{Q}}\!\left[g(S_T) \,\big|\, S_t = S\right]
\]
The expectation is taken under $\mathbb{Q}$ (where $S$ grows at $r$, not $\mu$). Since the payoff function $g$ does not depend on $\mu$, neither does $V$.
\end{definitionbox}

\textbf{Applying to the European call:} Under $\mathbb{Q}$, $S_T = S_t \exp\!\left[(r - \sigma^2/2)(T-t) + \sigma\sqrt{T-t}\,Z\right]$ with $Z \sim \mathcal{N}(0,1)$. Computing the expectation:
\[
  C = e^{-r(T-t)}\,\E^{\mathbb{Q}}[\max(S_T - K, 0)] = S\Phi(d_1) - Ke^{-r(T-t)}\Phi(d_2)
\]
This is the Black-Scholes formula, derived from first principles. Note: $\Phi(d_2) = \mathbb{Q}(S_T > K)$ is the risk-neutral probability of expiring ITM — not the real-world probability.

\begin{warningbox}{Real-World vs.\ Risk-Neutral Probability}
$\mathbb{Q}(S_T > K) = \Phi(d_2)$ uses the drift $r$ (risk-free rate). The \emph{actual} probability under $\mathbb{P}$ would use drift $\mu$ instead of $r$ in $d_1, d_2$. Since $\mu > r$ typically, the real-world probability of expiring ITM exceeds the risk-neutral one. This does \emph{not} mean Black-Scholes misprices the option — the two probabilities are under different measures and cannot be directly compared.
\end{warningbox}

\textbf{ML analogy:} The change of measure $\mathbb{P} \to \mathbb{Q}$ is exactly a \emph{likelihood ratio reweighting} of scenarios. You are reweighting the distribution of outcomes so that the stock grows at $r$ instead of $\mu$. It is the continuous-time analogue of importance sampling in Monte Carlo: instead of simulating under the true distribution and weighting by the Radon-Nikodym derivative, you simulate directly under $\mathbb{Q}$.

\begin{keyconceptbox}{The Core Insight of Risk-Neutral Pricing}
\begin{enumerate}[label=\arabic*.]
  \item \textbf{Delta-hedging eliminates first-order directional risk instantaneously:} The portfolio $\Pi = C - \Delta S$ satisfies $\partial\Pi/\partial S = 0$ at the current $S$, so $\mu$ drops out of the instantaneous P\&L. \emph{Residual risks remain}: (a) \textbf{gamma risk} $\frac{1}{2}\Gamma(\Delta S)^2$ from the quadratic term (non-zero whenever $\Gamma \neq 0$), (b) \textbf{vega risk} $\mathcal{V}\Delta\sigma$ from stochastic volatility (BS assumes $\sigma$ constant, which is false), and (c) \textbf{discrete rebalancing gap}: continuous hedging is not achievable in practice; each rebalancing interval leaves a residual $\propto \Gamma(\Delta S)^2$. In Prosperity, discrete rebalancing at every tick means gamma and vega exposures are real but small relative to the tick-level edge.
  \item \textbf{The risk-neutral measure absorbs the risk premium:} All assets grow at $r$ under $\mathbb{Q}$; the actual drift $\mu$ is embedded in the measure change, not in the pricing formula.
  \item \textbf{Practical implication for Prosperity:} With $r = 0$, the risk-neutral drift is zero. All vouchers price as if the underlying is a martingale (no drift). Your estimate of whether VOLCANIC ROCK will go up or down is irrelevant to fair value; only $\sigma$ matters.
\end{enumerate}
\end{keyconceptbox}

\section{The Greeks — Sensitivity Analysis}

The \textbf{Greeks} measure how the option price changes as each input changes. They are essential for hedging and for spotting mispricings in competition.

\begin{definitionbox}{The Option Greeks}
\begin{align}
  \Delta &= \frac{\partial C}{\partial S} = \Phi(d_1) \quad \text{(sensitivity to underlying price)} \\[4pt]
  \Gamma &= \frac{\partial^2 C}{\partial S^2} = \frac{\phi(d_1)}{S\sigma\sqrt{T-t}} \quad \text{(convexity; rate of change of $\Delta$)} \\[4pt]
  \Theta &= \frac{\partial C}{\partial t} = -\frac{S\phi(d_1)\sigma}{2\sqrt{T-t}} - rK e^{-r(T-t)}\Phi(d_2) \quad \text{(time decay)} \\[4pt]
  \mathcal{V} &= \frac{\partial C}{\partial \sigma} = S\phi(d_1)\sqrt{T-t} \quad \text{(sensitivity to volatility; ``vega'')} \\[4pt]
  \rho &= \frac{\partial C}{\partial r} = K(T-t)e^{-r(T-t)}\Phi(d_2) \quad \text{(sensitivity to interest rate)}
\end{align}
where $\phi(x) = \frac{1}{\sqrt{2\pi}}e^{-x^2/2}$ is the standard normal PDF.
\end{definitionbox}

\textbf{ML analogy:} The Greeks are exactly \emph{partial derivatives} (gradients) of the option price function with respect to each input — identical in concept to the gradient of a loss function with respect to model parameters. $\Delta$ is the first-order term in a Taylor expansion of $C$; $\Gamma$ is the second-order term.

\begin{keyconceptbox}{Delta as a Hedge Ratio}
If you hold one long call and want to be \emph{delta-neutral} (insensitive to small moves in $S$), you must short $\Delta = \Phi(d_1)$ units of the underlying per option. Since $0 \leq \Delta \leq 1$ for a call, a position of $N$ calls requires shorting between 0 and $N$ units of the underlying. This \emph{delta hedging} is the core insight that allows Black-Scholes to derive a unique price from no-arbitrage.
\end{keyconceptbox}

\section{Implied Volatility}

The Black-Scholes formula is invertible in $\sigma$: given the \emph{market price} of an option, we can solve for the volatility $\hat{\sigma}$ that makes the model price match the market price.

\begin{definitionbox}{Implied Volatility (IV)}
The \textbf{implied volatility} $\hat{\sigma}$ of an option trading at market price $C^*$ is:
\[
  \hat{\sigma} = \text{BS}^{-1}(C^*;\; S, K, T, r)
\]
i.e.\ the unique $\sigma$ satisfying $\text{BS}(S, K, T, r, \hat{\sigma}) = C^*$.
Since there is no closed form, it is found numerically (Newton-Raphson using vega $\mathcal{V}$).
\end{definitionbox}

\textbf{Trading use:} If the bot's option price implies $\hat{\sigma} = 0.30$ but your estimate of realized volatility is $\sigma = 0.20$, the option is overpriced — sell it (and delta-hedge if needed). If $\hat{\sigma} < \sigma$, the option is underpriced — buy it.

\begin{lstlisting}
def implied_vol(C_market: float, S: float, K: float,
                T: float, r: float = 0.0,
                tol: float = 1e-6, max_iter: int = 50) -> float:
    """Newton-Raphson implied volatility solver."""
    sigma = 0.3   # initial guess
    for _ in range(max_iter):
        price = black_scholes_call(S, K, T, r, sigma)
        sqrt_T = math.sqrt(T)
        d1 = (math.log(S / K) + (r + 0.5*sigma**2)*T) / (sigma*sqrt_T)
        vega = S * NormalDist().pdf(d1) * sqrt_T  # dC/d_sigma
        diff = price - C_market
        if abs(diff) < tol:
            break
        sigma -= diff / (vega + 1e-10)
        sigma = max(sigma, 1e-6)   # keep positive
    return sigma
\end{lstlisting}

\section{Vouchers in Prosperity — Practical Application}

In Prosperity~3, vouchers were European call options on a volcanic rock underlying (ticker \texttt{VOLCANIC\_ROCK}) with strikes 9500, 9750, 10000, 10250, 10500. The round ran 3 days; each day, $T$ decreased by $1/3$.

\begin{keyconceptbox}{Voucher Trading Strategy}
\begin{enumerate}[label=\arabic*.]
  \item \textbf{Estimate $\sigma$:} Compute the realized standard deviation of daily log-returns of the underlying from recent prices.
  \item \textbf{Compute fair value:} $C_{\text{fair}} = \text{BS}(S, K, T, r=0, \hat{\sigma})$.
  \item \textbf{Compare to market:} If $C_{\text{market}} < C_{\text{fair}} - \epsilon$, \textbf{buy} the voucher (underpriced). If $C_{\text{market}} > C_{\text{fair}} + \epsilon$, \textbf{sell} (overpriced).
  \item \textbf{Delta hedge:} To isolate the vol mispricing from directional risk, simultaneously trade $-\Delta$ units of the underlying per option.
  \item \textbf{Monitor IV:} Track implied vol across strikes. If ATM IV differs significantly from wing IVs (the ``volatility smile''), more sophisticated arbitrage is possible.
\end{enumerate}
\end{keyconceptbox}

% ── Multi-Leg Option Strategies ─────────────────────────────────────────
\section{Multi-Leg Option Strategies and Spread Arbitrage}

A single option has directional (delta) and volatility (vega) exposure simultaneously, making it hard to isolate a specific edge. \textbf{Multi-leg strategies} combine options to cancel unwanted exposures and target a specific risk premium.

\subsection{Vertical Spreads}

\begin{definitionbox}{Bull Call Spread}
Long one call at strike $K_1$, short one call at $K_2 > K_1$ on the same underlying and expiry:
\[
  \Pi_T = \max(S_T - K_1, 0) - \max(S_T - K_2, 0)
\]
\textbf{P\&L bounds:} $0 \leq \Pi_T \leq K_2 - K_1$.
The net cost is $C(K_1) - C(K_2) > 0$ (long call is always more expensive).
\textbf{Maximum profit:} $(K_2 - K_1) - [C(K_1) - C(K_2)]$ when $S_T \geq K_2$.
\textbf{Break-even:} $S_T^* = K_1 + C(K_1) - C(K_2)$.
\end{definitionbox}

The \textbf{bear put spread} is the mirror: long $P(K_2)$, short $P(K_1)$, profits when $S_T \leq K_1$.

\textbf{Greek profile of a bull call spread:}
\begin{center}
\begin{tabular}{lccc}
\toprule
\textbf{Greek} & \textbf{Long call $K_1$} & \textbf{Short call $K_2$} & \textbf{Net spread} \\
\midrule
Delta ($\Delta$) & $\Phi(d_1^{K_1}) > 0$ & $-\Phi(d_1^{K_2}) < 0$ & $\Phi(d_1^{K_1}) - \Phi(d_1^{K_2}) \in (0,1)$ \\
Vega ($\mathcal{V}$) & $S\phi(d_1^{K_1})\sqrt{\tau} > 0$ & $-S\phi(d_1^{K_2})\sqrt{\tau} < 0$ & Small (partially cancels) \\
Gamma ($\Gamma$) & $\phi(d_1^{K_1})/(S\sigma\sqrt{\tau}) > 0$ & $-\phi(d_1^{K_2})/(S\sigma\sqrt{\tau}) < 0$ & Small (partially cancels) \\
\bottomrule
\end{tabular}
\end{center}
The spread is primarily a \emph{directional bet with bounded risk}, much cheaper than an outright call but capped upside. Vega and Gamma are smaller than a single leg — the spread is less sensitive to volatility, which is useful when you have a view on direction but not on volatility.

\subsection{Butterfly Spread}

\begin{definitionbox}{Butterfly Spread}
Long one call at $K_1$, short two calls at $K_2$, long one call at $K_3$, with $K_1 < K_2 < K_3$ and $K_2 - K_1 = K_3 - K_2 = w$ (equal wings):
\[
  \Pi_T = \max(S_T - K_1, 0) - 2\max(S_T - K_2, 0) + \max(S_T - K_3, 0)
\]
\textbf{Maximum profit:} $w - \text{net debit}$ when $S_T = K_2$ at expiry.
\textbf{Profit region:} $K_1 < S_T < K_3$.
\end{definitionbox}

The butterfly is a bet that the underlying will stay close to $K_2$. Net delta $\approx 0$ when $S_T = K_2$; net vega $< 0$ (profits from falling volatility). It is a classic \emph{short volatility} structure: you sell the possibility of large moves.

\subsection{Smile Arbitrage in Prosperity}

When the voucher IV surface is non-flat across strikes (e.g.\ $\hat{\sigma}(9500) \neq \hat{\sigma}(10000)$), a \textbf{strike arbitrage} is possible without taking on directional or volatility risk:

\begin{keyconceptbox}{Smile Arbitrage Recipe}
\begin{enumerate}[label=\arabic*.]
  \item \textbf{Identify mispriced strikes:} Compute $\hat{\sigma}(K)$ for all available strikes. A strike where $\hat{\sigma}(K)$ is anomalously high (low) relative to neighbours is overpriced (underpriced).
  \item \textbf{Construct a spread:} Buy the underpriced strike, sell the overpriced strike. Example: buy $K_1 = 9500$ call, sell $K_2 = 9750$ call (a bull call spread).
  \item \textbf{Delta-hedge the spread:} The net delta of the spread is $\Delta_1 - \Delta_2$. Short $(\Delta_1 - \Delta_2)$ units of the underlying to make the portfolio delta-neutral. Rebalance each timestamp.
  \item \textbf{P\&L source:} The hedge makes the position insensitive to $S_T$ direction. Profit comes from the IV differential collapsing (i.e.\ the market corrects the relative mispricing).
\end{enumerate}

\textbf{Greeks of the delta-hedged spread:}
\[
  \Delta_{\text{net}} = 0,\quad
  \mathcal{V}_{\text{net}} = \mathcal{V}(K_1) - \mathcal{V}(K_2) \neq 0,\quad
  \Gamma_{\text{net}} = \Gamma(K_1) - \Gamma(K_2) \neq 0
\]
The residual vega and gamma remain. If $K_1 < K_2$ and both are near ATM, vega$_{\text{net}} \approx 0$ (vega is symmetric for calls near ATM) — the spread is nearly \emph{vega-neutral}. The remaining gamma exposure makes the strategy profitable if realized vol differs from the IV differential.
\end{keyconceptbox}

\begin{warningbox}{No-Arbitrage Constraint on the Smile}
Put-call parity and the absence of butterfly arbitrage impose constraints on any valid IV surface:
\begin{itemize}
  \item \textbf{Butterfly no-arbitrage:} $C(K_1) - 2C(K_2) + C(K_3) \geq 0$ for $K_1 < K_2 < K_3$ equally spaced. Equivalently, $\partial^2 C / \partial K^2 \geq 0$: the call price is convex in strike.
  \item \textbf{Vertical spread no-arbitrage:} $C(K_1) \geq C(K_2)$ for $K_1 < K_2$ (lower strike calls are worth at least as much).
\end{itemize}
If market prices violate these, an \emph{outright} arbitrage exists (not just relative mispricing). In Prosperity, bot prices may temporarily violate these constraints — check for them each timestamp.
\end{warningbox}

% ── Volatility Smile & Skew ─────────────────────────────────────────────
\section{The Volatility Smile and Skew}

Black-Scholes assumes a single constant $\sigma$ for all strikes and maturities. Empirically, this is false: inverting BS at different strikes produces different implied volatilities, tracing out a curve $K \mapsto \hat{\sigma}(K)$ — the \textbf{volatility smile} (or skew).

\begin{definitionbox}{Implied Volatility Surface}
The \textbf{implied volatility surface} $\hat{\sigma}(K, T)$ is the function that makes Black-Scholes match the observed market price at every strike $K$ and maturity $T$:
\[
  C_{\text{market}}(K, T) \;=\; \text{BS}\!\left(S,\, K,\, T,\, r,\; \hat{\sigma}(K,T)\right)
\]
If BS were exact, $\hat{\sigma}(K,T)$ would be constant. Non-flatness reveals BS model misspecification.
\end{definitionbox}

\subsection{Three Canonical Shapes}

\begin{enumerate}[label=\arabic*.]
  \item \textbf{Smile (symmetric):} IV is U-shaped in $K$ — higher at deep OTM and deep ITM strikes than at the money. Common in currency markets, where large moves in either direction are feared.
  \item \textbf{Skew / Smirk (left-skewed):} IV rises monotonically as $K$ falls. Dominant in equity index options. OTM puts (crash insurance) carry a premium: the market prices in the possibility of sharp downward moves that exceed what GBM predicts.
  \item \textbf{Term structure:} IV varies with $T$ at fixed $K$. Typically upward-sloping (more uncertainty over longer horizons), but can invert when a known event (earnings, macro announcement) falls within the near-term expiry window.
\end{enumerate}

\subsection{Why the Smile Exists: BS Lognormal Assumption Fails}

The root cause is that real asset returns are \emph{not} lognormal:
\begin{itemize}
  \item \textbf{Fat tails (leptokurtosis):} Extreme moves occur more often than Gaussian. OTM options are underpriced by BS $\Rightarrow$ market charges a premium $\Rightarrow$ higher IV at wings.
  \item \textbf{Negative skewness (equity):} Crashes are larger and faster than rallies. OTM puts are systematically overpriced $\Rightarrow$ higher IV at low strikes.
  \item \textbf{Stochastic volatility (Heston model):} $\sigma_t$ is itself a random process. When volatility and price are negatively correlated ($\rho < 0$), the IV surface shows left skew.
  \item \textbf{Jump risk (Lévy/Merton models):} Discrete price jumps fatten the tails and generate IV smiles without requiring stochastic vol.
\end{itemize}

\begin{warningbox}{Implication for Prosperity Vouchers}
When you compute IV for each voucher strike (\texttt{VOLCANIC\_ROCK\_VOUCHER\_9500} through \texttt{\_10500}), check whether $\hat{\sigma}(K)$ is flat. A flat smile means BS with your estimated $\sigma$ is consistent across all strikes and a single fair-value model suffices. A non-flat smile means:
\begin{itemize}
  \item Either the market is mispricing some strikes relative to others (tradeable), or
  \item Your $\sigma$ estimate is wrong for some strikes (e.g.\ realized vol is not constant), or
  \item The underlying distribution genuinely has skew/kurtosis (adjust your model).
\end{itemize}
In Prosperity, smile arbitrage is possible: buy the underpriced strike, sell the overpriced strike on the same underlying, and delta-hedge both legs. The net position is volatility-neutral and direction-neutral if Greeks are matched.
\end{warningbox}

\textbf{ML analogy:} The volatility smile is the \emph{calibration residual map} of the lognormal model. Just as a linear classifier's error varies across the input space when data is nonlinearly separable, BS's single-$\sigma$ model is misspecified when returns are fat-tailed. The smile makes the residual visible by reprojecting it back onto the model parameter space (one $\sigma$ per strike).

\begin{takeawaybox}{Chapter 34 Summary}
\begin{itemize}
  \item An option's payoff is a deterministic function of the underlying: $\max(S_T - K, 0)$ for a call.
  \item Black-Scholes prices European options under GBM: $C = S\Phi(d_1) - Ke^{-rT}\Phi(d_2)$. Set $r = 0$ in Prosperity.
  \item Greeks are partial derivatives of $C$: $\Delta$ (vs $S$), $\Gamma$ (convexity), $\Theta$ (time decay), $\mathcal{V}$ (vs $\sigma$).
  \item Implied volatility is the market's $\sigma$ estimate: trade when IV diverges from your realized vol estimate.
  \item As expiry approaches, time value decays to zero: $C \to \max(S-K, 0)$ at $T-t = 0$.
\end{itemize}
\end{takeawaybox}

% ──────────────────────────────────────────────────────────────────────────

\chapter{Cross-Product Arbitrage \& Pairs Trading}

When multiple related products trade simultaneously, their prices must satisfy \emph{no-arbitrage} constraints. Any persistent violation is a riskless profit opportunity. This chapter covers static basket arbitrage, cointegration-based pairs trading, and how to implement both in Prosperity.

\section{The No-Arbitrage Principle}

\begin{definitionbox}{Arbitrage}
An \textbf{arbitrage} is a portfolio with:
\begin{enumerate}[label=\arabic*.]
  \item Zero initial cost (self-financing).
  \item Non-negative payoff in all states of the world.
  \item Strictly positive payoff in at least one state.
\end{enumerate}
Equivalently: a strategy with zero risk that makes money with positive probability.
\end{definitionbox}

No-arbitrage is the \emph{only} assumption needed to derive most derivative pricing results. Competitive markets eliminate arbitrage quickly; mispricings in Prosperity are bounded by bot quote refresh speed and position limits.

\section{Static Basket Arbitrage}

\begin{definitionbox}{Basket Asset}
A \textbf{basket} is a portfolio containing $w_i$ units of asset $i$. Under no-arbitrage:
\[
  P_{\text{basket}} = \sum_{i} w_i P_i
\]
Any persistent deviation $\Delta = P_{\text{basket}} - \sum_i w_i P_i \ne 0$ is an arbitrage opportunity.
\end{definitionbox}

\begin{examplebox}{Prosperity Basket Trade}
Suppose \texttt{PICNIC\_BASKET1} contains 3 units of \texttt{CROISSANTS} and 1 unit of \texttt{DJEMBES}. No-arbitrage requires:
\[
  P_{\text{basket}} = 3\,P_{\text{croissant}} + P_{\text{djembe}}
\]
If the observed basket price is $P_{\text{basket}}^* = 3\,P_{\text{croissant}} + P_{\text{djembe}} + 20$ (basket overpriced by 20), the trade is:
\begin{itemize}
  \item \textbf{Sell} 1 basket at $P_{\text{basket}}^*$.
  \item \textbf{Buy} 3 croissants and 1 djembe at component prices.
  \item Net PnL $= +20$ per unit, risk-free (assuming simultaneous execution).
\end{itemize}
\end{examplebox}

\textbf{ML analogy:} Basket arbitrage is \emph{linear constraint satisfaction}. The relationship $\mathbf{w}^{\top}\mathbf{P} = P_{\text{basket}}$ is a linear equation in prices. You monitor the residual $\varepsilon_t = P_{\text{basket},t} - \mathbf{w}^{\top}\mathbf{P}_t$ and trade when $|\varepsilon_t| > \epsilon_{\text{threshold}}$.

\subsection{Estimating Basket Weights}

If the exact composition is unknown, estimate weights by OLS regression:
\[
  P_{\text{basket},t} = w_0 + w_1 P_{1,t} + w_2 P_{2,t} + \cdots + \varepsilon_t
\]
The in-sample residual $\hat{\varepsilon}_t$ is the \emph{spread} to trade. Beware: OLS on non-stationary prices can produce spurious regression (Chapter 27). Always check that the spread $\hat{\varepsilon}_t$ is stationary (ADF test, $p < 0.05$) before trading it.

\section{Cointegration — The Foundation of Pairs Trading}

\begin{definitionbox}{Cointegration}
Two non-stationary series $X_t$ and $Y_t$ (both $I(1)$, i.e.\ integrated of order~1) are \textbf{cointegrated} if there exists a coefficient $\beta$ such that:
\[
  Z_t = Y_t - \beta X_t \quad \text{is stationary (}I(0)\text{)}
\]
$Z_t$ is called the \textbf{cointegrating residual} or \textbf{spread}.
\end{definitionbox}

Cointegration means two assets share a common stochastic trend: they can drift apart in the short run but are ``tethered'' together in the long run. The spread $Z_t$ mean-reverts around zero (or a constant), making it tradeable.

\textbf{ML analogy:} Cointegration is the time-series version of \emph{finding a linear combination of features that reduces dimensionality to a stationary signal}. The cointegrating vector $[1, -\beta]$ is a single principal component that removes the common trend from both series.

\subsection{Testing for Cointegration: Engle-Granger}

\begin{enumerate}[label=\arabic*.]
  \item Regress $Y_t$ on $X_t$: $\hat{\beta} = \text{OLS}(Y_t \sim X_t)$.
  \item Compute the residual: $\hat{Z}_t = Y_t - \hat{\beta}X_t$.
  \item Apply the ADF test (Chapter 27) to $\hat{Z}_t$.
  \item If ADF $p$-value $< 0.05$, the series are cointegrated at 5\% confidence.
\end{enumerate}

\begin{warningbox}{Non-Standard Critical Values for the Cointegration Residual ADF}
A critical subtlety: the ADF test applied to the cointegration residual $\hat{Z}_t$ does \textbf{not} use standard ADF critical values. The standard critical values assume the mean is known; here $\hat{Z}_t$ was constructed by OLS regression, so its mean is forced to zero by construction — but the hedge ratio $\hat{\beta}$ was \emph{estimated from the data}. This introduces an additional degree of freedom, making the null distribution of the test statistic different (and less conservative).

The \textbf{MacKinnon (1991, 2010)} tables provide critical values for the cointegration residual ADF. For a bivariate system with no intercept in the ADF regression, the 5\% critical value is approximately $-3.34$ (vs.\ the standard $-2.86$ for an ADF with intercept). Using standard ADF $p$-values will over-reject the null (falsely conclude cointegration more often).

\textbf{Practical fix:} In Python, use \texttt{statsmodels.tsa.stattools.coint()} which implements the Engle-Granger test with correct MacKinnon critical values, rather than applying a raw ADF test to the residual.
\end{warningbox}

\subsection{The Johansen Test: Multiple Cointegrating Vectors}

The Engle-Granger procedure tests only a single cointegrating relationship between two series. When you have $k \geq 3$ assets, there can be multiple cointegrating vectors (a full cointegration space of rank $r$). The \textbf{Johansen (1988) test} handles this.

\begin{definitionbox}{Johansen Cointegration Test}
For a $k$-dimensional VAR(p) system $\Delta\mathbf{X}_t = \mathbf{\Pi}\mathbf{X}_{t-1} + \sum_{i=1}^{p-1}\mathbf{\Gamma}_i\Delta\mathbf{X}_{t-i} + \boldsymbol{\varepsilon}_t$, the matrix $\mathbf{\Pi}$ has rank $r$ where $0 \leq r \leq k$:
\begin{itemize}
  \item $r = 0$: no cointegration (all series are independent random walks).
  \item $0 < r < k$: exactly $r$ cointegrating vectors (cointegration subspace of dimension $r$).
  \item $r = k$: the system is stationary (all series are $I(0)$).
\end{itemize}
$\mathbf{\Pi} = \mathbf{\alpha}\mathbf{\beta}^{\top}$ where $\mathbf{\beta}$ is the $k \times r$ matrix of cointegrating vectors and $\mathbf{\alpha}$ is the $k \times r$ matrix of adjustment speeds.
\end{definitionbox}

Two test statistics are used to determine $r$:
\begin{align}
  \lambda_{\text{trace}}(r) &= -T\sum_{i=r+1}^{k}\ln(1-\hat{\lambda}_i) \quad \text{(trace statistic)} \\
  \lambda_{\max}(r, r+1) &= -T\ln(1-\hat{\lambda}_{r+1}) \quad \text{(max-eigenvalue statistic)}
\end{align}
where $\hat{\lambda}_1 \geq \hat{\lambda}_2 \geq \cdots \geq \hat{\lambda}_k$ are the eigenvalues of a specific reduced-rank regression.

\begin{examplebox}{Johansen Test for Basket Products}
Suppose you have three products: PICNIC\_BASKET1, CROISSANTS, and DJEMBES. The Johansen test might reveal $r = 1$ cointegrating vector $\mathbf{\beta} = [1, -3, -1]$, meaning $Z_t = \text{BASKET}_t - 3\,\text{CROISSANT}_t - \text{DJEMBE}_t$ is stationary. The adjustment speeds $\alpha_i$ tell you which product adjusts fastest when the spread deviates.

Python: \texttt{statsmodels.tsa.vector\_ar.vecm.coint\_johansen(data, det\_order=0, k\_ar\_diff=1)}.
\end{examplebox}

\begin{warningbox}{Sample Size Requirements for Cointegration Tests}
Both Engle-Granger and Johansen tests have low power in small samples. Simulation studies suggest:
\begin{itemize}
  \item \textbf{Engle-Granger (2 variables)}: reliable at $n \geq 500$; detects cointegration with $\geq 80\%$ power at $n = 250$ when the OU mean-reversion half-life is $\leq 20$ periods.
  \item \textbf{Johansen ($k$ variables)}: requires $n \geq 5k^2$ observations for the trace statistic to approximate its asymptotic distribution reliably. With $k = 3$ products, this means $n \geq 45$ — easily met. But power to correctly identify $r$ cointegrating vectors degrades when the adjustment speeds $\alpha_i$ are small.
\end{itemize}
\textbf{In Prosperity}: with $\approx 2{,}000$ ticks per product per day (and 2 tutorial days), $n \approx 4{,}000$ is sufficient for Johansen with $k \leq 5$ products. However, if the true half-life of the spread exceeds $\approx 200$ ticks, the test will have low power and may miss cointegration that is economically real. Always complement the statistical test with an out-of-sample z-score backtests to confirm the spread actually reverts in practice.
\end{warningbox}

\section{The Pairs Trading Signal}

\begin{definitionbox}{Pairs Trading}
\textbf{Pairs trading} exploits the mean-reversion of the spread $Z_t = Y_t - \beta X_t$ between two cointegrated assets. The standardized spread:
\[
  z_t = \frac{Z_t - \bar{Z}}{\sigma_Z}
\]
generates a trading signal:
\begin{itemize}
  \item $z_t > +z_{\text{entry}}$: spread is wide, $Y$ is relatively expensive. \textbf{Sell} $Y$, \textbf{buy} $\beta$ units of $X$.
  \item $z_t < -z_{\text{entry}}$: spread is narrow, $Y$ is relatively cheap. \textbf{Buy} $Y$, \textbf{sell} $\beta$ units of $X$.
  \item $|z_t| < z_{\text{exit}}$: close position (spread has reverted).
\end{itemize}
\end{definitionbox}

The entry threshold $z_{\text{entry}}$ (typically 1.5–2.0) and exit threshold $z_{\text{exit}}$ (typically 0.0–0.5) are grid-searched on in-sample data. Tighter entries increase edge per trade; looser entries increase trade frequency.

\begin{warningbox}{Non-Stationarity of $\sigma_Z$: the Fixed-$\sigma$ Trap}
The z-score formula $z_t = (Z_t - \bar{Z})/\sigma_Z$ assumes $\sigma_Z$ is constant. In practice, $\sigma_Z$ changes with:
\begin{itemize}
  \item \textbf{Regime transitions}: during a regime change, the spread $Z_t$ widens persistently — not as a mean-reversion opportunity but as a structural shift. A fixed $\sigma_Z$ computed in-sample will generate a false $z_t > z_{\text{entry}}$ signal precisely when no reversion will occur.
  \item \textbf{Cointegration breakdown}: if the cointegrating relationship changes ($\beta$ shifts), $\sigma_Z$ is misspecified for all future timestamps.
  \item \textbf{Volatility clustering}: GARCH effects in returns cause $\sigma_Z$ to vary over time even within a stable regime.
\end{itemize}

\textbf{Correct implementation}: use a \emph{rolling} $\sigma_Z$ estimated over a lookback window $W$:
\[
  \hat{\sigma}_{Z,t} = \sqrt{\frac{1}{W-1}\sum_{k=t-W+1}^{t}(Z_k - \bar{Z}_{t,W})^2}
\]
\begin{lstlisting}[caption={Rolling z-score for pairs trading}]
from collections import deque
import math

class PairsTradingSignal:
    def __init__(self, beta: float, window: int = 500,
                 z_entry: float = 2.0, z_exit: float = 0.5):
        self.beta = beta
        self.window = window
        self.z_entry = z_entry
        self.z_exit = z_exit
        self._spreads: deque = deque(maxlen=window)

    def update(self, price_y: float, price_x: float) -> float | None:
        """Add new prices; return z-score if window is full, else None."""
        self._spreads.append(price_y - self.beta * price_x)
        if len(self._spreads) < self.window:
            return None
        spreads = list(self._spreads)
        mu = sum(spreads) / len(spreads)
        var = sum((s - mu) ** 2 for s in spreads) / (len(spreads) - 1)
        sigma = math.sqrt(var) if var > 0 else 1e-9
        return (spreads[-1] - mu) / sigma

    def signal(self, z: float | None) -> int:
        """Return +1 (buy Y), -1 (sell Y), or 0 (flat)."""
        if z is None:
            return 0
        if z > self.z_entry:
            return -1   # spread wide: sell Y, buy X
        if z < -self.z_entry:
            return +1   # spread narrow: buy Y, sell X
        return 0
\end{lstlisting}

\textbf{Trade-off:} a shorter window reacts faster to regime changes but produces noisier z-scores; a longer window is more stable but slow to adapt. Typical choices for Prosperity-length series: $W \in [200, 500]$ timestamps.
\end{warningbox}

\begin{examplebox}{Pairs Trade PnL Calculation}
Suppose $\beta = 2$, $Z_t$ has mean 0 and $\sigma_Z = 5$. We enter at $z_t = +2.0$ (spread at $+10$):
\begin{itemize}
  \item Sell 1 unit of $Y$ at $P_Y$; buy 2 units of $X$ at $P_X$.
  \item Spread at entry: $P_Y - 2P_X = \bar{Z} + 2\sigma_Z = 0 + 10 = 10$.
  \item Spread reverts to 0 (exit): $P_Y' - 2P_X' = 0$.
  \item PnL: $(P_Y - P_Y') - 2(P_X' - P_X) = -(P_Y' - P_Y) + 2(P_X - P_X') = -(Z_{\text{exit}} - Z_{\text{entry}}) = -(0 - 10) = +10$.
\end{itemize}
Expected PnL per trade $\approx 2\sigma_Z = 10$ when entering at $z = 2$.
\end{examplebox}

\section{Implementation in Prosperity}

\begin{keyconceptbox}{Cross-Product Strategy Template}
\begin{enumerate}[label=\arabic*.]
  \item \textbf{Compute the spread:} $Z_t = P_{Y,t} - \hat{\beta}\,P_{X,t}$.
  \item \textbf{Maintain a rolling mean and std:} use EMA for $\bar{Z}$ and $\sigma_Z$ (update each tick).
  \item \textbf{Standardise:} $z_t = (Z_t - \bar{Z}) / \sigma_Z$.
  \item \textbf{Entry signal:} If $z_t > z_{\text{entry}}$, \emph{sell} $Y$ up to $L_Y$, \emph{buy} $\beta$ units of $X$ up to $L_X$.
  \item \textbf{Exit signal:} If $z_t < z_{\text{exit}}$ (and currently short spread), close both legs.
  \item \textbf{Position limits:} Each leg is constrained independently by its own position limit $L$.
  \item \textbf{Persist state:} Store $\bar{Z}$, $\sigma_Z$ (or EMA buffers) in \texttt{traderData}.
\end{enumerate}
\end{keyconceptbox}

\begin{takeawaybox}{Chapter 35 Summary}
\begin{itemize}
  \item Basket arbitrage exploits $P_{\text{basket}} \ne \sum_i w_i P_i$. Trade the mispriced leg.
  \item Cointegration: two $I(1)$ series share a common trend; their spread is $I(0)$ (stationary).
  \item Pairs trading: standardize the spread, enter when $|z_t| > z_{\text{entry}}$, exit at $z_{\text{exit}}$.
  \item Always verify stationarity of the spread (ADF test) before committing to a pairs strategy.
  \item OLS gives $\hat{\beta}$ and spread; use EMA to track moving mean and std in production.
\end{itemize}
\end{takeawaybox}

% ──────────────────────────────────────────────────────────────────────────

\chapter{Conversion Mechanics}

From Round~1 onward, the \texttt{run()} method can return a non-zero \texttt{conversions} integer. This allows you to convert between a domestic product and an equivalent foreign product at a known cost. Understanding when to convert is a bounded arbitrage problem.

\section{The Conversion API in Prosperity}

The \texttt{run()} method signature is:
\begin{lstlisting}
def run(self, state: TradingState):
    # ... generate orders ...
    conversions: int = 0   # positive = buy foreign / sell domestic
    return result, conversions, trader_data
\end{lstlisting}

A positive \texttt{conversions} value means: buy \texttt{conversions} units of the foreign product and receive the equivalent domestic units. A negative value does the reverse. The exchange rate is not 1:1 — it includes fees.

\section{Transport Costs and Tariffs}

The \texttt{TradingState.observations} object provides the cost components for any active conversion product:

\begin{center}
\renewcommand{\arraystretch}{1.4}
\begin{tabular}{lll}
\toprule
\textbf{Field} & \textbf{Type} & \textbf{Meaning} \\
\midrule
\texttt{bidPrice} & float & Foreign market bid (price if selling to foreign) \\
\texttt{askPrice} & float & Foreign market ask (price if buying from foreign) \\
\texttt{transportFees} & float & Fixed fee per unit converted \\
\texttt{exportTariff} & float & Fee per unit when exporting (selling to foreign) \\
\texttt{importTariff} & float & Fee per unit when importing (buying from foreign) \\
\bottomrule
\end{tabular}
\end{center}

\section{The Conversion PnL Formula}

When you \emph{import} (buy from foreign, receive domestic product):
\begin{align}
  \text{Cost per unit} &= \texttt{askPrice} + \texttt{transportFees} + \texttt{importTariff} \\
  \text{Net PnL per unit} &= P_{\text{domestic bid}} - (\texttt{askPrice} + \texttt{transportFees} + \texttt{importTariff})
\end{align}

When you \emph{export} (sell to foreign, give up domestic product):
\begin{align}
  \text{Revenue per unit} &= \texttt{bidPrice} - \texttt{transportFees} - \texttt{exportTariff} \\
  \text{Net PnL per unit} &= (\texttt{bidPrice} - \texttt{transportFees} - \texttt{exportTariff}) - P_{\text{domestic ask}}
\end{align}

\begin{definitionbox}{Conversion Threshold}
It is profitable to \textbf{import} (set \texttt{conversions} $> 0$) when:
\[
  P_{\text{dom}}^{\text{bid}} > P_{\text{for}}^{\text{ask}} + c_{\text{import}}, \quad c_{\text{import}} = \texttt{transportFees} + \texttt{importTariff}
\]
It is profitable to \textbf{export} (set \texttt{conversions} $< 0$) when:
\[
  P_{\text{for}}^{\text{bid}} - c_{\text{export}} > P_{\text{dom}}^{\text{ask}}, \quad c_{\text{export}} = \texttt{transportFees} + \texttt{exportTariff}
\]
\end{definitionbox}

\section{Optimal Conversion Threshold and Sizing}

Convert as many units as profitable (up to position limit). The expected PnL is linear in quantity:
\[
  \text{Total PnL} = n \cdot \left(P_{\text{dom}}^{\text{bid}} - P_{\text{for}}^{\text{ask}} - c_{\text{import}}\right)
\]
So convert the maximum $n^* = \min(\text{position limit}, \text{available capacity})$ whenever the per-unit edge $> 0$.

\begin{examplebox}{Conversion Decision}
Suppose observations show: \texttt{askPrice} $= 12{,}000$, \texttt{transportFees} $= 50$, \texttt{importTariff} $= 30$. Total cost per imported unit $= 12{,}080$. If the domestic best bid is $12{,}100$, the edge per unit is $+20$. Import as many units as your position limit and current inventory allow.
\end{examplebox}

\begin{keyconceptbox}{Conversion as Bounded Arbitrage}
Conversion exploits a \emph{cross-market mispricing} bounded by the total cost $c$. Unlike unlimited order-book arbitrage, the profit per unit is fixed (not determined by market depth). The main risk is that prices move between when you decide to convert and when the conversion is processed. In Prosperity, conversions are instantaneous, so the only risk is the spread between observation data and the live order book.
\end{keyconceptbox}

\begin{takeawaybox}{Chapter 36 Summary}
\begin{itemize}
  \item \texttt{conversions} $> 0$: import (buy foreign, receive domestic). Profitable when $P_{\text{dom}}^{\text{bid}} > P_{\text{for}}^{\text{ask}} + c_{\text{import}}$.
  \item \texttt{conversions} $< 0$: export (sell domestic, buy foreign). Profitable when $P_{\text{for}}^{\text{bid}} - c_{\text{export}} > P_{\text{dom}}^{\text{ask}}$.
  \item Total cost $c =$ \texttt{transportFees} $+$ applicable tariff.
  \item Convert maximum size whenever per-unit edge $> 0$ (linear PnL in quantity).
  \item Conversion PnL is independent of order book depth — only observation data and domestic mid matter.
\end{itemize}
\end{takeawaybox}

% ──────────────────────────────────────────────────────────────────────────

\chapter{Observation Data \& Exogenous Signals}

From Round~1 onward, \texttt{TradingState.observations} provides real-world data (sunlight index, sugar price, etc.). These \emph{exogenous features} may determine a product's fundamental value. Trading on them requires the same regression framework you already know from ML.

\section{What is Observation Data in Prosperity?}

The \texttt{ObservationRow} struct (from \texttt{prosperity3bt}) contains:

\begin{center}
\renewcommand{\arraystretch}{1.4}
\begin{tabular}{lll}
\toprule
\textbf{Field} & \textbf{Example} & \textbf{Interpretation} \\
\midrule
\texttt{sunlightIndex} & 48.7 & Hours of sunlight per day (affects crop yields) \\
\texttt{sugarPrice} & 330.0 & World sugar price (affects confectionery costs) \\
\texttt{bidPrice} / \texttt{askPrice} & — & Foreign exchange price (for conversion products) \\
\texttt{transportFees} & 3.5 & Per-unit shipping cost \\
\texttt{importTariff} / \texttt{exportTariff} & 2.1 & Per-unit border cost \\
\bottomrule
\end{tabular}
\end{center}

Access in trader:
\begin{lstlisting}
obs = state.observations.conversionObservations.get("MAGNIFICENT_MACARONS")
if obs:
    sunlight = obs.sunlightIndex
    sugar    = obs.sugarPrice
\end{lstlisting}

\section{Linear Regression on Exogenous Features}

If a product's fair value $\text{FV}_t$ depends linearly on observable $\mathbf{x}_t = [x_1, x_2, \ldots]$:
\[
  \text{FV}_t = \alpha + \boldsymbol{\beta}^{\top}\mathbf{x}_t + \varepsilon_t
\]
then regressing historical mid-prices on observations (EDA step) gives $\hat{\alpha},\hat{\boldsymbol{\beta}}$. In production, compute the predicted FV:
\[
  \widehat{\text{FV}}_t = \hat{\alpha} + \hat{\boldsymbol{\beta}}^{\top}\mathbf{x}_t
\]
and market-make around $\widehat{\text{FV}}_t$ exactly as in Chapters 15–17.

\begin{examplebox}{Magnificent Macarons (Prosperity 3)}
In Prosperity~3, the fair price of \texttt{MAGNIFICENT\_MACARONS} was a function of sunlight and sugar price:
\[
  \text{FV}_t \approx a_0 + a_1 \cdot \texttt{sunlightIndex}_t + a_2 \cdot \texttt{sugarPrice}_t
\]
EDA: regress historical macaron mid-prices on the two observation columns. The OLS coefficients $a_1, a_2$ directly indicate how to shift your market-making mid quote as observations change.
\end{examplebox}

\section{Feature Engineering for Non-Linear Relationships}

If the relationship is non-linear, use feature transforms before regression:

\begin{center}
\renewcommand{\arraystretch}{1.3}
\begin{tabular}{lll}
\toprule
\textbf{Suspected relationship} & \textbf{Transform} & \textbf{New feature} \\
\midrule
Diminishing returns & $\log(x)$ & $x' = \ln x$ \\
Exponential growth & $\exp$ & $x' = e^{\beta x}$ \\
Threshold effect & Indicator & $x' = \mathbf{1}[x > c]$ \\
Polynomial & Polynomial & $x' = x^2, x^3, \ldots$ \\
Interaction & Product & $x' = x_1 \cdot x_2$ \\
\bottomrule
\end{tabular}
\end{center}

\begin{warningbox}{Do Not Over-Engineer Features}
In Prosperity, you have $\sim$10{,}000 ticks per round. With many features, OLS will overfit unless you regularise (Ridge/Lasso). Prefer the simplest model (1–2 features) that passes a $p$-value test ($< 0.05$) and OOS residual stationarity check (ADF). Every extra feature adds parameter uncertainty.
\end{warningbox}

\section{Dynamic Updating vs.\ Fixed Coefficients}

Two approaches:

\begin{enumerate}[label=\arabic*.]
  \item \textbf{Fixed coefficients}: Run OLS on all available historical data before the round. Hard-code $\hat{\alpha}, \hat{\boldsymbol{\beta}}$ into \texttt{trader.py}. Simple but cannot adapt to regime shifts.
  \item \textbf{Recursive OLS (RLS) / EMA update}: Update $\hat{\boldsymbol{\beta}}$ every tick using a forgetting factor $\lambda \in (0.9, 1)$. The full RLS algorithm (Ljung \& S\"{o}derstr\"{o}m, \textit{Theory and Practice of Recursive Identification}, 1983) maintains a gain matrix $\mathbf{P}_t$ (covariance of the estimate):
  \begin{align}
    \mathbf{K}_t &= \frac{\mathbf{P}_{t-1}\,\mathbf{x}_t}{\lambda + \mathbf{x}_t^{\top}\mathbf{P}_{t-1}\,\mathbf{x}_t} \\[4pt]
    \hat{\boldsymbol{\beta}}_t &= \hat{\boldsymbol{\beta}}_{t-1} + \mathbf{K}_t\!\left(y_t - \hat{\boldsymbol{\beta}}_{t-1}^{\top}\mathbf{x}_t\right) \\[4pt]
    \mathbf{P}_t &= \frac{1}{\lambda}\!\left(\mathbf{P}_{t-1} - \mathbf{K}_t\,\mathbf{x}_t^{\top}\mathbf{P}_{t-1}\right)
  \end{align}
  The forgetting factor $\lambda$ exponentially down-weights older observations: data from $k$ steps ago carries weight $\lambda^k$. Setting $\lambda = 1$ recovers ordinary (non-forgetting) OLS. In Prosperity, a simplified scalar approximation suffices when features are roughly orthogonal:
  \[
    \hat{\boldsymbol{\beta}}_{t+1} \approx \hat{\boldsymbol{\beta}}_t + \frac{\mathbf{x}_t}{\lambda / (1-\lambda) + \|\mathbf{x}_t\|^2}\!\left(y_t - \hat{\boldsymbol{\beta}}_t^{\top}\mathbf{x}_t\right)
  \]
  \textbf{ML analogy:} RLS is the online version of ridge regression with exponential sample weights. The gain $\mathbf{K}_t$ plays the same role as the learning rate in stochastic gradient descent, but it is computed analytically rather than tuned as a hyperparameter. More adaptive than fixed OLS but requires persisting $\mathbf{P}_t$ in \texttt{traderData} (a $p \times p$ matrix for $p$ features).
\end{enumerate}

\begin{takeawaybox}{Chapter 37 Summary}
\begin{itemize}
  \item Observation data provides exogenous features $\mathbf{x}_t$ that may drive a product's fair value.
  \item Regress historical mid-prices on observations (EDA) to get $\hat{\alpha}, \hat{\boldsymbol{\beta}}$. Use $\widehat{\text{FV}}_t$ as your market-making anchor.
  \item Prefer simple 1–2 feature models; validate OOS before using (ADF test on residuals).
  \item Fixed coefficients work for stable regimes; recursive OLS adapts to slow drifts.
  \item Access observations via \texttt{state.observations.conversionObservations[product]}.
\end{itemize}
\end{takeawaybox}

% ──────────────────────────────────────────────────────────────────────────

\chapter{Sealed-Bid Auction Theory}

Round~2 of Prosperity requires a \texttt{bid()} method alongside the usual \texttt{run()} method. This is a \emph{first-price sealed-bid auction}: all participants submit simultaneously; the highest bidder wins and pays their own bid. Optimal bidding requires game theory, not market-making.

\section{Auction Formats}

\begin{definitionbox}{Four Standard Auction Formats}
\begin{enumerate}[label=\arabic*.]
  \item \textbf{English (ascending):} Bids rise openly; highest bidder wins and pays their final bid.
  \item \textbf{Dutch (descending):} Price starts high and falls; first bidder to accept wins and pays that price.
  \item \textbf{First-price sealed-bid:} All bid simultaneously; highest bidder wins and pays their \emph{own} bid.
  \item \textbf{Second-price (Vickrey):} All bid simultaneously; highest bidder wins and pays the \emph{second-highest} bid.
\end{enumerate}
\end{definitionbox}

\textbf{Revenue equivalence theorem} (Vickrey 1961): Under symmetric i.i.d.\ private values and risk-neutral bidders, all four formats yield the same expected revenue to the seller. Despite this, the optimal bidding strategy differs radically between formats.

\section{First-Price Sealed-Bid Auctions}

The key feature of a first-price auction: you pay your own bid. This creates an incentive to \emph{shade your bid below your true value} — you can win at a lower price and keep more surplus.

\begin{definitionbox}{Private Value}
Each bidder $i$ has a \textbf{private value} $v_i$: the maximum they are willing to pay. Values $v_i \sim F$ are drawn i.i.d.\ from the same distribution $F$ (symmetric case). Bidder $i$ only knows $v_i$, not $v_j$ for $j \ne i$.
\end{definitionbox}

The payoff to bidder $i$ if they win at bid $b_i$:
\[
  \pi_i = v_i - b_i \geq 0
\]
If they lose: $\pi_i = 0$. The bidder maximises expected payoff:
\[
  \max_{b_i} (v_i - b_i) \cdot \Pr(\text{win at } b_i)
\]

\section{Bayesian Nash Equilibrium Bidding Strategy}

\begin{definitionbox}{Bayesian Nash Equilibrium (BNE) Bid}
In a first-price sealed-bid auction with $n$ symmetric bidders and $v_i \sim F$ i.i.d., the unique symmetric Bayesian Nash Equilibrium (BNE) strategy is:
\[
  b^*(v) = v - \frac{\int_0^v F(x)^{n-1}\,\dd x}{F(v)^{n-1}}
\]
For the special case $v_i \sim \mathrm{Uniform}[0, \bar{v}]$:
\[
  b^*(v) = \frac{n-1}{n}\,v
\]
You bid a fraction $\frac{n-1}{n}$ of your true value, shading by $\frac{1}{n}\,v$.
\end{definitionbox}

\begin{examplebox}{BNE with 5 Bidders, Uniform Values}
Suppose you estimate your value as $v = 1{,}000$, and there are $n = 5$ bidders.
\[
  b^*(1000) = \frac{4}{5} \times 1000 = 800
\]
You should bid 800, not 1000. If you bid 1000 and win, your payoff is 0. If you bid 800 and win (which happens when all four opponents bid less than 800), your payoff is $1000 - 800 = 200$.
\end{examplebox}

\textbf{ML analogy:} BNE is the \emph{Nash equilibrium} of a game — a fixed point where no agent can unilaterally improve by changing strategy. It is analogous to finding the saddle point of a minimax loss (as in GANs). The shading formula is the closed-form solution to a first-order condition (derivative of expected payoff with respect to $b$, set to zero).

\section{When the Number of Bidders is Unknown}

In practice, $n$ is uncertain. Approaches:
\begin{enumerate}[label=\arabic*.]
  \item \textbf{Estimate from history:} If previous round leaderboards show $N_{\text{teams}} \approx 50$ teams, but only the top few are competitive, assume effective $n \approx 5$--$10$.
  \item \textbf{Robust bidding:} Bid conservatively (low shade) if uncertain — a winning bid at slightly less profit is better than losing because you shaded too aggressively.
  \item \textbf{Sensitivity analysis:} Compute $b^*(v)$ for $n \in \{3, 5, 10, 20\}$ and choose a value that wins with high probability while retaining meaningful surplus.
\end{enumerate}

\section{Estimating Your Value $v$ in Prosperity}

The ``value'' in Prosperity Round~2 is typically the fair profit from the item being auctioned:
\begin{itemize}
  \item Estimate the expected resale value of the auctioned item using your market-making strategy.
  \item Subtract expected costs (inventory risk, position limit usage).
  \item The net expected profit is your private value $v$.
\end{itemize}
Then apply the BNE formula: $b^* = \frac{n-1}{n}\,v$.

\begin{lstlisting}
def bid(self) -> int:
    """
    First-price sealed-bid auction.
    Estimate private value v from expected market-making PnL,
    then shade by (n-1)/n.
    """
    n_bidders = 8        # estimated number of competitors
    v = self._estimate_value()   # your expected PnL from winning item
    b_star = ((n_bidders - 1) / n_bidders) * v
    return int(round(b_star))

def _estimate_value(self) -> float:
    # Replace with product-specific fair value calculation
    return 1000.0
\end{lstlisting}

\begin{takeawaybox}{Chapter 38 Summary}
\begin{itemize}
  \item First-price sealed-bid auction: highest bid wins, winner pays own bid.
  \item Bidding your true value $v$ is suboptimal (zero surplus if you win). Always shade.
  \item BNE strategy: $b^*(v) = \frac{n-1}{n}\,v$ for $n$ symmetric uniform-value bidders.
  \item With unknown $n$: estimate from team count; use sensitivity analysis to choose shade factor.
  \item Estimate $v$ as expected market-making PnL from the auctioned item, then shade.
\end{itemize}
\end{takeawaybox}

% ──────────────────────────────────────────────────────────────────────────

\chapter{Summary — From Tutorial to Round~5}

\section{The Full Knowledge Map}

This chapter places every concept in the encyclopedia into the competitive context of Prosperity~4. The table below maps each strategy family to the round where it becomes relevant, the theoretical tool required, and the chapter that covers it.

\begin{center}
\renewcommand{\arraystretch}{1.5}
\begin{longtable}{p{3.5cm}p{2cm}p{4.5cm}p{2.5cm}}
\toprule
\textbf{Strategy family} & \textbf{First active} & \textbf{Core tool} & \textbf{Covered in} \\
\midrule
Market-making & Tutorial & FV estimation, passive quoting, inventory skew & Chs.\ 4, 15, 16, 23 \\
Mean-reversion (AR1) & Tutorial & AR(1) FV correction, regime switching & Chs.\ 27, 29, 32 \\
Options / vouchers & Round 1 & Black-Scholes, implied vol, Greeks & Ch.\ 34 \\
Basket arbitrage & Round 1 & No-arbitrage, static replication & Ch.\ 35 \\
Pairs trading & Round 1+ & Cointegration, spread z-score & Ch.\ 35 \\
Conversion arbitrage & Round 1 & Conversion cost formula, threshold rule & Ch.\ 36 \\
Exogenous signals & Round 1 & OLS regression on observations & Ch.\ 37 \\
Sealed-bid auction & Round 2 & BNE shading formula, value estimation & Ch.\ 38 \\
\bottomrule
\end{longtable}
\end{center}

\section{The Decision Tree for a New Product}

When a new product appears in a round manual, apply this decision tree:

\begin{keyconceptbox}{New Product Analysis Framework}
\begin{enumerate}[label=\arabic*.]
  \item \textbf{Read the manual.} What is the product? Does it have a known formula (option payoff)? Is it a basket of known components? Is it linked to observation data?
  \item \textbf{Run EDA.} Plot the price series, compute ACF/PACF, ADF, Hurst exponent. Classify: stationary (mean-revert), $I(1)$ (trend/RW), or regime-switching.
  \item \textbf{Choose FV model:}
    \begin{itemize}
      \item Strongly stationary $\Rightarrow$ constant FV (like EMERALDS).
      \item Slowly drifting $\Rightarrow$ EMA.
      \item Mean-reverting returns $\Rightarrow$ AR(1) correction (like TOMATOES).
      \item Driven by observations $\Rightarrow$ linear regression on $\mathbf{x}_t$.
      \item Option/voucher $\Rightarrow$ Black-Scholes.
      \item Basket $\Rightarrow$ component sum; trade spread.
    \end{itemize}
  \item \textbf{Identify cross-product links.} Run cointegration tests against all other products. If cointegrated, pairs-trade the spread.
  \item \textbf{Check conversion opportunity.} Compute per-unit conversion PnL at every tick. Convert when edge $> 0$.
  \item \textbf{Grid-search parameters.} Quote offset $\delta$, skew $S_{\text{range}}$, entry/exit thresholds. Evaluate OOS Sharpe on the later training day.
  \item \textbf{Persist minimal state.} Only store what \texttt{run()} cannot recompute: rolling buffers, last values, cointegration coefficients.
\end{enumerate}
\end{keyconceptbox}

\section{The Competitive Edge — Where PnL Actually Comes From}

\begin{center}
\begin{tikzpicture}
\begin{axis}[
  xbar stacked,
  width=12cm, height=7cm,
  bar width=14pt,
  ytick={1,2,3,4,5,6,7,8},
  yticklabels={Passive spread capture, Aggressive take (mispriced bots), Conversion arb, Basket arb, Pairs trading mean-reversion, Volatility arb (options), Exogenous signal directional, Auction surplus},
  xmin=0, xmax=100,
  xlabel={Approximate PnL contribution (\%)},
  title={Typical PnL source breakdown by round type},
  legend style={at={(0.98,0.05)},anchor=south east},
  grid=major, grid style={dotted,gray!40},
]
\addplot+[xbar, fill=DefBlue!80] coordinates
  {(50,1)(40,2)(10,3)(15,4)(20,5)(15,6)(10,7)(5,8)};
\addplot+[xbar, fill=ExGreen!70] coordinates
  {(20,1)(20,2)(30,3)(30,4)(30,5)(25,6)(20,7)(15,8)};
\addplot+[xbar, fill=TakeawayOrange!80] coordinates
  {(30,1)(40,2)(60,3)(55,4)(50,5)(60,6)(70,7)(80,8)};
\legend{Tutorial, Round 1--2, Round 3--5}
\end{axis}
\end{tikzpicture}
\end{center}

\begin{takeawaybox}{Part VIII Summary}
\begin{itemize}
  \item \textbf{Options:} Fair value $= S\Phi(d_1) - Ke^{-rT}\Phi(d_2)$ (Black-Scholes, $r=0$ in Prosperity). Trade when market price diverges from implied vol.
  \item \textbf{Basket arb:} Trade spread $\Delta = P_{\text{basket}} - \sum_i w_i P_i$ when $|\Delta| > \epsilon$.
  \item \textbf{Pairs trading:} Standardise spread $z_t$; enter at $|z_t| > z_{\text{entry}}$; exit at $|z_t| < z_{\text{exit}}$.
  \item \textbf{Conversions:} Convert when per-unit edge $= P_{\text{dom}} - P_{\text{for}} - c > 0$. No inventory risk — instantaneous execution.
  \item \textbf{Exogenous signals:} Regress mid-price on observation features; use fitted FV for market-making.
  \item \textbf{Auction:} BNE bid $= \frac{n-1}{n}\,v$. Estimate $v$ from expected downstream PnL.
  \item \textbf{Meta-rule:} Read the round manual first. Then apply the decision tree: EDA $\to$ FV model $\to$ cross-product links $\to$ conversions $\to$ grid-search $\to$ persist minimal state.
\end{itemize}
\end{takeawaybox}

% End of Part VIII

% ============================================================
%
%  APPENDIX — WORKED EXERCISES
%
% ============================================================
\appendix

\chapter{Worked Exercises}

This appendix contains worked problems covering every major concept in the encyclopedia. Each exercise is self-contained. Work through the solution yourself before reading the answer.

% ──────────────────────────────────────────────────────────────────────────────
\section*{Exercise 1: Order Book Arithmetic}

\textbf{Problem.} At timestamp $t$, the EMERALDS order book shows:
\[
  \text{Bids: } \{9998: 5,\; 9997: 10,\; 9993: 15\}, \quad \text{Asks: } \{10002: 5,\; 10007: 15,\; 10008: 20\}
\]
Your current position is $+12$ units. Position limit is 20 units.

\begin{enumerate}[label=(\alph*)]
  \item What is the best bid? Best ask? Mid-price? Bid-ask spread?
  \item You want to buy aggressively (market order). What is the effective price if you buy 8 units?
  \item Your strategy bids at 9993 and offers at 10007. If both passive orders are fully filled, what is the gross PnL (ignoring position risk)?
  \item After filling your passive bid (15 units at 9993), what is your new position? Does it violate the position limit?
\end{enumerate}

\textbf{Solution.}

\textbf{(a)} Best bid $= 9998$. Best ask $= 10002$. Mid $= (9998+10002)/2 = 10000$. Spread $= 10002 - 9998 = 4$ ticks.

\textbf{(b)} Aggressive buy of 8 units walks the ask side:
\begin{itemize}
  \item Fill 5 units at 10002, then 3 units at 10007.
  \item Effective price $= (5 \times 10002 + 3 \times 10007) / 8 = (50010 + 30021)/8 = 80031/8 = 10003.875$.
  \item Market impact: you paid 3.875 ticks above mid.
\end{itemize}

\textbf{(c)} Passive bid of 15 at 9993 filled; passive ask of 15 at 10007 filled. Gross PnL per round-trip $= (10007 - 9993) \times 15 = 14 \times 15 = 210$ XIREC.

\textbf{(d)} Starting position $= +12$. Passive bid fills $+15$ units. New position $= 12 + 15 = +27$. Limit is 20 → \textbf{violation}. The exchange clips the fill. You can only buy $\min(15,\; 20 - 12) = 8$ units. Your strategy must account for remaining capacity ($\text{buy\_cap} = \text{limit} - \text{pos}$) before placing the passive order.

% ──────────────────────────────────────────────────────────────────────────────
\section*{Exercise 2: Kelly Criterion}

\textbf{Problem.} A market-making strategy has the following statistics per trade:
\begin{itemize}
  \item Win probability: $p = 0.55$.
  \item Average win: $b = 7$ ticks (gain 7 XIREC per unit when right).
  \item Average loss: $\ell = 5$ ticks (lose 5 XIREC per unit when wrong).
\end{itemize}

\begin{enumerate}[label=(\alph*)]
  \item Compute the Kelly fraction $f^*$.
  \item Your capital is $C = 100{,}000$ XIREC and your position limit is 20 units. What trade size does Kelly recommend? What does the position limit cap it at?
  \item Compute the expected value $\mathbb{E}[\Delta C]$ per trade at position size 10 units. Then at size 20 units.
  \item Why do practitioners use \emph{fractional Kelly} (e.g.\ $0.5 f^*$)?
\end{enumerate}

\textbf{Solution.}

\textbf{(a)} Full Kelly formula:
\[
  f^* = \frac{p}{1} - \frac{1-p}{b/\ell} \cdot \frac{1}{\text{(see below)}}
\]
More precisely, with win amount $b$ and loss amount $\ell$ (in same units):
\[
  f^* = \frac{p\,b - (1-p)\,\ell}{b\,\ell} \cdot \ell = \frac{p}{1/b} - \frac{1-p}{1/\ell}
\]
Standard discrete Kelly: $f^* = p - (1-p)/\text{odds}$ where odds $= b/\ell = 7/5 = 1.4$:
\[
  f^* = 0.55 - \frac{0.45}{1.4} = 0.55 - 0.3214 = 0.2286
\]
Kelly recommends risking 22.86\% of capital per trade.

\textbf{(b)} Capital-based size: $f^* \times C / \ell = 0.2286 \times 100{,}000 / 5 = 4{,}571$ units. Position limit caps this at 20 units. The position limit is the binding constraint — you are far below Kelly.

\textbf{(c)} $\mathbb{E}[\Delta C]$ at size $q$:
\[
  \mathbb{E}[\Delta C] = q\,(p \cdot b - (1-p) \cdot \ell) = q\,(0.55 \times 7 - 0.45 \times 5) = q\,(3.85 - 2.25) = 1.6\,q
\]
At $q = 10$: $\mathbb{E} = 16$ XIREC/trade. At $q = 20$: $\mathbb{E} = 32$ XIREC/trade. The maximum-size trade is always preferred here since we are operating well below the Kelly fraction. At the Kelly fraction, increasing size beyond $f^* C$ reduces expected log-growth.

\textbf{(d)} Fractional Kelly reduces variance of outcomes and protects against model estimation error. If the true win probability $p$ was estimated with error, the full Kelly bet could be catastrophic. At $0.5f^*$, the expected log-growth is approximately $0.5^2 \times$ the Kelly growth (quadratic loss), while drawdown risk is halved (linear reduction). This is the standard risk-adjusted Kelly approach used by professional traders.

% ──────────────────────────────────────────────────────────────────────────────
\section*{Exercise 3: Black-Scholes Pricing and Greeks}

\textbf{Problem.} A voucher is a European call on VOLCANIC ROCK with:
\[
  S = 10000, \quad K = 10250, \quad T = 2/3 \text{ day (day 1 of 3)}, \quad r = 0, \quad \sigma = 0.15 \text{ (daily)}
\]

\begin{enumerate}[label=(\alph*)]
  \item Compute $d_1$, $d_2$, and the call price $C$.
  \item Compute $\Delta$ and $\mathcal{V}$ (vega).
  \item The market quotes the voucher at $C^* = 120$. Is the voucher overpriced or underpriced? What trade do you place, and how do you delta-hedge it?
  \item If $S$ moves from 10000 to 10050, estimate the change in $C$ using the delta approximation. Then use the delta-gamma approximation.
\end{enumerate}

\textbf{Solution.}

\textbf{(a)} With $r = 0$:
\[
  d_1 = \frac{\ln(10000/10250) + 0.5 \times 0.15^2 \times (2/3)}{0.15\sqrt{2/3}} = \frac{-0.02469 + 0.00750}{0.12247} = \frac{-0.01719}{0.12247} \approx -0.1404
\]
\[
  d_2 = d_1 - \sigma\sqrt{T} = -0.1404 - 0.1225 = -0.2629
\]
\[
  C = 10000\,\Phi(-0.1404) - 10250\,\Phi(-0.2629) = 10000 \times 0.4441 - 10250 \times 0.3963 = 4441 - 4062 \approx 379 \text{ XIREC}
\]

\textbf{(b)} $\Delta = \Phi(d_1) = \Phi(-0.1404) = 0.4441$. Vega $= S\phi(d_1)\sqrt{T} = 10000 \times 0.3948 \times 0.8165 \approx 3223$ XIREC per unit of $\sigma$.

\textbf{(c)} Market price $C^* = 120 < C_{\text{fair}} = 379$: the voucher is severely \textbf{underpriced}. Trade: \textbf{buy the voucher} at 120. To delta-hedge: simultaneously \textbf{sell $\Delta = 0.4441$ units of VOLCANIC ROCK} per voucher. The hedge removes directional exposure and isolates the vol mispricing.

\textbf{(d)} Delta approximation: $\Delta C \approx \Delta \times \Delta S = 0.4441 \times 50 = 22.2$ XIREC. Delta-gamma: $\Delta C \approx \Delta \times \Delta S + \frac{1}{2}\Gamma\,(\Delta S)^2$. Gamma $= \phi(d_1)/(S\sigma\sqrt{T}) = 0.3948/(10000 \times 0.1225) = 3.22 \times 10^{-5}$. Gamma correction: $\frac{1}{2} \times 3.22\times10^{-5} \times 50^2 = 0.040$. Total: $22.2 + 0.04 = 22.24$ XIREC (gamma is negligible at this moneyness/horizon).

% ──────────────────────────────────────────────────────────────────────────────
\section*{Exercise 4: AR(1) Fair Value Correction}

\textbf{Problem.} The TOMATOES mid-price at timestamps $t-1, t, t+1$ is:
\[
  m_{t-2} = 5000,\quad m_{t-1} = 5008,\quad m_t = 5003
\]
The AR(1) coefficient in the tight-spread regime is $\phi = 0.05$.

\begin{enumerate}[label=(\alph*)]
  \item Using the AR(1) fair value formula $FV_t = m_t - \phi(m_t - m_{t-1})$, compute $FV_t$.
  \item What is the economic interpretation of this correction?
  \item With $\text{QUOTE\_OFFSET} = 5$ and zero inventory skew, where do you post your passive bid and ask?
  \item Now suppose inventory is $+40$ units with $\text{position\_limit} = 80$ and $\text{SKEW\_RANGE} = 1$. Recompute the passive quotes.
\end{enumerate}

\textbf{Solution.}

\textbf{(a)} $m_t - m_{t-1} = 5003 - 5008 = -5$. $FV_t = 5003 - 0.05 \times (-5) = 5003 + 0.25 = 5003.25$.

\textbf{(b)} The mid fell by 5 ticks from $t-1$ to $t$. AR(1) $\phi = -0.303$ in the statistical data means returns are negatively autocorrelated: a down-move is slightly likely to be followed by an up-move. The correction $+0.25$ ticks nudges the FV slightly above the current mid to reflect this expected mean-reversion. With $\phi = 0.05$ (grid-optimised), the correction is mild — only 0.25 ticks — because too-strong corrections misalign our quotes with bot crossing thresholds.

\textbf{(c)} $\text{passive\_bid} = \lfloor FV_t - 5 \rfloor = \lfloor 5003.25 - 5 \rfloor = \lfloor -1.75 \rfloor$... wait, we need the floor of $5003.25 - 5 = 4998.25$, so $\text{passive\_bid} = 4998$. $\text{passive\_ask} = \lceil 5003.25 + 5 \rceil = \lceil 5008.25 \rceil = 5009$.

\textbf{(d)} Inventory skew: $\text{skew} = (40/80) \times 1 = 0.5$ ticks. $\text{passive\_bid} = \lfloor 5003.25 - 5 - 0.5\rfloor = \lfloor 4997.75 \rfloor = 4997$. $\text{passive\_ask} = \lceil 5003.25 + 5 - 0.5\rceil = \lceil 5007.75\rceil = 5008$. The positive inventory tilts both quotes downward (making the sell more competitive, the buy less competitive) to incentivise reducing the long position.

% ──────────────────────────────────────────────────────────────────────────────
\section*{Exercise 5: Sealed-Bid Auction BNE}

\textbf{Problem.} In a first-price sealed-bid auction with $n = 5$ bidders and private values uniformly distributed on $[0, V_{\max}]$:

\begin{enumerate}[label=(\alph*)]
  \item Derive the symmetric BNE bidding function $b^*(v)$ from first principles.
  \item Compute the optimal bid if your private value is $v = 800$ and $V_{\max} = 1000$.
  \item What is the expected revenue for the seller?
  \item If you estimate $V_{\max} = 1000$ but are uncertain (std $= 200$), how should you adjust your bid?
\end{enumerate}

\textbf{Solution.}

\textbf{(a)} At BNE, all bidders use the same strategy $b(v)$. Bidder with value $v$ solves:
\[
  \max_{b}\;(v - b)\,\Pr(\text{win}) = (v - b)\,\Pr(b(v_j) < b \;\forall j \neq i) = (v - b)\,\left[\frac{b^{-1}(b)}{V_{\max}}\right]^{n-1}
\]
where $b^{-1}$ is the inverse bidding function. With symmetric strategy, $\Pr(\text{win}) = [v/V_{\max}]^{n-1}$ at the optimum (since all rivals with value $< v$ bid less). First-order condition:
\[
  \frac{\partial}{\partial b}[(v-b)\cdot(v/V_{\max})^{n-1}] = 0 \implies b = v - \frac{v}{n} = \frac{n-1}{n}\,v
\]
(This is only a sketch; the full derivation integrates the FOC as a first-order ODE.) The BNE is:
\[
  b^*(v) = \frac{n-1}{n}\,v
\]

\textbf{(b)} $n = 5$, $v = 800$: $b^* = \frac{4}{5} \times 800 = 640$.

\textbf{(c)} Expected revenue by the Revenue Equivalence Theorem equals $\mathbb{E}[\text{second-highest value among }n]$. For $U[0,1000]$ with $n=5$: $\mathbb{E}[\text{2nd order statistic}] = \frac{n-1}{n+1}V_{\max} = \frac{4}{6} \times 1000 = 666.67$ XIREC.

\textbf{(d)} Under uncertainty about $V_{\max}$, the bid function becomes $b^*(v) = \frac{n-1}{n}\,v$, which depends only on your own value $v$ and $n$ — not on $V_{\max}$. The BNE strategy is robust to uncertainty about the range of opponents' values, as long as the uniform distribution assumption holds. If you have asymmetric information about $V_{\max}$, adjust the estimate of $n$ (opponents who have high-value information may bid more aggressively).

% ──────────────────────────────────────────────────────────────────────────────
\section*{Exercise 6: Cointegration and Pairs Trading}

\textbf{Problem.} Over 100 timestamps, the prices of two products are:
\[
  X_t \sim I(1), \quad Y_t \sim I(1), \quad \hat{Z}_t = Y_t - 2\,X_t \sim I(0)
\]
The spread $\hat{Z}_t$ has sample statistics: $\bar{Z} = 10$, $\sigma_Z = 4$.

At time $t$, $Y_t = 520$, $X_t = 255$, so $Z_t = 520 - 2 \times 255 = 10 = \bar{Z}$.
\vspace{3pt}

At time $t+5$, $Y_{t+5} = 532$, $X_{t+5} = 256$, so $Z_{t+5} = 532 - 512 = 20$.

\begin{enumerate}[label=(\alph*)]
  \item Compute the z-score at time $t$ and $t+5$.
  \item At $t+5$, entry threshold is $z^* = 2.0$. Do you enter? If yes, describe the position.
  \item The spread at $t+10$ reverts to $Z_{t+10} = 12$. Compute the PnL if you entered at $t+5$ with 1 unit of $Y$ and 2 units of $X$.
  \item You run an ADF test on $\hat{Z}_t$ and get test statistic $-2.50$. Standard ADF 5\% critical value is $-2.86$. MacKinnon EG 5\% critical value is $-3.34$. What do you conclude?
\end{enumerate}

\textbf{Solution.}

\textbf{(a)} $z_t = (10 - 10)/4 = 0$. At $t+5$: $z_{t+5} = (20 - 10)/4 = 2.5$.

\textbf{(b)} $z_{t+5} = 2.5 > z^* = 2.0$: yes, enter. Direction: spread is wide ($Y$ relatively overpriced vs. $X$). Position: \textbf{sell} 1 unit of $Y$ at 532, \textbf{buy} 2 units of $X$ at 256. Net cash outflow: $-532 + 2 \times 256 = -532 + 512 = -20$ (short the spread at value 20).

\textbf{(c)} At entry: sold $Y$ at 532, bought $2X$ at 512 (combined cost of $X$ leg). At exit: buy back $Y$ (price: $Z_{t+10} + 2X_{t+10}$, but we just use spread logic): spread fell from 20 to 12. PnL $= Z_{\text{entry}} - Z_{\text{exit}} = 20 - 12 = 8$ XIREC (you shorted the spread when it was wide and it narrowed).

\textbf{(d)} The test statistic $-2.50$ fails both critical values. Under the standard ADF critical value ($-2.86$), you would fail to reject at 5\% — no cointegration. Under the MacKinnon EG critical value ($-3.34$), you also fail. The correct conclusion is: \textbf{no evidence of cointegration at 5\%}. Using standard ADF tables would have given the same conclusion here, but if the statistic were between $-3.34$ and $-2.86$ (e.g.\ $-3.10$), using standard tables would \emph{incorrectly} conclude cointegration while the correct MacKinnon tables would not.

% ──────────────────────────────────────────────────────────────────────────────
\section*{Exercise 7: Walk-Forward Grid Search}

\textbf{Problem.} You want to optimise the QUOTE\_OFFSET parameter $\delta \in \{3, 4, 5, 6, 7\}$ for TOMATOES. You have 2 days of data (Day $-2$: timestamps 0--99900; Day $-1$: timestamps 0--99900).

\begin{enumerate}[label=(\alph*)]
  \item Describe the correct walk-forward optimisation procedure. Why is in-sample optimisation alone invalid?
  \item You test all 5 values of $\delta$ on Day $-2$ and get PnLs: \{3: 45k, 4: 68k, 5: 82k, 6: 79k, 7: 61k\}. What is the best in-sample $\delta$?
  \item You apply $\delta = 5$ to Day $-1$ and get 74k. You apply $\delta = 4$ (second-best IS) and get 81k. What does this tell you about overfitting?
  \item How would you combine both days into a single decision about which $\delta$ to submit?
\end{enumerate}

\textbf{Solution.}

\textbf{(a)} Walk-forward procedure: (1) Fit/select parameters on the in-sample period (Day $-2$). (2) Evaluate the chosen parameters on the out-of-sample period (Day $-1$, never touched during selection). (3) Report the out-of-sample PnL as the unbiased performance estimate. In-sample optimisation alone is invalid due to \textbf{overfitting}: with 5 parameter values, the best in-sample result includes approximately $\log_2(5) \approx 2.3$ bits of noise from the data. The IS result is biased upward; OOS reveals the true edge.

\textbf{(b)} Best IS $\delta = 5$ with PnL 82k.

\textbf{(c)} On OOS, $\delta = 5$ yields 74k but $\delta = 4$ yields 81k. The optimal IS parameter did not transfer to OOS — this is mild overfitting. The OOS gap ($82k \to 74k$) is approximately 10\%, suggesting the parameter selection is somewhat noisy but not catastrophically overfit. The fact that $\delta = 4$ (a nearby value) performs best OOS suggests the true optimum is in the range $[4, 5]$ — both are reasonable choices.

\textbf{(d)} Combined decision framework: weight OOS performance more heavily than IS performance (as it is unbiased). $\delta = 4$: IS 68k + OOS 81k = average 74.5k. $\delta = 5$: IS 82k + OOS 74k = average 78k. On a harmonic or geometric mean basis, $\delta = 5$ may still dominate. However, the \textbf{risk-adjusted} choice depends also on variance: if $\delta = 4$ is more robust (smaller spread between IS and OOS), it may be preferable. In practice, submit $\delta = 5$ for the live round (best combined evidence) but note that $\delta = 4$ is nearly as good — if you can test a third day, do so.

\begin{takeawaybox}{Appendix A Summary}
\begin{itemize}
  \item Order book arithmetic: always track \texttt{buy\_cap} and \texttt{sell\_cap} before placing orders to avoid position limit violations.
  \item Kelly at position limits: in Prosperity, position limits bind before Kelly. Max out size; Kelly is not the binding constraint.
  \item Black-Scholes: with $r = 0$ and OTM options, $\Delta < 0.5$ and vega is large. Implied vol arbitrage dominates directional bets.
  \item AR(1) FV correction: the correction is small ($\phi = 0.05$ gives $< 1$ tick shift). Its value is in aggregate over thousands of trades.
  \item Auction: $b^* = \frac{n-1}{n}v$ regardless of $V_{\max}$. Your job is to estimate $v$ (your downstream PnL from winning).
  \item Cointegration: always use MacKinnon critical values (not standard ADF) for the Engle-Granger residual test.
  \item Walk-forward: OOS PnL is the only unbiased estimate. IS PnL is optimistic by construction.
\end{itemize}
\end{takeawaybox}

% End of Appendix

\end{document}
